\documentclass[11pt,a4paper]{article}
\makeatletter\def\input@path{{sections/paper_arithmetic/}{sections/paper/}}\makeatother
% LuaLaTeX. The optional font file is only a local build convenience.
\usepackage{luatexja-fontspec}
\IfFileExists{../local-fonts.tex}{\input{../local-fonts.tex}}{}
\usepackage[margin=24mm,headheight=15pt]{geometry}
\usepackage{amsmath,amssymb,amsthm,mathtools}
\usepackage{booktabs,tabularx,array,longtable,enumitem}
\usepackage{xcolor,tikz}
\usetikzlibrary{arrows.meta,calc,positioning}
\usepackage{microtype}
\usepackage[unicode,colorlinks=true,linkcolor=blue!45!black,citecolor=blue!45!black,urlcolor=blue!45!black]{hyperref}
\usepackage{fancyhdr}
\pagestyle{fancy}
\fancyhf{}
\fancyhead[L]{\small General type surfaces and Lagrangian fibrations}
\fancyhead[R]{\small v25}
\fancyfoot[C]{\thepage}
\renewcommand{\headrulewidth}{0.3pt}
\setlength{\parindent}{1em}
\setlength{\parskip}{0.25em}
\setlength{\emergencystretch}{2em}
\setlist{leftmargin=1.8em,itemsep=0.2em,topsep=0.3em}
\theoremstyle{plain}
\newtheorem{theorem}{定理}[section]
\newtheorem{proposition}[theorem]{命題}
\newtheorem{lemma}[theorem]{補題}
\newtheorem{corollary}[theorem]{系}
\theoremstyle{definition}
\newtheorem{definition}[theorem]{定義}
\newtheorem{remark}[theorem]{注意}
\newtheorem{problem}[theorem]{問題}
\newtheorem{example}[theorem]{例}
\renewcommand{\proofname}{証明}
\renewcommand{\refname}{参考文献}
\renewcommand{\contentsname}{目次}
\newcommand{\PP}{\mathbb P}
\newcommand{\CC}{\mathbb C}
\newcommand{\RR}{\mathbb R}
\newcommand{\ZZ}{\mathbb Z}
\newcommand{\QQ}{\mathbb Q}
\newcommand{\TT}{\mathbb T}
\newcommand{\cX}{\mathcal X}
\newcommand{\eps}{\varepsilon}
\newcommand{\dd}{\,\mathrm d}
\newcommand{\figeight}{(S^1\vee S^1)}
\newcommand{\Bl}{\operatorname{Bl}}
\newcommand{\Cone}{\operatorname{Cone}}
\newcommand{\Sk}{\operatorname{Skel}}
\newcommand{\PD}{\operatorname{PD}}
\newcommand{\Diff}{\operatorname{Diff}}
\newcommand{\id}{\operatorname{id}}

\makeatletter
\def\input@path{{sections/paper_arithmetic/}{sections/paper/}}
\makeatother
\hypersetup{pdftitle={五次曲面の整数アフィン接着と高分母支持グラフ 第25稿},pdfauthor={}}
\title{五次曲面の整数アフィン接着と高分母支持グラフ\\[.4em]\large 四rayの完全符号分類，二単位枝の接着障害，三高係数の有限復元\\第25稿}
\author{}
\date{2026年10月28日}
\begin{document}
\maketitle
\begin{abstract}
五つの平面からなる五次曲面の退化に由来する標準 $K_5$ branch graph を固定し，
四直線 core の作用適合整数アフィン charts の接着を研究する．
局所の整数行和・pair products・正の basin weights と，整数 clutching，
対向分岐係数の相殺，非負 node 配分を同時に課す．
全ての pair product が狭義正である四ray整数データでは，正符号 graph は
完全 matching または三角形と孤立頂点に限られ，正 weights の存在は完全 matching と同値である．
道型 $P_4$ を正値恒等式で排除することにより，三高・四高を統一して扱う．
局所尺度の復元と大域 node 条件から，非空高分母支持を三型・二十六標識に帰着する．
三高 core 間の単位枝の node 配分は $\{5,0,0\}$ または $\{4,1,0\}$ であり，
二単位枝を持つ $G_{5,14}$ では両方とも $\{4,1,0\}$ に限られ，中央の四高 core の正 matching も固定される．
三高係数の同時非負性では非対称量 $\delta=\pm1$ を排除し，比例領域 $\delta=0$ を
二つの高係数が零の軸へ帰着する．この零軸の端点は低行の一整数 $s$ だけで完全有限復元できる．
退化面では元 pair 残余の正整数積表示により互いに素な領域を排除し，
固定 $s$ に加えて $x\ge1$ の固定 $b$ からの有限復元も得る．
これらの一般定理は任意の探索上限を用いない．
残る三支持型の大域接着，一般分母剛性，比例零軸・退化面・一般非対称領域の全体は未解決である．
対象は指定した cores，cones，marked punctures を保つ接着であり，
任意の Lagrangian torus fibration の分類を主張しない．
\end{abstract}
\tableofcontents
% Current arithmetic introduction, v25. Prior literature and stable ar17 labels retained.
\section{序論}
\label{ar17:introduction}

一般型曲面の Lagrangian torus fibration を退化から構成するとき，
局所の作用周期が整数格子を定めることに加え，その格子を異なる成分の間で接続する必要がある．
本稿は，五つの平面へ退化する五次曲面に現れる標準 $K_5$ branch graph と，
五つの四直線 $K_4$ cores を固定し，この整数接着問題を扱う．
局所の行和と pair products，正の basin weights，marked punctures の整数 clutching，
有限な node 予算を同時に追うことで，高分母 support を三つの五頂点 graph 型へ帰着する．

局所正値性の出発点は，三高・四高に共通する四ray整数データの全符号分類である．
全ての pair product が正なら，整数の行和と pair equations だけで，正符号 graph は
完全 matching または三角形と孤立頂点に限られる．
道 $P_4$ は元の残る pair 残余を正値恒等式に変形して排除する．
三角形型には gain が一以上の閉路があるので，正の weights が存在することは
完全 matching 型であることと同値となる．

大域の支持障害には，二高分母の正負枝の整数剛性と全尺度分類，
および外側 core の一つの pair による尺度の線形復元を用いる．
これらは book 型と，低辺の補 graph が $P_3$ となる $G_{5,15}$ を完全に排除する．
従って非空高分母支持は三型・二十六標識，最小次数三以上，三点連結に限られる．
同じ core に二本の単位枝が接すれば全分母が一となり，高分母全体の最大公約数も一となる．

残る支持型に単位枝があるとき，その両端はともに三高 core である．
三角形の必要条件を完全に解くと，その node 配分は
$\{5,0,0\}$ または $\{4,1,0\}$ に限られ，隣接高分母の少なくとも一つが三または五となる．
対向低係数の相殺は，少なくとも一つの外側 core に正の pair を要求する．
二本の単位枝を同時に使うと，$G_{5,14}$ の $\{5,0,0\}$ 配分は全尺度で排除され，
両単位枝はともに $\{4,1,0\}$ に限られる．
この制約は，残る接着問題を node の位置，小分母の位置，局所符号の整合から調べる入口となる．

三高係数の同時非負性については，非対称量 $\delta=sD_b-tD_a$ と
尺度の整数格子を用いる．$\delta=\pm1$ は全域で排除され，比例領域 $\delta=0$ は
二つの高係数が零の軸へ縮む．その零軸では二整数を固定した真の四次式による有限復元を得る．
端点 $n=0,z$ では，整数根の位置をさらに制限して低行の一整数 $s$ だけで有限復元する．
尺度残余の退化面でも固定 $s$ による有限復元があり，元残余の積表示から
互いに素な領域の排除と，$x\ge1$ に対する固定 $b$ の有限復元を得る．
固定した入力ごとの完全性を，入力自体への一様上界や一般不存在とは区別する．
残る三型の大域接着，一般分母剛性，三高係数の一般の非対称問題は未解決である．

\subsection{整数データと中心となる結果}

cores を $V=\{0,1,2,3,4\}$ で標識し，共有 branch $ij$ の分母を
$q_{ij}=q_{ji}>0$，core $i$ の rays $j,k$ の間の整数係数を $p^i_{jk}$ とする．
局所の基本方程式は
\[
 \sum_{k\ne i,j}p^i_{jk}=-q_{ij},\qquad
 p^i_{jk}p^i_{kj}=q_{ij}q_{ik}-1.
\]
さらに正の実数 $w^i_j$ が存在し，
\[
 w^i_k-\frac{p^i_{jk}}{q_{ij}}w^i_j>0
\]
となることを要求する．この不等式は puncture basins の正面積に対応する．
共有 branch では
\[
 q_{ij}\mid p^i_{jk}+p^j_{ik},\qquad
 n_{ij;k}=1-\frac{p^i_{jk}+p^j_{ik}}{q_{ij}}\in\ZZ_{\ge0},
\]
\[
 D^i_j+D^j_i=0,\qquad
 D^i_j=-q_{ij}-\sum_{k\ne i,j}p^i_{kj}
\]
を課す．正確な定義は節~\ref{ar17:setup}に置く．
局所方程式は分母の奇数性を含意し，両端の row sums から各 branch の node 総数は五となる．
$q_{ij}>1$ の辺からなる graph を高分母 support $H$ と呼ぶ．

\begin{theorem}[四ray整数データの符号と正値性]
\label{ar24:four-ray-positive}\label{ar18:main-matching}
四つの正整数 $q_j$ と整数 $p_{jk}$ が
\[
 \sum_{k\ne j}p_{jk}=-q_j,\qquad p_{jk}p_{kj}=q_jq_k-1>0
\]
を満たすとする．正の pair からなる無向 graph $G^+$ は
$2K_2$ または $K_3\sqcup K_1$ に限られる．
正の basin weights が存在すること，$G^+$ が孤立頂点を持たないこと，
$G^+=2K_2$ であることは同値である．
三高の場合には，weights を仮定しない三角形型は低 ray を含む．
\end{theorem}

定理~\ref{fh25:no-path}，\ref{fh25:all-signs} と系~\ref{fh25:matching} がこれを示す．
孤立頂点からは，各行の二整数の和が一となる gap を作る．
非正の gap を選ぶと gain 一以上の有向閉路を得るが，pair product は長さ二を禁じるため，
残る三頂点は三角形を作る．孤立頂点がなければ，負の行和により候補は
$2K_2,P_4,C_4$ であり，$C_4$ は負の補辺の pair product で排除される．
$P_4$ では五つの pair と四行を解いた後，第六 pair の残余を
非負項の和へ変形し，狭義正の一項によって矛盾を得る．
この排除は低分母一を使わず，三高・四高を同時に扱う．
完全 matching の各辺では，非空の有理開区間から weight 比を独立に選べる．

正値性を仮定しない三角形型は実在する
（注意~\ref{mh18:positivity-essential}）．従って整数の行和・pair equations と
正 weights の役割を区別する．整数 clutching，node 配分，対向 $D$ 相殺は
この局所分類の仮定に含まれない．

\begin{theorem}[高分母 support の三型への帰着]
\label{ar17:main-support}
定義~\ref{ar17:admissible}の五次の算術的接着データを考える．
$H\ne\varnothing$ なら，$H$ は五頂点全てを使う三点連結 graph であり，
最小次数三以上，少なくとも八辺を持ち，$b_1(H)\ge4$ となる．
補 graph は高々二辺の matching であり，残る必要候補は
$G_{5,14},G_{5,16},G_{5,17}$ の三型，標識を保って二十六個である．
高分母全体の最大公約数は一である．
一つの core に接する二本の枝が単位分母なら，全分母が一となる．
三本以上の相異なる枝が単位分母である場合にも同じ結論が成り立つ．
\end{theorem}

この定理は，本稿の graph 障害を合わせた結論である
（系~\ref{sf23:remaining-support}，\ref{gp23:primitive}）．
三型という数は整数接着データや幾何的 atlas の存在数ではなく，
証明した必要条件を通過する support の数を表す．
本稿の対象は指定した cores，標準 cones，marked punctures を保つ整数 atlas の接着であり，
一般型曲面上の任意の Lagrangian torus fibration の分類ではない．

\begin{theorem}[単位枝の node 配分と小分母]
\label{ar24:unit-edge-boundary}
五次の算術的接着データにおいて，単位枝 $ij$ の両端がともに正の三高 core であるとする．
三つの方向に割り当てる node 数の多重集合は $\{5,0,0\}$ または $\{4,1,0\}$ に限られる．
最大 node の位置 $c$ は，両端の低 row がともに負となる唯一の第三 core である．
\begin{enumerate}
\item $n_{ij;c}=5$ なら，$(q_{ic},q_{jc})$ は $(3,7),(5,5),(7,3)$ のいずれかである．
この core の pair $ij$ は負，他の二つの外側 pair は正となる．
\item $n_{ij;c}=4$ なら $\min(q_{ic},q_{jc})=3$ である．
この core の pair が負なら両分母は三である．
\item いずれの場合も，三つの外側 cores の少なくとも一つで pair $ij$ は正である．
\end{enumerate}
さらに $G_{5,14}$ では五node配分自体が排除され，二つの単位枝の配分はともに
$\{4,1,0\}$ となる．$G_{5,16}$ の五node配分までを排除したわけではない．
\end{theorem}

定理~\ref{ut24:unit-edge-boundary}，命題~\ref{ut24:five-node-positive}，
\ref{ut24:positive-witness} と系~\ref{ut24:g514-five-node} がこの主張を与える．
二つの低 row の正位置は node 非負性によって異なる．共通負位置の絶対値を $s,s'$ とすると
$n_{ij;c}=1+s+s'\le5$ であり，第三 core の pair を整数として解くと $s=s'=1$ が除かれる．
反対符号位置も四つの通常尺度族と全ての小例外へ分類できる．
五nodeの場合は対向低 $D$ の相殺が二つの column 差を一に固定する．
三つの外側 pair が全て負と仮定した場合にも，完全表は一方の低データを
\[
 \text{row}=(-2,-1,2),\qquad \text{column}=(-1,-2,2r),\qquad r\ge1
\]
に強制する．その局所完成方程式は
$(3A+6r+1)(A+6r+1)=0$ となり，整数性または非零 entry の条件に反する．
$G_{5,14}$ の五node排除は定理~\ref{ug25:no-five} による．
第二単位枝と matching の伝播により一つの core の二高分母を $(3,3)$ または $(3,7)$ に固定し，
正 pair の全因子 $LM=8,20$ を調べる．残る低 column の線形式が恒等零になる二場合も，
元の逆 entry の非整数性によって排除する．
さらに定理~\ref{cm25:central} は，単位枝を $01,23$ としたとき，
唯一の四高 core $4$ の正 matching が $\{01,23\}$ に限られることを示す．
交差する二 matching は四つの正規符号配置へ帰着し，元の clutching と行和が
分母による一の整除または零 entry を強制するため排除される．
小分母の出現は，全分母の一様上界や残る支持型の排除を意味しない．

\subsection{matching 性が大域の接着に与える制約}

三高分母 core の低 row の符号が定まることにより，その row を隣接する core の
node 数から復元する方法が有効になる．二高分母の標準形と共通低 puncture の比較も合わせ，
次の三型を排除する：
\[
 \begin{split}
 G_{5,5}&=\{01,02,03,12,14,34\},\\
 G_{5,10}&=\{01,02,03,12,13,24,34\},\\
 G_{5,12}&=K_4\text{ に一本の葉辺を加えた graph}.
 \end{split}
\]
$G_{5,5}$ は三角形と四角形が一辺を共有する house graph である
（定理~\ref{hs18:no-house}）．この型では二価 cores の比較と二本の低辺の $D$ 相殺だけで十分であり，
三高分母の完全 matching 定理そのものは必要としない．
$G_{5,10}$ の排除では，唯一の二価 core の標準形を，
共通低 puncture における通常接着と正負の例外接着に分ける
（定理~\ref{gten18:no-support}）．
通常型は一つの中央模型に強制されるが，その反対 corner の signed divisor が存在しない．
正例外型は $D$ 相殺と奇分母の parity に反する．
負例外型では，固定分母十三の無限共鳴族を全て残したまま，
node 予算と三本の辺の $D$ 相殺から矛盾を導く．
この無限族を全て残して扱うことが，局所整数模型の分類から大域の接着障害へ進む際の要点となる．

$G_{5,12}$ では，第五 core が一高分母の葉となるため，
全低分母 hub の $K_4$ 証明をそのまま適用することはできない．
一高分母標準形の自由度を保ち，隣接する三高分母 cores の matching 性と
共有 punctures の整数条件を用いて排除する
（定理~\ref{lh18:no-K4-leaf}）．
この結果は support 全体が $K_4$ と一本の葉辺からなる場合に関するものであり，
$K_4$ を部分 graph として含む全ての support を禁止する主張ではない．

これらに先立つ四使用頂点 $K_4$ の排除では，全低分母の第五 core を hub とし，
その node 条件と $D$ 相殺を通常辺上の正の流れへ書き換える
（定理~\ref{kg17:no-K4}）．
二例外は大きい負 reverse の位置に応じた場合分けで排除される．
四例外は，node 予算によって符号が安定する無限領域を有限個の型へまとめ，
残る六候補を判別式 $-47$ の局所二次式で排除する．
この $K_4$ 証明自体は，一般の三高分母 matching 性を使わずに成立する．
各証明で用いる局所方程式と大域条件を分けることで，
他の support に再利用できる部分を明確にする．

\subsection{葉と二つの三価頂点を持つ七高辺型の排除}

葉を持つ六高辺型 $G_{5,7}$ と $G_{5,9}$ は，それぞれ三価，四価の頂点に一葉を持つ．
両型とも葉の完全標準形を用いるが，隣接 core の高次数が異なるため，
$K_4$ 一葉型の matching 条件をそのまま流用することはできない．
定理~\ref{g719:no-G5-7} では，葉の分母 $Q$ に $Q\le21$ を証明し，
全自由整数の partner 条件から得る $212$ 型を，中心辺の $D$ 相殺で排除する．
定理~\ref{dl19:no-G5-9} では，三高 core の全配置に対して $Q\le13$ を証明する．
自由整数の全半直線十八型を canonical 第二高分母三の矛盾で除き，
残る点九百八型からの $2596$ 個の $D$ 線形系を厳密に排除する．
各自由分母は正整数変数による等差数列のままであり，
有理解の不整合または $3a+3b=-6$，$a,b\ge1$ という符号矛盾で全族を除く．
これにより，高分母支持の葉は全て排除される．

葉を持たない七高辺型のうち
\[
 G_{5,11}=\{01,02,03,04,12,13,24\}
\]
も排除される（定理~\ref{g51120:no-support}）．
二つの二価 cores の標準形を四つの向きに分け，
三高 cores の matching 性と node 予算から reverse data に
$r\le9,S\le7,T\le13$ を証明する．
これは分母と scale を一律に打ち切る条件ではない．
双一次式の二種類の退化を保って全 scale を因数分解し，
九十個の点型と，七十八個の剰余類付き全半直線型を得る．
node 条件・低行の符号と共通低 puncture の接着型が残す $1395$ 配置のうち，
$1309$ 配置は分母の等式に反し，残りは低辺の $D$ 残差が奇数または非零定数となる．
最後の表は，交換と剰余類による重複を除くと十組の模型へ圧縮できる．
七・十三の共鳴を含む全無限族を残した上での排除である．

\subsection{二高分母の整数剛性と book の完全排除}

Book 型 $G_{5,13}$ の三つの二価 cores の間には低辺の三角形がある．
固定 core の全相手を復元する方法，低 $D$ の非零周期，正負枝の node 制約により，
全正・全負・二負枝を除き，一負枝と二正枝の二つの sector へ帰着できる
（定理~\ref{bs21:two-sectors}）．この帰着は，先に解析的に導いた $w\le134$ の範囲を
全尺度で検証する有限証明書を含む．残る二 sector では負枝の $w\ge135$ が必要である．

この残余を閉じる局所整数剛性は，book の接合そのものには依存しない．
負枝を $S=T=w$，$r=w+H$ と書くと，正の整数尺度の条件は
\[
 H=w-1\quad\hbox{または}\quad H=w
\]
を強制する．例えば $E=1-w^2/[H(H+1)]>0$ の側では
$c=sH/(w+H)\in\ZZ_{\ge2}$ と $Ec<1$ を得る．
$H>w$ なら $s-c=cw/H$ は $1$ 以上 $c-1$ 以下の整数であり，
$Ec>2-1/c>1$ となるので矛盾する．負の $E$ 側も対称な議論で決まる．

正枝では $L=3$ または四，$F>0$ とすると
\[
 (r,S,T)=(2R,3R,3R+2)\quad(L=3),\qquad
 (S,T)=(2r,2r+3)\quad(L=4)
\]
に限られる．これは尺度方程式の $0<E_+\alpha<1$ と，その分子に現れる整数の
非負性から従う（補題~\ref{bc22:positive-rigidity}）．
正枝の $F<0$ 側も同様に剛性を持ち，$L=3$ では
$(r,S,T)=(2R-1,3R-2,3R)$，$L=4$ では $(S,T)=(2r-1,2r+2)$ となる．
$F=0$ は整除性により不可能である．
これにより $L=1,3,4$ の全模型を六つの尺度族へ分類する
（命題~\ref{rs22:positive-scales}，\ref{rs22:negative-F-scales}，\ref{rs22:negative-scales}）．
分母七・十三の共鳴自由尺度を含み，全族で奇数高分母は自動成立する．

Book の残余型 I では $H=w$ と node offset から
$w=-\kappa\le18$ となり，$w\ge135$ に反する．
型 II では負枝の $r_N=2w$ が偶数である一方，正枝の $L=4$ と node 条件は
$r_N=2r_1+1$ を要求する．従って型 II も存在しない．
定理~\ref{bc22:book-obstruction} が book 全体を排除し，
系~\ref{sf22:three-unit-branches} が三本の単位分母からの伝播を与える．
この二 sector の排除には，追加の有限探索や四高中心二 core の復元を要しない．

\subsection{二価 core を一つ持つ \texorpdfstring{$G_{5,15}$}{G5,15} の完全排除}

八高辺型 $G_{5,15}$ の低辺を $23,24$ とし，二高 core を $2$，三高 cores を $3,4$ とする．
共有高辺 $34$ の共通低 puncture の比較は，通常型，正例外型，負例外型の三択を与える．
定理~\ref{g51521:reduction} は，正例外型を排除し，通常型を四中心・十四配置へ，
負例外型を共有分母 $Q=3,5$ の二つの template へ完全に還元する．
この段階だけでは未知尺度はなお無限に動き得る．

補題~\ref{oc23:linear-transfer} は，低辺の両端から外側 core に入る一つの pair を使う．
既知の分母 $a$，entry $p$，隣接低 entry $\varepsilon j$ と尺度 $S$ に対して
\[
 \{(au-p)(jv-\varepsilon)-aj\}S=a-1-(au-p)v,
 \qquad u,v\in\{-4,-3,-2,-1,0,1\}
\]
を得る．係数が非零なら尺度は各 node 組で高々一つであり，
係数と定数がともに零になる場合も整数整除で完全に分類する．
通常型では $a-1>j$ により共鳴がなく，十四配置の全尺度は六つの小さい表に限られる．
唯一残る候補は元の行和と二つの pair が要求する
$HK=62,(H+1)M=14,M+N=6$ に反する
（定理~\ref{g51523:no-ordinary}）．

負例外では，外側 pair は隣接低 columns を中心尺度の定数シフトに制限する．
小さい境界と共鳴を全て保持した後，四つの中心族で尺度変数 $\ell$ が奇数，$t$ が偶数となる．
共有高 row の二つの変数は node による定数差を持つため，元の二つの局所二次式の
resultant が必要条件になる．四十個の resultant は，三つの無限領域における
全係数の符号と四つの境界点で排除される．唯一の零評価は正の column 条件に反する
（定理~\ref{ne23:no-negative-exception}）．
先頭係数で割らずに resultant の必要性を証明しているので，一次退化も取りこぼさない．
これは明示的な有限個の整数多項式の全係数による証明であり，分母探索範囲の外挿ではない．
通常型・正例外・負例外を合わせて $G_{5,15}$ は全て除かれ，
残る高分母支持の最小次数は三以上になる．

\subsection{幾何的な先行研究と本稿の範囲}

複素超曲面の高次元 pairs-of-pants による分解と，多面体的な底上の torus fibers は，
Mikhalkin\cite[Theorems~1 and~1$'$]{Mikhalkin}によって研究された．
Leung--Wan\cite{LeungWan}は一般型超曲面の退化に tropical geometry を用い，
正則最高次形式の Calabi--Yau 型開部分への集中と，対応する fibers の
漸近的な special Lagrangian 性を調べた．
Zhang\cite[Theorem~1.1]{Zhang}は，次数 $d\ge5$ の超曲面に微分同相な四次元多様体を，
$d(d-4)^2$ 個の四次元 pair-of-pants と K3 surfaces の部分集合に分解する結果を与えた．
これらは一般型における分解と fibration の幾何的な背景を与える．
本稿が固定する五つの四直線 cores の整数分母配置は，それらの分解に追加する問題である．

Evans--Mauri\cite{EM}は contact link と Liouville cone を通じて
局所 Lagrangian torus fibrations を構成し，条件付きの SNC divisor 近傍の模型を与えた．
本稿の四直線 core の解釈はこの方法を背景とする．
また Fernández de Bobadilla--Pełka\cite[Theorem~6.14]{FernandezBobadillaPelka}は，
与えられた Kähler form の exact modification を用いて A'Campo space の境界で
torus fibrations を構成し，symplectic connection によって近傍 fibers の
codimension-zero 部分へ運ぶ一般的方法を与えた．
maximal Calabi--Yau degenerations では，generic region の torus fibration と
体積充填・漸近的な幾何の結果を得ている
\cite[Theorems~1.1--1.2]{FernandezBobadillaPelka}．
本稿で問うのは，こうした幾何的構成と整合する特定の標準整数 cone charts の接着であり，
それぞれの存在定理だけから本稿の分母の剛性は従わない．
SNC 境界に沿う相対構成，作用切片を含む neck の接続，
曲面全体の fibration の存在は伴論文\cite{SurfaceGeometry}で扱う．

標準 $K_4$ cone を用いることは，その link の抽象 graph を指定することより強い．
整数 fan は一様 matroid $U_{3,4}$ の Bergman fan であり，
線形 tropicalization，matroid，hyperplane arrangement の関係は
Feichtner--Sturmfels\cite{FeichtnerSturmfels}，
Ardila--Klivans\cite{ArdilaKlivans}で体系的に扱われている．
本稿では，さらに face 上の作用微分から来る格子との同定を課す．
tropical atlas の重なりを ambient 整数アフィン写像の制限として定める条件は，
Shaw\cite[Definition~3.1]{ShawSurfaces}，
Jell--Rau--Shaw\cite[Definitions~2.1 and~5.1]{JellRauShaw}に明示されている．
この条件と，torus cycles の作用微分との適合性を同時に要求する点が，
本稿の整数 clutching の出発点である．

Cartwright\cite[Definitions~2.5 and~2.8]{CartwrightComplexes}の tropical complex は，
dual complex に交点数のデータを加え，局所交叉行列を用いて
線形関数と divisor theory を定める．
Gross--Siebert\cite{GrossSiebert}は，特異点を持つ実アフィン幾何と
追加の接着データから Calabi--Yau 退化を再構成する．
本稿の $p,q$ は実際の作用周期の格子から定まるので，
これらとの比較には交点数，アフィン層，特異集合，作用切片の対応を指定する必要がある．
本稿の整数条件が各理論の仮定を満たすことを示さずに，
それらの存在定理や有限性を分母へ適用することはしない．

\subsection{principal minors，cycle gains，整数性}

一つの core の行列を $A=\operatorname{diag}(q_j)+P$ と書けば，
\[
 A\mathbf1=0,\qquad \det A_{\{j,k\}}=1\quad(j\ne k)
\]
となる．非対称行列の principal-minor assignment と
principal minors 間の多項式関係は古典的な問題である．
Lin--Sturmfels\cite{LinSturmfels}は $4\times4$ 行列の関係式を求め，
同論文 Proposition~4 では集合分割による principal minors と cycle-sums の相互変換を与えた．
本稿の三角形の cycle products の消去や Schur 補行列も，
この一般的な代数的枠組みの中に位置付けられる．
本稿で追加して追跡するのは，固定した kernel vector $\mathbf1$，
個々の entry の整数性と整除性，strict basin inequalities，
および異なる五行列の marked punctures での整合である．

正の weights の存在判定は，既知の cycle-gain 条件の特殊化である．
$p_{jk}>0$ の有向辺に gain $p_{jk}/q_j$ を与えると，
$w_k>(p_{jk}/q_j)w_j$ が全て同時に成立することは，
全ての directed cycles の gain が一未満であることと同値になる．
max-times algebra では，これは maximum cycle geometric mean，
Kleene star，subeigenvector の理論で説明できる
\cite[\S2, Theorem~2.2 and Proposition~2.5]{SergeevSchneiderButkovic}．
strict な場合は，十分小さい $\varepsilon>0$ を取り，全ての gains を共通因子 $1+\varepsilon$ で拡大しても
cycle 条件を保てるため，通常の非 strict な subeigenvector 条件から回収できる．
本稿ではこの判定を初等的に示し，有理 weights と整数積による厳密検算に用いる．

一方，cycle 条件そのものは正符号 graph を matching に限定しない．
道上の gain inequalities や，十分小さい gains を持つ三角形には
正の weights が存在し得る．
定理~\ref{ar18:main-matching}の内容は，これらに
row sums と全ての pair products を整数として同時に課すと，
三高・四高のいずれでも完全 matching だけが残るという点にある．
本稿の四ray判定では，孤立頂点から作る二つの整数 $g_{ij}$ の和が一であることを使い，
少なくとも一つが非正となることから禁じられた閉路を作る．
道型を排除する正値恒等式では，二つの正整数積が一以上になることも使う．
この整数の間隙を実係数の一般的な principal-minor relation や
数値的な tolerance に置き換えることはできない．
既存の一般理論との関係をこのように区別し，本稿の各整数結果は明記した条件から直接証明する．

関連する整数行列の理論として，
Corrales--Valencia\cite[\S\S2--3]{CorralesValencia}の graph 上の arithmetical structures がある．
同理論は非正の非対角成分を持つ $Z$-matrices，$M$-matrices を用い，
固定した graph に対する有限性を扱う．
本稿の $A\mathbf1=0$ は形式上似た条件であるが，
高分母の局所行列には正の非対角成分が必要となる
（補題~\ref{lit10:nonpositive-q1}）．
従ってその有限性を，ここでの分母の一様上界として用いることはできない．

determinant 一の整数式は frieze patterns にも現れる．
Cuntz--Holm\cite[\S\S2--3 and~7]{CuntzHolmFriezes}は，
零を許す tame integral friezes の quiddity cycles の表示と，
離散集合上の零を持たない固定高さの friezes の有限性を与えた．
Cuntz--Holm--Pagano\cite{CuntzHolmPagano}は代数的数の環上での有限性を研究した．
frieze では隣接 minors の条件に加え，境界行，周期性，quiddity cycle の積条件を課す．
本稿の六つの face determinants からそれらの構造への写像は与えていないため，
同じ unimodular な式を持つことだけから分類や有限性を適用しない．

\subsection{局所の無限族，完全復元，残る問題}

局所には，任意の奇数 $q>1$ に対して分母 $(q,1,1,1)$ の正の chart が存在する
（系~\ref{q9:local-all-odd}）．
二高分母の全整数行列には signed-divisor 表示があり，
正値性は二本の高低 pairing の符号を決める
（命題~\ref{kl10:two-high-parametrization}，定理~\ref{ls11:two-high-signs}）．
固定した一方の奇分母に対する partner の集合は，七と十三の共鳴を除けば有限であるが，
この二つでは明示的な無限族が残る（定理~\ref{fp12:partner-structure}）．
従って，局所整数性から分母全体の有限性を結論することはできない．

三高分母についても，matching 性は全行列の有限性を意味しない．
matching 型から得られる Schur 三元一次系は特異になり得る．
実際，階数二の系から
\[
 (q_a,q_b,q_c,q_0)
 =(409+600n,\ 3493+5100n,\ 1303+1904n,\ 1),
 \qquad n\in\ZZ_{\ge0}
\]
という正の整数局所族が得られる（定理~\ref{sh14:singular-family}）．
matching 定理により，この Schur 表示と $\beta\le-1$ の結果
（定理~\ref{sm14:beta-negative}）は全ての正の三高分母模型へ適用できる．
三つの高 $D$ が全て非負なら，$b=-\beta$ として
\[
 d=-D_0\ge1,\qquad \alpha\ge k\ge b+1\ge3,
 \qquad \Delta\ge(-B)(-B-1)>0
\]
が必要である（定理~\ref{hd20:nonsingular}）．
全 $D$ 零の場合は Cramer 恒等式と整数ベクトルの norm の整除性で排除し，
二つの境界 $k=b-1,b$ は元の線形系の正値性から除く．
従って，全ての特異 Schur 系をこの同時非負性の問題から外せる．
特異な正整数局所模型自体の存在は，この帰着と両立する．

残る非特異領域については，高係数を
$(x,y,z)=(D_a,D_b,D_c)\ge0$，$x+y+z=d\ge1$ と指定し，低 row の二つの負成分の大きさを $s,t$，
$u=s+t-1$ とする．命題~\ref{hd21:reconstruction} は新しい必要条件 $x\le s-1$ を与え，
三尺度 $S,T,U$ の整数性を
\[
 F=b-1+ux+sz,\qquad m=F^2+F+1,\qquad m\mid FE-J
\]
という一つの合同条件にまとめる．ここで $E,J$ は固定した五整数 $(s,t,x,y,z)$ だけで定まる
整数である（式~\eqref{hd21:fixed-data}）．
この合同条件に，$\alpha\ge k\ge b+1$ と元の一つの pair-product 方程式，
および尺度・entry の正値性を加えると，
三高係数が全て非負である正の整数模型の必要十分な復元判定になる．
正の有理 weights も明示的に構成できる．

さらに定理~\ref{hd21:norm-bound} は，$N=E^2+EJ+J^2$ に対して
\[
 3\le m\le N,\qquad m\mid N
\]
を証明する．基礎となるのは，尺度の整数差を用いる正定値二次形式の恒等式であり，
$N=0$ となる二つの候補も元の pair-product 方程式で排除する．
従って $(s,t,x,y,z)$ を固定すれば $b$ の全範囲を有限に復元できる．
さらに本稿では $K=sT-b>0$ の分子を上から評価し，
\[
 F\le st+(2u+1)(s-x)-2
\]
を導く．従って $s,t$ を固定すれば $x,b,z$ は明示的な有限範囲に入る．
残る元の pair 残余式は $y$ の二次式であり，二次係数が消える場合も一次係数が
厳密に負となるため，恒等零にはならない．これにより固定 $(s,t)$ ごとの完全有限復元を得る．
さらに $s=t,x=y$ の対称領域は無上界で排除できる．
この場合は $E+J=0$ と尺度の整数性から $E/m$ が整数となり，
三尺度の正値性が $E=J=0$ を強制して既知の零ノルム排除に反する．
尺度の整数性を $n=\ell-U\in\ZZ$ として取り出す．
ここで $B=b-1$，$h=s-x$，$j=t-y$，$\ell=u+z$，$\delta=sy-tx$ と置くと
\[
 S=j-(F+1)n-\ell\delta,\qquad
 T=h+Fn+\ell\delta,\qquad U=\ell-n
\]
となる．$n$ は格子変数であって node の個数ではない．
この表示と元の pair-product 残余から，次の排除と有限復元を得る．

\begin{theorem}[非対称量の排除，比例領域と退化面の有限復元]
\label{ar24:high-D-reductions}
三高係数が同時に非負である正の canonical 整数模型を考え，直前の記号を用いる．
\begin{enumerate}
\item $\delta=-1$ および $\delta=1$ は起こらない．
\item $\delta=0$ なら $x=y=0$，$z\ge1$，$s<t$，$0\le n\le z$ である．
$s,z$ を固定したとき，各 $n=0,\ldots,z$ の候補 $b$ は真の四次式の整数根となる．
従って原方程式を戻す前の候補数は高々 $4(z+1)$ である．
\item 退化係数 $a=u(s-x)-(b-1)(F+1)$ が零なら，$s\ge3$ であり，
\[
 2\le b\le1+s^2(s-1)^2
\]
が成り立つ．$s$ だけを固定して全ての残りの整数を有限に復元できる．
\end{enumerate}
いずれも分母や残りの尺度に探索上限を仮定しない．
\end{theorem}

第一項は定理~\ref{hd23:main} と定理~\ref{ax24:positive-axis} による．
正側の境界 $x=0$ では $s=y=1$ となるが，元の pair 残余と尺度差方程式を合わせると，
正整数の整除性とその大きさの比較が矛盾する．
差方程式だけで元の残余が零になると仮定してはならない．

第二項は定理~\ref{hp24:main} と系~\ref{hp24:finite-axis} による．
$x>0$ の比例領域では尺度の正値性が $n=0$ を強制し，
元の残余の厳密な係数符号と整数の間隙から矛盾を導く．
残る零軸では $0\le n\le z$ を証明した上で，$b$ の四次式を得る．
その先頭係数は $n^2-zn-1<0$ であり，恒等零への退化はない．
さらに
\[
 C_n=n\bigl((sz)^2-3sz+3\bigr)+(s^2-1)z-3s+1
\]
と置くと，模型に対して $C_n>0$ であり，
\[
 b+1\mid C_n,\qquad
 b+1\le C_n\le s^2z^3-3sz^2+(s^2+2)z-3s+1
\]
という明示的な約数上界も得られる
（補題~\ref{hp24:divisor-bound}）．零の余りを持つ二つの共鳴は，
それぞれ既知の対称排除と判別式九十三の平方障害によって除かれる．

第三項は定理~\ref{ax24:degenerate-R1} と系~\ref{ax24:degenerate-finite} による．
退化面で残る $x=0,R_0=b^2-sT=1$ の例外枝を，resultant の約数条件と
連続する二平方の間隔で排除する．残る枝では元の残余から
\[
 b-1\mid(s-x)(s-1)(R_0-1)
\]
が従い，上の一変数の明示上界を得る．
候補ごとの残余は $y$ の真の一次式となるため，その復元も有限である．
端点については，定理~\ref{hpe25:main} が
\[
 n=0:\quad s\ge3,\ s<b<3s/2,\ 1\le z<s/2,
\]
\[
 n=z:\quad s\ge2,\ 1\le z\le s,\ 1\le C_z/(b+1)\le s-1
\]
を与える．前者は元残余の符号による三つの領域排除，後者は正の整数商を変数とする
モニック四次式の単調性と隣接整数間の根の位置による．
両端点は固定 $s$ だけで完全に有限化される．

退化面では，定理~\ref{dg25:product} の正整数 $M,L$ により
$M(t-y)=tT$，$hM=(b-1)L$ となる．$h\ge b$ なら $1\le L\le b-1$ なので
$\gcd(h,b-1)\ge2$ が必要である．この条件は $x\ge1$ の退化面全体に適用でき，
系~\ref{dg25:fixed-b} は固定 $b$ から全残りの整数を有限復元する．

これらの結果は比例零軸や退化面全体の不存在を結論せず，
$s,z$ または $s$ 自体の一様上界も与えていない．

別の有限性として，低 ray の行と列を両方固定すれば，
三高分母の全整数 completion は高々二つとなる
（命題~\ref{fc17:at-most-two}）．
同符号の low-row entries を選ぶと，二次係数が消えない一変数二次式が得られ，
その整数根と整除条件から全行列を復元できる．
同じ固定データに対する二つの正の標識付き completion も存在する
（注意~\ref{fc17:sharp}）．
ここでは低 column を動かす自由度を残していないため，上の無限族と矛盾しない．
また二つの標識付き解が，幾何学的に非同値な二つの charts を与えると主張するものではない．

大域の解析には素数ごとの合同条件も併用できる．
$\ell\mid q_{ij}$ の辺からなる素数 support $G_\ell$ は，
任意の奇素数 $\ell$ に対して $K_4$ と book $B_3$ を含まず，$G_3$ は matching となる
（定理~\ref{ps10:forbidden}）．
七辺以上の素数 support は $\ell=5$，$G_5=K_{2,2,1}$ の場合に限られ，
任意の七つの分母の gcd は五を割る
（定理~\ref{ps11:dense}，系~\ref{ps11:dense-gcd}）．
これらは高分母 support $H$ と異なる graph に対する制約であり，
その剰余解の存在は整数 lifting や atlas の存在を保証しない．

証明中の有限計算では，先に方程式から有限範囲または有限個の無限型代表を導く．
node 予算から得られる有限な場合分け，符号が安定する全無限領域，
および明示的な共鳴族の不等式による排除を区別して記述する．
全分母が百一以下なら全て一となる検算
（定理~\ref{ex10:bounded-exclusion}）は，指定した分母範囲についての結果である．
これを，本稿の上限を仮定しない graph 障害と混同しない．

中心的な未解決問題は，残る三型のいずれかが
異なる高分母を混在させた算術的接着データを持つかどうかである．
一般の五次分母剛性，すなわち全ての分母が一に強制されるかという問題も未解決である．
二つの互いに素な低辺を持つ $G_{5,14}$ では，両低辺に強制される $\{4,1,0\}$ 配分，
隣接分母三，固定された中央 matching を同時に合わせる必要がある．
$G_{5,16}$ の唯一の単位枝には一般の二配分 $\{5,0,0\},\{4,1,0\}$ と小分母の制約が適用されるが，
五node配分の排除までは得ていない．$G_{5,17}$ は単位枝を持たない．
四高 core でも完全 matching に限定し，元の行和・pair products と大域接着を同時に調べる．
高分母全体の原始性と，二本の単位枝があるときの任意七分母の互いに素性は，
素数 support の合同条件を整数 lifting へ適用する追加の必要条件となる
（系~\ref{gp23:primitive}，\ref{gp23:two-units}）．
三高同時非負性では，$\delta=0$ の零軸，$|\delta|\ge2$ の領域，
および退化面の一整数有限復元から一般的な構造へ進む問題が残る．

今後は，これらの還元を matching 性，Schur 表示，素数 support と組み合わせ，
残る無限領域に対する構造定理または整数接着の構成を求める必要がある．
整数候補が得られた場合には，局所の実面積と作用切片を同時に整合させる幾何的な課題が残る．
本稿の結果は，その実現可能性を調べるための局所符号分類，
特定の入力に対する完全な整数復元，および標準五次模型に対する厳密な必要条件を与える．

% Standalone arithmetic setup for the third paper, v17.
% SurfaceGeometry is the companion geometric manuscript (root supplies its entry).
\section{標準五次模型と整数データ}
\label{ar17:setup}

\subsection{五つの cores と四直線の作用座標}

五つの一般位置の平面の和へ退化する五次曲面を考え，成分を
$V=\{0,1,2,3,4\}$ で標識する．各成分の blow-down 上の境界は
一般位置の四直線であり，その局所模型を \emph{core} と呼ぶ．
core $i$ の四つの境界直線，従って四つの branch rays は
$V\setminus\{i\}$ で標識される．共有二重曲線 $C_{ij}$ に対応する
branch を $ij$ と記すと，大域 branch graph は $K_5$ である．
その branch の三つの punctures は $k\in V\setminus\{i,j\}$ に対応する．
順序付き小解消では，この二重曲線の一方の側に五つの blow-up patches がある．
これらの幾何的な準備，SNC 境界に沿う相対構成，および実際の smoothing との比較は
伴論文\cite{SurfaceGeometry}で扱う．以下の整数問題は，これらの構成を再実行せずに定義できる．

一つの四直線 core では標識を $0,1,2,3$ と取り直すことがある．
標準 fan の ambient tangent lattice を
\[
 N_4=\ZZ^4/\ZZ(1,1,1,1),\qquad r_j=[e_j]
\]
とする．$\sum_jr_j=0$ であり，fan は四本の rays $\RR_{\ge0}r_j$ と
六つの二次元 cones $\RR_{\ge0}r_j+\RR_{\ge0}r_k$，$j<k$，からなる．
これは一様 matroid $U_{3,4}$ の Bergman fan の標準的な粗い fan 構造で，link は $K_4$ である
\cite{FeichtnerSturmfels,ArdilaKlivans}．

\begin{definition}[作用座標に適合する cone chart]
\label{per:chart-definition}
Liouville cone 型 torus fibration の底の core 頂点を $o$ とする．
標準 $K_4$ fan の原点近傍への，strata を保つ同相写像 $F$ を考える．
正則 face $e$ の torus fiber $T_e$ に対して
$L_e=H_1(T_e;\ZZ)$ と置く．作用一形式 $\dd I_\gamma$ は $L_e$ を整数余接線格子に同定する．
各 face 上で $F$ が作用座標の整数線形変換であり，その微分が
\[
 L_e^*=\operatorname{Hom}(L_e,\ZZ)
 \longrightarrow N_4\cap\operatorname{span}_{\RR}F(e)
\]
の格子同型となり，さらに各 branch ray の接線格子が全 incident faces で飽和しているとき，
$F$ を\emph{作用適合 cone chart} と呼ぶ．core の像は原点に正規化する．
\end{definition}

この定義で同定するのは fan の整数接線格子と $L_e^*$ である．
補集合の全ての一次ホモロジー類に ambient linear function を割り当てる追加条件は課さない．
大域の作用適合 atlas では，重なりの変換は ambient 整数アフィン変換の制限であり，
共有 branch に追加の特異点を入れずに接続するものを対象とする．
単に各 face 上の格子を別々に同定する条件とは区別する．

\subsection{局所の行和，pair products，面積正値性}

一つの core の四つの branch denominators を $q_j\in\ZZ_{>0}$ とし，
$P=(p_{jk})$ を対角成分零の整数行列とする．
本稿の符号では，crossing face $jk$ 上の作用座標は
\begin{equation}\label{ar17:action-matrix}
 \begin{pmatrix}I_{\mu_j}\\ I_{\mu_k}\end{pmatrix}
 =\begin{pmatrix}-q_j&-p_{kj}\\-p_{jk}&-q_k\end{pmatrix}
   \begin{pmatrix}t_j\\t_k\end{pmatrix},\qquad
 F=t_jr_j+t_kr_k.
\end{equation}
ここで $\mu_j,\mu_k$ は正の meridians であり，作用積分の単位は $1/(2\pi)$ とする．
ray $j$ 上の normal action は $-q_jt_j$ なので，$q_j$ は
primitive branch 座標に対する normal action の整数係数の絶対値である．

\begin{definition}[局所整数データと正値性]
\label{ar17:local-data}
局所整数データとは，全ての $j\ne k$ について
\begin{align}
 \sum_{k\ne j}p_{jk}&=-q_j,\label{ar17:local-row}\\
 p_{jk}p_{kj}&=q_jq_k-1\label{ar17:local-pair}
\end{align}
を満たす $(q,P)$ である．さらに，正実数 $w_j$ が存在して
\begin{equation}\label{ar17:local-positive}
 B_{jk}:=w_k-\frac{p_{jk}}{q_j}w_j>0\qquad(j\ne k)
\end{equation}
となるとき，このデータを\emph{正}と呼ぶ．
$q_j=1$ の ray を低分母，$q_j>1$ の ray を高分母と呼ぶ．
\end{definition}

式~\eqref{ar17:local-row} は三つの punctures の period relation，
式~\eqref{ar17:local-pair} は作用行列~\eqref{ar17:action-matrix} の determinant 一を表す．
後者から $\gcd(q_j,p_{jk})=1$ も従う．
また，もし $q_j$ が偶数なら三つの $p_{jk}$ が全て奇数となり，
その和が $-q_j$ に等しいことと矛盾する．従って局所整数データの分母は自動的に奇数である．

$W=\sum_jw_j$ と置くと，行和から $\sum_{k\ne j}B_{jk}=W$ となる．
幾何的には $2\pi B_{jk}$ は直線 $j$ の $k$-puncture basin の面積であり，
三つの basins が面積 $2\pi W$ の全直線を尽くす．
命題~\ref{q9:four-line-criterion} は，この整数条件と正値性から，
指定した marked Morse foliation の相対 isotopy class 内で実際の局所 chart を得ることを示す．
正の weights の存在は，$p_{jk}>0$ のときだけ有向辺 $j\to k$ を置き，
全ての simple directed cycles $C$ に対し
\begin{equation}\label{ar17:positive-cycles}
 \prod_{(j\to k)\in C}p_{jk}<\prod_{j\in C}q_j
\end{equation}
を要求することと同値である（補題~\ref{kl10:cycle-positive}）．
四頂点なので長さ二，三，四だけを検査すればよく，成立する場合には weights を有理数に取れる．

局所の normal-action 分岐係数を
\begin{equation}\label{ar17:D-local}
 D_j=-q_j-\sum_{k\ne j}p_{kj}
\end{equation}
と定める．これは標準 branch star の relation
$\sum_{k\ne j}r_k=-r_j$ に対し，normal action の transverse slopes の和を取った係数である．
行和を全て足すと
\begin{equation}\label{ar17:D-sum}
 \sum_jD_j=0
\end{equation}
が従う．文中の $\mathfrak d_j$ と $D_j$ は，この同じ正規化の係数を表す．

\subsection{五つの core 間の整数 clutching と node 条件}

大域標識 $V$ に戻る．共有 branch の分母は両端で共通の
$q_{ij}=q_{ji}\in\ZZ_{>0}$ とする．core $i$ の行列成分を
$p^i_{jk}$，$j,k\in V\setminus\{i\}$，$j\ne k$，と書く．
局所の条件は，その core において $q_j=q_{ij}$ とした
式~\eqref{ar17:local-row}--\eqref{ar17:local-positive} である．
特に，core $i$ の ray $ij$ における係数は
\begin{equation}\label{ar17:D-global}
 D^i_j=-q_{ij}-\sum_{k\in V\setminus\{i,j\}}p^i_{kj}
\end{equation}
となる．

共有 puncture $ijk$ の整数 clutching は
\begin{equation}\label{ar17:clutching}
 q_{ij}\mid p^i_{jk}+p^j_{ik}
\end{equation}
を要求する．node 数は，この整数を用いて
\begin{equation}\label{ar17:node}
 n_{ij;k}=n_{ji;k}
 :=1-\frac{p^i_{jk}+p^j_{ik}}{q_{ij}}
 \in\ZZ_{\ge0}
\end{equation}
と定める．この式は marked puncture における normal framing の比較から来る．
両端の三つの row entries の和がそれぞれ $-q_{ij}$ であるため，
\begin{equation}\label{ar17:node-budget}
 \sum_{k\in V\setminus\{i,j\}}n_{ij;k}=5,
 \qquad 0\le n_{ij;k}\le5
\end{equation}
が従う．五という総数を，三つの node 数から独立の自由変数として扱わない．
さらに，共有 branch の対向端で normal action の向きが逆になることから
\begin{equation}\label{ar17:D-opposition}
 D^i_j+D^j_i=0
\end{equation}
を課す．

\begin{definition}[五次の算術的接着データ]
\label{ar17:admissible}
$K_5$ の十辺上の正整数 $q_{ij}$ と，対角成分が零の五つの整数行列 $P^i$ が，
各 core の行和，pair products，正値性，および全共有辺の
整数 clutching，node 非負性，対向 $D$ 相殺を満たすとき，
\emph{五次の算術的接着データ}と呼ぶ．
その\emph{高分母 support} は，頂点集合 $V$ と辺集合
\begin{equation}\label{ar17:high-support}
 E(H)=\{\{i,j\}:q_{ij}>1\}
\end{equation}
を持つ単純 graph $H\subset K_5$ とする．
support の使用頂点とは少なくとも一本の高辺に接する頂点であり，
core $i$ の高分母次数は $\deg_H(i)$ である．
\end{definition}

定理ごとに，この定義のどの部分を使うかを明記する．
例えば素数 support の合同障害は行和・pair products・整数 clutching だけを使い，
いくつかの小さい support の障害は正値性や対向 $D$ 相殺を必要としない．
一方，幾何的な大域 atlas を構成する際は，局所 weights による実面積と
有限 smoothing parameter における作用切片の一致も必要になる．
定義~\ref{ar17:admissible} は大域 atlas に対する必要な算術条件をまとめたものであり，
これらの実条件まで含む十分条件ではない．

\subsection{三つの異なる有向・無向 graph}

support $H$ の vertices は五つの cores である．これに対し，
式~\eqref{ar17:positive-cycles} の有向 graph の vertices は
一つの core に属する四つの rays であり，別の graph である．
さらに奇素数 $\ell$ に対して
\[
 E(G_\ell)=\{\{i,j\}:\ell\mid q_{ij}\}
\]
と定める素数 support は，$H$ の部分 graph である．
高分母の辺が同じ素因数を共有するとは限らないので，
$G_\ell$ に対する禁止部分 graph の結果を $H$ にそのまま移すことはできない．
以下ではこの三者を区別して局所正値性，合同障害，大域接着を組み合わせる．

% v9. A genuinely geometric local high-denominator example for four lines.
% Uses existing bibliography key EM and per:chart-definition.

\section{四直線 core の奇数分母と五次の追加制約}
\label{q9:quintic-local-high-q}

四 puncture neck の剛性は両端の比較から生じる．
三 puncture 曲線だけを用いる四直線 core では，高分母の作用適合 chart を実際に構成できる．
次の基準では，局所の整数条件と正の basin areas が十分であることを，sector 全体で確認する．

\begin{proposition}[四直線 core の十分条件]\label{q9:four-line-criterion}
$L_1,\ldots,L_4\subset\PP^2$ を一般位置の四直線とし，正の weights $w_i$ と
\[
 [\omega]=2\pi\sum_{i=1}^4 w_i[L_i],\qquad W=\sum_iw_i
\]
を固定する．各直線に三 puncture の marked Morse foliation の相対 isotopy class を指定する．
正の整数 $q_i$ と整数 $p_{ij}$，$i\ne j$，が
\begin{equation}\label{q9:four-line-arithmetic}
 \sum_{j\ne i}p_{ij}=-q_i,
 \qquad q_iq_j-p_{ij}p_{ji}=1
\end{equation}
を満たし，さらに
\begin{equation}\label{q9:four-line-basin-positive}
 B_{ij}:=w_j-w_i\frac{p_{ij}}{q_i}>0
 \qquad(i\ne j)
\end{equation}
とする．このとき，指定した isotopy classes 内で foliations の代表を選び，
本稿の比例曲率 collars と Liouville cone construction による core に，
標準 $K_4$ fan への作用適合 cone chart を構成できる．
その四本の branch denominators は $q_i$ である．
\end{proposition}

\begin{proof}
各直線の面積は $2\pi W$ であり，\eqref{q9:four-line-arithmetic} から
\[
 \sum_{j\ne i}B_{ij}=\sum_{j\ne i}w_j+w_i=W
\]
を得る．従って $L_i$ の三つの critical basins の面積を $2\pi B_{ij}$ と指定できる．
具体的には，既存の marked foliation の各 puncture disk を十分小さく固定し，
その外で正の density を調整して三 basin の面積を指定する．
全積分が変わらないので，puncture disks に相対的な Moser isotopy によって，
元の面積形式の下で所期の foliation の代表を得る．
比例曲率 collar での critical slope は
\[
 h_{ij}=\frac{w_j-B_{ij}}{w_i}=\frac{p_{ij}}{q_i}.
\]
graph-period correction は fiber tangent を消す basic form で行うので，この作用値を変えない．

$N_4=\ZZ^4/\ZZ(1,1,1,1)$ の標準 rays を $r_i$ とする．
crossing face $ij$ の regular torus の基底を，二直線の正の meridians $(\mu_i,\mu_j)$ とする．
その作用座標において
\begin{equation}\label{q9:four-line-action-matrix}
 \begin{pmatrix}I_{\mu_i}\\I_{\mu_j}\end{pmatrix}
 =\begin{pmatrix}-q_i&-p_{ji}\\-p_{ij}&-q_j\end{pmatrix}
   \begin{pmatrix}x\\y\end{pmatrix},
 \qquad F=xr_i+yr_j
\end{equation}
と定める．行列の determinant は一である．
$i$ 側の critical ray の action direction は $(-q_i,-p_{ij})$ であり，
$j$ 側は $(-p_{ji},-q_j)$ なので，両境界 rays は所期の $r_i,r_j$ へ写る．

sector 全体で $x,y>0$ であることを確認する．
$i$ 側の比例曲率 collar では $I_{\mu_i}<0$ とし，
$h=I_{\mu_j}/I_{\mu_i}$ を用いると
\[
 x=(-I_{\mu_i})(q_j-p_{ji}h),\qquad
 y=(-I_{\mu_i})(q_i h-p_{ij})
\]
である．partial basin を critical level から puncture に向かって縮めると，
$h$ は $p_{ij}/q_i$ から $w_j/w_i$ へ向かって厳密単調に増える．
従って正則部分では $y>0$ となる．
$x$ の括弧内は $h$ の affine function であり，両端でそれぞれ
\[
 q_j-p_{ji}\frac{p_{ij}}{q_i}=\frac1{q_i}>0,
 \qquad
 q_j-p_{ji}\frac{w_j}{w_i}
   =\frac{q_j}{w_i}B_{ji}>0
\]
を取るため，全 interval で正である．$j$ 側も同様である．
product crossing の極限作用点 $(-w_i,-w_j)$ における二係数は
\[
 x=q_jw_i-p_{ji}w_j=q_jB_{ji}>0,
 \qquad
 y=q_iw_j-p_{ij}w_i=q_iB_{ij}>0.
\]
transition caps は critical leaves から離れて固定できる．
そこでは上の正値性に一様な余裕があり，cap 内の normal-radius に比例する有界な補正は，
normal radius を十分小さくすることでこの余裕より小さくできる．
従って caps，product crossings，それらの丸めを通じても正の sector 内に留まる．

各 $i$ に incident な三つの maps は共通 ray 上で
$F=(-I_{\mu_i}/q_i)r_i$ に一致する．
\eqref{q9:four-line-arithmetic} の determinant 条件から
$\gcd(q_i,p_{ij})=1$ であるため，それぞれの action direction は primitive であり，
共通 branch lattice は三枚の sheets 全てで飽和している．
三 arm の period relation は，row-sum 条件 $\sum_{j\ne i}p_{ij}=-q_i$ と整合する．
標準 fan の各 cone の span lattice は $\ZZ r_i+\ZZ r_j$ なので，
\eqref{q9:four-line-action-matrix} は作用の整数余接線格子と fan の整数接線格子の
双対関係を通じて必要な格子同型を与える．

最後に，各 contact-link edge の action vector を $a(t)$ とすると，
Liouville cone の作用 map は $(r,t)\mapsto r a(t)$ である．
正則 torus fibration の作用は局所座標なので $\det(a(t),a'(t))\ne0$ であり，
正の sector 内の projective action ratio は厳密単調である\cite[Remarks 2.11--2.12]{EM}．
両 endpoint rays の間を一度だけ走るので，各 face は対応する cone sector へ同相に写る．
四つの rays 上での一致と compact link を用いて六つの faces を core まで貼れば，
定義~\ref{per:chart-definition} の作用適合 cone chart を得る．
\end{proof}

\begin{corollary}[任意の奇数分母を持つ局所 chart]\label{q9:local-all-odd}
任意の整数 $h\geq1$ に対し，$q=2h+1$ とする．
四直線を $A,B,C,D$ と標識し，
\[
 (q_A,q_B,q_C,q_D)=(q,1,1,1),\qquad
 (p_{ij})=
 \begin{pmatrix}
 0&-2h&-2&1\\
 -1&0&0&0\\
 -h&h-1&0&0\\
 2h&0&-q&0
 \end{pmatrix}
\]
とする．正の weights
\begin{equation}\label{q9:local-all-odd-weights}
 (w_A,w_B,w_C,w_D)=\left(q,h,1,1+\frac1{4h}\right)
\end{equation}
に対し，これらの分母と critical numerators を持つ局所作用適合 $K_4$ chart が存在する．
全 weights の共通の正の倍率は任意なので，任意の正の hyperplane Kähler class に対して例が得られる．
\end{corollary}

\begin{proof}
各 row の和は $-q_i$ である．$A$ に接する三辺では
$p_{ij}p_{ji}=2h=q_iq_j-1$，残る三辺では零となる．
従って全六 determinant が一である．
normalized basin areas は
\[
 (B_{ij})=
 \begin{pmatrix}
 -&3h&3&1/(4h)\\
 3h+1&-&1&1+1/(4h)\\
 3h+1&1&-&1+1/(4h)\\
 1/2&h&2h+5/2+1/(4h)&-
 \end{pmatrix}.
\]
全て正であり，各 row の和は $W=3h+3+1/(4h)$ である．
命題~\ref{q9:four-line-criterion} を適用する．
\end{proof}

この族の normal-action branch coefficients は，$A,B,C,D$ の順に
$(-3h,h,2h+2,-2)$ であり，総和は零である．
各 normal Hamiltonian や全補集合の $H_1$ が ambient linear forms に延長する条件は課していない．
三価 branch の必要条件である奇数性に対し，局所の四直線 core では任意に大きい奇数分母が生じる．
これは一つの四直線 core の幾何学的な存在族であり，五つの quintic cores の同時接着を与えるものではない．
共有 branch の分母と係数の相殺に加え，作用切片，面積，node 配分を同時に合わせる問題が残る．

% Reusable v9 lemma: can be placed after the odd-denominator K4 family.
\subsection{孤立した高分母 branch の接着障害}
\label{qh9:isolated-section}

全平面五次の一つの four-line core を考え，四 vertices の分母を
$(q,1,1,1)$，$q>1$ とする．outgoing numerators を $p_{ij}$ と書く．
各 line の三 basins が全曲線を尽くすので
\[
 \sum_{j\ne i}p_{ij}=-q_i,
 \qquad q_iq_j-p_{ij}p_{ji}=1.
\]
正の weights に対して全 basin areas が正，すなわち
$w_j-w_i p_{ij}/q_i>0$ であることも仮定する．

\begin{proposition}[唯一の高分母 vertex の係数]\label{qh9:isolated-local}
高分母 vertex から出る三 numerators は，順序を除いて
\[
 \bigl(-(q-1),-(b+1),b\bigr)
\]
である．ここで $b$ は正整数で，$b$ と $b+1$ はともに $q-1$ を割る．
$q-1=t b(b+1)$ と書くと，その vertex の normal-action 分岐係数は
\[
 \mathfrak d_{\rm high}=-t\bigl(b(b+1)+1\bigr)<0
\]
となる．
\end{proposition}

\begin{proof}
$m=q-1$ と置く．高分母 vertex を $0$ とすれば，各 edge で
$p_{0j}p_{j0}=m>0$ なので，三つの $p_{0j}$ は零でなく，絶対値は $m$ の約数である．
和は $-q=-(m+1)$ なので，二つ以上が正ではあり得ない．

全てが負で $p_{0j}=-a_j$ と仮定する．$a_1+a_2+a_3=m+1$ であり，
各 $a_j$ は正なので，いずれも $m$ には等しくない．
従って $m$ の真の約数として $a_j\le m/2$ であり，
低分母 vertex $j$ から他の二つの低分母 vertices へ出る numerators の和は
\[
 -1-p_{j0}=\frac{m}{a_j}-1>0.
\]
各低分母 vertex から少なくとも一つの正の directed edge が出る．
低分母同士では $p_{jk}p_{kj}=0$ なので，これらの正の edges は
三 vertices 上に directed cycle を持つ．正の整数 $p_{jk}$ に対して basin positivity は
$w_k>p_{jk}w_j\ge w_j$ を要求するため，cycle を一周して矛盾する．

よって正の entry がちょうど一つあり，三 entry を $(-a,-c,b)$ と書ける．
$a+c-b=m+1$ である．$a,c$ の両方が $m$ の真の約数なら
$a+c\le m$ なので不可能であり，順序を替えて $a=m$ とできる．
すると $c=b+1$ となる．$b$ と $b+1$ は互いに素でともに $m$ を割るので，
$m=t b(b+1)$，$t\in\ZZ_{>0}$ と書ける．

高分母 vertex へ入る三 numerators は
\[
 \left(-1,-\frac{m}{b+1},\frac{m}{b}\right)
\]
であり，その和は $-1+t$ である．$K_4$ fan では $\kappa=-1$ なので，
分岐係数は $-q-(-1+t)=-m-t=-t(b(b+1)+1)$ となる．
\end{proof}

\begin{corollary}[高分母 support の孤立 edge はない]\label{qh9:isolated-edge}
全平面五次の既存の branch graph $K_5$ 上に，大域作用適合 atlas があるとする．
$q>1$ の branches だけからなる部分 graph は，ただ一本の edge からなる
連結成分を持たない．
\end{corollary}

\begin{proof}
そのような edge の両端の core は，どちらも分母 pattern $(q,1,1,1)$ を持つ．
命題~\ref{qh9:isolated-local} により，その branch の両端の normal-action 係数は
ともに負となり，必要な対向端の相殺と矛盾する．
\end{proof}

これは五次の全分母が一であることまでは示さない．
二本以上の高分母 branches が隣接する場合の相殺条件と，
全共有曲線における整数 framing・実面積の同時整合が，次の問題として残る．

% v10 local K4 structure: exact positivity criterion, two-high parametrization,
% and geometric counterexamples to proposed local maximum principles.
% Uses v9 q9:four-line-criterion and existing theorem environments/macros.

\subsection{正の局所 weights の有限な判定}
\label{kl10:positivity-section}

四直線 core について，分母 $q_i>0$ と整数 numerators $p_{ij}$ が
\[
 \sum_{j\ne i}p_{ij}=-q_i,
 \qquad p_{ij}p_{ji}=q_iq_j-1
\]
を満たすとする．この節での係数は
$D_i=-q_i-\sum_{j\ne i}p_{ji}$ である．
これらの row sums から $\sum_iD_i=0$ が従う．
次の閉路判定は，対数を取った差分不等式の標準的な判定であり，
max-times algebra の Kleene star の存在条件の特殊化でもある
\cite[\S2, Theorem~2.2 (arXiv version)]{SergeevSchneiderButkovic}．
局所の幾何学的な実現に残る正値性は
$w_j>w_i p_{ij}/q_i$ であり，次の有限条件と同値である．

\begin{lemma}[有向 cycle による正値性判定]\label{kl10:cycle-positive}
$p_{ij}>0$ の場合にだけ有向辺 $i\to j$ を置き，その gain を
$c_{ij}=p_{ij}/q_i>0$ とする．
全ての basin areas が正となる weights $w_i>0$ が存在することと，
この graph の全ての simple directed cycles $C$ について
\begin{equation}\label{kl10:cycle-test}
 \prod_{(i\to j)\in C}p_{ij}<\prod_{i\in C}q_i
\end{equation}
が成立することは同値である．存在する場合には，全 weights を有理数に取れる．
\end{lemma}

\begin{proof}
正の weights があれば，cycle に沿って $w_j>c_{ij}w_i$ を掛けることで必要条件を得る．
逆に，全 cycle gains が一未満とする．simple cycles は有限個なので，
有理数 $\theta>1$ を十分一に近く選べば，辺の gains を
$\widetilde c_{ij}=\theta c_{ij}$ に替えても全 cycle gains は一未満である．
頂点 $j$ で終わる全ての simple directed paths の gains と，empty path の gain 一の最大値を $w_j$ とする．
これは有限個の正の有理数の最大なので，有理数である．
$w_i$ を与える path に $i\to j$ を付け加える．
繰返し頂点が生じたら cycle を削除することで，gain は小さくならず，終点 $j$ の simple path を得る．
従って $w_j\geq\widetilde c_{ij}w_i>c_{ij}w_i$ となる．
$p_{ij}\leq0$ の場合の basin positivity は $w_i,w_j>0$ から自動的に従う．
\end{proof}

四頂点なので検査する cycle の長さは二，三，四だけである．
長さ二の条件は $p_{ij}p_{ji}=q_iq_j-1$ から従うが，
長さ三，四の条件は row sums と pair products だけから省略しない．
検査対象は最大八個の有向三角形と六個の有向四角形である．
この判定は有限の整数積の比較であり，数値的な線形計画法の tolerance を用いない．

\subsection{二つの高分母を持つ局所模型の算術的な表示}
\label{kl10:two-high-section}

\begin{proposition}[二高分母の factorization 表示]\label{kl10:two-high-parametrization}
分母が $(a,b,1,1)$，$a,b>1$，である局所四直線模型を考える．
二つの高分母頂点および二つの分母一の頂点の標識を必要に応じて交換すれば，
その全ての整数 numerator matrices は次の形に表せる．
正の整数 $r$ と非零整数 $u$ を取り，
\[
 r\mid a-1,\qquad r+1\mid b-1,\qquad u\mid a-1,
 \qquad x=\frac{a-1}{r},\qquad y=\frac{b-1}{r+1}
\]
と置く．さらに
\[
 A=a+x+u,\qquad A\mid ab-1,\qquad
 B=\frac{ab-1}{A},\qquad v=B-b+y
\]
が $v\ne0$，$v\mid b-1$ を満たすとする．このとき
\begin{equation}\label{kl10:two-high-matrix}
 (p_{ij})=
 \begin{pmatrix}
 0&-A&x&u\\
 -B&0&-y&v\\
 r&-r-1&0&0\\
 (a-1)/u&(b-1)/v&d&0
 \end{pmatrix},\qquad
 d=-1-\frac{a-1}{u}-\frac{b-1}{v}
\end{equation}
となる．ここで $A,B$ は自動的に正である．
逆に，以上の条件を満たすデータは全 row sums と pair products を満たす．

正の weights が存在するための必要十分条件は，次の単一条件である：
\begin{equation}\label{kl10:two-high-positive}
 u<0\quad\text{または}\quad rud<a.
\end{equation}
これを満たすデータは命題~\ref{q9:four-line-criterion} により実際の局所作用 chart を与える．
\end{proposition}

\begin{proof}
分母一の二頂点を $C,D$ とすると，$p_{CD}p_{DC}=0$ である．
交換によって $p_{CD}=0$ としてよい．$C$ の row sum から
$p_{CA}+p_{CB}=-1$ であり，$a,b>1$ なので両者は非零整数である．
従って一方が正，一方が負であり，高分母頂点を交換して
$p_{CA}=r>0$，$p_{CB}=-r-1$ とできる．
pair products により $p_{AC}=x=(a-1)/r$，$p_{BC}=-y=-(b-1)/(r+1)$ となる．
$p_{AD}=u$ を選ぶと $u\mid a-1$ であり，$A$ の row sum は $p_{AB}=-a-x-u=-A$ を与える．
$|u|\leq a-1$ なので $A\geq x+1>0$ である．
残る pair product と $B$ の row sum から，$B$ と $v$ が指定通りに決まり，
$D$ の row sum が $d$ を与える．逆方向の検査は直接代入による．

正の辺は，常に存在する $A\leftrightarrow C$ に加え，
$u>0$ なら $A\leftrightarrow D$，$v>0$ なら $B\leftrightarrow D$，
$d>0$ なら $D\to C$ だけである．$A\leftrightarrow B$ と $B\leftrightarrow C$ の entries は負である．
$u<0$ の場合には長さ三以上の正の directed cycle はなく，正値性は自動的に成立する．
$u>0$ の場合には
\[
 A v=y(1-ru)-(x+u+1)<0
\]
なので $v<0$ である．長さ三以上の正の cycle が存在するとすれば
$A\to D\to C\to A$ のみで，その gain は $rud/a$ となる．
$d\leq0$ ならこの cycle は存在せず，同時に $rud<a$ は自動的に成立する．
補題~\ref{kl10:cycle-positive} により \eqref{kl10:two-high-positive} が必要十分である．
\end{proof}

この表示は任意の固定された $(a,b)$ に対し有限である．
$r$ と $u$ はそれぞれ $a-1$ の正の divisor と signed divisor から選び，
残りは整除性で決まる．この意味で，二高分母の局所問題は正値性も含めて有限な factorization 問題に帰着した．
これは全 $(a,b)$ にわたる可能な分母対を有限個に分類する主張ではない．

\begin{corollary}[分母三を含む二高分母の分類]\label{kl10:three-pair-classification}
$b>1$ を奇数とする．分母 $(3,b,1,1)$ を持つ正の局所四直線 chart が存在するための
必要十分条件は
\[
 b\in\{7,11\}
\]
である．この主張では $b$ の上限を仮定しない．
\end{corollary}

\begin{proof}
命題~\ref{kl10:two-high-parametrization} の標準表示における高分母の順序は
$(a,b)=(k,3)$ または $(3,k)$ である．まず $(k,3)$ の場合には $r=1$，
$x=k-1$，$y=1$ となる．$u\mid k-1$ より
\[
 A=2k-1+u\geq k,\qquad B=\frac{3k-1}{A}<3.
\]
$B=1$ は $u=k$ を強制し，$B=2$ は $v=B-3+1=0$ を強制するので，いずれも不可能である．

$(3,k)$ の場合には $r\in\{1,2\}$，$u\in\{-2,-1,1,2\}$ である．
式~\eqref{kl10:two-high-matrix} の整除条件を代入した全八場合は次の通りである．
表の最終列は，正値性をまだ課さずに全ての整数条件を満たす $k>1$ を示す．
\begin{center}
\small
\begin{tabular}{ccccc}
$r$ & $u$ & $A$ & $v$ & 可能な $k$ \\
\hline
$1$ & $-2$ & $3$ & --- & $\varnothing$ \\
$1$ & $-1$ & $4$ & $(k-3)/4$ & $7,11$ \\
$1$ & $1$ & $6$ & --- & $\varnothing$ \\
$1$ & $2$ & $7$ & $-(k+9)/14$ & $5,19,61,131$ \\
$2$ & $-2$ & $2$ & $5(k-1)/6$ & $\varnothing$ \\
$2$ & $-1$ & $3$ & --- & $\varnothing$ \\
$2$ & $1$ & $5$ & $-(k+8)/15$ & $7,37,127$ \\
$2$ & $2$ & $6$ & --- & $\varnothing$
\end{tabular}
\end{center}
$A=3,6$ の場合は $A\nmid3k-1$ であり，$(r,u)=(2,-2)$ の場合は
$(k-1)/v=6/5\notin\mathbb Z$ である．$(r,u)=(1,-1)$ では
$k=4v+3$ と $v\mid k-1$ から $v\mid2$ となり，$v>0$ より $v=1,2$ である．
これが $k=7,11$ を与え，$u<0$ なので両者は正の局所 chart を与える．

残る二行は正値性によって排除される．$(r,u)=(1,2)$ では
$v=-z$ と置くと $k=14z-9$，$z\mid10$ であり，
\[
 rud=24-\frac{20}{z}\geq4>3.
\]
$(r,u)=(2,1)$ では $k=15z-8$，$z\mid9$ であり，
\[
 rud=24-\frac{18}{z}\geq6>3.
\]
いずれも条件~\eqref{kl10:two-high-positive} に反する．
\end{proof}

\subsection{複数高分母と局所最大原理への反例}
\label{kl10:counterexamples}

複数の高分母を持つ正の局所 chart は存在する．例えば
\[
 q=(7,3,1,1),\quad
 P=\begin{pmatrix}
 0&-5&-3&1\\-4&0&2&-1\\-2&1&0&0\\6&-2&-5&0
 \end{pmatrix},\qquad
 w=(13/2,5/4,1,1)
\]
は全条件を満たし，$D=(-7,3,5,-1)$ を与える．
同じ素数を共有する二分母についても
\[
 q=(7,7,1,1),\quad
 P=\begin{pmatrix}
 0&-8&-2&3\\-6&0&1&-2\\-3&6&0&-4\\2&-3&0&0
 \end{pmatrix},\qquad
 w=(13/6,13/2,1,1)
\]
が正の局所 chart を与え，$D=(0,-2,0,2)$ となる．
従って「一つの core の高分母頂点は高々一つ」という局所命題は成立しない．

最大分母の頂点の係数が負であるという命題は，最大分母が一意の場合にも成立しない．
具体的には
\begin{equation}\label{kl10:unique-max-counterexample}
 q=(13,7,7,1),\quad
 P=\begin{pmatrix}
 0&-10&-9&6\\-9&0&8&-6\\-10&6&0&-3\\2&-1&-2&0
 \end{pmatrix},\qquad
 w=(25/12,1,23/20,1)
\end{equation}
について $D=(4,-2,-4,2)$ であり，一意の最大分母十三の頂点で $D=4>0$ となる．
正の numerators に対応する四つの basin areas は
$1/26,1/12,1/140,1/70$ で，残りは numerators が非正なので正である．
全 row sums と pair products は整数の直接計算で確認できる．

全四分母が同じである局所模型もある．
\begin{equation}\label{kl10:uniform-eleven}
 q=(11,11,11,11),\quad
 P=\begin{pmatrix}
 0&-20&-15&24\\-6&0&10&-15\\-8&12&0&-15\\5&-8&-8&0
 \end{pmatrix},\qquad
 w=(1,23/21,1,35/16)
\end{equation}
では，正の numerators に対応する basin areas は $1/176,1/231,1/231,1/176$ であり，
$D=(-2,5,2,-5)$ となる．
いずれの例も命題~\ref{q9:four-line-criterion} による幾何学的な存在例であるが，
五つの cores を接続する大域 atlas の存在を主張するものではない．
大域の不存在を示す場合には，局所の分母や係数の大小だけに依拠せず，
対向端の clutching と全体の incidence を使う必要がある．

\subsection{非正の numerators だけでは高分母は生じない}
\label{lit10:nonpositive}

\begin{lemma}\label{lit10:nonpositive-q1}
四つの正整数 $q_i$ と非正整数 $p_{ij}$ が
$\sum_{j\ne i}p_{ij}=-q_i$ と $q_iq_j-p_{ij}p_{ji}=1$ を満たすなら，
すべての $q_i$ は一である．
\end{lemma}

\begin{proof}
$a_{ij}=-p_{ij}\ge0$ と置く．$q_i=1$ の vertex があり，
異なる二つの vertices の分母が一より大きいなら，その二方向の
$a_{ij}$ がともに正となり $\sum_j a_{ij}=1$ に反する．
従って分母一の vertex がある場合，高分母 vertex は高々一つである．
唯一の高分母を $q>1$ とすると，他の各 vertex の outgoing row は
高分母 vertex への一を除いて零となる．
高分母 vertex の三つの outgoing entries はすべて $q-1$ なので，
$q=3(q-1)$ となり整数性に反する．

残る場合はすべての $q_i>1$ である．
このときすべての $a_{ij}\ge1$ なので $q_i\ge3$ および
$a_{ij}\le q_i-2$ となる．しかし
\[
 1=q_iq_j-a_{ij}a_{ji}
   \ge q_iq_j-(q_i-2)(q_j-2)
   =2q_i+2q_j-4\ge8
\]
は矛盾である．
\end{proof}

従って高分母の局所模型は必ず正の numerator を含む．
通常の arithmetical structures の $Z$-matrix 仮定\cite[\S\S1--3]{CorralesValencia}へ直接帰着する方法では
高分母模型の全体を扱えない．

% v11 staging module. General two-high sign theorem and triangular-support obstruction.
% Assumes kl10:two-high-parametrization, kl10:two-high-positive,
% and the already established opposing-end cancellation of D.

\section{二高分母の符号と三角形成分の接着障害}
\label{ls11:section}

二高分母の factorization 表示の正値性条件は，分母が異なる場合にも
単純な符号条件まで強められる．これを共有 puncture の整数 clutching と組み合わせると，
高分母 support の一つの連結成分が三角形になる場合を，分母の上限なしで排除できる．

\begin{theorem}[二高分母の正値性と符号]\label{ls11:two-high-signs}
命題~\ref{kl10:two-high-parametrization} の標準表示に従う
分母 $(a,b,1,1)$，$a,b>1$，の整数 numerator matrix を考える．
正の局所 weights が存在することは
\[
 u<0<v
\]
と同値である．従って，正の局所模型の各高分母 row は，もう一つの高分母へ負の entry を持ち，
二つの低分母へは正と負の entry を一つずつ持つ．
この結論は $a=b$ を仮定しない．
\end{theorem}

\begin{proof}
標準表示の記号を用いる．$A=a+x+u\geq x+1\geq2$ であり，
pair-product 条件と $x=(a-1)/r$，$y=(b-1)/(r+1)$ から
\begin{equation}\label{ls11:linear-v}
 A(r+1)v=(1-ru)b-(A+r+1)
\end{equation}
を得る．$u<0$ なら $1-ru>0$ である．もし $v\leq-1$ なら
\[
 (1-ru)b=A(r+1)v+A+r+1
 \leq r+1-Ar\leq0
\]
となり，$b>1$ に反する．$v\ne0$ なので $v>0$ である．
また，$u<0$ なら既知の正値性条件~\eqref{kl10:two-high-positive} が成立する．

残る $u>0$ の場合が正の weights を持たないことを示す．
この場合は $v<0$ なので $v=-z$，$z>0$ と書く．
\[
 S=\frac{a-1}{u},\qquad T=\frac{b-1}{z},\qquad
 d=T-S-1
\]
と置くと $S,T$ は正の整数である．正値性は
$rud<a=uS+1$ を要求し，整数性から
\begin{equation}\label{ls11:cycle-bound}
 rd\leq S
\end{equation}
となる．次の正の量と整数を導入する：
\[
 \begin{split}
 F&=(2r+1)S+r-r^2d,\qquad G=(r+1)(S+r),\\
 H&=r(r+1),\qquad K=r(r+T+1).
 \end{split}
\]
pair-product 条件を $a=uS+1$，$b=zT+1$，$T=S+d+1$ によって書き直すと
\begin{equation}\label{ls11:bilinear-positive}
 u(zF-G)=H-zK
\end{equation}
を得る．式~\eqref{ls11:cycle-bound} より
\[
 F\geq(r+1)S+r>0.
\]
一方，$z\geq1$，$T\geq1$ より
$H-zK\leq H-K=-rT<0$ である．従って $zF<G$ であり，特に
\[
 0<G-F=r\{r(d+1)-S\}.
\]
従って
\begin{equation}\label{ls11:short-interval}
 rd\leq S<r(d+1).
\end{equation}
$S>0$ なので $d\geq0$ である．

$d=0$ なら，二つの低分母頂点間の entries は両方とも零である．
これらの頂点を交換すると，再び同じ標準表示になり，新しい $r,S$ はそれぞれ $S,r$ となる．
式~\eqref{ls11:short-interval} は元の標識で $S<r$，交換後には $r<S$ を与えるので矛盾する．

$d\geq1$ なら $S\geq r$ であり，
\[
 1\leq z<\frac GF
 \leq\frac{(r+1)(S+r)}{(r+1)S+r}
 =1+\frac{r^2}{(r+1)S+r}<2.
\]
従って $z=1$ である．$S=rd+e$，$0\leq e<r$ と書けば
\[
 T=S+d+1=(r+1)d+e+1.
\]
標準表示の整除条件 $r+1\mid b-1=zT$ は，$z=1$ より $r+1\mid e+1$ を要求する．
しかし $1\leq e+1\leq r$ なので矛盾する．これで $u>0$ は排除された．
\end{proof}

この定理は，既知の factorization 表示における cycle 不等式を削除するものではなく，
その整数条件を全て用いて，正値な場合を $u<0<v$ にちょうど絞り込む結果である．
局所の正の chart の実現は命題~\ref{q9:four-line-criterion} に従う．

\begin{lemma}[両端が低分母に接する puncture]\label{ls11:low-puncture}
共有 branch の分母を $q>1$ とし，その同じ marked puncture が，
両側の core において分母一の ray に向かうとする．
共有 branch からその ray へ向かう二つの numerators を $p,p'$ と書く．
整数 clutching $q\mid p+p'$ が成立するなら，次のいずれかに限られる：
\[
 p'=-p,\qquad
 \{p,p'\}=\{1,q-1\},\qquad
 \{p,p'\}=\{-1,-(q-1)\}.
\]
この三場合のいずれでも，逆向きの numerators の和は元の和と等しい：
\[
 \frac{q-1}{p}+\frac{q-1}{p'}=p+p'.
\]
meridian node 表示では，その puncture の node 数はそれぞれ $1,0,2$ である．
\end{lemma}

\begin{proof}
crossing determinant により $p,p'$ は $q-1$ の非零 signed divisors である．
従って $|p+p'|<2q$ であり，その和は $-q,0,q$ のいずれかとなる．
和が零なら第一の場合である．和が $q$ なら両者は正であり，
両者とも $q-1$ より小さい divisor なら各々 $(q-1)/2$ 以下なので矛盾する．
従って一方が $q-1$，他方が一である．負の場合も同様である．
逆向きの和についての等式と node 数は，各場合への直接代入で従う．
\end{proof}

\begin{theorem}[高分母の三角形成分の不存在]\label{ls11:no-triangle-component}
現在の全平面五次の構成条件の下で，大域作用適合 atlas が存在するとする．
その高分母 support
\[
 H=\bigl\{\{i,j\}:q_{ij}>1\bigr\}\subset K_5
\]
は，三角形を一つの連結成分として持たない．
これは分母の大きさによらない一般障害である．
証明に用いるのは局所の正値性，row sums，crossing determinants，
marked puncture の整数 clutching，および対向端の $D$ の相殺である．
node 数の非負性や各直線の全 node 数五は，この定理には必要ない．
\end{theorem}

\begin{proof}
三角形の core labels を $0,1,2$ とし，添字は三を法として読む．
辺 $\{i,i+1\}$ の分母を $q_i>1$ と書く．
三角形成分という仮定から，この三つの cores はそれぞれ二高分母を持ち，
残る二つの core labels へ向かう分母は一である．
各共有 branch の二つの外側 punctures は，それら二つの固定された labels に対応する．

定理~\ref{ls11:two-high-signs} より，その二つの outgoing entries は各端で一正一負である．
同じ marked punctures を接続するとき，符号が対応するなら補題~\ref{ls11:low-puncture} により
二つの和は $q_i,-q_i$，符号が逆に対応するなら零，零となる．
従って二つの外側 punctures の outgoing sums の総和は零であり，
同補題により incoming sums の総和も零である．

core $i$ における二つの高分母 rays の numerators を
\[
 p_{i;i-1,i+1}=-\alpha_i,\qquad
 p_{i;i+1,i-1}=-\beta_i,
 \qquad \alpha_i,\beta_i>0
\]
と書く．局所 crossing determinant は
\begin{equation}\label{ls11:triangle-products}
 \alpha_i\beta_i=q_{i-1}q_i-1
\end{equation}
である．両端の outgoing row sums から，辺 $\{i,i+1\}$ について
\[
 \beta_i+\alpha_{i+1}=2q_i
\]
となる．また対向端の $D$ の相殺と incoming sums の総和零から
\[
 \alpha_i+\beta_{i+1}=2q_i
\]
を得る．差を取ると $\alpha_i-\beta_i$ は $i$ によらない整数 $\delta$ となる．
meridian node 表示を用いる場合には，最初の等式は三角形に向かう puncture の node 数三を意味するが，
ここではこの個数や全 node budget を別の仮定として用いていない．

$T_i=(\alpha_i+\beta_i)/2>0$ と置くと，
\[
 q_i=\frac{T_i+T_{i+1}}2,
 \qquad
 (T_{i-1}+T_i)(T_i+T_{i+1})=4T_i^2-\delta^2+4.
\]
$T_i$ が最大となる頂点では左辺は $4T_i^2$ 以下なので $\delta^2\geq4$，
最小となる頂点では逆に $\delta^2\leq4$ である．
従って $\delta^2=4$ であり，最大頂点での等号から三つの $T_i$ は全て等しい．
従って三辺の分母も全て同じ $q$ となり，各 core の二つの high--high numerators は
$-q-1,-q+1$ となる．

この特別な二高分母模型を標準表示で調べる．
$m=q-1=t r(r+1)$，$x=t(r+1)$，$y=tr$ とし，$u=-s<0$，$v=z>0$ と置く．
high--high entries は $-A,-B$，
\[
 A=q+x-s,\qquad B=q-y+z,\qquad \{A,B\}=\{q-1,q+1\}
\]
である．もし $A=q-1$ なら $s=x+1$ であり，$s\mid m=rx$ から $x+1\mid r$ となる．
しかし $x=t(r+1)\geq r+1$ なので不可能である．
従って $A=q+1$ であり，$s=x-1$，$z=y-1$ となる．
$s\mid rx$ より $x-1\mid r$，一方 $x-1\geq r$ なので $t=1$，$s=r$ である．
さらに $z=r-1>0$，$z\mid r(r+1)$ より $r-1\mid2$ となる．
従って $r=2,3$，すなわち $q=7,13$ に限られる．

直接計算すると，その二高分母における $D$ の組は，順序を除き
\[
 q=7:\quad (0,-2),\qquad
 q=13:\quad (0,-1)
\]
である．どちらも高分母二頂点の係数和は負である．
三つの cores にわたりこの和を足すと負になるが，各三角形辺の対向端では相殺するため零でなければならない．
この矛盾が結論を与える．
\end{proof}

本定理は，高分母 support が三角形を部分 graph として含む場合を全て排除するものではない．
三つの cores のいずれかにさらに高分母 branch が接続する場合には，
二高分母の符号定理をその core へ適用できない．
また，四角形・五角形や一般の path 成分の排除は，本定理だけからは従わない．

% v12 staging module; do not edit the immutable v11 bundle.
% Inputs: the standard two-high factorization (kl10:two-high-parametrization)
% and its positivity/sign theorem (ls11:two-high-signs).

\section{異なる二高分母に対する負の係数}
\label{td12:section}

二つの高分母が異なる場合，その二頂点の係数 $D$ の和は正にもなり得る．
しかし，二つの係数がともに非負になることはない．
以下では分母の上限を置かず，実際には奇数性も用いずに，この局所的な障害を証明する．

\begin{lemma}[標準形に残る低分母間 entry の負性]
\label{td12:low-entry}
正の局所 weights を持つ二高分母標準形において
\[
 u=-s<0,\qquad v=z>0,\qquad
 S=\frac{a-1}{s},\qquad T=\frac{b-1}{z}
\]
と置く．このとき $S\leq T$ であり，残る低分母間 entry
$d=S-T-1$ は負である．
\end{lemma}

\begin{proof}
標準形の整数条件から $r,s,z,S,T$ は正の整数であり，
$r\mid a-1$，$r+1\mid b-1$ である．
$d\geq0$，従って $S\geq T+1$ と仮定する．
\[
 R=r(r+1),\qquad
 K=(r+1)^2S-r^2T-R=(2r+1)S+r^2d-r>0
\]
と置く．高分母間の pair-product 条件に $a=sS+1$，$b=zT+1$ を代入すると
\begin{equation}\label{td12:low-entry-product}
 0=R+s(r+1)(S-r)+zr(r+1-T)+szK
\end{equation}
となる．$S\geq r$，$T\leq r+1$ なら右辺の全項は非負で，$R>0$ である．
$T>r+1$ なら第二項は正であり，
\[
 K-r(T-r-1)=(r+1)\{(r+1)S-rT\}>0
\]
から第三項と第四項の和も正となる．
残る $S<r$ の場合には $T<S<r$ で第三項は正である．
整除性 $r+1\mid zT$ より $zT\geq r+1$ であり，
\[
 K-T(r-S)=(r+1)\{(r+1)S-rT-r\}+ST>0
\]
なので $zK>(r+1)(r-S)$ となり，第二項と第四項の和も正である．
最後の括弧内は $S\geq T+1$ により $T+1$ 以上である．
全場合で式~\eqref{td12:low-entry-product} と矛盾する．
\end{proof}

\begin{corollary}[二高分母の明示的な正 weights]
\label{td12:explicit-weights}
二高分母標準形の整数条件と $u=-s<0<v=z$ が成立するとき，
\[
 (w_A,w_B,w_C,w_D)
 =\left(a,b,\frac{a-1}{r}+\frac1{2r},z+\frac1{2T}\right)
\]
は正の有理 weights を与える．
\end{corollary}

\begin{proof}
定理~\ref{ls11:two-high-signs} と補題~\ref{td12:low-entry} により，
正の numerators は $A\leftrightarrow C$ と $B\leftrightarrow D$ の四つだけである．
$B_{ij}=w_j-w_i p_{ij}/q_i$ と書くと，表示した weights に対して
\[
 B_{AC}=\frac1{2r},\qquad B_{CA}=\frac12,\qquad
 B_{BD}=\frac1{2T},\qquad B_{DB}=\frac12.
\]
他の numerators は非正であり，それぞれの basin も正となる．
\end{proof}

\begin{theorem}[二高分母の少なくとも一方の係数は負]
\label{td12:one-negative}
分母 $(a,b,1,1)$，$a,b>1$，を持つ整数局所模型が正の weights を持つとする．
二つの高分母頂点の係数を $D_A,D_B$ とすると
\[
 \min(D_A,D_B)<0.
\]
この結果は $a=b$ を仮定せず，分母の上限，分母の奇数性，
puncture の clutching，node 個数のいずれも用いない．
\end{theorem}

\begin{proof}
定理~\ref{ls11:two-high-signs} と標準形により，頂点を適切に標識すると
\[
 P=\begin{pmatrix}
 0&-A&x&-s\\
 -B&0&-y&z\\
 r&-r-1&0&0\\
 -S&T&S-T-1&0
 \end{pmatrix},
\quad
 x=\frac{a-1}{r},\quad y=\frac{b-1}{r+1},
\]
\[
 a=sS+1,\quad b=zT+1,\quad
 A=a+x-s>0,\quad B=b-y+z>0,\quad AB=ab-1
\]
と書ける．補題~\ref{td12:low-entry} により $S\leq T$ である．
次の整数を導入する：
\begin{equation}\label{td12:auxiliaries}
 h=S-r,\qquad k=T-r-1,\qquad
 \alpha=x-s=\frac{sh}{r},\qquad
 \beta=y-z=\frac{zk}{r+1}.
\end{equation}
$A=a+\alpha$，$B=b-\beta$ なので pair product は
\begin{equation}\label{td12:pair-alpha}
 a\beta-b\alpha+\alpha\beta=1
\end{equation}
となる．$\alpha=0$ または $\beta=0$ は $a,b>1$ と整数性に反する．
$\alpha>0>\beta$ なら式~\eqref{td12:pair-alpha} の左辺は負である．
$\alpha<0<\beta$ なら同左辺は
$A\beta-b\alpha\geq A+b>1$ となる．
従って $\alpha,\beta,h,k$ はすべて非零で同符号である．

\[
 M=rs+1,\qquad N=(r+1)z+1,\qquad
 \delta=\alpha-\beta,\qquad L=M-N+r\delta
\]
と置く．式~\eqref{td12:auxiliaries} と~\eqref{td12:pair-alpha} から
\begin{equation}\label{td12:bezout}
 a=M+r\alpha,\qquad b=N+(r+1)\beta,\qquad
 M\beta-N\alpha=1
\end{equation}
であり，特に $\gcd(\alpha,\beta)=1$ である．また
\begin{equation}\label{td12:defect-identities}
 \alpha L=1+a\delta,\qquad
 D_A=h-L,\qquad D_B=L+\delta-k,
 \qquad D_A+D_B=h-k+\delta.
\end{equation}
以下の三場合で結論を示す．

\paragraph{第一の場合：$\delta<0$．}
整数性と $a>1$ より $1+a\delta<0$ である．
$\alpha>0$ なら $L<0$ であり，$k>0$ より $D_B=L+\delta-k<0$ となる．
$\alpha<0$ なら $L>0$ であり，$h<0$ より $D_A=h-L<0$ となる．

\paragraph{第二の場合：$\delta\geq0$，$\alpha,\beta>0$．}
両方の係数が非負と仮定する．式~\eqref{td12:defect-identities} より
\[
 \delta\geq k-h,\qquad \alpha h\geq1+a\delta.
\]
式~\eqref{td12:bezout} を用いると
\begin{equation}\label{td12:sign-ratio}
 \alpha\beta
 \left(\frac{r^2}{h}-\frac{(r+1)^2}{k}\right)=\delta+1>0
\end{equation}
なので
\[
 k-h>\frac{(2r+1)h}{r^2}>\frac hr.
\]
一方，$a=s(r+h)+1$ と $\alpha=sh/r$ から
\[
 \delta<\frac{\alpha h}{a}
 <\frac{h^2}{r(r+h)}<\frac hr.
\]
従って $\delta\geq k-h>h/r>\delta$ となり矛盾する．
この議論は $\delta=0$ も含む．

\paragraph{第三の場合：$\delta\geq0$，$\alpha,\beta<0$．}
\[
 \alpha=-c,\quad\beta=-e,\quad h=-H,\quad k=-K
\]
と書く．$1\leq c\leq e$，$1\leq H\leq r-1$，$1\leq K\leq r$ であり，
$S\leq T$ は $K\leq H+1$ と同値である．
式~\eqref{td12:sign-ratio} は
\[
 Ece=e-c+1,\qquad
 E=\frac{(r+1)^2}{K}-\frac{r^2}{H}>0
\]
となる．右辺が $e$ 以下なので $Ec\leq1$ である．
もし $K\leq H$ なら
\[
 E\geq\frac{2r+1}{H}>2
\]
となり矛盾する．従って $K=H+1$ である．
\[
 w=S=T=r-H\geq1
\]
と置くと
\begin{equation}\label{td12:negative-reduction}
 E=1-\frac{w^2}{H(H+1)}\in(0,1),\qquad
 E=\frac{\delta+1}{c(c+\delta)},\quad \delta=e-c.
\end{equation}
$c=1$ なら $E=1$ となるので $c\geq2$ である．
また $\delta=0$ なら $c=e$ と $\gcd(c,e)=1$ より $c=1$ なので，$\delta\geq1$ である．

$\alpha L=1+a\delta$ より $U=-L=(1+a\delta)/c>0$ と書ける．
$s=rc/H$，$r=H+w$，$a=sw+1$ を代入すると
\[
 U=w\delta+\frac{w^2\delta}{H}+\frac{\delta+1}{c},\qquad
 D_B=H+1+\delta-U.
\]
式~\eqref{td12:negative-reduction} と $c\geq2$ から $E<1/2$ である．
$\delta\geq2$ なら
\[
 D_B=(H+1)(1-\delta+\delta E)
       +\delta(1-w)-\frac{\delta+1}{c}<0,
\]
なぜなら $1-\delta+\delta E<1-\delta/2\leq0$，$w\geq1$ だからである．

残る $\delta=1$ では
\[
 D_B=1-w+\frac{2(H-c)}{c(c+1)}.
\]
$H<c$ なら直ちに $D_B<0$ である．$H\geq c\geq2$ なら
\[
 D_B\leq\frac{H+1}{3}-w<0.
\]
最後の厳密な不等式は
\[
 w^2=H(H+1)\left(1-\frac2{c(c+1)}\right)
 \geq\frac23H(H+1)>\frac{(H+1)^2}{9}
\]
から従う．これで全場合が尽くされ，結論を得る．
\end{proof}

第三の場合の証明は，その枝が必ず $S=T$，従って $d=-1$ に属することも示す．
ただし本定理は，高分母二頂点の和が負であるとは主張しない．
異なる分母ではその和が正になる既知の正局所模型とも整合する．

\subsection{道成分の排除と閉路上の係数の向き}
\label{pc12:section}

高分母 support の次数は，その core における高分母 rays の本数に等しい．
従って，一高分母の負係数と定理~\ref{td12:one-negative} を組み合わせると，
異なる高分母を許した場合にも，次数二以下の成分を制限できる．

\begin{theorem}[高分母の道成分の不存在]\label{pc12:no-path}
現在の全平面五次の構成条件の下で，大域作用適合 atlas が存在するとする．
高分母 support
\[
 H=\{ij:q_{ij}>1\}\subset K_5
\]
は，辺を持つ path を一つの連結成分として持たない．
分母の一致や上限を仮定せず，node 数の非負性と全 node budget も用いない．
\end{theorem}

\begin{proof}
path 成分の頂点を順に $i_0,i_1,\ldots,i_m$，$m\geq1$ とする．
各葉は一高分母 core なので，命題~\ref{qh9:isolated-local} により
\[
 D^{i_0}_{i_1}<0,\qquad D^{i_m}_{i_{m-1}}<0.
\]
対向端相殺から $D^{i_1}_{i_0}>0$ となる．
$m=1$ なら既に矛盾する．$m\geq2$ のとき，$i_1$ は二高分母 core であるから，
定理~\ref{td12:one-negative} により，もう一つの高分母係数は
$D^{i_1}_{i_2}<0$ でなければならない．
同じ議論を順に繰り返すと，最後の辺で $D^{i_m}_{i_{m-1}}>0$ を得る．
これは最後の葉の負係数に反する．
\end{proof}

\begin{proposition}[閉路に沿う非零係数]\label{pc12:cycle-orientation}
同じ条件の下で，$H$ の一つの連結成分が cycle であるとする．
その全ての高分母係数は非零であり，各 core には正と負の係数が一つずつある．
各辺を負の係数を持つ端から正の係数を持つ端へ向けると，
全辺の向きは一つの有向閉路となる．
特に，$H$ の最大次数が二以下で $H$ が辺を持つなら，
辺を持つ成分は一つだけで，それは四角形または五角形である．
\end{proposition}

\begin{proof}
cycle の頂点数と辺数を $m$ とする．各 core は二高分母を持つので，
定理~\ref{td12:one-negative} により，負の高分母係数が少なくとも一つある．
従って負の係数の総数は少なくとも $m$ である．
一方，各辺では対向端相殺により，両係数が零であるか，
一方が負でもう一方が正である．従って負の係数は各辺に高々一つであり，
総数は高々 $m$ である．等号が成立するため，各 core に負の係数はちょうど一つで，
零の係数を持つ辺はない．各頂点の入次数と出次数は共に一であるから，向きは閉路に沿って揃う．

最大次数二以下の有限単純 graph の，辺を持つ連結成分は path または cycle である．
path は定理~\ref{pc12:no-path}，三角形は定理~\ref{ls11:no-triangle-component} で排除される．
$H\subset K_5$ なので残るのは四角形または五角形であり，二つのそのような成分は共存できない．
\end{proof}

\begin{proposition}[木成分の中心への予備的な還元]\label{pc12:tree-reduction}
同じ条件の下で，$H$ の辺を持つ連結成分が tree であるなら，
その graph は $K_{1,3}$，または $K_{1,3}$ の一辺を一度細分した五頂点の tree に限られる．
いずれの場合にも，唯一の三価 core における三つの高分母係数は全て正でなければならない．
\end{proposition}

\begin{proof}
五頂点以下の tree で最大次数二以下のものは path なので排除済みである．
最大次数四のものは $K_{1,4}$ に限られるが，各葉の負係数と対向端相殺により，
中心の四つの係数は全て正となる．これは row sums から従う
$\sum_jD^i_j=0$ に反する．
最大次数三のものは，記載した二種類に限られる．
葉から中心へ進む各道において，定理~\ref{pc12:no-path} の証明と同じ符号伝播を用いると，
中心に到達する全ての高分母係数が正になる．
\end{proof}

これらは連結成分についての主張である．別の高分母辺が接続する path 部分 graph や，
三高分母以上を持つ core の任意の係数配置を排除してはいない．
後の定理~\ref{fo13:no-tree} では，葉からの係数の大きさを用いて全ての木成分を排除する．
一般三高分母の全高係数の正性については，独立した未解決の局所問題として残る．

% v12: an unbounded leaf-to-degree-two denominator escalation, using neither
% node positivity nor a finite node budget.

\subsection{葉から二価頂点へ進むときの分母の増大}
\label{pe12:section}

\begin{proposition}[葉に隣接する二価頂点]\label{pe12:leaf-escalation}
全平面五次の高分母 support の辺 $ij$ において，$i$ の次数は一，$j$ の次数は二であるとする．
$j$ のもう一つの高分母の隣接頂点を $k$，
$q_{ij}=a>1$，$q_{jk}=b>1$ と書く．
正の局所 charts，marked puncture の整数 clutching，対向端の $D$ の相殺が成立するなら
\[
 b>a.
\]
より正確には，$i$ の chart における $ij$ から $ik$ への numerator を $x$ とすると，
$x$ は $a-1$ の非零 signed divisor であり，
\begin{equation}\label{pe12:exact-escalation}
 b=c^2a+c\left(x+\frac{a-1}{x}\right)+1,
 \qquad
 c=\begin{cases}1,&x>0,\\2,&x<0.\end{cases}
\end{equation}
node 数の非負性や全 node budget は仮定しない．
\end{proposition}

\begin{proof}
残る二つの labels を $l,m$ とする．$il,im,jl,jm$ の分母は全て一である．
$i$ の一高分母 row は一正二負である一方，$j$ の二高分母 row は
$jk$ へ負，$jl,jm$ へ一正一負である．
$j$ における $ij\to jk$ と $jk\to ij$ の entries をそれぞれ
$-\alpha,-\beta$ と書くと，$\alpha,\beta>0$ かつ
\[
 \alpha\beta=ab-1
\]
である．

$l,m$ の二つの共通 low-low punctures における，両端の outgoing numerators の総和を $\Lambda a$ と書く．
低低 puncture の clutching 補題により，符号が逆の一組は和零，
同符号の一組は和 $a$ または $-a$ である．
$x>0$ なら $i$ の残る二 entries は両方負なので，$\Lambda=-1$ である．
$x<0$ なら両端で一正一負となり，同符号同士を合わせても逆符号同士を合わせても
二組の総和は零なので，$\Lambda=0$ である．
同じ補題により，これら二つの punctures における incoming numerators の総和も $\Lambda a$ である．

両端の outgoing row sums を加えると
\[
 -\alpha+x+\Lambda a=-2a.
\]
また対向端の $D$ の相殺は
\[
 -2a+\beta-\frac{a-1}{x}-\Lambda a=0
\]
を与える．従って $c=2+\Lambda$ と置けば
\[
 \alpha=ca+x,\qquad \beta=ca+\frac{a-1}{x}.
\]
この二式を $\alpha\beta=ab-1$ に代入して $a>0$ で割ると
\eqref{pe12:exact-escalation} を得る．

$x>0$ なら $c=1$ なので $b-a=x+(a-1)/x+1>0$ である．
$x=-s<0$ なら $s>0$，$S=(a-1)/s>0$ と置けて
\[
 b=4a-2(s+S)+1,
 \qquad b-(2a+1)=2(s-1)(S-1)\geq0.
\]
この場合も $b\geq2a+1>a$ である．
\end{proof}


% v12 staging module. Depends on quintic_two_high_defect_v12.tex.

\section{低分母間の最小 entry と分母の下降}
\label{ml12:section}

二高分母の正の標準形では，低分母間の二つの entries のうち一方が零で，他方が負である．
この非零 entry が $-1$ の場合には，二つの高分母の大小と係数 $D$ の符号が完全に対応する．
この結果は，非零 entry がさらに小さい場合へ拡張しない．

\begin{lemma}[低分母間 entry の符号による二分]
\label{ml12:dichotomy}
定理~\ref{td12:one-negative} の標準形と補助変数を用いる．
非零の低分母間 entry を $d=S-T-1<0$ とすると
\[
 \alpha,\beta<0\quad\Longleftrightarrow\quad d=-1.
\]
他方，$\alpha,\beta>0$ なら
\[
 rT>(r+1)S,\qquad d\leq-3.
\]
従って，この標準形の非零の低分母間 entry が $-2$ となることはない．
\end{lemma}

\begin{proof}
$\alpha,\beta<0$ のとき
\[
 \alpha=-c,\quad\beta=-e,\quad h=-H,\quad k=-J
\]
と書く．$c,e>0$，$1\leq H\leq r-1$，$1\leq J\leq r$ であり，
$S\leq T$ より $J\leq H+1$ である．
式~\eqref{td12:sign-ratio} は，$c,e$ の大小によらず
\[
 Ece=e-c+1,\qquad E=\frac{(r+1)^2}{J}-\frac{r^2}{H}
\]
を与える．従って $Ec\leq1$ である．
$J\leq H$ なら $E\geq(2r+1)/H>2$ となって矛盾するため，
$J=H+1$，すなわち $S=T$，$d=-1$ となる．
この段階では $E$ が正であるとは仮定していない．

$\alpha,\beta>0$ なら $S\geq r+1$，$T\geq r+2$ である．
\[
 R=r(r+1),\qquad F=(r+1)^2S-r^2T-R
\]
と置くと，式~\eqref{td12:low-entry-product} は
\[
 0=R+s(r+1)(S-r)+z\{sF-r(T-r-1)\}
\]
となる．もし $rT\leq(r+1)S$ なら
\[
 F-r(T-r-1)=(r+1)\{(r+1)S-rT\}\geq0.
\]
$T\geq r+2$ より $F>0$ であり，$s\geq1$ から最後の波括弧も非負となるので矛盾する．
従って $r(T-S)>S>r$，すなわち $T-S\geq2$，$d\leq-3$ となる．
$\alpha,\beta$ は非零で同符号であるため，これで二分と同値が従う．
\end{proof}

\begin{theorem}[最小の低分母間 entry における分母と係数の符号]
\label{ml12:denominator-descent}
二高分母 $(a,b,1,1)$ の正の整数局所模型で，非零の低分母間 entry が $-1$ であるとする．
このとき $a\ne b$ であり，大きい分母の係数 $D$ は負，小さい分母の係数は正である．
標準形の標識では
\[
 (a-b)D_A<0,\qquad (a-b)D_B>0.
\]
この結論にも分母の上限や奇数性は必要ない．
\end{theorem}

\begin{proof}
補題~\ref{ml12:dichotomy} より $\alpha,\beta<0$ であり，
\[
 \alpha=-c,\quad \beta=-e,\quad h=-H,\quad k=-H-1,
 \quad w=S=T=r-H\geq1
\]
と書ける．前証明の式は
\begin{equation}\label{ml12:E-equation}
 Ece=e-c+1,\qquad
 E=1-\frac{w^2}{H(H+1)}
\end{equation}
となる．また
\begin{equation}\label{ml12:s-z}
 s=c+\frac{cw}{H},\qquad z=e+\frac{ew}{H+1},\qquad
 a=sw+1,\quad b=zw+1.
\end{equation}
$H^2<H(H+1)<(H+1)^2$ なので $E=0$ は不可能である．

$w\geq H+1$ なら $E<0$ である．式~\eqref{ml12:E-equation} から
$c\geq e+2$，従って $\delta=e-c\leq-2$ となる．
式~\eqref{ml12:s-z} より $s>z$，従って $a>b$ である．
$\alpha<0$ と $1+a\delta<0$ により
$\alpha L=1+a\delta$ は $L>0$ を与えるので $D_A=-H-L<0$ となる．
さらに $j=c-e\geq2$ と書けば
\[
 D_B
 =j\left(w+\frac{w^2}{H}-1\right)+H+1+\frac{j-1}{c}>0.
\]
ここでは $a/c=w+w^2/H+1/c$ を
$D_B=L+\delta+H+1$ に代入した．

$w\leq H$ なら $1/(H+1)\leq E<1$ である．
式~\eqref{ml12:E-equation} から $e\geq c$ となる．
$c=1$ なら $E=1$ で矛盾し，$c=e$ なら
$\gcd(c,e)=1$ より $c=1$ となるため，$c\geq2$，
$\delta=e-c\geq1$ である．
\[
 Ec=\frac{e-c+1}{e}<1
\]
と $E\geq1/(H+1)$ から $c\leq H$ を得る．
式~\eqref{ml12:s-z} を引くと
\[
 s-z=-\delta+\frac{w(c-H\delta)}{H(H+1)}<0
\]
なので $a<b$ である．
これは定理~\ref{td12:one-negative} の証明の第三の場合に属し，
そこで示した厳密な不等式より $D_B<0$ となる．
また $D_A+D_B=\delta+1>0$ だから $D_A>0$ である．
\end{proof}

\begin{corollary}[全頂点で低分母間 entry が最小の閉路は存在しない]
\label{ml12:no-minimal-cycle}
高分母 support の一つの連結成分が閉路であり，その各 core で非零の低分母間 entry が $-1$ であるとする．
正の局所 weights と共有高分母 branch における対向 $D$ の相殺を同時に満たす atlas は存在しない．
\end{corollary}

\begin{proof}
各 core の二つの高分母端のうち，負の $D$ を持つ端から正の $D$ を持つ端へ向きを付ける．
定理~\ref{ml12:denominator-descent} より，この向きに沿って分母は厳密に減少する．
共有 branch の両端では $D$ が相殺するため，向きは閉路全体で一貫する．
閉路を一周して分母が厳密に減少することは不可能である．
\end{proof}

ここで仮定する entry は，各 core の二つの低分母頂点の間の\emph{非零} entry である．
大域の固定標識に従って一方向だけを選んだ entry は零である場合があり，
そのような entries の和から，この系の仮定を直ちに結論してはならない．

% v12 staging module: complete fixed-denominator factorization and resonances.

\section{固定した分母の partner と二つの無限共鳴}
\label{fp12:section}

二高分母の正局所模型で，一方の分母を固定したときの相手分母を調べる．
標準形の二つの高分母の役割は対称ではないため，固定した分母が
canonical な $A$ 側にある場合と $B$ 側にある場合を両方調べる必要がある．

\begin{theorem}[固定分母の有限な factorization と共鳴]
\label{fp12:partner-structure}
奇整数 $n>1$ に対し，正の整数局所模型 $(n,m,1,1)$ を持つ奇整数 $m>1$ の集合を
$\mathcal P(n)$ とする．この集合は $n\notin\{7,13\}$ なら有限である．
$n=7,13$ の場合にはそれぞれ
\[
 \{6k+1:k\geq1\}\subset\mathcal P(7),\qquad
 \{12k+1:k\geq1\}\subset\mathcal P(13)
\]
となり，これらの表示された無限集合を除く partners は有限個である．
従って $\mathcal P(n)$ が無限であることと $n\in\{7,13\}$ は同値である．
全ての有限な候補も，以下の二方向の正約数計算で上限なしに列挙できる．
\end{theorem}

\begin{proof}
正の二高分母標準形を
\[
 x=\frac{a-1}{r},\quad y=\frac{b-1}{r+1},\quad
 u=-s<0<v=z,\quad
 S=\frac{a-1}{s},\quad T=\frac{b-1}{z},
\]
\[
 A=a+x-s,\qquad B=b-y+z,\qquad AB=ab-1
\]
と書く．

\paragraph{固定分母が canonical な $A$ 側にある場合．}
$a=n$ とし，まず正約数 $r,s\mid n-1$ を選ぶ．
\[
 M=1+rs,\qquad C_A=(r+1)(s-x-1)
\]
と置けば，pair product は
\begin{equation}\label{fp12:A-factorization}
 z\{A(r+1)-MT\}=C_A
\end{equation}
と同値である．$C_A\ne0$ なら $z$ は $|C_A|$ の正約数なので，
\[
 T=\frac{A(r+1)-C_A/z}{M},\qquad b=zT+1
\]
を計算し，$T\in\mathbb Z_{>0}$，$r+1\mid zT$，$b$ が奇数という条件を検査すればよい．
この計算は有限である．

$C_A=0$ なら $s=x+1\geq2$ であり，$s\mid a-1=r(s-1)$ より $r=ts$，$t\geq1$ となる．
このとき
\[
 A=a-1=M-r-1,\qquad M=ts^2+1.
\]
式~\eqref{fp12:A-factorization} から $M\mid A(r+1)$，従って $M\mid(r+1)^2$ である．
しかし
\[
 (r+1)^2=tM+\{t(2s-1)+1\},\qquad
 0<t(2s-1)+1<M,
\]
最後の差が $t(s-1)^2>0$ なので矛盾する．
従ってこの向きに無限の枝はない．

\paragraph{固定分母が canonical な $B$ 側にある場合．}
$b=n$ とし，$r+1\mid n-1$，$r\geq1$ と正約数 $z\mid n-1$ を選ぶ．
\[
 N=1+(r+1)z,\qquad C_B=r(y-z-1)
\]
と置けば，pair product は
\begin{equation}\label{fp12:B-factorization}
 s\{NS-rB\}=C_B
\end{equation}
と同値である．$C_B\ne0$ なら $s$ は $|C_B|$ の正約数であり，
\[
 S=\frac{rB+C_B/s}{N},\qquad a=sS+1
\]
に対して $S\in\mathbb Z_{>0}$，$r\mid sS$，$a$ が奇数という条件を検査する．
この計算も有限である．

残る $C_B=0$ は $y=z+1$ と同値である．
$z\mid b-1=(r+1)(z+1)$ より $r+1=tz$，$t\geq1$ であり，$tz\geq2$ である．
\[
 b=tz(z+1)+1,\qquad N=tz^2+1,\qquad B=N+r.
\]
$S=rB/N$ の整数性は $N\mid r^2$，すなわち
\begin{equation}\label{fp12:resonance-divisibility}
 tz^2+1\mid t(2z+1)-1
\end{equation}
と同値である．$z\geq3$ なら右辺は正で左辺より小さいため不可能である．
$z=1$ なら $t\geq2$ であり，$t+1\mid3t-1$ は $t+1\mid4$ と同値なので $t=3$ となる．
$z=2$ なら $4t+1\mid5t-1$ である．$t=1$ は不可能であり，$t\geq2$ では
\[
 4t+1\leq5t-1<2(4t+1)
\]
なので $t=2$ に限られる．従って，共鳴する固定データはちょうど
\[
 (b,r,z,S,T)=(7,2,1,3,6),\qquad (13,3,2,4,6)
\]
である．

前者では $r\mid sS$ により $s=2k$，後者では $s=3k$ となる．
従って全ての共鳴模型は，それぞれ
\[
 \begin{array}{c|cccccc}
  b&a&r&s&z&S&T\\\hline
  7&6k+1&2&2k&1&3&6\\
  13&12k+1&3&3k&2&4&6
 \end{array}\qquad(k\geq1)
\]
にちょうど一致する．どちらも標準形の全整数条件を満たし，
$u=-s<0<v=z$ なので正の weights を持つ．
各固定分母の他の全枝は先の有限計算に含まれる．
これは $b=7,13$ の全ての局所模型がこの二族に入るという主張ではなく，
有限 factorization の右辺が零となる枝を全て分類したものである．
\end{proof}

\begin{proposition}[三と五の全 partners と負の低分母 entry]
\label{fp12:small-partners}
分母に上限を置かずに
\[
 \mathcal P(3)=\{7,11\},\qquad
 \mathcal P(5)=\{13,21,31,37\}
\]
である．これらの小さい分母は canonical な $A$ 側にしか現れず，
対応する標準形の全データは次の表で尽きる．
\begin{center}
\small
\begin{tabular}{ccccccc}
 $a$&$b$&$r$&$s$&$z$&$(S,T)$&$d$\\\hline
 3&7&1&1&1&$(2,6)$&$-5$\\
 3&11&1&1&2&$(2,5)$&$-4$\\
 5&13&1&1&1&$(4,12)$&$-9$\\
 5&21&1&1&2&$(4,10)$&$-7$\\
 5&31&1&2&6&$(2,5)$&$-4$\\
 5&37&1&1&4&$(4,9)$&$-6$
\end{tabular}
\end{center}
従って，小さい分母三の row から低分母へ向かう負の entry は常に $-1$ であり，
分母五の row では $-1$ または $-2$ に限られる．
\end{proposition}

\begin{proof}
定理~\ref{fp12:partner-structure} の有限計算を両方向で行う．
$a=3$ の $A$ 側では $r,s\in\{1,2\}$，$a=5$ では $r,s\in\{1,2,4\}$ を調べるだけである．
$T$ の正整数性と $r+1\mid zT$ を満たすものは，三の場合に
$(r,s,z)=(1,1,1),(1,1,2)$，五の場合に
$(1,1,1),(1,1,2),(1,1,4),(1,2,6)$ となる．
$B$ 側では，三の場合に $r+1=2$，五の場合に $r+1\in\{2,4\}$ であるが，
式~\eqref{fp12:B-factorization} の候補で正整数 $S$ と $r\mid sS$ を同時に満たすものはない．
表の全模型は正の weights を持つ．標準形の負の低分母 entry が $-s$ であることから，
最後の主張も従う．これは $b$ の探索上限を用いた列挙ではなく，
固定分母に対して完全であることが証明された正約数計算である．
\end{proof}

% v13. Complete fixed-denominator divisor calculation; no denominator cutoff.
\section{小さい分母を持つ二価 branch の局所伝播}
\label{sp13:section}

二高分母の局所模型を共有高 branch の両端で比較するとき，
分母だけでなく係数 $D$ の値を保持することが有効である．
ここでは小さい四つの分母について，可能な係数を上限なしに決定する．
結果は閉路に限らず，より高い次数の vertices に接続する二価部分にも適用できる．

\begin{proposition}[固定した小分母の係数集合]
\label{sp13:coefficient-sets}
奇数の二高分母を持つ正の整数局所模型において，
分母 $n$ の高 ray が取り得る係数 $D$ の集合を $\mathcal D(n)$ と書く．
このとき
\[
 \mathcal D(3)=\{3,6\},\qquad
 \mathcal D(5)=\{6,11,18,21\},
\]
\[
 \mathcal D(7)=F_7\cup\{7k-9:k\geq1\},\qquad
 F_7=\{-7,0,3,7,9,16,21,23,44\},
\]
\[
 \mathcal D(13)=F_{13}\cup\{13k-14:k\geq1\},
\]
ただし
\begin{align*}
 F_{13}={}&\{-15,-13,-7,-6,-3,-2,0,13,14,15,18,22,23,26,\notag\\
          &\hspace{13mm}27,28,31,41,44,52,57,58,80,83,92,161\}.
\end{align*}
特に
\[
 \mathcal D(7)\cap(-\mathcal D(7))=\{-7,0,7\},\qquad
 \mathcal D(13)\cap(-\mathcal D(13))=\{-15,-13,0,13,15\}.
\]
これらは相手の分母に上限を置かない等式である．
\end{proposition}

\begin{proof}
定理~\ref{fp12:partner-structure} の二方向の factorization を用いる．
固定した分母が canonical な $A$ 側の場合には，
$r,s\mid n-1$ と $z\mid|(r+1)(s-(n-1)/r-1)|$ を調べる．
同定理で証明した通り，最後の絶対値の中が零となる枝に解はない．
$B$ 側の場合には，$r+1\mid n-1$ と $z\mid n-1$ を選び，
$C_B=r((n-1)/(r+1)-z-1)\ne0$ なら $s\mid|C_B|$ を調べる．
各候補に正整数性，残る整除条件，奇数性を課す．

非共鳴の固定 ray 付き模型の数は $n=3,5,7,13$ の順に $2,4,13,28$ である．
分母三と五のデータは命題~\ref{fp12:small-partners} に表示されている．
各模型で
\[
 D_A=-a+B-r+S,\qquad D_B=-b+A+r+1-T
\]
を評価すると，それぞれ表示した有限集合を得る．
これは証明された有限個の約数候補を全て評価する計算であり，
相手分母や numerator の打切りを含まない．
全ての候補，棄却理由，受理された標準形と係数は付属の
\texttt{small\_denominator\_propagation\_verification\_v13.json}
で再生成できる．

残る共鳴は既知の二族だけである．
$(a,b,r,s,z,S,T)=(6k+1,7,2,2k,1,3,6)$ では固定 ray は $B$ であり，
$D_B=7k-9$ となる．
$(a,b,r,s,z,S,T)=(12k+1,13,3,3k,2,4,6)$ では $D_B=13k-14$ である．
これで全ての模型が尽くされる．
両集合とその符号反転との共通部分を取ると最後の二式を得る．
\end{proof}

\begin{theorem}[小分母の二価接続に必要な隣接分母]
\label{sp13:propagation}
高 branch の両端の core が二高分母型であり，
正の整数局所模型と対向係数の相殺 $D+D'=0$ を仮定する．
共有分母を $n$ とし，両端の core におけるもう一つの高分母を $u,v$ とする．
このとき次が成り立つ．
\begin{enumerate}
\item $n=3$ または $5$ は不可能である．
\item $n=7$ で $D\ne0$ なら，
      $\{u,v\}=\{3,13\}$ または $\{3,19\}$ であり，係数は $\pm7$ である．
      $D=0$ なら $u=v=7$ である．
\item $n=13$ で $D\ne0$ なら，
      $\{u,v\}=\{5,25\}$ で係数 $\pm15$，または
      $\{u,v\}=\{7,21\}$ で係数 $\pm13$ となる．
      $D=0$ なら $u,v\in\{13,17\}$ である．
\end{enumerate}
ここで零係数の可能性は，閉路などで既に非零性が証明されている場合を除き，
取り除いてはならない．この定理は接着の必要条件であり，
列挙された分母の組に対する global atlas の存在を主張しない．
\end{theorem}

\begin{proof}
命題~\ref{sp13:coefficient-sets} により，三と五では係数が常に正なので相殺できない．
七と十三では同命題の共通部分を用い，残る係数を持つ模型だけを取り出す．
その全データは次の通りである．表の ``側'' は固定分母の canonical な位置を表す．
\begin{center}
\small
\begin{tabular}{ccrrrrrr}
 $n$ & 側 & $a$ & $b$ & $r$ & $s$ & $z$ & $D_n$\\\hline
 7&B&3&7&1&1&1&$-7$\\
 7&A&7&13&2&1&1&7\\
 7&A&7&19&1&2&2&7\\
 7&A&7&7&2&2&1&0\\\hline
 13&B&5&13&1&1&1&$-15$\\
 13&A&13&25&2&1&1&15\\
 13&B&7&13&2&1&1&$-13$\\
 13&A&13&21&3&1&1&13\\
 13&A&13&13&3&3&2&0\\
 13&A&13&17&1&3&1&0
\end{tabular}
\end{center}
共鳴の係数列に新たな相殺可能値はない．
表の同じ $n$ に対して正負が相殺する行を組にすれば結論となる．
\end{proof}

\begin{corollary}[固定七・十三の共鳴枝は二価接続を延長できない]
\label{sp13:no-resonant-continuation}
定理~\ref{fp12:partner-structure} の共鳴模型の固定 ray $b=7$ または $13$ を，
もう一つの二高分母 core に接続し，対向 $D$ を相殺することはできない．
この結論に node 非負性も閉路の仮定も必要ない．
\end{corollary}

\begin{proof}
固定 ray の係数はそれぞれ $7k-9$，$13k-14$，$k\geq1$ である．
これらは命題~\ref{sp13:coefficient-sets} の各共通部分に属さない．
\end{proof}

% v18: the matching form covers every positive integer three-high model.
\section{三高分母の線形表示と特異な整数局所族}
\label{sh14:section}

二高分母の理論だけでは，三価以上の core を含む高分母 support を扱いきれない．
後の定理~\ref{mh18:matching} により，正の weights を持つ整数三高分母模型では，
正の numerators の辺が必ず完全マッチングとなる．ここではその標準表示を先に与える．
この型では pair equations を三元一次方程式へ帰着できるが，
その係数行列が可逆であるとは限らない．後半では，階数二の系から得られる
正の整数局所模型の無限族を与え，指定した slice の全解を決定する．

\subsection{完全マッチング型の表示}
\label{sh14:matching-setup}

分母を $(q_a,q_b,q_c,q_0)$，$q_a,q_b,q_c>1$，$q_0=1$ とし，
正の辺が $ab$ と $c0$ となるように標識を選ぶ．
正の weights を持つ整数模型では，この標識は定理~\ref{mh18:matching} により常に選べる．
正値性を課さない任意の整数 row/pair 模型にまで拡張することはできない．
正整数 $s,t,S,T,U$ と $u=s+t-1$ を用いて
\[
 (p_{0a},p_{0b},p_{0c})=(-s,-t,u),\qquad
 (p_{a0},p_{b0},p_{c0})=(-S,-T,U)
\]
と書く．pair products により
\[
 q_a=1+sS,\qquad q_b=1+tT,\qquad q_c=1+uU.
\]
$a,b$ を必要に応じて交換すれば，ある整数 $k\geq0$ に対し
$p_{ca}=-(sU-k)$，$p_{cb}=-(tU+k+1)$ とできる．
実際，それらの絶対値から $sU,tU$ をそれぞれ引いた二つの整数の和は一なので，
少なくとも一方は非正である．
さらに整数 $\alpha,\beta$ を導入すると，row sums を満たす行列は
\begin{equation}\label{sh14:matching-P}
 P=\begin{pmatrix}
 0&tS+\alpha&-(uS+\alpha+1)&-S\\
 sT+\beta&0&-(uT+\beta+1)&-T\\
 -(sU-k)&-(tU+k+1)&0&U\\
 -s&-t&u&0
 \end{pmatrix}
\end{equation}
となる．この表示では
\[
 H=tS+\alpha,\quad K=sT+\beta,\quad
 L=uS+\alpha+1,\quad N=uT+\beta+1,\quad
 M=sU-k,\quad O=tU+k+1
\]
が全て正であることも保持する．

\begin{proposition}[三つの pair equations の線形化]
\label{sh14:linear-reduction}
\[
 \mathcal A=s\alpha-t,\qquad
 \mathcal B=t\beta-s,\qquad
 \mathcal C=s+(s+t-1)k
\]
と置く．行列~\eqref{sh14:matching-P} が残る三つの pair equations を満たすことは
\begin{equation}\label{sh14:matching-linear}
 \begin{pmatrix}
  \mathcal B&\mathcal A&0\\
  -\mathcal C&0&\mathcal A+1\\
  0&\mathcal C-1&\mathcal B+1
 \end{pmatrix}
 \begin{pmatrix}S\\T\\U\end{pmatrix}
 =\begin{pmatrix}
 -\alpha\beta\\k(\alpha+1)\\-(\beta+1)(k+1)
 \end{pmatrix}
\end{equation}
と同値である．このとき正の weights は
\begin{equation}\label{sh14:matching-weights}
 w=\left(q_a,\ H+\frac1{2K},\ u+\frac1{2U},\ 1\right)
\end{equation}
で与えられる．分岐係数は $d=U+1-S-T=-D_0$ として
\begin{equation}\label{sh14:matching-D}
 D_a=sd-\beta-k-1,\qquad
 D_b=td-\alpha+k,\qquad
 D_c=\alpha+\beta+1-ud
\end{equation}
となる．
\end{proposition}

\begin{proof}
$HK=q_aq_b-1$，$LM=q_aq_c-1$，$NO=q_bq_c-1$ を展開し，
$ST,SU,TU$ の項をそれぞれ相殺すると，式~\eqref{sh14:matching-linear} の三行を得る．
Schur 補行列として見るなら，$\operatorname{diag}(q)+P$ の低分母 diagonal entry 一を pivot とすると
\[
 B=\begin{pmatrix}
 1&\alpha&-\alpha-1\\
 \beta&1&-\beta-1\\
 k&-k-1&1
 \end{pmatrix}
\]
が残る．その diagonal entries は一，row sums は零である．
正の辺に対応する四つの basin gaps は，式~\eqref{sh14:matching-weights} に対して
\[
 \frac1{2K},\qquad \frac1{2q_b},\qquad
 \frac1{2U},\qquad \frac1{2q_c}
\]
となる．他の entries は負なので，全ての basin areas は正である．
最後に $D_i=-q_i-\sum_{j\ne i}p_{ji}$ に代入すると係数公式を得る．
\end{proof}

この Schur 帰着は二次 principal minors を一に指定した非対称整数行列の問題である．
一般の principal-minor relations と cycle sums の関係は Lin--Sturmfels\cite[\S2, Proposition~4]{LinSturmfels}に従う．
ここで必要な row sums，正の weights，整数性の条件は，その一般的な関係式に加えて課す条件である．

\subsection{階数二の系から生じる全解}

\begin{theorem}[特異な三高分母 Schur 系の整数無限族]
\label{sh14:singular-family}
式~\eqref{sh14:matching-P} において
\[
 (s,t,\alpha,\beta,k)=(3,12,12,-4,1)
\]
と固定する．正の整数 $S,T,U$ を持つ全ての解は
\begin{equation}\label{sh14:singular-parameters}
 S=136+200n,\qquad T=291+425n,\qquad U=93+136n,
 \qquad n\in\mathbb Z_{\geq0}
\end{equation}
で尽くされる．その分母は
\begin{equation}\label{sh14:singular-q}
 (q_a,q_b,q_c,q_0)
 =(409+600n,\ 3493+5100n,\ 1303+1904n,\ 1)
\end{equation}
であり，三つの高分母は全て奇数である．
全ての模型は正の weights を持ち，対応する線形系の係数行列は階数二である．
分岐係数は
\begin{align*}
 D_a&=-997-1467n,&D_b&=-4007-5868n,\\
 D_c&=4671+6846n,&D_0&=333+489n.
\end{align*}
\end{theorem}

\begin{proof}
$u=14$，$(\mathcal A,\mathcal B,\mathcal C)=(24,-51,17)$ なので，線形系は
\begin{equation}\label{sh14:rank-two-system}
 \begin{pmatrix}
 -51&24&0\\
 -17&0&25\\
 0&16&-50
 \end{pmatrix}
 \begin{pmatrix}S\\T\\U\end{pmatrix}
 =\begin{pmatrix}48\\13\\6\end{pmatrix}
\end{equation}
となる．第一行は第二行の三倍と第三行の $3/2$ 倍の和であり，右辺も同じ関係を満たす．
また左上の二次 minor は $408\ne0$ であるから，係数行列の階数はちょうど二である．
後二行を解くと
\[
 S=\frac{25U-13}{17},\qquad T=\frac{25U+3}{8}.
\]
整数性は $U\equiv8\pmod{17}$ および $U\equiv5\pmod8$ と同値であり，
従って $U=93+136n$ となる．$U>0$ は $n\geq0$ と同値であり，
この範囲では $S,T$ も正なので，式~\eqref{sh14:singular-parameters} が全解である．

実際の行列は
\[
 P=\begin{pmatrix}
 0&1644+2400n&-1917-2800n&-136-200n\\
 869+1275n&0&-4071-5950n&-291-425n\\
 -278-408n&-1118-1632n&0&93+136n\\
 -3&-12&14&0
 \end{pmatrix}.
\]
表示した符号は全ての $n\geq0$ で保たれる．
前命題により pair equations と正値性が従い，row sums は構成から成立する．
分母と係数の式は直接代入で得られる．
\end{proof}

この族は，特異な線形系を捨てて Cramer の公式だけで三高分母を分類することができないことを示す．
同時に，これは固定した五つの補助整数に対する完全な slice の分類であり，
全ての特異系や全ての三高分母模型の分類ではない．
この族では二つの高係数が負であるから，全ての高係数が正となる例でもない．
また，局所模型の存在から五つの core の同時接着や大域 atlas の存在は結論しない．

\subsection{四rayの道型を排除する正値恒等式}
\label{fh25:section}

本節では，三高・四高に共通する整数局所模型の符号を完全に決定する．
全ての pair products が非零である四ray core について，
これまで必要候補として残していた道型 $P_4$ は実際には存在しない．
この排除自体は正の basin weights を仮定せず，低分母一の ray も必要としない．

\begin{theorem}[整数四ray模型における道型の非存在]
\label{fh25:no-path}
正整数 $q_1,q_2,q_3,q_4$ と整数 $p_{ij}$ が
\begin{equation}
 \sum_{j\ne i}p_{ij}=-q_i,\qquad
 p_{ij}p_{ji}=q_iq_j-1>0
 \label{fh25:raw-equations}
\end{equation}
を満たすとする．$p_{ij}>0$ である pair を辺とする無向グラフ $G^+$ は，
道 $P_4$ ではあり得ない．
\end{theorem}

\begin{proof}
道を $1-2-3-4$ と標識すると，$p_{ii}=0$ を補った行列は
\begin{equation}
 P=\begin{pmatrix}
  0&a&-d&-e\\
  A&0&b&-f\\
  -D&B&0&c\\
  -E&-F&C&0
 \end{pmatrix},
 \qquad a,A,b,B,c,C,d,D,e,E,f,F\in\mathbb Z_{>0}
 \label{fh25:path-matrix}
\end{equation}
と書ける．ここで大文字と小文字は独立な entries である．
行和と六つの pair equations はそれぞれ
\begin{gather}
 e=q_1+a-d,\quad f=q_2+A+b,\quad
 D=q_3+B+c,\quad E=q_4+C-F,
 \label{fh25:rows}\\
 aA=q_1q_2-1,\quad bB=q_2q_3-1,\quad cC=q_3q_4-1,
 \label{fh25:positive-pairs}\\
 dD=q_1q_3-1,\quad eE=q_1q_4-1,\quad fF=q_2q_4-1
 \label{fh25:negative-pairs}
\end{gather}
である．次の量を定める．
\begin{gather}
 q=q_3,\quad t=q_1,\quad s=bB+1,\quad r=cC+1,
 \quad L=B+c,\quad D=L+q,\label{fh25:abbreviations}\\
 h=ab-1\ge0,\qquad K=BC-1\ge0.
 \label{fh25:integer-inequalities}
\end{gather}
最後の二つの不等式こそ，この証明で正の整数 entries を使う箇所である．
正の pair equations から $q_2=s/q$，$q_4=r/q$ を得る．さらに
\begin{align}
 T&=(a+t)s+qh=aqf,\label{fh25:T}\\
 U&=tL+aD+1=De,\label{fh25:U}\\
 V&=ts(r+Cq)+q\{aCs+h(r+Cq)+aq\}=qTE
 \label{fh25:V}
\end{align}
が成り立つ．実際，最初の等式は
\[
 aqf=aq\left(\frac{s}{q}+\frac{ts/q-1}{a}+b\right)
      =(a+t)s+q(ab-1)
\]
であり，第二は $De=D(t+a)-(tq-1)$ である．第三については
\[
 F=\frac{q_2q_4-1}{f}=\frac{a(sr-q^2)}{qT},\qquad
 qTE=(r+Cq)T-a(sr-q^2)
\]
を展開すればよい．全ての分母 $a,q,D,T$ は正である．

残る pair equation $eE=tq_4-1$ は
\begin{equation}
 \mathcal P:=UV-DT(tr-q)=0
 \label{fh25:obstruction}
\end{equation}
を要求する．一方，恒等式
\begin{equation}
 L(r+Cq)-Dr=q(BC-1)=qK
 \label{fh25:key-identity}
\end{equation}
を用いて $t$ の次数別に整理すると
\begin{align}
 \mathcal P
 ={}&qsK\,t^2\notag\\
 &+\bigl\{aqCs(L+D)+q^2hK+aLq^2+s(r+Cq)+Dsq\bigr\}t\notag\\
 &+(aD+1)q\bigl\{aCs+h(r+Cq)+aq\bigr\}
       +Dq(as+qh).
 \label{fh25:positive-certificate}
\end{align}
右辺は非負の項の和であり，例えば $s(r+Cq)t$ は狭義正である．
従って $\mathcal P>0$ で，\eqref{fh25:obstruction} と矛盾する．
これはパラメータの上界を置かない恒等式による排除である．
\end{proof}

\begin{remark}[排除に必要な離散性]
\label{fh25:real-scope}
証明の正値部分は $q,t,a,b,c,B,C>0$ と $ab\ge1$, $BC\ge1$ の下で成立する．
従って全ての変数が整数であることよりも弱い仮定で道型は排除される．
しかし任意の正の実数 entries から $ab\ge1$ と $BC\ge1$ は従わないので，
整数性を外して無条件の実数非存在を主張してはならない．
また $q_i=1$ の有無は証明に用いていない．
三高と四高の両方に同じ証明を適用できるが，二つの低 ray が作る
$q_iq_j-1=0$ の pair は定理の仮定の外にある．
\end{remark}

\begin{theorem}[正 weights を仮定しない全符号分類]
\label{fh25:all-signs}
\eqref{fh25:raw-equations} を満たす整数 data の正符号グラフは，
完全マッチング $2K_2$ または三角形と孤立頂点 $K_3\sqcup K_1$ である．
後者では，正三角形のある向きの gain は一以上である．
三高の場合，その三角形は分母一の ray を含む．
\end{theorem}

\begin{proof}
孤立頂点 $h$ がある場合から考える．残る三頂点に対して
$c_i=-p_{hi}\in\mathbb Z_{>0}$ と置くと，$\sum_{i\ne h}c_i=q_h$ である．
異なる $i,j\ne h$ に対し，整数 $g_{ij}=q_i c_j-c_i p_{ij}$ を定める．
行和と pair equations によって
\[
 \sum_{j\ne i,h}g_{ij}
 =q_i(q_h-c_i)-c_i(-q_i-p_{ih})
 =q_iq_h+c_i p_{ih}=1
\]
となる．左辺は二つの整数の和なので，各 $i$ に対して $g_{ij}\le0$ となる
$j$ を一つ選べる．選んだ辺は $p_{ij}>0$ であり，
\[
 \frac{p_{ij}}{q_i}\ge\frac{c_j}{c_i}
\]
を満たす．三頂点上で各頂点から一本の辺を選んだ有向グラフには閉路があり，
その gain は一以上である．長さ二の閉路の gain は
$(q_iq_j-1)/(q_iq_j)<1$ なので，選ばれた閉路の長さは三でなければならない．
従って残る三頂点の全 pair は正であり，$G^+=K_3\sqcup K_1$ である．

孤立頂点がない場合，行和が負なので全ての次数は一または二である．
四頂点上の候補は $2K_2$, $P_4$, $C_4$ となる．
定理~\ref{fh25:no-path} は $P_4$ を排除する．
$C_4$ の補グラフは二本の負の辺からなる．その一辺 $ij$ について，
両端の残る二 entries は正であるから $-p_{ij}>q_i$，$-p_{ji}>q_j$ となり，
$p_{ij}p_{ji}=q_iq_j-1$ に反する．従って $G^+=2K_2$ である．

最後に，三高で分母一の ray が孤立すると，その行の三つの負の整数 entries
の和が $-1$ となる必要があり，不可能である．従って三角形は分母一の ray を含む．
\end{proof}

\begin{corollary}[三高・四高の統一された完全マッチング判定]
\label{fh25:matching}
\eqref{fh25:raw-equations} を仮定する．次の三条件は同値である．
\begin{enumerate}
 \item 正の basin weights，すなわち全ての $i\ne j$ に対して
 $w_j-(p_{ij}/q_i)w_i>0$ を満たす $w_i>0$ が存在する．
 \item $G^+$ は孤立頂点を持たない．
 \item $G^+$ は完全マッチング $2K_2$ である．
\end{enumerate}
第三条件の下では正の有理 basin weights が存在する．
特に，三高でのみ利用できた低分母一に依存する道型排除は，
本定理によって四高と共通の一つの恒等式に置き換えられる．
\end{corollary}

\begin{proof}
正 weights があれば，正の有向閉路に沿って basin 不等式を掛けることにより，
その gain は一より小さい．従って定理~\ref{fh25:all-signs} の三角形型は
第一条件と両立しない．これで第一条件から第二条件を得る．
第二条件から第三条件も同定理の全符号分類から従う．

第三条件の下では，マッチングの各辺 $ij$ ごとに一方の weight を一とし，
\[
 \frac{p_{ij}}{q_i}<\frac{w_j}{w_i}<\frac{q_j}{p_{ji}}
\]
を満たす正の有理数を選ぶ．区間は pair product によって非空である．
辺は互いに独立であり，負の entries の basin 不等式は正の weights から
自動的に従う．これで第三条件から第一条件を得る．
\end{proof}

\begin{remark}[大域問題との区別]
この結果は四ray core の局所的な符号と正値性の判定である．
完全マッチング型であることだけから，整数 clutching，node 配分，
対向 $D$ 相殺，あるいは閉じた整数 affine atlas の存在は従わない．
四高の残る研究対象は，$P_4$ の探索ではなく完全マッチング型の大域整合性である．
\end{remark}

% v25. Unified four-ray signs replace the old three-high path cases.
% Stable legacy labels are retained.
% No denominator cutoff is used.

\subsection{三高分母の完全マッチング定理と Schur 境界}
\label{sm14:section}

\begin{proposition}[正値性を仮定しない三高分母の全符号型]
\label{sm14:four-sign-types}
$q_0=1,q_1,q_2,q_3>1$ の整数行和・pair-product 方程式を満たす局所模型では，
正符号 graph は完全 matching，または低 ray を含む三角形と孤立頂点の和である．
標識付き候補は各三個である．
\end{proposition}
\begin{proof}
定理~\ref{fh25:all-signs} の直接の特殊化である．
\end{proof}

\input{three_high_matching_v25.tex}

定理~\ref{mh18:matching} により，正の weights を持つ三高分母の整数模型は
全て完全マッチング型である．以下の Schur 表示は，したがって全ての正の三高分母模型を覆う．
行列の整数 row/pair equations と正値性の役割は区別され，
正値性を除いた三角形型は実際に存在する．

以下では完全マッチングに合わせて標識を選び，\ref{sh14:matching-setup} の記号
$s,t,S,T,U\geq1$，$u=s+t-1$，$k\geq0$，$\alpha,\beta\in\mathbb Z$ を用いる．
\[
 A=s\alpha-t,\qquad B=t\beta-s,\qquad C=s+uk
\]
と略記する．式~\eqref{sh14:matching-linear} は
\begin{equation}\label{sm14:linear}
 BS+AT=-\alpha\beta,\qquad
 -CS+(A+1)U=k(\alpha+1),\qquad
 (C-1)T+(B+1)U=-(\beta+1)(k+1)
\end{equation}
である．

\begin{lemma}[基本的な整数 bound]
\label{sm14:basic-bounds}
この表示では $\alpha\geq1$，$A\geq1$，$B\leq-2$，$C\geq2$ が成り立つ．
\end{lemma}

\begin{proof}
$\alpha+1\leq0$ なら，正の entries の絶対値 $L=uS+\alpha+1$，$M=sU-k$ は
$LM\leq suSU=(q_a-1)(q_c-1)<q_aq_c-1$ を満たして矛盾する．
従って $\alpha\geq0$ であり，第二式から $A\geq0$，さらに $s\alpha\geq t\geq1$ から
$\alpha\geq1$ を得る．$\beta\leq-1$ なら $B\leq-s-t\leq-2$，
$\beta\geq0$ なら第三式と $C\geq1$ から $B+1<0$ を得る．
$A=0$ なら第一式から $\beta>0$ だが，$t=s\alpha$ より
$B=s(\alpha\beta-1)\geq0$ となって矛盾する．
$C=1$ なら $s=1,k=0$ となり，第三式は $t\beta U=-\beta-1$ となる．
整数 $\beta\geq0$，$\beta\leq-1$ のどちらでも両辺の符号が矛盾する．
\end{proof}

\begin{theorem}[Schur 補助変数の符号]
\label{sm14:beta-negative}
完全マッチング型の正の整数模型では，canonical な向き $k\geq0$ において
\[
 \beta\leq-1
\]
が成り立つ．分母の上界，奇数性，接着および node 条件は不要である．
\end{theorem}

\begin{proof}
$\beta\geq0$ と仮定し，$b=-B=s-t\beta\geq2$ と置く．
$C=s+uk\geq s\geq b$ なので
\[
 \Delta=b(A+1)-C(A+b)\leq b(1-b)<0.
\]
三つの線形式から $S,U$ を消去すると，$\Delta T=N$，ただし
\[
 N=b(A+1)(\beta+1)(k+1)
   -(b-1)\bigl(C\alpha\beta+bk(\alpha+1)\bigr)
\]
となる．$A=s\alpha-t$，$b=s-t\beta$ を代入すると，恒等式
\begin{align*}
 N-\Delta
  &=(s+\beta)A+s+b(b+\beta-1)\\
  &\quad+k\bigl\{uA+\alpha[s+\beta(s+b-1)]+b(\beta+1)\bigr\}
\end{align*}
を得る．右辺は狭義正である．従って $T=N/\Delta<1$ となり，
$T$ が正整数であることに反する．
\end{proof}

\begin{theorem}[$\beta=-1$ の境界と高係数の負性]
\label{sm14:beta-boundary}
同じ完全マッチング型で $\beta=-1$ なら，必ず $k=0$ であり，$D_0>0$ となる．
従って三つの高分母係数の少なくとも一つは負である．
この結論にも分母上限，奇数性，接着および node 条件を用いない．
\end{theorem}

\begin{proof}
ここでは $b=s+t$，$u=b-1$，$B=-b$ と書く．$k\geq1$ と仮定すると
$C=s+uk\geq b$ であるから
\[
 F=C(A+b)-b(A+1)\geq b(b-1)>0.
\]
線形式~\eqref{sm14:linear} を解くと
\[
 T=\frac{u[(\alpha+b)k-s\alpha]}{F},\qquad
 S=\frac{u(A-\alpha)k-\alpha(s-1)(A+1)}{F}
\]
となる．$S>0$ から $A>\alpha$ が必要である．ところが
\[
 F-u[(\alpha+b)k-s\alpha]
 =(s-1)(s\alpha+b)+t^2+uk(A-\alpha)>0
\]
なので $T<1$ となり，整数性に反する．従って $k=0$ である．

このとき $C=s\geq2$ であり，第三式は $uU=(s-1)T$ となる．
$G=tA-b(s-1)$ と置き，第二式と第一式も合わせると
$GT=\alpha su>0$ を得る．従って $G>0$ であり，特異な場合 $G=0$ は存在しない．
その上で解を表示すると
\[
 T=\frac{\alpha su}{G},\quad
 U=\frac{\alpha s(s-1)}{G},\quad
 S=\frac{\alpha(s-1)(A+1)}{G}.
\]
$G>0$ から
$A+1>s+s(s-1)/t>s$ であり，$sS=(A+1)U$ により $S>U$ を得る．
従って $D_0=S+T-U-1>0$ である．全係数の和は零だから，高係数の和は負となる．
\end{proof}

\begin{corollary}[$\beta=-1$ の全整数模型を与える有限な約数手順]
\label{sm14:beta-boundary-factorization}
任意の固定した $s\geq2,t\geq1$ に対し，$\beta=-1$ の完全マッチング型は，
\[
 K_0=s^2+st+t^2-s-t,\qquad g=\gcd(s-1,t)
\]
を用いて次の有限手順で全て得られる．正の約数 $G\mid K_0g$ のうち
$\alpha=(G+K_0)/(st)$ が正整数になるものを取り，
\[
 S=\frac{\alpha(s-1)(s\alpha-t+1)}G,\qquad
 T=\frac{\alpha s(s+t-1)}G,\qquad
 U=\frac{\alpha s(s-1)}G
\]
が全て整数となるものだけを残す．$k=0$ として canonical 行列を復元する．
奇分母を要求する場合は，さらに $sS,tT,(s+t-1)U$ が全て偶数である条件を課す．
\end{corollary}

\begin{proof}
前定理の $G$ は $st\alpha-K_0$ である．その表示から
\[
 tT-(s+t-1)=\frac{K_0(s+t-1)}G,
 \qquad tU-(s-1)=\frac{K_0(s-1)}G
\]
となる．左辺は整数なので
$G\mid K_0\gcd(s+t-1,s-1)=K_0g$ である．
逆に，記載した検査を通ったデータは線形式を満たし，$S,T,U>0$ である．
canonical entries は
\[
 H=tS+\alpha,\quad K=sT-1,\quad
 L=uS+\alpha+1,\quad N=uT,\quad M=sU,\quad O=tU+1
\]
で全て正であり，row sums と pair equations が成り立つ．
完全マッチングの二辺に中点 weights を取れば正値性も従う．
よって候補生成にも棄却にも人工的な分母上限を置いていない．
\end{proof}

以上により，完全マッチング型で三つの高 $D$ が全て非負となる可能性を調べる際は
$\beta\leq-2$ に制限できる．完全マッチング性は定理~\ref{mh18:matching} で証明されたが，
残る $\beta\leq-2$ における三つの高 $D$ の同時非負性は未解決である．

% v20. General three-high D sign problem: reduction to a nonsingular sector.
% No search bound, parity hypothesis, clutching or node condition is used.
\section{高係数の同時非負性に対する非特異領域への帰着}
\label{hd20:section}

三高分母の完全 matching 定理により，正の weights を持つ整数模型は
\ref{sh14:matching-setup} の canonical 表示で尽くされる．
この節では，三つの高係数が全て非負となる可能性を調べる．
一般の非特異模型についての不存在は主張せず，
その候補から全ての特異 Schur 系と二つの境界を除く．

記号を
\[
 b=-\beta,\quad u=s+t-1,\quad
 A=s\alpha-t,\quad B=-s-tb,\quad C=s+uk,
 \qquad d=U+1-S-T=-D_0
\]
とする．$s,t,S,T,U$ は正整数，$k\ge0$ である．
既知の $\beta=-1$ の境界定理を除いた範囲では $b\ge2$ であり，
$A\ge1$，$C\ge2$ が成り立つ．三つの高係数は
\begin{equation}\label{hd20:D}
 D_a=sd+b-k-1,\qquad
 D_b=td-\alpha+k,\qquad
 D_c=\alpha-b+1-ud
\end{equation}
である．その和は $d$ である．

\begin{theorem}[同時非負性が要求する非特異領域]
\label{hd20:nonsingular}
$q_0=1$，$q_a,q_b,q_c>1$ とし，整数 row sums，pair products，
正の weights を満たす局所模型を考える．
$D_a,D_b,D_c\ge0$ なら，canonical な向き $k\ge0$ において
\[
 b=-\beta\ge2,\qquad d\ge1,\qquad
 \alpha\ge k\ge b+1
\]
であり，Schur 線形系の determinant は
\begin{equation}\label{hd20:positive-det}
 \Delta=C(A-B)+B(A+1)
 \ge(-B)(-B-1)>0
\end{equation}
を満たす．従って，特異な Schur 線形系から三高係数が全て非負である
整数模型が生じることはない．
分母の上限，分母の奇数性，clutching，node 配分，対向 $D$ の相殺は仮定しない．
\end{theorem}

証明を，全係数零の排除と $k=b-1,b$ の境界の排除に分ける．
境界では $d$ の整数性と $S\ge1$ を用いる．

\subsection{全係数零の無上界排除}

\begin{proposition}[三高模型の $D$ ベクトルは零ではない]
\label{hd20:no-zero-D}
上の局所模型において $(D_a,D_b,D_c,D_0)=(0,0,0,0)$ は起こらない．
特に，$D_0=0$ なら三つの高係数には正のものと負のものがともに現れる．
\end{proposition}

\begin{proof}
全係数が零であると仮定する．既知の $\beta=-1$ の境界では $D_0>0$ なので，
$b\ge2$ である．$d=0$ を式~\eqref{hd20:D}へ代入すると
\[
 \alpha=k=b-1
\]
となる．以下
\[
 m=b^2-b+1,\qquad
 K=s^2+st+t^2-s-t,\qquad
 r=(-s,-t,u)^{\mathsf T},\quad x=(-S,-T,U)^{\mathsf T}
\]
と置く．$m$ は三以上の奇整数，$K\ge1$ である．
このとき三次 Schur 補行列は
\[
 B_0=\begin{pmatrix}
 1&b-1&-b\\
 -b&1&b-1\\
 b-1&-b&1
 \end{pmatrix},\qquad
 B_0\mathbf1=B_0^{\mathsf T}\mathbf1=0
\]
である．ここで $B_0$ は三次行列を，$B=-s-tb$ は線形系の scalar を表す．

線形系~\eqref{sh14:matching-linear} の determinant を展開すると
\[
 \Delta=mK>0
\]
となる．$N_S,N_T,N_U$ をそれぞれ $S,T,U$ の Cramer 分子とすると，
直接の多項式恒等式により
\begin{align}
 N_S&=(m-1)\bigl(ms+(B_0r)_a\bigr),\notag\\
 N_T&=(m-1)\bigl(mt+(B_0r)_b\bigr),\notag\\
 N_U&=(m-1)\bigl(mu-(B_0r)_c\bigr),\label{hd20:zero-cramer}\\
 N_U+\Delta-N_S-N_T&=m\bigl(K-(m-1)\bigr).\notag
\end{align}
$U+1-S-T=d=0$ なので最後の式から $K=m-1$ が従う．
残る三式は
\begin{equation}\label{hd20:zero-rotation}
 x=r-\frac1m B_0r
\end{equation}
と書ける．$r,x$ が整数ベクトルであることから，各成分について $m\mid B_0r$ となる．

一方，$J$ を全成分が一の三次行列とすると，直接計算により
\[
 B_0^{\mathsf T}B_0=m(3I-J),\qquad
 \mathbf1^{\mathsf T}r=-1,\qquad
 \lVert r\rVert^2=2K+1=2m-1
\]
である．従って
\[
 \lVert B_0r\rVert^2
 =m\bigl(3\lVert r\rVert^2-1\bigr)=m(6m-4).
\]
左辺は $m^2$ で割り切れるので，$m\mid4$ である．
これは $m\ge3$ が奇整数であることに反する．
最後の主張は $D_a+D_b+D_c=-D_0=0$ と，三つが同時に零とはならないことから従う．
\end{proof}

\subsection{二つの境界の排除}

\begin{lemma}[境界からの同時非負性は生じない]
\label{hd20:k-boundary}
$b\ge2$，$k=b-1$ の canonical 整数模型において
$D_a,D_b,D_c\ge0$ は起こらない．
\end{lemma}

\begin{proof}
非負性と前命題から $d\ge1$ である．$k=b-1$ なら $D_a=sd$ なので，
$D_a\le d$ と $s\ge1$ によって $s=1$ が従う．このとき $u=t$ であり，
Schur 線形系~\eqref{sh14:matching-linear} の第三式は
\[
 t(b-1)T-tbU=b(b-1)>0
\]
となる．ところが $U=S+T+d-1\ge S+T>T$ なので左辺は負となり，矛盾する．
\end{proof}

\begin{lemma}[隣接する境界の排除]
\label{hd20:k-next-boundary}
$b\ge2$，$k=b$ の canonical 整数模型においても
$D_a,D_b,D_c\ge0$ は起こらない．
\end{lemma}

\begin{proof}
前命題により $d\ge1$ である．$D_a=sd-1\le d$ なので
$(s-1)d\le1$ であり，$s=1$ または $(s,d)=(2,1)$ に限られる．
$s=1$ なら線形系の第三式は
\[
 tb(T-U)=b^2-1>0
\]
となるが，やはり $U\ge S+T>T$ なので矛盾する．

$(s,d)=(2,1)$ の場合は $D_a=1,D_b=D_c=0$ であり，
$\alpha=b+t$，$U=S+T$ となる．線形系の第一式と第三式は，それぞれ
\[
 -(2+tb)S+(2b+t)T=b(b+t),
 \qquad bT-(1+tb)S=b^2-1
\]
である．第一式の $b$ 倍から第三式の $(2b+t)$ 倍を引くと
\[
 t(1+b^2+tb)S=t+2b-b^3<t
\]
を得る．$b\ge2,t\ge1$ なので，これは正整数 $S\ge1$ に反する．
両境界の証明に determinant による除算はなく，特異系を落としていない．
\end{proof}

\begin{proof}[定理~\ref{hd20:nonsingular}の証明]
$\beta=-1$ の境界は既知の定理で除かれるので，$b\ge2$ としてよい．
非負性と命題~\ref{hd20:no-zero-D}から $d\ge1$ である．
さらに $0\le D_a\le d$ より
\[
 (s-1)d+b-1\le k\le sd+b-1.
\]
特に $k\ge b-1$ である．補題~\ref{hd20:k-boundary}と
補題~\ref{hd20:k-next-boundary}によって $k=b-1,b$ は起こらず，
整数性から $k\ge b+1$ を得る．また $D_c\ge0$ から
$\alpha\ge ud+b-1$ なので
\[
 \alpha-k\ge(u-s)d=(t-1)d\ge0.
\]
$\bar b=-B=s+tb$ と置くと
\[
 C-\bar b=(s-1)b+u(k-b)\ge0.
\]
従って determinant の恒等式
\[
 \Delta=\bar b(\bar b-1)+(C-\bar b)(A+\bar b)
\]
から式~\eqref{hd20:positive-det}が従う．$\bar b\ge3$ なので右辺は狭義正である．
\end{proof}

\begin{remark}[残る非特異領域の量化]
\label{hd20:remaining}
本定理は全ての特異 Schur 系を，高係数の同時非負性という問題から除くものである．
特異な正整数局所模型自体の不存在を意味しない．
定理~\ref{sh14:singular-family}の無限族はそのまま存在し，
少なくとも一つの高係数が負となる．
また $\Delta>0$ の領域では一般の同時非負性は未解決である．
Cramer 分子の整除性，$S,T,U$ の正整数性，canonical entries の正値性を
引き続き確認する必要がある．
\end{remark}

\begin{example}[残る領域には正整数局所模型がある]
\label{hd20:remaining-example}
$(s,t,\alpha,\beta,k)=(1,2,10,-2,5)$，$(S,T,U)=(4,5,11)$ は
\[
 q=(5,11,23,1),\qquad
 P=\begin{pmatrix}
 0&18&-19&-4\\
 3&0&-9&-5\\
 -6&-28&0&11\\
 -1&-2&2&0
 \end{pmatrix},\qquad
 D=(-1,1,3,-3)
\]
を与える．正の weights は $(5,109/6,45/22,1)$ と取れる．
この例は $k\ge b+1$ や $\Delta>0$ だけを理由に局所模型を除外できないことを示す．
三つの高係数の一つは負であり，同時非負性の反例ではない．
\end{example}

本節の Cramer 恒等式と norm 恒等式は，
\texttt{verify\_high\_d\_boundary\_v20.py} により
多変数整数・有理係数多項式として確認する．
有限範囲の代入から一般定理を推論していない．

% v21. Exact reconstruction in the remaining three-high nonnegative-D sector.
% All algebra is over Z or Q.  No numerical cutoff is a theorem hypothesis.
\section{三高係数の非負領域における単一整除条件と有限復元}
\label{hd21:section}

定理~\ref{hd20:nonsingular} は，三高係数の同時非負性を
$b\ge2$，$d\ge1$，$\alpha\ge k\ge b+1$，$\Delta>0$ の領域へ帰着した．
ここではこの残余領域を，$D$ 自身を固定した復元問題として書き直す．
得られるのは，単一の合同条件，一つの元の pair-product 方程式，
および正値性の条件からなる必要十分判定である．
さらに，固定した五つの整数に対して残る $b$ の全範囲を有限化する．
三高係数が全て非負である模型の一般不存在を主張するものではない．

\subsection{係数を指定した二変数表示}

三高 canonical 表示を，$b=-\beta$ として再掲する：
\begin{equation}\label{hd21:canonical}
\begin{split}
 &u=s+t-1,\qquad
 q=(1+sS,\,1+tT,\,1+uU,\,1),\\
 &P=\begin{pmatrix}
 0&tS+\alpha&-(uS+\alpha+1)&-S\\
 sT-b&0&-(uT-b+1)&-T\\
 -(sU-k)&-(tU+k+1)&0&U\\
 -s&-t&u&0
 \end{pmatrix}.
\end{split}
\end{equation}
ここで $s,t,S,T,U$ は正整数であり，表示された各 off-diagonal entry の
符号はそのまま保つ．従って，負号の後の量を含む次の六つの量は全て正である：
\begin{equation}\label{hd21:sign-factors}
 tS+\alpha,\quad sT-b,\quad
 uS+\alpha+1,\quad uT-b+1,\quad
 sU-k,\quad tU+k+1.
\end{equation}

以下
\begin{equation}\label{hd21:D-data}
 x=D_a\ge0,\quad y=D_b\ge0,\quad z=D_c\ge0,
 \qquad d=x+y+z\ge1
\end{equation}
と置く．式~\eqref{hd20:D} を解くと
\begin{equation}\label{hd21:ak}
 k=sd+b-1-x,\qquad
 \alpha=ud+b-1+z,\qquad U=S+T+d-1.
\end{equation}
このとき $\alpha-k=td-y\ge0$ である．
第20版で既に除外した $k=b-1,b$ の境界は再び探索しない．

次の記号は，$s,t,x,y,z$ と $b$ の依存関係を分離するためのものである：
\begin{equation}\label{hd21:fixed-data}
\begin{gathered}
 h=s-x,\qquad j=t-y,\qquad \ell=u+z,\\
 \lambda=ux+sz,\qquad \mu=uy+tz,\\
 F=b-1+\lambda,\qquad m=F^2+F+1,\\
 E=\mu h-\ell-j,\qquad J=h+\ell-\lambda j.
\end{gathered}
\end{equation}
$h,j,\ell,\lambda,\mu,E,J$ は $b$ に依存しない．
$F\ge1$，$m\ge3$ である．$j$ は正と仮定しない．

\begin{proposition}[単一合同条件を持つ復元表示]
\label{hd21:reconstruction}
式~\eqref{hd21:canonical} の正の整数模型が
式~\eqref{hd21:D-data} を満たすなら $h=s-x\ge1$ であり，
\begin{equation}\label{hd21:scales}
\begin{split}
 S&=h+\frac{FE-J}{m},\\
 T&=j+\frac{E+(F+1)J}{m},\\
 U&=\ell+\frac{(F+1)E+FJ}{m}
\end{split}
\end{equation}
である．これら三尺度が全て整数となることは，一つの合同条件
\begin{equation}\label{hd21:congruence}
 m\mid FE-J
\end{equation}
と同値である．

より正確には，任意の整数
\[
 s,t\ge1,\qquad x,y,z\ge0,\quad d=x+y+z\ge1,
 \qquad b\ge2
\]
を指定し，式~\eqref{hd21:ak}，\eqref{hd21:fixed-data}，
\eqref{hd21:scales} で全ての量を定める．
第20版の残余領域で求める模型は，次の条件を全て満たすものとちょうど一致する：
\begin{enumerate}
 \item $h\ge1$，$\alpha\ge k\ge b+1$，および合同条件~\eqref{hd21:congruence}；
 \item $S,T,U>0$，および六つの符号条件~\eqref{hd21:sign-factors}；
 \item $A=s\alpha-t$，$B=-s-tb$ として，元の第一方程式
 \begin{equation}\label{hd21:one-residue}
 BS+AT=\alpha b.
 \end{equation}
\end{enumerate}
特に，復元に際して正の weights の存在を別の存在問題として解く必要はない．
上の条件が満たされれば，その weights は明示的に構成できる．
\end{proposition}

\begin{proof}
$C=s+uk$ と置く．元の三つの高--高 pair-product 方程式は
\begin{equation}\label{hd21:original-linear}
\begin{split}
 r_1&:=BS+AT-\alpha b=0,\\
 r_2&:=-CS+(A+1)U-k(\alpha+1)=0,\\
 r_3&:=(C-1)T+(B+1)U-(b-1)(k+1)=0
\end{split}
\end{equation}
である．これは単なる必要条件ではない．例えば
\[
 (tS+\alpha)(sT-b)-(q_aq_b-1)=r_1,
\]
であり，残りの二つの積からもそれぞれ $r_2,r_3$ が得られる．
低頂点 $0$ に接する三つの pair products と，全ての row sums は，
式~\eqref{hd21:canonical} によって恒等的に成り立つ．

$U=S+T+d-1$ を代入し，式~\eqref{hd21:ak} を用いる．
\[
 R=(b-1)j+\ell,
 \qquad W=(\alpha+1)h-\ell
\]
と置けば，差を取るだけで
\begin{equation}\label{hd21:row-differences}
 r_2-r_1=(F+1)S+T-W,
 \qquad r_1-r_3=-S+FT-R
\end{equation}
となる．従って三式が零なら
\begin{equation}\label{hd21:two-relations}
 S=FT-R,\qquad (F+1)S+T=W.
\end{equation}
逆に，この二式と $r_1=0$ を満たせば，式~\eqref{hd21:row-differences} により
$r_2=r_3=0$ となる．ここで方程式を捨ててはいない．

さらに
\[
 W=\alpha(s-x)+1-t-x-z
\]
である．$F\ge1$，$S,T>0$，$\alpha>0$ であり，
$1-t-x-z\le0$ だから，式~\eqref{hd21:two-relations} の第二式は
$s-x>0$ を強制する．整数性から $h\ge1$ を得る．
これは $D_a\le s-1$ という新しい必要条件である．

二式の determinant は $m=F^2+F+1>0$ であり，解は一意である．
\[
 \alpha=F+\mu,\quad
 R=Fj+(\ell-\lambda j),\quad
 W=(F+1)h+(\mu h-\ell)
\]
を用いて解くと，式~\eqref{hd21:scales} の最初の二式を得る．
$\ell=h+j+d-1$ なので第三式は $U=S+T+d-1$ から従う．

整数性については，$S$ が整数であることが
式~\eqref{hd21:congruence} と同値である．恒等式
\begin{equation}\label{hd21:one-divisibility}
 E+(F+1)J=mE-(F+1)(FE-J)
\end{equation}
により，この一つの合同条件から $T$ も整数となる．
$U=S+T+d-1$ も整数である．逆向きは $S$ だけから明らかである．

最後に，復元した整数が命題の全ての条件を満たすとする．
既に三つの元の pair products と全 row sums を確認した．
式~\eqref{hd20:D} への代入により，$D$ は指定した
$(x,y,z,-d)$ と厳密に一致する．$S,T,U>0$ から $q_a,q_b,q_c>1$ である．
\[
 H_*=tS+\alpha>0,\qquad K_*=sT-b>0
\]
と置き，weights を
\begin{equation}\label{hd21:weights}
 w=\left(q_a,\ H_*+\frac{1}{2K_*},\ u+\frac{1}{2U},\ 1\right)
\end{equation}
で定める．正の二辺 $ab,c0$ における四つの basin gaps は，順に
\[
 \frac{1}{2K_*},\quad \frac{1}{2q_b},\quad
 \frac{1}{2q_c},\quad \frac{1}{2U}
\]
である．他の全ての off-diagonal entries は負であり，weights は正なので，
残る gaps も正である．これで十分性が示された．
必要性は canonical 完全 matching 定理，定理~\ref{hd20:nonsingular}，
および上の恒等変形から従う．
\end{proof}

\subsection{Eisenstein 型ノルムと \texorpdfstring{$b$}{b} の全範囲の有限化}

\begin{theorem}[固定した五つの整数データに対する有限復元]
\label{hd21:norm-bound}
命題~\ref{hd21:reconstruction} の正の整数模型では
$(E,J)\ne(0,0)$ であり，
\begin{equation}\label{hd21:norm}
 E^2+EJ+J^2
 =m\bigl((S-h)^2+(S-h)(T-j)+(T-j)^2\bigr).
\end{equation}
従って $N=E^2+EJ+J^2$ と置くと
\begin{equation}\label{hd21:finite-bound}
 m\mid N,\qquad
 3\le m\le N,\qquad
 2\le b\le
 1-\lambda+
 \left\lfloor\frac{\sqrt{4N-3}-1}{2}\right\rfloor.
\end{equation}
特に，$s,t,x,y,z$ を固定すれば，$b\ge2$ の全範囲を有限個の厳密演算で
完全に復元または排除できる．この上限は分母の探索上限を仮定したものではなく，
元の整数方程式の必要条件である．
\end{theorem}

\begin{proof}
$v=S-h$，$w=T-j$ と置く．式~\eqref{hd21:scales} の二式は，
逆に
\begin{equation}\label{hd21:norm-linear}
 E=(F+1)v+w,\qquad J=-v+Fw
\end{equation}
と書ける．両辺の二次式を展開すれば
\[
 ((F+1)v+w)^2+((F+1)v+w)(-v+Fw)+(-v+Fw)^2
 =(F^2+F+1)(v^2+vw+w^2)
\]
であり，式~\eqref{hd21:norm} が従う．
この正定値二次形式を Eisenstein 型ノルムと呼ぶが，
この証明には数体や整数環の外部定理を用いない．

$(E,J)=(0,0)$ の可能性を完全に調べる．
$J=0$ と $h\ge1$，$\ell\ge1$ より
$\lambda j=h+\ell>0$ であり，$\lambda>0$，$j>0$ である．
恒等式
\begin{equation}\label{hd21:zero-norm-identity}
 E+J=\ell(sy-tx)+h-j
\end{equation}
と $\ell=h+j+d-1\ge h+j$ を用いると，
\[
 |sy-tx|=\frac{|j-h|}{\ell}<1.
\]
左辺は整数なので $sy-tx=0$，従って $h=j$ である．
$s=x+h$，$t=y+h$ と書けば
$sy-tx=h(y-x)$ なので $x=y$，$s=t$ も従う．

この場合 $d=2x+z$，$s=h+x$，
$\lambda=\mu=sd-x$，$\ell=2h+d-1$ である．
$E=0$ を整理すると
\begin{equation}\label{hd21:zero-norm-classification}
 (d-1)\bigl(h(h+x)-1\bigr)=h(3-h)
\end{equation}
となる．$d\ge1$ と $h\ge1$ から $h\le3$ である．
なお $h(h+x)-1=0$ の唯一の場合 $h=1,x=0$ は右辺が二であり不可能である．
$h=3$ では $d=1$，従って $x=y=0,z=1,s=t=3$ となる．
$h=2$ では左辺が零または三以上なのに右辺は二であり，不可能である．
$h=1$ では $(d-1)x=2$ である．
$d\ge2x$ を満たす解は $x=y=1,d=3,z=1,s=t=2$ のみである．
従って零ノルム候補は正確に
\[
 (s,t,x,y,z)=(3,3,0,0,1),\qquad (2,2,1,1,1)
\]
の二組である．

第一の候補では $(S,T,U)=(3,3,6)$，$\alpha=b+5$ であり，
元の式~\eqref{hd21:one-residue} は
\[
 b(b+5)=27
\]
を要求する．その discriminant は $133$ で平方数ではないので，整数解はない．
第二の候補では $(S,T,U)=(1,1,4)$，$\alpha=b+9$ であり，
同じ元の式は $b(b+9)=14$ を要求する．$b\ge2$ では左辺が二十二以上なので不可能である．
よって正の整数模型では $(E,J)\ne(0,0)$ である．

従って $(v,w)\ne(0,0)$ であり，$v,w$ の整数性と二次形式の正定値性から
$v^2+vw+w^2\ge1$ を得る．式~\eqref{hd21:norm} は
$m\mid N$ と $m\le N$ を与える．
$m=F^2+F+1$ および $F=b-1+\lambda\ge1$ を解けば，
式~\eqref{hd21:finite-bound} の最後の不等式を得る．
\end{proof}

\subsection{復元手順，検算，および残る量化}

固定した五つの整数 $(s,t,x,y,z)$ に対する完全な手順は次の通りである．
$s,t\ge1$，$x,y,z\ge0$，$d\ge1$ の入力条件を確認し，
$h\le0$ なら命題~\ref{hd21:reconstruction} により排除する．
$E=J=0$ なら上の二候補の分類と元の第一式により排除する．
それ以外は $N>0$ であり，式~\eqref{hd21:finite-bound} の有限範囲内の $b$ を調べる．
各 $b$ で合同条件~\eqref{hd21:congruence}，尺度と entry の正値性，
$\alpha\ge k\ge b+1$，最後に元の式~\eqref{hd21:one-residue} を厳密に検査する．
残れば式~\eqref{hd21:canonical} と~\eqref{hd21:weights} から模型が復元される．
これらの判定は全て必要十分であり，元の三元一次系への復帰を含む．

付属の Python 標準ライブラリのみの検証器
\texttt{verify\_high\_d\_reconstruction\_v21.py} は，
二十四個の多項式恒等式を係数ごとに検証する．
特に，復元尺度の分母を払った元の三つの残差が全て等しいことを直接検証する．
また，既知の整数模型
$q=(5,11,23,1)$，$D=(-1,1,3,-3)$，
既知の singular rank-two 族の $n=0$，および既知の有理零 $D$ 模型について，
row sums，六つの pair products，四つの $D$，十二の basin gaps を
有理数のまま再計算する．前二者の高係数は全て非負ではなく，
第三の模型の尺度は整数ではない．従っていずれも本節の不存在候補の反例ではない．
零ノルムの二候補についても，元の第一式の非零残差を回帰検証する．

この結果が有限化するのは，指定した $(s,t,x,y,z)$ ごとの $b$ である．
$s,t,x,y,z$ 自体にはまだ一様上限を得ていない．
従って，この五つの変数を有限範囲で調べて模型が現れなくても，
一般不存在は従わない．次の課題は，単一合同条件と元の残余式の併用により，
この五変数側にも無上界の排除またはさらなる構造を与えることである．

% v22. Exact finite reconstruction after fixing only the two entries of the low row.
% No denominator cutoff, parity, clutching, or node-count hypothesis is used.
\section{低行を固定した三高非負係数模型の有限性}
\label{hd22:section}

前節の復元は五整数 $(s,t,x,y,z)$ ごとに $b$ の全範囲を有限化した．
ここでは正の entry $sT-b$ と，残していた元の pair-product 方程式を併用する．
これにより，低行 $(-s,-t,s+t-1,0)$ の二整数 $(s,t)$ だけを固定した段階で，
三高係数が同時に非負である全整数模型の候補が有限となる．
特に，$D_b=y$ と $D_c=z$ に外から上限を与える必要はない．
低行自体の二変数には依然として上限を得ておらず，一般不存在とは区別する．

\subsection{一つの正 entry が与える一様上限}

第21版の記号を用い，とくに
\[
 u=s+t-1,\quad x=D_a,\quad y=D_b,\quad z=D_c,\quad h=s-x
\]
とする．定理~\ref{hd20:nonsingular}と命題~\ref{hd21:reconstruction}から，
問題の全ての模型では
\[
 s,t\ge1,\quad 0\le x\le s-1,\quad y,z\ge0,\quad
 d=x+y+z\ge1,\quad b\ge2
\]
である．次の補助記号を追加する：
\begin{equation}\label{hd22:basic}
 \begin{gathered}
 B=b-1\ge1,\qquad \lambda=ux+sz,\qquad F=B+\lambda,\qquad
 m=F^2+F+1,\\
 a=uh-B(F+1),\qquad
 n_0=(F+1)(Bt+h)+F(u+z)+htz.
 \end{gathered}
\end{equation}
本節の $B=b-1$ は正整数であり，前節の scalar $-s-tb$ とは異なる．
後者を用いる箇所では $H=s+tb$ と書く．

\begin{lemma}[第二尺度の係数分離]
\label{hd22:T-affine}
第21版の復元尺度は，$s,t,x,z,b$ を固定すると
\begin{equation}\label{hd22:T-affine-eq}
 mT=n_0+ay
\end{equation}
を満たす．さらに $K=sT-b$ と置けば
\begin{equation}\label{hd22:K-numerator}
 mK=Q_0+say,
\end{equation}
ただし
\begin{equation}\label{hd22:Q0}
 \begin{aligned}
 Q_0={}&-BF^2+
 \bigl(B(st-2)+(2u+1)h-1\bigr)F\\
 &+B(tx-1)+h(s-utx)-1
 \end{aligned}
\end{equation}
である．
\end{lemma}

\begin{proof}
復元に用いた量を $y$ に関して分離すると
\[
 \begin{aligned}
 E&=E_0+(uh+1)y,& E_0&=thz-u-z-t,\\
 J&=J_0+\lambda y,& J_0&=h+u+z-\lambda t.
 \end{aligned}
\]
$mT=m(t-y)+E+(F+1)J$ の $y$ の係数は
\[
 -m+uh+1+(F+1)\lambda=uh-B(F+1)=a
\]
である．定数項を整理すると $n_0$ となる．
次に $mK=sn_0-bm+say$ とし，$sz=F-B-ux$ と $h=s-x$ を代入すれば
$sn_0-bm=Q_0$ を得る．これらは整数係数多項式の恒等式である．
\end{proof}

\begin{theorem}[低行ごとの $F$ 上限]
\label{hd22:F-bound}
三高係数が同時に非負である正の整数模型は
\begin{equation}\label{hd22:M-bound}
 F\le M(s,t,x):=st+(2u+1)(s-x)-2
\end{equation}
を満たす．従って
\begin{equation}\label{hd22:bz-bound}
 1\le b-1+sz\le M(s,t,x)-ux.
\end{equation}
右端が一未満なら，その $(s,t,x)$ に候補はない．
この上限には $y=D_b$ が現れない．
\end{theorem}

\begin{proof}
$F\ge M+1$ と仮定する．$B\ge1$ だから
\[
 a=uh-B(F+1)\le uh-(F+1)
 \le -st-(u+1)h<0.
\]
$y\ge0$ より，補題~\ref{hd22:T-affine}は $mK\le Q_0$ を与える．
一方，模型の符号条件は $K=sT-b>0$ である．
この二つを矛盾させるため，
\[
 \begin{aligned}
 Q_1={}&-BF^2+
 \bigl(B(st-1)+(2u+1)h-1\bigr)F
 -B(B+1)+sh-1,\\
 Q_2={}&-F^2+MF+sh-3
 \end{aligned}
\]
と置く．直接の恒等式
\begin{align}
 Q_1-Q_0&=B\bigl((s-1)x+sz\bigr)+uhtx\ge0,
 \label{hd22:bound-gap1}\\
 Q_2-Q_1&=(B-1)\bigl(F(F-st+1)+B+2\bigr)\ge0
 \label{hd22:bound-gap2}
\end{align}
を用いる．第二式の不等号は，$F\ge M+1>st-1$ と $B\ge1$ による．
従って
\[
 \begin{aligned}
 0<mK&\le Q_0\le Q_1\le Q_2\\
 &\le -F+sh-3\\
 &\le -st-(s+2t-1)h-2<0,
 \end{aligned}
\]
となる．最後の行は $F\ge M+1$ と $2u+1-s=s+2t-1$ を用いた．
矛盾と $F$ の整数性により $F\le M$ であり，式~\eqref{hd22:bz-bound} は
$F=b-1+ux+sz$ から従う．
\end{proof}

この証明は $T$ の整数性を使わず，$K>0$ と $y\ge0$ を直接使っている．
ただし $h\ge1$，$b\ge2$ などの入力範囲は，前節までの整数模型の帰着から
受け取っている．従って正の実数模型一般に対する不存在は導かれない．

\subsection{残った係数に関する二次方程式}

固定した $s,t,x,z,b$ に対して，以下は全て整数である：
\begin{equation}\label{hd22:quadratic-data}
 \begin{gathered}
 \alpha_0=F+tz,\quad H=s+tb,\quad A_0=s\alpha_0-t,\\
 E_0=thz-u-z-t,\quad J_0=h+u+z-\lambda t,\\
 n_{S,0}=mh+FE_0-J_0,\quad c=Fuh+B,\\
 L=(2s-1)F+tb+2stz+s-2t,\\
 C=F\bigl(b(t-1)+s\bigr)+b^2t+bs-b,\\
 C_0=-Hn_{S,0}+A_0n_0-b\alpha_0m.
 \end{gathered}
\end{equation}
ここで $L>0$，$C>0$ である．実際，
$L=(2s-1)F+t(b-2)+2stz+s$，
$C=F\bigl(b(t-1)+s\bigr)+b(bt+s-1)$ と書ける．

\begin{proposition}[恒等零とならない残余多項式]
\label{hd22:residual-quadratic}
復元尺度を元の第一 pair-product 方程式へ代入した残余を
\[
 r=(-s-tb)S+(s\alpha-t)T-\alpha b
\]
とする．このとき $\alpha=\alpha_0+uy$ であり，
\begin{equation}\label{hd22:quadratic}
 mr=\mathcal R(y),\qquad
 \mathcal R(Y)=su\,aY^2+(aL-C)Y+C_0.
\end{equation}
$a\ne0$ なら二次式であり，$a=0$ なら傾き $-C<0$ の一次式である．
従って，どの固定五組 $(s,t,x,z,b)$ に対しても $\mathcal R$ は恒等零ではなく，
その非負整数根は高々二個である．
$a\le0$ の場合は $Y\ge0$ で狭義単調減少なので，高々一個である．
\end{proposition}

\begin{proof}
前補題の計算と同様に
\[
 mS=n_{S,0}+cy,\qquad mT=n_0+ay,\qquad
 c=Bm+Fa
\]
である．従って分母 $m$ を払った元の残余は
\[
 -H(n_{S,0}+cy)+(A_0+suy)(n_0+ay)
 -b(\alpha_0+uy)m.
\]
二次係数は $sua$，定数項は $C_0$ である．一次係数についての恒等式
\begin{equation}\label{hd22:linear-coefficient}
 -Hc+A_0a+sun_0-bum=aL-C
\end{equation}
が式~\eqref{hd22:quadratic} を与える．
これは $F=B+ux+sz$，$u=s+t-1$，$h=s-x$ による直接展開で確かめられる．
特に $a=0$ の場合にも一次係数が負なので，一次式まで消えて恒等零となる
例外は存在しない．$a<0$ なら二次係数も一次係数も負であり，
$a=0$ の場合と合わせて最後の主張も従う．
\end{proof}

\begin{theorem}[二整数を固定した有限・必要十分復元]
\label{hd22:fixed-row-finite}
低行の正整数 $(s,t)$ を固定する．三高係数が全て非負である canonical 整数模型は，
次の有限手順で漏れなく復元される：
\begin{enumerate}
 \item $x=0,\ldots,s-1$ を取り，$M=st+(2u+1)(s-x)-2$ とする．
 \item $z\ge0$，$b\ge2$ のうち $b-1+sz\le M-ux$ を満たすものを全て取る．
 \item 式~\eqref{hd22:quadratic} の非負整数根 $y$ を全て取る．
 各五組に高々二個であり，$a\le0$ なら高々一個である．
 \item $d=x+y+z\ge1$，$\alpha\ge k\ge b+1$ を検査する．
 第21版の単一合同条件 $m\mid FE-J$ を検査し，
 式~\eqref{hd21:scales} で $S,T,U$ を復元する．
 \item $S,T,U>0$ と式~\eqref{hd21:sign-factors} の六つの正値性を検査する．
 通過した候補から式~\eqref{hd21:canonical} と~\eqref{hd21:weights} を用いて
 $q,P,w$ を構成する．
\end{enumerate}
最後に残るものは正確に所望の模型であり，別の存在判定は必要ない．
特に，固定 $(s,t)$ に対して $x,y,z,b,S,T,U$ の全範囲が有限化される．
\end{theorem}

\begin{proof}
全ての模型が最初の三段階に現れることは，
$h\ge1$，定理~\ref{hd22:F-bound}，命題~\ref{hd22:residual-quadratic} による．
第二段階の集合は明示的に有限である．例えば $M-ux\ge1$ の場合，
\[
 0\le z\le\left\lfloor\frac{M-ux-1}{s}\right\rfloor,
 \qquad 2\le b\le M-ux-sz+1
\]
と列挙すればよい．
第三段階は恒等零でない一次または二次式の整数根を求めるだけである．
一次式では割り切れを調べ，二次式では判別式が非負の平方数であることと，
二つの二次公式の分子が $2sua$ で割り切れることを調べる．
浮動小数点近似も外からの $y$ 上限も不要である．

必要性は既に得られている．十分性では，第21版の復元が
三つの元の pair-product 残余を等しくすることに注意する．
ここでは $\mathcal R(y)=0$ も検査しているので，第一残余そのものが零となり，
三つ全てが零である．単一合同条件は三尺度の整数性を与え，
正値性検査は canonical な符号を保証する．
命題~\ref{hd21:reconstruction} の十分性から，
全 row sums，六つの pair products，指定した $D=(x,y,z,-d)$，
および正の weights が従う．
\end{proof}

\begin{remark}[有限化の範囲と計算の意味]
\label{hd22:scope}
本定理に分母上限，奇偶条件，clutching，node 配分，対向する $D$ の相殺条件はない．
固定しているのは低行の二整数 $(s,t)$ だけである．
したがって，有限範囲の $(s,t)$ について本手順を実行した結果は，
その低行の範囲で全ての高係数・全ての分母を尽くした判定となる．
ただし $(s,t)$ が範囲外の模型を除外せず，
三高非負係数模型の一般不存在も，その存在も，本節では結論しない．
\end{remark}

\subsection{対称な低行・高係数に対する無上界の排除}

\begin{proposition}[対称 sector の不存在]
\label{hd22:symmetric-sector}
三高係数が全て非負である正の整数模型において，
$s=t$ と $x=y$ は同時には成立しない．
この主張に $s,x,z,b$ の上限はない．
\end{proposition}

\begin{proof}
$s=t$，$x=y$ と仮定する．$h=j$ であり，
式~\eqref{hd21:zero-norm-identity} は $E+J=0$ を与える．
単一合同条件は $m\mid(F+1)E$ となる．
$m-F(F+1)=1$ だから $\gcd(m,F+1)=1$ であり，$n=E/m$ は整数である．
復元式は
\[
 S=h+(F+1)n,\qquad T=h-Fn
\]
となる．一方 $d=2x+z\ge1$ であり，
\[
 F=b-1+sd-x=b-1+h+s(d-1)\ge h+1.
\]
$n\ge1$ なら $T<0$，$n\le-1$ なら $S<0$ なので，
$S,T>0$ は $n=0$ を強制する．従って $E=J=0$ であるが，
これは定理~\ref{hd21:norm-bound}の零ノルム候補の排除に反する．
\end{proof}

検証器 \nolinkurl{verify_high_d_low_row_finiteness_v22.py} は，
上限証明の二つの差，残余二次式，退化時にも非零となる一次係数，
および復元後の三つの元の残余の一致を，整数係数多項式として係数ごとに確認する．
付属の計算結果では，指定した $(s,t)$ の範囲だけを有限に区切り，
各低行に対する $x,y,z,b$ は本定理の全範囲を処理している．
具体的には $1\le s,t\le12$ の $144$ 低行について，
$(s,t,x,z,b)$ の候補 $2{,}004{,}015$ 組を処理した．
そのうち $219$ 組は $a=0$ の一次式であり，恒等零の例外として捨てていない．
非負整数残余根は計六個で，全て $x=y=z=0$，従って $d=0$ であり対象外だった．
この低行の範囲では三高非負係数模型は存在しない．
この結論は，定理の有限化を用いた当該低行範囲についての厳密計算結果である．
$s$ または $t$ が十二を超える場合の一般結論へは外挿しない．

% v23: unbounded asymmetric lattice obstructions. All proofs are exact.
\section{三高非負係数模型の非対称格子障害}
\label{hd23:section}

低行を固定した有限復元の後に残った二整数 $(s,t)$ 自体には，依然として一様上限を得ていない．
本節は探索範囲を拡大せず，尺度の整数性を一つの格子変数として取り出す．
これにより，完全対称の場合に加えて，以下の四領域を無上界で除外する．
特に，非対称量 $sy-tx$ の符号を入れ替えたときに証明は同一ではなく，
正側では元の pair-product 残余そのものが必要になる．

\subsection{記号と格子表示}

前節までと同じ canonical 表示を用いる．必要な入力と記号を全て再掲する：
\begin{equation}\label{hd23:data}
 \begin{gathered}
 s,t\ge1,\quad u=s+t-1,\quad b\ge2,\quad B=b-1,\\
 x=D_a\ge0,\quad y=D_b\ge0,\quad z=D_c\ge0,
 \quad d=x+y+z\ge1,\\
 h=s-x\ge1,\quad j=t-y,\quad \ell=u+z,\\
 \lambda=ux+sz,\quad \mu=uy+tz,\quad F=B+\lambda,
 \quad m=F^2+F+1,\\
 E=\mu h-\ell-j,\quad J=h+\ell-\lambda j,\\
 \alpha=ud+B+z,\qquad k=sd+B-x.
 \end{gathered}
\end{equation}
これらは整数であり，$j$ の非負性は仮定しない．
対象とする模型の三尺度 $S,T,U$ は正整数であり，
\begin{equation}\label{hd23:reconstruction}
 \begin{gathered}
 S=h+\frac{FE-J}{m},\qquad
 T=j+\frac{E+(F+1)J}{m},\qquad U=S+T+d-1,\\
 E=(F+1)(S-h)+(T-j),\qquad
 J=-(S-h)+F(T-j)
 \end{gathered}
\end{equation}
を満たす．さらに $sT-b>0$ を含む六つの符号条件と
\begin{equation}\label{hd23:original-residual}
 \mathfrak r:=-(s+tb)S+(s\alpha-t)T-\alpha b=0
\end{equation}
が必要である．後者は元の第一 pair-product 方程式であり，
三つの残余の差が零であるという条件だけからは従わない．

\begin{lemma}[非対称量と一つの整数格子変数]
\label{hd23:lattice}
\[
 \delta=sy-tx,\qquad n=\ell-U\in\mathbb Z
\]
と置く．$n$ は本節の格子変数であり，node の個数ではない．このとき
\begin{equation}\label{hd23:lattice-form}
 \boxed{\quad
 S=j-(F+1)n-\ell\delta,\qquad
 T=h+Fn+\ell\delta,\qquad U=\ell-n.
 \quad}
\end{equation}
また
\begin{equation}\label{hd23:F-delta}
 F+\delta=B+h+s(d-1),\qquad
 E+J=\ell\delta+h-j,\qquad \ell=h+j+d-1.
\end{equation}
\end{lemma}

\begin{proof}
$\mu h-\lambda j=\ell(sy-tx)$ を展開すれば第二の恒等式を得る．
式~\eqref{hd23:reconstruction} から
\[
 E+J=F\bigl((S-h)+(T-j)\bigr)+(T-j)
       =-Fn+(T-j)
\]
である．これを $E+J=\ell\delta+h-j$ と比較して $T$ を求め，
$S+T=h+j-n$ を使えば $S$ も得られる．
$F+\delta=B+sd-x$ なので最初の恒等式も従う．
尺度の整数性，またはそれと同値な単一合同 $m\mid FE-J$ を
使って $n\in\mathbb Z$ としていることが重要である．
\end{proof}

\begin{theorem}[四つの無上界排除領域]
\label{hd23:main}
三高係数が同時に非負である正の canonical 整数模型では，
次のいずれも起こらない：
\begin{enumerate}
 \item $\delta=0$ かつ $s\ge t$；
 \item $\delta=-1$；
 \item $\delta=1$ かつ $x\ge1$；
 \item $a:=uh-B(F+1)=0$ かつ $s\le2$．
\end{enumerate}
各主張に，分母，$D$，$s,t,b$ の探索上限を課していない．
ただし，例えば $s=t$ だけを仮定した全領域や，$\delta=1,x=0$ は
この定理の排除対象に含めていない．
\end{theorem}

\subsection{比例する高係数と低行の一方向の順序}

\begin{proof}[定理~\ref{hd23:main} の第一項]
$sy=tx$ から $j=th/s>0$ である．$s\ge t$ なら $h\ge j\ge1$ となる．
式~\eqref{hd23:F-delta} は $F\ge h+1$ を与える．
式~\eqref{hd23:lattice-form} で $n\le-1$ なら $T\le h-F<0$，
$n\ge1$ なら $S\le j-F-1<0$ である．従って
\[
 n=0,\qquad S=j,\quad T=h,\quad U=\ell.
\]
$\delta=0$ の下では $\lambda=h+s(d-1)$，$\mu h=\lambda j$ である．
式~\eqref{hd23:reconstruction} の $E$ を計算すると，必要条件は
\begin{equation}\label{hd23:delta-zero-equation}
 0=(d-1)(sh-1)+h(h-3)+(b+1)(h-j)
\end{equation}
となる．

$h\ge4$ なら右辺は正である．$h=3$ なら全項が非負なので，
$d=1$ かつ $j=h=3$ が必要である．$h=2$ では，$j=1$ なら最後の項が
三以上となって矛盾する．$j=2$ なら $(d-1)(2s-1)=2$ が必要だが，
$s\ge2$ より左辺は零または三以上である．
$h=1$ では $j=1$ であり，$(d-1)(s-1)=2$ が必要である．

残った場合は全て $h=j$ なので，$sy=tx$ と $h,j>0$ から $s=t$，$x=y$ も従う．
$h=3,d=1$ は $(s,t,x,y,z)=(3,3,0,0,1)$ だけである．
$h=1$ では $x=y=s-1$，$d\ge2(s-1)$ である．
$(d-1)(s-1)=2$ の候補のうちこの不等式を満たすものは
$(s,t,x,y,z)=(2,2,1,1,1)$ だけである．

第一候補では $S=T=3$，$\alpha=b+5$ であり，式~\eqref{hd23:original-residual} は
$27-b(b+5)=0$ を要求する．$3(3+5)<27<4(4+5)$ と単調性により，整数解はない．
第二候補では $S=T=1$，$\alpha=b+9$ であり，残余は
$14-b(b+9)<0$ となる．これで両候補も排除された．
\end{proof}

\subsection{非対称量が負の一の場合}

\begin{proof}[定理~\ref{hd23:main} の第二項]
$\delta=-1$ なら $tx=sy+1$ なので $x\ge1$，$s\ge2$ である．
また $x<s$ より $y<t$，従って $j>0$ である．
$\ell=h+j+d-1\ge h+j>h$ なので，式~\eqref{hd23:lattice-form} の
$T=h+Fn-\ell>0$ は $n\ge1$ を強制する．
そのとき $S>0$ は $F+1<j+\ell$ を要求する．
もし $x\ge2$ なら
\[
 F+1-j-\ell
 =B+u(x-1)+(s-1)z+1-t+y
 \ge B+s+y>0,
\]
となって矛盾する．従って $x=1$，$t=sy+1$ である．

$g=B+(s-1)z\ge1$ と置くと $F=\ell+g$ である．
$n\ge2$ なら
\[
 S\le j-2(F+1)+\ell=j-\ell-2g-2<0
\]
となるため，$n=1$ であり，
\[
 S=j-g-1,\qquad T=s-1+g,\qquad U=\ell-1.
\]
この三尺度を $E=(F+1)(S-h)+(T-j)$ に代入する．左辺から右辺を引くと
\begin{equation}\label{hd23:delta-minus-positive}
 \begin{aligned}
 Q_-={}&\bigl((s-1)(s-2)+(b-2)+s(s-1)z\bigr)y\\
 &+s^2+(b-2)(b+2s-1)\\
 &+\bigl((2s-1)b+2s^2-3s-1\bigr)z+s(s-1)z^2.
 \end{aligned}
\end{equation}
$s\ge2,b\ge2,y,z\ge0$ である．各項は非負であり，$s^2>0$ なので
$Q_->0$ である．これは必要な差方程式 $Q_-=0$ に反する．
この領域は元の残余式を課す前に除外される．
\end{proof}

\subsection{非対称量が正の一の場合}

\begin{proof}[定理~\ref{hd23:main} の第三項]
$\delta=1$，$x\ge1$ と仮定する．$\ell\ge u\ge t$，$j\le t$ なので，
$n\ge0$ なら $S=j-(F+1)n-\ell\le0$ となる．
従って $p=-n\ge1$ として
\[
 S=j+(F+1)p-\ell,\qquad T=h-Fp+\ell
\]
と書ける．$T>0$ は $F<h+\ell$ を要求する．もし $x\ge2$ なら
\[
 F-h-\ell=B+(u+1)x-s-u+(s-1)z\ge B+t+1>0
\]
で矛盾する．ゆえに $x=1$，$t=sy-1$，特に $s\ge2,y\ge1$ である．

再び $g=B+(s-1)z$ と置くと $F=\ell+g$ であり，$F>h=s-1$ である．
$p\ge2$ なら $T\le h-2F+\ell<0$ なので $p=1$ である．
従って $T=s-1-g>0$ は
\[
 z=0,\qquad \nu:=s-b\ge1
\]
を強制する．尺度と補助量は
\begin{equation}\label{hd23:delta-plus-data}
 \begin{gathered}
 s=b+\nu,\quad t=sy-1,\quad
 S=(s-1)y+b-1,\quad T=\nu,\quad U=s+t,\\
 d=1+y,\qquad \alpha=u(1+y)+b-1
 \end{gathered}
\end{equation}
となる．ここまで元の残余式はまだ使っていない．

$E=(F+1)(S-h)+(T-j)$ の左辺と右辺の差は
\begin{equation}\label{hd23:delta-plus-Q}
 Q_+=\nu^2+2b\nu-4\nu-b+2
       -\bigl(b^2-\nu(\nu-1)\bigr)y
\end{equation}
であり，模型では $Q_+=0$ が必要である．一方，元の残余
$\mathfrak r$ には次の整数係数恒等式がある：
\begin{equation}\label{hd23:delta-plus-residual}
 y\mathfrak r=\bigl(sy(y+1)+1\bigr)Q_++N_+,
\end{equation}
ただし
\begin{equation}\label{hd23:N-positive}
 \begin{aligned}
 N_+={}&b\nu(y^3+y^2-2)+b(y+1)+\nu^2(y^3-1)\\
       &+\nu(\nu-1)y^2+\nu y+4\nu-2.
 \end{aligned}
\end{equation}
$y\ge1$，$\nu\ge1$，$b\ge2$ により $N_+>0$ である．
従って $Q_+=0$ と $\mathfrak r=0$ は両立しない．
式~\eqref{hd23:delta-plus-residual} は元の残余を保持することで得た障害であり，
差方程式だけの候補を模型と取り違えていない．
\end{proof}

\subsection{二次係数の退化面と整数尺度の平方障害}

\begin{lemma}[退化面における有界整数]
\label{hd23:degenerate-integer}
$a=uh-B(F+1)=0$ とする．このとき
\begin{equation}\label{hd23:degenerate-relations}
 s\ell=bF,\qquad u(s-bx)=B(b+sz),\qquad s-bx>0.
\end{equation}
さらに $K_0=h(s-1)+1$ と置くと，整数
\[
 R_0=b^2-sT
\]
は
\begin{equation}\label{hd23:degenerate-bound}
 1\le R_0\le K_0,\qquad
 R_0m=(bK_0-sh)F+b^2-sh
\end{equation}
を満たす．この補題は元の残余 $\mathfrak r=0$ を使わない．
\end{lemma}

\begin{proof}
一般の恒等式
\[
 a=s\ell-bF=u(s-bx)-B(b+sz)
\]
が式~\eqref{hd23:degenerate-relations} を与える．
式~\eqref{hd23:reconstruction} を展開し，$B(F+1)=uh$ を使うと
\[
 \begin{aligned}
 mT&=B(F+1)j+(F+1)h+F\ell+\mu h\\
   &=\ell(ht+F)+(F+1)h
 \end{aligned}
\]
となる．$s\ell=bF$ と $ht=B(F+1)-sh+h$ により
\[
 m(b^2-sT)=(bK_0-sh)F+b^2-sh
\]
を得る．$A_0=bK_0-sh$ と書くと
\[
 A_0-B=h(Bs-b)+1\ge0.
\]
$s=1$ では $h=1$ により等号であり，$s\ge2$ では $Bs-b\ge0$ である．
$BF=uh-B\ge sh-B$ も成り立つので，
\[
 A_0F+b^2-sh\ge BF+b^2-sh\ge b^2-B>0.
\]
従って $R_0>0$ である．

$w=F-B\ge0$ とすれば，直接の展開により
\begin{equation}\label{hd23:degenerate-upper-gap}
 \begin{aligned}
 &(K_0+1)m-\bigl(A_0F+b^2-sh\bigr)\\
 &\quad=(K_0+1)w^2+
 \bigl((b-1)K_0+2b-1+sh\bigr)w+K_0+1+b(sh-1)>0.
 \end{aligned}
\end{equation}
$R_0$ の整数性から $1\le R_0\le K_0$ が従う．
\end{proof}

\begin{proof}[定理~\ref{hd23:main} の第四項]
$s-bx>0$ と $b\ge2$ により，$s\le2$ なら $x=0$ である．従って $h=s$，
$F=b-1+sz$ となる．

$s=1$ では $K_0=1$ なので $R_0=1$ であり，式~\eqref{hd23:degenerate-bound} は
\[
 F^2+(2-b)F+2-b^2=0
\]
となる．その判別式 $D=5b^2-4b-4$ は平方数でなければならない．
$D=q^2$ と仮定すると
\[
 (5b-2)^2-5q^2=24.
\]
法三では二つの平方の和が零となるため，$5b-2$ と $q$ はともに三で割り切れる．
すると左辺は九で割り切れるが右辺は割り切れず，矛盾する．

$s=2$ では $K_0=3$ なので $R_0\in\{1,2,3\}$ である．
式~\eqref{hd23:degenerate-bound} は
\[
 R_0F^2+(R_0-3b+4)F+R_0-b^2+4=0
\]
となる．各 $R_0$ における判別式は，順に
\[
 D_1=13b^2-30b+5,\quad
 D_2=17b^2-36b-12,\quad
 D_3=21b^2-42b-35.
\]
$D_1$ が平方数なら，法五の $D_1\equiv3b^2$ により $5\mid b$ と
平方根の五整除性が必要だが，その場合 $D_1\equiv5\pmod{25}$ となって不可能である．
$D_2$ が平方数なら，法三の $D_2\equiv2b^2$ により $3\mid b$ と平方根の
三整除性が必要だが，$D_2\equiv6\pmod9$ となる．
最後に $D_3=7\bigl(3(b-1)^2-8\bigr)$ である．平方数なら括弧内も七で
割り切れ，$(b-1)^2\equiv5\pmod7$ を要求するが，五は法七の平方剰余ではない．
従って全ての場合が排除された．
\end{proof}

\subsection{保持する必要十分条件と未解決範囲}

新しい障害は第22版の有限復元への必要条件として追加できる．
残る候補では引き続き $\alpha\ge k\ge b+1$，$m\mid FE-J$，正整数尺度，
六つの符号条件，および元の残余 $\mathfrak r=0$ を全て確認する．
通過後の $q,P$ と正の weights の構成は第21版の必要十分復元をそのまま用いる．
本節の格子表示だけを十分条件としない．

$\delta=1,x=0$ なら $sy=1$ より $s=y=1$ である．
これは第三項とは別の残余であり，本節で全域排除したとはしない．
$\delta=0,s<t$，$|\delta|\ge2$，および一般の $a=0,s\ge3$ も残る．
したがって三高同時非負模型の一般不存在は依然として未解決である．

標準 Python の
\nolinkurl{verify_high_d_asymmetry_v23.py}
は本節の恒等式を多項式の係数として再計算し，
四つの平方障害を法 $9,25,9,49$ の完全な剰余類集合で独立に確認する．
これらの剰余類検査は探索上限ではなく，全整数を覆う合同式の検証である．
既存の $144$ 低行の完全復元を繰り返すことは，新定理の証明にも検算にも必要ない．

% v24. Unbounded exclusion of the positive axis and exceptional degenerate branch.
\section{正側残余軸の排除と退化面の一整数有限性}
\label{ax24:section}

三高係数が同時に非負である canonical 整数模型の記号
\eqref{hd23:data} を用いる．本節では，前節で残った
$\delta=1,x=0$ と $a=0,x=0,R_0=1$ を全域で排除する．
その結果，$\delta=1$ の排除が完成し，退化面 $a=0$ の全体は
低行の一整数 $s$ だけを固定して有限に復元できる．
一般の $\delta$ や，退化面全体の不存在を主張するものではない．

\subsection{正側残余軸の整除性障害}

\begin{theorem}[非対称量一の完全排除]
\label{ax24:positive-axis}
三高係数が同時に非負である正の canonical 整数模型では
$\delta=sy-tx=1$ は起こらない．
この主張に分母，低行，係数，尺度の探索上限はない．
\end{theorem}

\begin{proof}
$x\ge1$ は定理~\ref{hd23:main} により排除済みなので，$x=0$ とする．
$sy=1$ より $s=y=1$ である．$k=b+z\ge b+1$ から $z\ge1$ を得る．
$B=b-1\ge1$，$F=B+z$，$m=F^2+F+1$ と置く．
格子表示~\eqref{hd23:lattice-form} で $S>0$ は $p=-n\ge1$ を強制し，
\begin{equation}\label{ax24:axis-data}
 \begin{gathered}
 S=(F+1)p-z-1,\qquad T=t+z+1-Fp,\qquad U=t+z+p,\\
 \alpha=t(z+1)+F,\\
 Q:=(z-1)(t-1)-mp+F(z+2)=0
 \end{gathered}
\end{equation}
となる．最後の式は元の二つの尺度差方程式の一方である．
$z=1$ なら $mp=3F$ となるが，$F=b\ge2$ と
$m-3F=(F-1)^2>0$ より不可能である．従って $z\ge2$ とする．

次の整数を定義する：
\begin{equation}\label{ax24:axis-auxiliary}
 \begin{gathered}
 C=F+2=B+z+2,\qquad R=t-B(p+z+1),\\
 N=b^2+bz+z^2,\qquad K=(z+1)F+1.
 \end{gathered}
\end{equation}
元の第一 pair-product 残余を $\mathfrak r$ とすると，直接の展開により
\begin{equation}\label{ax24:two-identities}
 \mathfrak r=(t+1)Q+tR-BC,\qquad
 BQ=mR+BCK-Nt
\end{equation}
が成り立つ．模型では $Q=\mathfrak r=0$ なので $tR=BC>0$ であり，$R>0$ である．
$t=B(p+z+1)+R$ を代入すると
\[
 R^2=B\bigl(C-R(p+z+1)\bigr).
\]
従って正整数 $k_*=R^2/B$ が存在し，
\begin{equation}\label{ax24:small-divisor}
 C=k_*+R(p+z+1)
\end{equation}
となる．もし $R\ge B+1$ なら
$C>R(z+2)\ge(B+1)(z+2)>B+z+2=C$ となるので，$R\le B$ である．
ゆえに $k_*\le R$ であり，$p\ge1,z\ge2$ と合わせて
\begin{equation}\label{ax24:large-C}
 C\ge k_*+(z+2)R\ge5k_*
\end{equation}
を得る．

一方，式~\eqref{ax24:two-identities} の第二式に $R$ を掛け，
$tR=BC$ と $R^2=Bk_*$ を使うと
\[
 mk_*=C(N-RK).
\]
$m=C(F-1)+3$ なので $C\mid3k_*$ である．
これは $0<3k_*<C$ と矛盾する．
\end{proof}

この証明は元の残余 $\mathfrak r=0$ を本質的に用いる．
尺度差方程式と整数尺度だけから $tR=BC$ を結論してはいけない．

\subsection{退化面に残った例外枝の有界な約数}

\begin{theorem}[退化面の例外枝の排除]
\label{ax24:degenerate-R1}
三高係数が同時に非負である正の canonical 整数模型において，
\[
 a=uh-(b-1)(F+1)=0,\qquad x=0,\qquad R_0=b^2-sT=1
\]
は同時には成立しない．$s,b$ に上限を課さない．
\end{theorem}

\begin{proof}
$s=1,2$ は定理~\ref{hd23:main} で既に除かれているので $s\ge3$ とする．
$B=b-1$，$h=s$ と置く．$F=B+sz$，$us=B(F+1)$，$sT=b^2-1$ より
\[
 s\mid B(B+1),\qquad s\mid B(B+2)
\]
である．差を取れば $s\mid B$ となる．従って
\[
 B=s\kappa,\qquad F=sL,\qquad
 \kappa\ge1,\quad L=\kappa+z\ge\kappa
\]
と書ける．以下
\begin{equation}\label{ax24:degenerate-data}
 \begin{gathered}
 K_0=s^2-s+1,\qquad
 w=L+1-\kappa K_0,\qquad D=s(w+\kappa)+1,\\
 u=\kappa(sL+1),\qquad
 t=\kappa sL+\kappa-s+1,\qquad
 T=\kappa(s\kappa+2)
 \end{gathered}
\end{equation}
と置く．退化面の恒等式~\eqref{hd23:degenerate-bound} は
\begin{equation}\label{ax24:degenerate-Q}
 Q_s:=sLw-s\kappa^2+s-2\kappa=0
\end{equation}
を与える．$L>0$ と
$Lw=\kappa^2-1+2\kappa/s>0$ より，整数 $w$ は正である．

$S=FT-B(t-y)-(u+z)$ を使うと
\[
 S=(s\kappa-1)L+s\kappa(s-\kappa-1+y)
\]
であり，元の第一残余は
\begin{equation}\label{ax24:degenerate-original}
 \mathfrak r=\kappa\bigl[t(sw-1)-Dy\bigr]
\end{equation}
と因数分解される．従って整数 $y$ の存在には $D\mid t(sw-1)$ が必要である．

この整除性と式~\eqref{ax24:degenerate-Q} から，
\begin{equation}\label{ax24:divisor-bound}
 D\mid C_s:=K_0\bigl((s-1)^4+3\bigr)
\end{equation}
を導く．そのために，変数 $X$ の整数多項式を
\begin{equation}\label{ax24:resultant-polynomials}
 \begin{aligned}
 A_s(X)&=s^2K_0X^2+s(s-1)^2X-(s^2+s+1),\\
 G_s(X)&=s^2(s-1)X^2-sX-s+1,\\
 H_s(X)&=(sX+2)G_s(X)
 \end{aligned}
\end{equation}
と定める．直接展開による二恒等式
\begin{align*}
 sQ_s+A_s(\kappa)
 &=D\bigl[\kappa s^2(s-1)+sw-s-1\bigr],\\
 t(sw-1)+H_s(\kappa)
 &=D\bigl[\kappa^2s^2(s-1)+\kappa sw-\kappa s-\kappa-s+1\bigr]
\end{align*}
から $D\mid A_s(\kappa),H_s(\kappa)$ となる．整数係数の resultant の
B\'ezout 恒等式により，$D$ は両多項式の resultant を割る．
ここで
\begin{equation}\label{ax24:resultants}
 \begin{aligned}
 \operatorname{Res}_X(A_s,sX+2)&=s^2K_0,\\
 \operatorname{Res}_X(A_s,G_s)&=-s^6\bigl((s-1)^4+3\bigr),\\
 \operatorname{Res}_X(A_s,H_s)&=-s^8C_s.
 \end{aligned}
\end{equation}
第一式は $s^2A_s(-2/s)$ の計算であり，第二式は二次式の公式
$(af-cd)^2-(ae-bd)(bf-ce)$ から従う．実際，$J=(s-1)^2+1$ とすれば
第二式の右辺は $s^6\bigl(4(s-1)^2-J(J+2)\bigr)$ である．
第三式は resultant の積性による．$\gcd(D,s)=1$ なので
式~\eqref{ax24:divisor-bound} を得る．

次に $r=\kappa-K_0w$ と置くと，式~\eqref{ax24:degenerate-Q} は
\begin{equation}\label{ax24:small-cofactor}
 \kappa(sr+2)=s(w^2-w+1)
\end{equation}
となる．$s\ge3$ と右辺の正値性から $r\ge0$ である．
$r=0$ なら $\kappa=K_0w$，$2K_0w=s(w^2-w+1)$ なので $w\mid s$ である．
従って $w\le s$ だが，同じ式は $2K_0\le sw\le s^2$ を要求する．
$2K_0-s^2=(s-1)^2+1>0$ より矛盾する．

ゆえに $r\ge1$ である．式~\eqref{ax24:small-cofactor} は
$w^2-w+1>K_0w$ を与えるので $w\ge K_0+1$ となり，
\begin{equation}\label{ax24:D-lower}
 D=s(\kappa+w)+1\ge s\bigl((K_0+1)^2+1\bigr)+1
\end{equation}
を得る．また
\[
 s^2(K_0+1)^2-C_s
 =3s^4(s-2)+10s^2(s-1)+8s-4>0.
\]
正整数 $e=C_s/D$ は従って $e<s$ を満たす．
$C_s\equiv4\pmod s$，$D\equiv1\pmod s$ より $e\equiv4\pmod s$ でもある．
$s=3$ では $C_s=133$ に対し式~\eqref{ax24:D-lower} が $D\ge196$ を要求する．
$s=4$ では $0<e<4$ と $e\equiv0\pmod4$ が両立しない．
従って残る $s\ge5$ では
\begin{equation}\label{ax24:cofactor-four}
 e=4,\qquad D=C_s/4
\end{equation}
に強制される．

最後にこの唯一の約数も平方障害で除く．
$M=(D-1)/s=\kappa+w$ とすると，式~\eqref{ax24:degenerate-Q} は
\begin{equation}\label{ax24:kappa-quadratic}
 sK_0\kappa^2+
 \bigl(2-s[M(K_0-2)+1]\bigr)\kappa
 -s\bigl(M(M-1)+1\bigr)=0
\end{equation}
となる．$D=C_s/4$ を代入したこの二次式の判別式を $\Delta$ とする．
直接計算により
\begin{equation}\label{ax24:square-gap}
 16\Delta=P(s)^2-4W(s)
\end{equation}
であり，ここで
\begin{align*}
 P(s)={}&s^8-6s^7+17s^6-30s^5+41s^4\\
        &-44s^3+34s^2-20s-2,\\
 W(s)={}&2s^7-13s^6+34s^5-73s^4+100s^3-106s^2+68s-15.
\end{align*}
$v=s-4\ge0$ における $W(s)$ と $G(s)=2P(s)-1-4W(s)$ の係数は次の通りである：
\begin{center}\small
\begin{tabular}{c|r|r}
次数 $j$ & $[v^j]W(v+4)$ & $[v^j]G(v+4)$\\\hline
0&609&26135\\
1&6324&42152\\
2&8934&34636\\
3&5652&19704\\
4&1967&8134\\
5&394&2316\\
6&43&422\\
7&2&44\\
8&0&2
\end{tabular}
\end{center}
従って $W>0$，$G>0$ である．$P$ は正整数であり，
\[
 (P-1)^2<16\Delta<P^2
\]
となる．しかし整数根 $\kappa$ が存在すれば $\Delta$ は整数の平方であり，
$16\Delta$ も平方となる．連続する二平方の間にあることと矛盾する．
\end{proof}

\subsection{退化面全体の一整数を固定した有限復元}

\begin{corollary}[退化面の低行一整数有限性]
\label{ax24:degenerate-finite}
$a=0$ を満たす三高非負係数模型は，$s\ge3$ に対して
\begin{equation}\label{ax24:uniform-b-bound}
 2\le b\le1+s^2(s-1)^2
\end{equation}
を満たす．従って $s$ だけを固定したとき，全ての残りの整数を有限に復元できる．
この系は $s$ 自体に一様上界を与えず，退化面全体の不存在も結論しない．
\end{corollary}

\begin{proof}
$x\ge1$ の場合は式~\eqref{hd23:degenerate-relations} の $s-bx>0$ により
$b\le\lfloor(s-1)/x\rfloor$ である．$x=0$ の場合は前定理から $R_0\ge2$ である．
退化面の元残余による整除条件は
\begin{equation}\label{ax24:general-degenerate-divisibility}
 B=b-1\mid h(s-1)(R_0-1)
\end{equation}
であり，$x=0$ では $h=s$，$R_0\le s(s-1)+1$ なので
$b\le1+s^2(s-1)^2$ となる．

式~\eqref{ax24:general-degenerate-divisibility} の導出も記す．
$S_0=FT-Bt-\ell$，$\alpha_0=F+tz$，$H=s+tb$，$K=sT-b$ と置く．
$a=0$ では $S=S_0+By$，$\alpha=\alpha_0+uy$ であり，
\[
 \mathfrak r=\mathfrak r_0-(BH-uK)y,
 \qquad \mathfrak r_0=-HS_0+(s\alpha_0-t)T-\alpha_0b.
\]
$uh=B(F+1)$ より $B\mid h(BH-uK)$ なので $B\mid h\mathfrak r_0$ である．
法 $B$ での恒等式
\[
 h(\mathfrak r_0+tT)
 =h(tT+1)(s\ell-F)+t(1-sT)uh
 \pmod B
\]
と $s\ell=bF$ から $B\mid htT$ を得る．
$ht\equiv h(1-s)\pmod B$，$sT=b^2-R_0\equiv1-R_0\pmod B$ を順に使えば
式~\eqref{ax24:general-degenerate-divisibility} が従う．

具体的な有限復元では，$x=0,\ldots,s-1$，$h=s-x$，上の有限範囲の $b$，
および $R_0=1,\ldots,h(s-1)+1$ を取る．$x=0,R_0=1$ は除外する．
各組に対し，二次式
\[
 R_0F^2+(R_0-b[h(s-1)+1]+sh)F+R_0-b^2+sh=0
\]
の正整数根 $F$ は高々二個である．次に
\[
 u=\frac{B(F+1)}h,\quad \ell=\frac{bF}s,\quad
 t=u-s+1,\quad z=\ell-u,\quad T=\frac{b^2-R_0}s
\]
を復元し，整数性と必要な正値性を確認する．
元残余は $y$ の一次式で，その傾きは命題~\ref{hd22:residual-quadratic} により
厳密に負である．従って $y$ は高々一個に定まる．
非負整数性，単一合同，三尺度，六符号，元残余を検査し，
命題~\ref{hd21:reconstruction} の行列と正 weights を構成すれば必要十分な復元となる．
\end{proof}

\nolinkurl{verify_high_d_axis_degenerate_v24.py} は，新しい多項式恒等式，
二つの resultant，平方間隔と全係数符号を厳密整数演算で確認する．
有限範囲の回帰確認を上の無上界証明の代用としていない．

% v24: proportional high-D sector, unbounded exclusion off the zero axis.
\section{比例する高係数の排除と残る零軸の有限復元}
\label{hp24:section}

本節では前節の canonical 表示と正の整数模型の必要十分復元を用いる．
記号を再掲すると，
\[
\begin{gathered}
s,t\ge1,\quad b\ge2,\quad B=b-1,\quad u=s+t-1,\\
x,y,z\ge0,\quad d=x+y+z\ge1,\quad h=s-x\ge1,\quad j=t-y,\\
\ell=u+z,\quad F=B+ux+sz,\quad m=F^2+F+1,\quad
\alpha=ud+B+z
\end{gathered}
\]
である．本節の対象は $\delta=sy-tx=0$ の領域である．
尺度は格子変数 $n=\ell-U\in\mathbb Z$ を用いて
\begin{equation}\label{hp24:lattice}
 S=j-(F+1)n,\qquad T=h+Fn,\qquad U=\ell-n
\end{equation}
と書ける．差方程式に加えて，元の pair-product 残余
\begin{equation}\label{hp24:residual}
 \mathfrak r=-(s+tb)S+(s\alpha-t)T-\alpha b=0
\end{equation}
を必ず保持する．

\begin{theorem}[比例領域の零軸への帰着]
\label{hp24:main}
三高 $D$ が同時に非負である正の canonical 整数模型で $sy=tx$ なら，
\[
 x=y=0,\qquad z\ge1,\qquad s<t,\qquad 0\le n\le z
\]
が必要である．従って $sy=tx$ かつ $x>0$ の領域は，低行整数，分母，
尺度，高係数に上限を置かずに排除される．
残る $x=y=0$ の零軸の一般不存在を，本定理では主張しない．
\end{theorem}

\subsection{正の比例係数では格子変数が零になる}

$sy=tx$ と $h=s-x>0$ より $j=th/s>0$ である．
また $F=B+h+s(d-1)\ge h+1$ なので，$T>0$ は $n\ge0$ を要求する．
$x\ge1$ のときは
\[
 F+1-j=B+(s+t-1)(x-1)+s+sz+y>0.
\]
従って $n\ge1$ なら $S<0$ となり，$n=0$ である．
この議論では $s$ と $t$ の順序を交換していない．
$s\ge t$ は定理~\ref{hd23:main} で排除済みなので，以下では $s<t$ とする．

このとき $S=j,T=h$ であり，差方程式は
\begin{equation}\label{hp24:Q}
 Q=(d-1)(sh-1)+h(h-3)+(b+1)(h-j)=0
\end{equation}
となる．また $w=j-h\ge1$ と置けば
\[
 s=h+x,\qquad j=h+w,\qquad t=\frac{sj}{h},\qquad
 d-1=\frac{x(h+j)}h+z-1.
\]
これらの式は，元の整数模型から得られる有理式として用いる．
任意の実数解を模型とみなすことはしない．

\subsection{二以上の正係数を持つ領域の整数係数証明}

$x\ge2$ とする．次の整数係数多項式を定義する：
\begin{equation}\label{hp24:positive-certificate-data}
\begin{aligned}
s&=h+x,&j&=h+w,&q&=sh-1,\\
D_*&=x(h+j)+h(z-1),& L&=h^2+sD_*,\\
C&=D_*q+h^2(h-3)-hw,&
A&=hC-h^2w+w(h+j)L,\\
\mathcal N&=-A(sh^2w-C)+2sjh^3w^2+Csj^2hw.
\end{aligned}
\end{equation}
差方程式~\eqref{hp24:Q} は $bhw=C$ を与え，
$\lambda=ux+sz=L/h$，$\alpha=A/(h^2w)$ となる．
従って元の残余を代入すると，恒等的に
\begin{equation}\label{hp24:negative-residual-certificate}
 \mathcal N=-h^3w^2\mathfrak r
\end{equation}
である．

ここで式~\eqref{hp24:positive-certificate-data} に
\[
 h=H+1,\qquad x=X+2,\qquad w=W+1
\]
を代入し，$H,X,W,z$ の多項式として展開する．
非零係数は $281$ 個で，全て正であり，定数項は $312$ である．
これは明示的な整数多項式の有限係数命題であり，
同梱の \nolinkurl{verify_high_d_proportional_v24.py} が整数辞書演算で全係数を再計算する．
係数一覧と，変数の各次数ごとの係数数・最小係数は結果JSONに保持する．
検算は他の検算器を import せず，元の残余から分母を払って得た多項式との
恒等式~\eqref{hp24:negative-residual-certificate} も独立に係数比較する．

$H,X,W,z\ge0$ より $\mathcal N>0$ である．
しかし模型では $\mathfrak r=0$ が必要なので矛盾する．
以上は有界なパラメータ探索ではなく，表示された全多項式の係数による証明である．

\subsection{係数一の領域の解析的排除}

$x=1$ とする．$sy=t$ より，$s=h+1$，$y=v+1$ と書けば
\[
 v\ge1,\qquad t=(h+1)(v+1),\qquad j=h(v+1),\qquad d-1=v+1+z.
\]
式~\eqref{hp24:Q} は
\begin{equation}\label{hp24:x-one-v}
 (hb-h^2+1)v=(h^2+h-1)z+2h^2-2h-1
\end{equation}
となる．$h\ge2$ では右辺が正なので $hb-h^2+1>0$ であり，
整数 $b$ は $b\ge h$ を満たす．$k=b-h\ge0$ と置くと
\[
 v=\frac{(h^2+h-1)z+2h^2-2h-1}{hk+1}.
\]
この式を元の残余へ代入し，
$\mathcal R_h(k,z)=(hk+1)^2\mathfrak r$ と置く．
直接展開により
\begin{equation}\label{hp24:R0}
\begin{aligned}
\mathcal R_h(0,z)={}&h^2(h+1)(h^2+h-1)z^2\\
&+h(4h^4+h^3-6h^2+2)z+h^2(4h^3-6h^2+3)>0
\end{aligned}
\end{equation}
である．一方，
\begin{equation}\label{hp24:R1}
\mathcal R_h(1,z)=-h(h+1)^2\bigl((h^2+h-1)z^2+h(3h+1)z+2h^2+1\bigr)<0.
\end{equation}
さらに $\mathcal R_h(k,z)=\sum_{i=0}^4 R_i(h,z)k^i$ の正次数係数は
\begin{align*}
-R_1={}&2h^2(h^3+2h^2-1)z^2+h(7h^4-10h^2+h+2)z\\
 &+6h^4(h-2)+3h^3+3h^2-1,\\
-R_2={}&h(h^3+2h^2-1)z^2+h^2(8h^2+7h-2)z
 +10h^4-6h^3-h+1,\\
-R_3={}&h\bigl(2h(h+1)z+6h^2+h+2\bigr),\qquad -R_4=h^2.
\end{align*}
$h\ge2,z\ge0$ で全て厳密に正である．従って $k\ge1$ では
$\mathcal R_h(k,z)\le \mathcal R_h(1,z)<0$ となる．整数 $k\ge0$ に零点はない．

$h=1$ では $b\ge2$ より $k=b-1\ge1$ である．$k=K+1$ と書けば
\begin{equation}\label{hp24:h-one}
\begin{aligned}
-\mathcal R_1(K+1,z)={}&K^4+13K^3+37K^2+38K+12\\
&+(4K^3+25K^2+38K+16)z+(2K^2+8K+4)z^2>0.
\end{aligned}
\end{equation}
従ってこの場合も $\mathfrak r=0$ と両立しない．
これで $x>0$ の全域が排除された．

\subsection{零軸の格子幅と四次方程式}

$x=0$ なら $sy=tx$ から $y=0$ であり，$d=z\ge1$ である．
ここでは
\[
 F=b-1+sz,\quad m=F^2+F+1,\quad K=s(F-1)-z+1
\]
と置く．差方程式は
\begin{equation}\label{hp24:axis-t}
 (b+1)t=mn+K
\end{equation}
であり，元の残余は
\begin{equation}\label{hp24:axis-r}
 \mathfrak r=-bt^2+(s^2z-2s-bz)t+F(s^2-b)
       +n\bigl(sm+(F^2+b)t\bigr)
\end{equation}
となる．従って
\begin{equation}\label{hp24:axis-polynomial}
\begin{aligned}
\mathcal P_{s,z,n}(b)={}&-b(mn+K)^2
 +(b+1)(s^2z-2s-bz)(mn+K)\\
 &+(b+1)^2F(s^2-b)\\
 &+n\bigl(s m(b+1)^2+(F^2+b)(mn+K)(b+1)\bigr)
\end{aligned}
\end{equation}
は $(b+1)^2\mathfrak r$ に等しい．

$n\ge z+1$ と仮定し，
\[
 s=\sigma+1,\quad b=\beta+2,\quad z=\zeta+1,
 \quad n=\zeta+\nu+2
\]
を代入する．式~\eqref{hp24:axis-polynomial} の非零係数は $209$ 個で，全て正で，
定数項は $63$ である．この有限係数命題も同じ独立な整数辞書演算で確認できる．
$\sigma,\beta,\zeta,\nu\ge0$ より $\mathcal P>0$ となり，模型に必要な
$\mathfrak r=0$ と矛盾する．従って $0\le n\le z$ である．
これで定理~\ref{hp24:main} が証明された．

\begin{lemma}[整除余りの共鳴と明示上界]
\label{hp24:divisor-bound}
零軸の模型では
\[
 C_n=n\bigl((sz)^2-3sz+3\bigr)+(s^2-1)z-3s+1
\]
が正であり，
\begin{equation}\label{hp24:divisor-bound-formula}
 b+1\mid C_n,\qquad b+1\le C_n\le
 s^2z^3-3sz^2+(s^2+2)z-3s+1
\end{equation}
を満たす．
\end{lemma}

\begin{proof}
次の恒等式を用いる：
\[
 mn+K=(b+1)\bigl(s+n(b+2sz-2)\bigr)+C_n.
\]
従って $b+1\mid C_n$ である．$M_0=(sz)^2-3sz+3\ge1$ なので，
$C_n$ は $n$ とともに増加する．
$C_0=(s^2-1)z-3s+1$ は，$s\ge4$ なら正，
$s=3$ なら $8(z-1)$，$s=2$ なら $3z-5$，$s=1$ なら $-2$ である．
$0\le n\le z$ を合わせると，$C_n=0$ は
\[
 (s,z,n)=(3,1,0),\qquad (1,2,2)
\]
だけである．その他の $C_n\le0$ は $-1,-2$ であり，$b+1\ge3$ が割ることはない．
第一の零例外は式~\eqref{hp24:axis-t} から $t=s=3$ となり既に排除済みである．
第二では $t=2b+5$，$S=1,T=2b+3,U=2b+5$ となるが，元残余は
\[
 \mathfrak r=-b^2+5b+17
\]
である．判別式 $93$ は $9^2<93<10^2$ により平方数でなく，整数根はない．
従って模型では $C_n>0$ である．最後の上界は $n\le z$ と $M_0>0$ から従う．
\end{proof}

\begin{corollary}[二整数を固定した零軸の完全有限復元]
\label{hp24:finite-axis}
$s,z\ge1$ を固定する．比例領域の正整数模型は次の有限手順で全て復元される．
\begin{enumerate}
\item $n=0,\ldots,z$ ごとに，四次式 $\mathcal P_{s,z,n}(b)=0$ の整数根 $b\ge2$ を全て求める．
\item 式~\eqref{hp24:axis-t} から $t$ を求め，整数性と $t>s$ を課す．
\item $S=t-(F+1)n,T=s+Fn,U=s+t+z-1-n$ について，正整数性，
六つの符号条件，元の全 pair-products と正 weights を検査する．
\end{enumerate}
各 $n$ の候補 $b$ は高々四個であり，検査前の候補数は高々 $4(z+1)$ である．
この手順は $b,t$ や分母の探索上限を含まない．
\end{corollary}

\begin{proof}
式~\eqref{hp24:axis-polynomial} の $b^4$ の係数は
\[
 n^2-zn-1<0\qquad(0\le n\le z)
\]
なので，各式は恒等的に零にならない真の四次式である．
定数項が零なら $b$ の因子を取り去り，残った整数多項式の正の整数根を
整数根定理から定数項の約数として全て列挙すればよい．
候補を得た後の必要十分復元は従前の canonical 復元を用いる．
模型から出発すれば全手順を通過し，通過したデータは元の全方程式と正 weights を持つ．
\end{proof}

この系は固定した $s,z$ ごとの完全性であり，$s,z$ 自体の一様上界や
零軸の一般不存在を意味しない．\nolinkurl{verify_high_d_proportional_v24.py} \texttt{--axis s z}
はこの復元を実行する．

% v25: unbounded endpoint exclusions on the proportional zero axis.
\section{比例零軸の両端点における無上界排除と固定一整数での有限復元}
\label{hpe25:section}

本節は比例零軸 $x=y=0$ の端点 $n=0,z$ を扱う．
定理~\ref{hp24:main} の必要条件
\[
 s<t,\qquad b\ge2,\qquad z\ge1,\qquad 0\le n\le z
\]
を前提とする．記号を
\[
 F=b-1+sz,\quad m=F^2+F+1,\quad K=s(F-1)-z+1,\quad
 C_n=n((sz)^2-3sz+3)+(s^2-1)z-3s+1
\]
とし，差方程式と元の pair-product 残余を
\begin{align}
 (b+1)t&=mn+K,\label{hpe25:Q}\\
 \mathfrak r&=-bt^2+(s^2z-2s-bz)t+F(s^2-b)
 +n\{sm+(F^2+b)t\}=0\label{hpe25:residual}
\end{align}
と書く．既知の $C_n=0$ の二つの例外は模型を与えないので，
以下の模型では必ず $C_n>0$ かつ $b+1\mid C_n$ である．

\begin{theorem}[比例零軸の端点有限化]\label{hpe25:main}
正の canonical 整数模型に対し，次が必要である．
\begin{enumerate}
\item $n=0$ なら
\[
 s\ge3,\qquad s<b<\frac32s,\qquad 1\le z<\frac s2.
\]
\item $n=z$ なら
\[
 s\ge2,\qquad 1\le z\le s,\qquad
 1\le e:=\frac{C_z}{b+1}\le s-1.
\]
\end{enumerate}
従って両端点は，低行の一整数 $s$ だけを固定すれば全て有限に復元できる．
$n=0$ では各 $b$ について $z$ の二次方程式を解くことで，
完全検査前の $(b,z)$ の候補数は高々
$2\lfloor(s-1)/2\rfloor$ である．
$n=z$ では $(z,e)$ の候補数は高々 $s(s-1)$ である．
いずれも $t,S,T,U$ に事前の上限を置かない．
本定理は全ての $s$ を通じた両端点の不存在を主張せず，
内部 $0<n<z$ の一般不存在も主張しない．
\end{theorem}

\subsection{整数商と元の残余}

正整数 $e=C_n/(b+1)$ を導入すると，式~\eqref{hpe25:Q} は
\begin{equation}\label{hpe25:e-coordinates}
 t=s+n(b+2sz-2)+e,\qquad S=s+n(sz-2)+e
\end{equation}
となる．ここで $T=s+Fn$，$U=s+t+z-1-n$ である．
$e$ は独立に選べる正実数ではなく，$b+1$ による正整数商である．
以下の四次式は式~\eqref{hpe25:residual} を代入して作り，
差方程式だけから候補を模型として認めることはしない．

\subsection{端点 \texorpdfstring{$n=0$}{n=0} の三つの無上界排除}

$n=0$ では $S=t,T=s$ なので，正の matching entry
$p_{ba}=s^2-b>0$ より $s\ge2$ かつ $2\le b\le s^2-1$ である．
$t=K/(b+1)$ を元の残余に代入し，
\[
 P_0(s,b,z):=(b+1)^2\mathfrak r
            =A_2z^2+A_1z+A_0
\]
と置くと
\begin{align}
 A_2={}&(s^2-1)(s^2-b^2),\label{hpe25:P0coeff}\\
 A_1={}&(b-1)(-2b^2s-bs-b+3s^3-s^2-2s),\notag\\
 A_0={}&-b^4-b^3+3b^2s^2-2b^2s+b^2-3bs^2+2bs+3s^2-2s.\notag
\end{align}
$s=2$ では $b=2,3$ のみであり，それぞれ
\[
 P_0(2,2,z)=-2(3z-2),\qquad
 P_0(2,3,z)=-(z+1)(15z+43)
\]
となるので，$z\ge1$ に零点はない．従って $s\ge3$ としてよい．

まず $b\le s$ を排除する．$s=b+a$ と書くと
\begin{align*}
P_0={}&z^2a^4+z(4bz+3b-3)a^3\\
 &+\{(5b^2-1)z^2+(9b^2-10b+1)z+3b^2-3b+3\}a^2\\
 &+\{2b(b^2-1)z^2+(7b^3-10b^2+b+2)z
             +6b^3-8b^2+8b-2\}a\\
 &+b(b-1)\{(b-3)(b+1)z+2(b-1)^2\}.
\end{align*}
$b\ge3$ では全ての項が非負で，定数項は正である．
$b=2$ では $s\ge3$ より $a\ge1$ であり，
一次項だけでも負の定数項 $-6z+4$ を上回る：
\[
 (12z^2+20z+30)a-6z+4\ge12z^2+14z+34>0.
\]
従って $P_0=0$ は $b>s$ を必要とする．

次に $b\ge3s/2$ を排除する．$b=3s/2+a$，$a\ge0$ と置くと
\[
 -16P_0=\sum_{i=0}^4 L_i(s,z)a^i
\]
であり，係数は
\begin{align*}
 L_4={}&16,\qquad L_3=16(2sz+6s+1),\\
 L_2={}&8\{2(s^2-1)z^2+(18s^2-2s+2)z+21s^2+13s-2\},\\
 L_1={}&4\{12s(s^2-1)z^2+(42s^3-8s^2+16s-4)z
                    +18s^3+63s^2-20s\},\\
 L_0={}&s\{20s(s^2-1)(z-1)^2
 +(76s^3+36s^2+4s-56)(z-1)\\
 &\hspace{30mm}+29s^3+234s^2-108s-24\}.
\end{align*}
これらは $s\ge3,z\ge1$ で全て正なので $P_0<0$ となる．

最後に $z\ge s/2$ を排除する．既に $b>s$ なので整数
$a=b-s-1\ge0$ を用い，$w=z-s/2\ge0$ と置く．直接展開により
\begin{equation}\label{hpe25:wide-z}
 -4P_0(s,s+1+a,s/2+w)=J_2w^2+J_1w+J_0
\end{equation}
であり，
\begin{align*}
 J_2={}&4(a+1)(s^2-1)(a+2s+1),\\
 J_1={}&4\{2sa^3+(s^3+6s^2+4s+1)a^2
 +(2s^4+5s^3+9s^2+5s+1)a+s^2(s^2+7s+4)\},\\
 J_0={}&4a^4+4(s^2+4s+5)a^3
 +(s^4+12s^3+21s^2+70s+32)a^2\\
 &+2(s^5+4s^4+6s^3+38s^2+37s+10)a
 +5s^4+26s^3+29s^2+28s+4.
\end{align*}
$s\ge3$ では右辺は正なので，この領域も零点を持たない．
これで第一の必要条件が証明された．さらに $b>s$ より
式~\eqref{hpe25:P0coeff} の $A_2$ は負であり，固定 $(s,b)$ ごとの
$z$ の候補は高々二つである．整数 $b$ の個数は
$\lfloor(s-1)/2\rfloor$ なので，定理の候補数が従う．

\subsection{端点 \texorpdfstring{$n=z$}{n=z} の単調四次式}

$n=z$ では
\begin{align*}
 C&:=C_z=s^2z(z^2+1)-3s(z^2+1)+2z+1,\\
 A&:=3s(z^2+1)-4z-1,\\
 B&:=3s^2(z^2+1)^2-2s(z^2+1)(4z+1)+5z^2+4z-2,\\
 H&:=sz(z^2+1)-z^2-z+3
\end{align*}
と置く．式~\eqref{hpe25:e-coordinates} と $b=C/e-1$ を元の残余へ
代入すると，恒等的に
\begin{equation}\label{hpe25:R}
 R(e):=e^2\mathfrak r=e^4+Ae^3+Be^2+CHe-C^2
\end{equation}
となる．従って $e$ はこのモニック四次式の正整数零点でなければならない．

まず $s\ge2$ とする．$a=s-2\ge0$ と置くと
\[
 A=3a(z^2+1)+6z^2-4z+5>0,
\]
また
\begin{align*}
 B={}&3a^2(z^2+1)^2+2a(z^2+1)(6z^2-4z+5)\\
 &+12(z-1)^4+32(z-1)^3+49(z-1)^2+38(z-1)+15>0.
\end{align*}
さらに $H\ge2z^3-z^2+z+3>0$ である．$C>0$ と合わせて
$R'(e)>0$ が $e\ge0$ の全域で成立する．$R(0)=-C^2<0$ なので，
$R$ は正実零点をちょうど一つ持つ．

この零点が $s$ より小さいことを示す．$a=s-2$ に関して
$R(s)=\sum_{i=0}^4 Q_i(z)a^i$ と展開すると
\begin{align*}
Q_0={}&16z^5+8z^4+28z^3+38z^2+6z+25,\\
Q_1={}&24z^5+52z^4+18z^3+171z^2-31z+117,\\
Q_2={}&12z^5+58z^4+2z^3+177z^2-30z+130,\\
Q_3={}&2z^5+23z^4+69z^2-6z+53,\\
Q_4={}&3z^4+9z^2+7.
\end{align*}
$z\ge1$ で各係数が正であるから $R(s)>0$ である．
従って正整数零点は $1\le e\le s-1$ を満たす．

次に $z\ge s+1$ を排除する．$w=z-s-1\ge0$ とすると
\begin{equation}\label{hpe25:root-gap}
 -R(s-1)=\sum_{i=0}^6 E_i(s)w^i,
\end{equation}
ここで
\begin{align*}
 E_6={}&s^3,\\
 E_5={}&2s^2(3s^2+2s-2),\\
 E_4={}&s(15s^4+17s^3-6s^2-16s+5),\\
 E_3={}&2(10s^6+14s^5+2s^4-14s^3-13s^2+11s-1),\\
 E_2={}&15s^7+22s^6+8s^5-15s^4-34s^3-15s^2+38s-10,\\
 E_1={}&(s-1)(6s^7+14s^6+14s^5+4s^4+4s^3-4s^2-24s+13),\\
 E_0={}&(s-1)^2(s^3+2s^2+s-1)(s^4+s^3-4s+1).
\end{align*}
これらは全て $s\ge2$ で正である．例えば $E_2$ では
$15s^7-15s^4$，$22s^6-34s^3$，$8s^5-15s^2$，$38s-10$
をそれぞれ正の組として評価できる．他の係数も同様に，または
$s=a+2$ に展開して全係数が正であることから従う．
従って $R(s-1)<0<R(s)$ であり，唯一の正実零点は開区間
$(s-1,s)$ にある．正整数零点は存在しない．これで $z\le s$ を得る．

残る $s=1$ では $C=(z-1)^3-1$ なので $C>0$ は $z\ge3$ を要求する．
$z\ge3$ で $A,B,H$ は正であり，$B$ は $w=z-2$ と書けば
$3w^4+16w^3+33w^2+32w+11$ である．また
\[
 R(1)=2w^5+14w^4+44w^3+72w^2+61w+18>0.
\]
従って唯一の正実零点は $(0,1)$ にあり，正整数解はない．
定理第二項と候補数 $s(s-1)$ が従う．

\subsection{完全復元と検証の境界}

$n=0$ では定理の有限範囲内で $b$ を列挙し，
式~\eqref{hpe25:P0coeff} の整数零点 $z$ を判別式の完全平方検査で求める．
その後 $b+1\mid K$，$t=K/(b+1)>s$ と元の残余を検査する．
$n=z$ では定理の $(z,e)$ を列挙し，$e\mid C$，$b=C/e-1\ge2$ として
$t=s+z(b+2sz-2)+e$ を求め，式~\eqref{hpe25:R} と元の残余を検査する．

いずれも通過した候補に対し，全ての $S,T,U$ の正値性，符号，
四つの行和，六つの pair-product，欠損
$(D_a,D_b,D_c,D_0)=(0,0,z,-z)$，正の重みの全有向不等式を
有理数演算で再検査して初めて模型と認める．
同梱の \nolinkurl{verify_high_d_axis_endpoints_v25.py} は標準ライブラリのみを用い，
元の残余から恒等式を独立に再構成し，上記の正係数証明を係数比較する．
これは有界探索による無上界主張ではない．
付随する $1\le s\le400$ の端点復元で模型が零個であったことは，
この新しい定理を使った有限検査結果としてのみ記録する．
一般の $s$，および内部 $0<n<z$ は依然として開いている．

% v25. Exact factorization of the original degenerate pair residual.
\section{退化面の積表示と互いに素な領域の排除}
\label{dg25:section}

三高係数が同時に非負である canonical 整数模型の記号
\eqref{hd23:data} を用いる．退化面 $a=uh-B(F+1)=0$ では，
元の pair 残余を単なる一次式として解く以上に，正整数の積として扱える．
これは $x\ge1$ の残余に対する新しい互いに素な場合の障害と，
Schur 整数 $b$ だけを固定した有限復元を与える．
退化面全体の不存在は結論しない．

\subsection{元の残余の正整数積への変換}

\begin{theorem}[退化面の積表示]
\label{dg25:product}
$a=0$ とし，
\begin{equation}\label{dg25:data}
 \begin{gathered}
 A=B(s-1)-1=Bs-b,\quad R=b^2-sT,\qquad
 L=(F+1)R-hA,\quad M=uR-BA
 \end{gathered}
\end{equation}
と置く．このとき $L,M$ は正整数であり，
\begin{equation}\label{dg25:identities}
 hM=BL,\qquad FL=bF+s(T-h),\qquad
 \mathfrak r=M(t-y)-tT.
\end{equation}
従って元の pair 方程式 $\mathfrak r=0$ は
\begin{equation}\label{dg25:positive-product}
 Mj=tT,\qquad j=t-y\in\{1,\ldots,t\}
\end{equation}
を要求する．特に $y<t$，$M\ge T$，$M\mid tT$ である．
\end{theorem}

\begin{proof}
既知の退化関係より $s\ell=bF$ であり，
$K_0=h(s-1)+1$ と置いたとき
\[
 R(F^2+F+1)=(bK_0-sh)F+b^2-sh.
\]
ここで $bK_0-sh=hA+b$ である．従って
\begin{equation}\label{dg25:Q}
 Q:=R(F^2+F+1)-(hA+b)F-b^2+sh=0.
\end{equation}
定義から $hM-BL=Ra=0$ であり，$R=b^2-sT$ を用いると
$FL-bF-s(T-h)=Q=0$ となる．

$M$ の正値性は恒等式として確認できる．$m=F^2+F+1$，
\[
 C=F\{b(t-1)+s\}+b^2t+bs-b>0
\]
と置けば，$t=u-s+1$ を使った直接展開により
\begin{equation}\label{dg25:positive-M}
 mM-C=uQ+(FA-s)a
\end{equation}
である．ゆえに $M=C/m>0$ であり，$hM=BL$ と $h,B>0$ から $L>0$ でもある．

退化面上では $T$ は $y$ に依存せず，
\[
 S=FT-B(t-y)-\ell,\quad
 \alpha=F+tz+uy,\quad z=\ell-u
\]
である．これらを元の残余に代入すると，より一般の恒等式
\begin{equation}\label{dg25:original-identity}
 -(s+tb)S+(s\alpha-t)T-\alpha b
 =M(t-y)-tT+(s\ell-bF)(tT+1)
\end{equation}
を得る．$s\ell=bF$ により主張の残余式になる．
模型ではその残余が零であり $M,t,T>0$ なので $j>0$ である．
$y\ge0$ より $j\le t$，従って $M=tT/j\ge T$ となる．
\end{proof}

差方程式だけから式~\eqref{dg25:positive-product} を結論してはいけない．
最後に $\mathfrak r=0$ を用いることが不可欠である．
一般の三高模型に $j>0$ を仮定するのも正しくない．ここでは退化面上で証明した．

\begin{corollary}[互いに素な領域の無上界排除]
\label{dg25:coprime}
同じ退化模型で $h\ge b$ なら，
\begin{equation}\label{dg25:small-L}
 1\le L\le b-1,\qquad
 g:=\gcd(h,b-1)\ge2,\qquad
 \frac hg\le b-1,\qquad
 T\le M\le b-2
\end{equation}
である．従って $h\ge b$ かつ $\gcd(h,b-1)=1$ の領域は存在しない．
特に $x\ge1$ の全ての退化模型にこの結論が適用される．
\end{corollary}

\begin{proof}
退化面の既知の範囲は $R\ge1$ である．従って
$T=(b^2-R)/s<b^2/s\le h$ となる．最後は $s\ge h\ge b$ による．
式~\eqref{dg25:identities} から $L<b$ であり，正整数性と合わせて最初の範囲を得る．
$h\mid BL$ は $h/g\mid L$ と同値だから，$h/g\le b-1$ である．
$h\ge b$ の下では $g=1$ は不可能である．
さらに $M=BL/h\le B^2/h<B$ なので，整数性から $M\le b-2$ を得る．

$x\ge1$ の場合には，退化関係 $s-bx>0$ から $s\ge b+1$ と
$x\le(s-1)/b$ が従う．
\[
 s-b-\frac{s-1}{b}=\frac{(b-1)(s-b-1)}b\ge0
\]
なので $x\le s-b$，従って $h=s-x\ge b$ である．
\end{proof}

\subsection{正の第一係数を持つ退化面の固定 \texorpdfstring{$b$}{b} 復元}

\begin{corollary}[固定 $b$ に対する完全有限手順]
\label{dg25:fixed-b}
$a=0,x\ge1$ の三高非負係数模型は $b\ge3$ を満たす．
固定した $b$ ごとに，次の有限手順で全て復元できる．
\begin{enumerate}
 \item $s=b+1,\ldots,b^2-1$，$x=1,\ldots,\lfloor(s-1)/b\rfloor$ を取り，
 $h=s-x$，$B=b-1$，$A=Bs-b$，$g=\gcd(h,B)$ とする．
 \item $T=1,\ldots,\min\{b-2,\lfloor(b^2-1)/s\rfloor\}$ と，
 $h/g$ の正倍数 $L\le B$ を取る．$R=b^2-sT$，$M=BL/h$ とし，$M\ge T$ を検査する．
 \item
 \[
 F=\frac{s(h-T)}{b-L},\quad u=\frac{B(F+1)}h,\quad
 \ell=\frac{bF}s,\quad t=u-s+1,\quad z=\ell-u
 \]
 が整数で，$t\ge1,z\ge0$，$F=B+ux+sz$，
 $L=(F+1)R-hA$ を満たすものを残す．
 \item $M\mid tT$ を検査し，$y=t-tT/M\ge0$ とする．
 元の尺度復元~\eqref{hd23:reconstruction} で $S,T,U$ を計算し，
 三尺度の正整数性，$d\ge1$，$\alpha\ge k\ge b+1$，六つの符号条件，
 全ての行和・pair products と指定 $D$ を確認する．
 \item 通過した模型に完全 matching の有理正 weights を与える．
\end{enumerate}
この手順では $b$ 自体に上界を証明していない．$x=0$ の全退化面には
系~\ref{ax24:degenerate-finite} の固定 $s$ 手順を使う．
\end{corollary}

\begin{proof}
$b\ge3$ は $1\le T\le b-2$ から従う．$s\ge b+1$ は前の証明による．
$R\ge1,T\ge1$ より $s\le b^2-1$ である．
残りの全範囲と復元式は前定理と前系から従い，$b-L>0$ なので零分母はない．
従って全ての所望の模型が列挙される．
最後に元の全方程式と正値性を検査するため，逆向きの十分性も明白である．
これは必要条件だけを満たす整数を模型と認定する手順ではない．
\end{proof}

\paragraph{検算の射程．}
\nolinkurl{verify_degenerate_divisor_v25.py} は式~\eqref{dg25:positive-M}，
\eqref{dg25:original-identity} と二つの積恒等式を整数多項式として係数ごとに検査する．
\texttt{--fixed-b b} は上の証明された有限範囲を全て処理し，任意の分母や $t,z,y$ の
探索上限を追加しない．有限個の $b$ に対する実行結果を，退化面全体の不存在へ外挿しない．

% v17. Complete three-high reconstruction with a fixed low row and column.
\subsection{低分母の行と列を固定した三高分母行列の二次復元}
\label{fc17:section}

三高分母の完全な分類を先に行わなくても，低 ray に接するデータが固定されると，
残る高高 entries は高々二つの候補から復元できる．この手順では正符号 graph の
matching 性も，Schur 系の非特異性も仮定しない．

\begin{proposition}[固定低行列データの完全復元]
\label{fc17:at-most-two}
分母を $q_0=1$，$q_1,q_2,q_3>1$ とし，低 ray の行と列を
\[
 r_i=p_{0i},\qquad x_i=p_{i0},\qquad
 r_i,x_i\in\ZZ,\quad \sum_{i=1}^3r_i=-1,\quad r_ix_i=q_i-1
\]
と固定する．このデータを持ち，行和と全 pair products を満たす整数行列 $P$ は
高々二つである．正の weights を持つものも，この候補の部分集合として完全に決まる．

具体的には，同符号の二つの $r_i$ を $r_2,r_3$ とするよう高標識を交換し，
\[
 s_i=-q_i-x_i,\qquad Q_{ij}=q_iq_j-1
\]
と置く．$A=p_{12}$ は次の真の二次方程式を満たす：
\begin{equation}\label{fc17:quadratic}
 \Phi(A):=(s_2A-Q_{12})\{s_3(s_1-A)-Q_{13}\}
             -Q_{23}A(s_1-A)=0.
\end{equation}
その整数根のうち
\begin{equation}\label{fc17:divisibility}
 A\ne0,s_1,\qquad A\mid Q_{12},\qquad s_1-A\mid Q_{13}
\end{equation}
を満たすものに対して，行列は
\begin{equation}\label{fc17:complete-P}
 P=\begin{pmatrix}
 0&r_1&r_2&r_3\\
 x_1&0&A&s_1-A\\
 x_2&Q_{12}/A&0&s_2-Q_{12}/A\\
 x_3&Q_{13}/(s_1-A)&s_3-Q_{13}/(s_1-A)&0
 \end{pmatrix}
\end{equation}
により一意に定まる．これが全整数解である．
\end{proposition}

\begin{proof}
$r_ix_i=q_i-1>0$ なので，$r_i,x_i$ は非零で同符号である．
$r_i>0$ なら $q_i+x_i\ge q_i+1$，$r_i<0$ なら
$1\le q_i+x_i\le q_i-1$ である．従って $s_i<0$ である．
三つの $r_i$ には同符号の対があるため，指定した標識を選べる．
二次係数は
\begin{equation}\label{fc17:leading-coefficient}
 a=Q_{23}-s_2s_3.
\end{equation}
$r_2,r_3>0$ なら $s_2s_3>q_2q_3$ なので $a<0$，
$r_2,r_3<0$ なら $s_2s_3\le(q_2-1)(q_3-1)<q_2q_3-1$ なので $a>0$ である．
従って二次係数が消える場合はない．

高 row $1$ と pair $12,13$ から，式~\eqref{fc17:complete-P} の対応する entries が定まる．
pair products は正なので，$A$ と $s_1-A$ は非零であり，整数性は
式~\eqref{fc17:divisibility} である．残る二つの row sums が最後の二 entries を決め，
残った pair $23$ は，分母 $A(s_1-A)$ を払うと式~\eqref{fc17:quadratic} となる．
逆に，これらの条件から全ての元の row sums と pair products が成立する．
真の二次式なので候補は高々二つであり，正値性は
補題~\ref{kl10:cycle-positive} の全正 cycle 条件を検査すればよい．
\end{proof}

$\Phi(A)=aA^2+bA+c$ と書けば，
\begin{align*}
 b&=s_2(s_3s_1-Q_{13})+Q_{12}s_3-Q_{23}s_1,\\
 c&=-Q_{12}(s_3s_1-Q_{13}).
\end{align*}
その判別式は対称な式
\begin{equation}\label{fc17:discriminant}
 \Delta=
 (s_1s_2s_3-s_1Q_{23}-s_2Q_{13}-s_3Q_{12})^2
 -4Q_{12}Q_{13}Q_{23}
\end{equation}
を持つ．従って $\Delta$ が非負の平方数であること，
$(-b\pm\sqrt\Delta)/(2a)$ が整数であること，
式~\eqref{fc17:divisibility} を順に確認すればよい．
\nolinkurl{exact_three_high_completion_v17.py} は整数平方根と整数演算で候補を復元し，
元の全行和・pair products・正 cycle を再検査する．
分母や entry を有限範囲で走査する必要はない．

ここで固定しているのは低 row と低 column の両方である．
定理~\ref{sh14:singular-family} の無限族では低 column 自体が動くので，
この高々二つという結論とは両立する．

\begin{remark}[標識を固定した場合の上界の達成]
\label{fc17:sharp}
上界二は，正の weights を要求しても達成される．例えば
\[
 r=(-2,-2,3),\qquad x=(-3,-3,2),\qquad q=(1,7,7,7)
\]
を固定すると，
\[
 P=\begin{pmatrix}
 0&-2&-2&3\\
 -3&0&4&-8\\
 -3&12&0&-16\\
 2&-6&-3&0
 \end{pmatrix}
\]
と，高 ray $1,2$ を交換した行列が，異なる二つの completion となる．
その正の simple cycles は二辺 cycles だけで，gain は $6/7$ と $48/49$ である．
従って両者は正の weights を持ち，命題~\ref{q9:four-line-criterion} による局所 chart を与える．
この二つは標識交換で対応するので，幾何学的な非同値性を主張する例ではない．
\end{remark}

\begin{corollary}[小さい三高分母データの実数障害]
\label{fc17:335}
低 row と column が，標識交換後
\[
 (r_1,r_2,r_3)=(-4,1,2),\qquad
 (x_1,x_2,x_3)=(-1,2,1)
\]
である整数局所 chart は存在しない．実際，分母は $(1,5,3,3)$ であり，
残る高高 entries は実数としても元の row sums と pair products を満たせない．
\end{corollary}

\begin{proof}
この場合 $s=(-4,-5,-4)$，$Q_{12}=Q_{13}=14$，$Q_{23}=8$ であり，
式~\eqref{fc17:quadratic} は
\[
 -12A^2-34A-28=0
\]
となる．しかし
\[
 24(6A^2+17A+14)=(12A+17)^2+47>0
\]
なので実数根はない．この矛盾には正の weights や別の core との $D$ 相殺を用いない．
\end{proof}

\section{四・五角形成分の接着障害}\label{cy12:section}

% v12: rigorous necessary conditions for C4/C5 support components.
% These statements use nonnegative node counts, unlike the path obstruction.

\subsection{四・五角形成分に対する node と符号の制約}
\label{cn12:section}

この節では正の局所 charts と整数 clutching に加え，
\begin{equation}\label{cn12:node-convention}
 n_{ij;k}=1-\frac{p^i_{jk}+p^j_{ik}}{q_{ij}}\geq0
\end{equation}
を仮定する．両端の row sums から $\sum_{k\ne i,j}n_{ij;k}=5$ なので，各 $n_{ij;k}\leq5$ も従う．
まず一般的な必要条件を導く．四・五角形成分の完全排除は後続の定理で与える．

\begin{lemma}[cycle edge の三 punctures]\label{cn12:three-punctures}
高分母 support の連結成分が $C_4$ または $C_5$ であるとする．
その高分母 edge $ij$ の三 punctures のうち，二つは一方の core で高分母へ，
他方で低分母へ向かい，残る一つは両端で低分母へ向かう．
後者の outgoing sum を $\lambda_{ij}q_{ij}$ と書くと
\[
 \lambda_{ij}\in\{-1,0,1\}.
\]
二つの mixed punctures の node 数はそれぞれ一以上であり，その和は $4+\lambda_{ij}$ である．
mixed node 数の下限一は，node 非負性を追加しなくても整数 clutching から従う．
\end{lemma}

\begin{proof}
cycle の次数二と長さが四以上であることから，三 punctures の型が指定通りになる．
共通 low-low puncture には既知の signed-divisor clutching 補題を適用する．
mixed puncture では一端の entry は二高分母間の負の整数 $-\alpha$，
他端の entry は $q-1$ の非零 signed divisor $x$ である．
従って $(x-\alpha)/q<1$ であり，この値が整数なので零以下である．
node 数は $1-(x-\alpha)/q\geq1$ となる．
共通 puncture の node 数が $1-\lambda_{ij}$ なので，全数五から和の式を得る．
\end{proof}

\begin{proposition}[正の pairing の有限な符号型]\label{cn12:cycle-signs}
$C_4$ の vertices を $i\in\mathbb Z/4$，残る label を $*$ とし，
\[
 \varepsilon_i=\operatorname{sign}(p^i_{i+1,*})\in\{1,-1\}
\]
と置く．このとき
\[
 \varepsilon_i=\varepsilon_{i+2}.
\]
従って四つの符号は全て等しいか，交互である．共通 puncture の係数は
\[
 \lambda_{i,i+1}=\frac{\varepsilon_i-\varepsilon_{i+1}}2.
\]

$C_5$ の vertices を $i\in\mathbb Z/5$ とし，
\[
 \varepsilon_i=\operatorname{sign}(p^i_{i+1,i+2})
\]
と置く．$\varepsilon_i=1$ となる vertices は高々二つで，二つなら cycle 上で隣接する．
この場合
\[
 \lambda_{i,i+1}=-\frac{\varepsilon_i+\varepsilon_{i+1}}2.
\]
従って $\lambda$ の個数は，次の三つのいずれかである：
\[
 (\#\{\lambda=1\},\#\{\lambda=0\},\#\{\lambda=-1\})
 \in\{(5,0,0),(3,2,0),(2,2,1)\}.
\]
\end{proposition}

\begin{proof}
二高分母の標準符号から，各高分母 row の低分母二 entries は一正一負であり，
同じ低分母 column に対する二つの高分母 entries も異符号である．
二つの高 edges $i-1,i$ と $i,i+1$ が作る wedge を考える．
外側の低分母 edge $i-1,i+1$ に沿う puncture $i$ では，二つの incoming numerators が同時に正だと，
それぞれ一以上なので node 数が $1-(\text{二つの和})<0$ となる．

$C_4$ ではそれらの符号は $-\varepsilon_{i-1}$ と $\varepsilon_{i+1}$ である．
もう一つの center $i+2$ に対しても同じ議論を行えば，
$\varepsilon_{i-1}$ と $\varepsilon_{i+1}$ が異なる両場合が排除される．
$C_5$ では二つの符号は $\varepsilon_{i-1}$ と $\varepsilon_{i+1}$ である．
従って正の vertices は距離二の pair を含まない．五角形でそのような集合は高々二点で，
二点なら隣接する．

共通 low-low puncture の二つの符号は $C_4$ では
$\varepsilon_i,-\varepsilon_{i+1}$，$C_5$ では
$-\varepsilon_i,-\varepsilon_{i+1}$ である．
同符号なら和は $\pm q$，逆符号なら零なので，表示した $\lambda$ の式が従う．
正の vertices が零，一，二の場合を数えると最後の三組を得る．
\end{proof}

\begin{lemma}[一高分母 triangle の小さい reverse divisor]\label{cn12:reverse-bound}
triangle $i,j,k$ について $q_{ij}=q>1$，$q_{ik}=q_{jk}=1$ とする．
共通 puncture の outgoing numerators が逆符号なら，その reverse numerators は $c,-c$ と書け，
\[
 1\leq |c|\leq4.
\]
これは第三 core の低低 entries の符号を仮定しない．
\end{lemma}

\begin{proof}
逆符号の outgoing numerators は $p,-p$ であり，reverse は $(q-1)/p,-(q-1)/p$ である．
第三 core の対応する二つの entries の積は零なので，少なくとも一つが零である．
それに対応する低 edge の node 数は $1-c$ または $1+c$ である．
これが零以上五以下という条件から $|c|\leq4$ を得る．
\end{proof}

\begin{proposition}[$C_5$ の一高分母 triangle の node 条件]\label{cn12:C5-table}
$C_5$ の高 edge $ij$ と，両端で共通の低 label $k$ を考える．
第三 core $k$ の対応する低低 entries は，一方が零，他方が $-L$，$L\geq1$ である．
$i,j$ を必要に応じて交換して，零を $ki\to kj$ 側に置く．
共通 puncture の reverse numerators を $r,r'$ とすると，三つの node 数は
\[
 (n_{ij;k},n_{ik;j},n_{jk;i})=(1-\lambda,1-r,1-r'+L).
\]
可能性は次の表で尽きる．
\begin{center}
\small
\begin{tabular}{ccll}
$\lambda$ & $(r,r')$ & 三つの node 数 & 必要条件 \\
\hline
$1$ & $(1,q-1)$ & $(0,0,L-q+2)$ & $q-2\leq L\leq q+3$ \\
$0$ & $(1,-1)$ & $(1,0,L+2)$ & $1\leq L\leq3$ \\
$0$ & $(-c,c)$ & $(1,1+c,1-c+L)$ & $1\leq c\leq4$, $\max(1,c-1)\leq L\leq c+4$ \\
$-1$ & $(-(q-1),-1)$ & $(2,q,L+2)$ & $q\leq5$, $1\leq L\leq3$ \\
$-1$ & $(-1,-(q-1))$ & $(2,2,q+L)$ & $q+L\leq5$
\end{tabular}
\end{center}
表では $r,r'$ が $q-1$ の signed divisors であるという元の条件も保持する．
特に第三行には $c\mid q-1$ が必要である．
さらに補題~\ref{ml12:dichotomy} により，非零低低 entry の大きさは
$L\in\{1\}\cup\mathbb Z_{\geq3}$ であり，$L=2$ は除外する．
この表は node 式の必要条件を整理したもので，局所整数模型や幾何学的 chart の十分条件ではない．
特に $\lambda=-1$ の高 edge は $q\in\{3,5\}$ に限られる．
従って，全ての高分母が七以上の $C_5$ 成分では，命題~\ref{cn12:cycle-signs} の
三つの型のうち最後の型は存在しない．
\end{proposition}

\begin{proof}
第三 core は二高分母であり，その二つの低分母が $ki,kj$ である．
低低 entry の負性から，二つの entries はちょうど $0,-L$ である．
signed-divisor clutching の三場合を逆向きに書くと，
$\lambda=1$ では $\{r,r'\}=\{1,q-1\}$，
$\lambda=0$ では $r'=-r$，
$\lambda=-1$ では $\{r,r'\}=\{-1,-(q-1)\}$ である．
各 node 数を区間 $[0,5]$ に入れるだけで表を得る．
例えば $\lambda=1$ で $r=q-1$ とすると $1-r=2-q<0$ なので，表示した順序だけが可能である．
最後に高分母は奇数で一より大きいので，$q\leq5$ は $q=3,5$ を意味する．
\end{proof}


% v13: the constant-pairing C4 obstruction no longer uses D opposition.
% The new tau=1 argument uses the fifth core and the unique S=T=2 model.

\subsection{四角形成分の同符号 pairing の排除}
\label{cq12:section}

\begin{lemma}[二つの reverse divisors が二である最小 entry の模型]
\label{cq13:reverse-two}
正の二高分母標準形で $S=T=2$ なら，canonical な整数データは
\[
 (r,s,z,S,T)=(4,4,5,2,2),\qquad (a,b)=(9,11)
\]
に限られる．
\end{lemma}

\begin{proof}
$d=S-T-1=-1$ なので，補題~\ref{ml12:dichotomy} より
$\alpha=-u$，$\beta=-v$，$H=r-2\geq1$ と書ける．ここで $u,v$ は正の整数である．
式~\eqref{ml12:E-equation} と式~\eqref{ml12:s-z} は
\[
 Euv=v-u+1,\qquad E=1-\frac4{H(H+1)},\qquad
 s=u\frac{H+2}{H},\quad z=v\frac{H+3}{H+1}
\]
を与える．$H=1$ なら $E=-1$ で，最初の式は $u(1-v)=v+1$ となる．
左辺は非正，右辺は正なので不可能である．
$H\geq2$ では $0<E<1$ である．$u=1$ なら $E=1$ となるため $u\geq2$ であり，
\[
 Eu=1-\frac{u-1}{v}<1
\]
となる．$H\geq3$ なら $E\geq2/3$ なので $Eu\geq4/3$ となって矛盾する．
従って $H=2$，$E=1/3$，$u=2$，$v=3$ である．
表示式から $r=4$，$s=4$，$z=5$，$(a,b)=(sS+1,zT+1)=(9,11)$ を得る．
このデータが標準形の各条件を満たすことは直接確認できる．
\end{proof}

\begin{theorem}[四角形成分に残る pairing]\label{cq12:no-constant-pairing}
\label{cq13:no-constant-pairing}
高分母 support の連結成分が $C_4$ であり，その vertices を
$i\in\mathbb Z/4$，残る第五の label を $*$ とする．
正の局所 charts，marked puncture の整数 clutching，
および全ての node 数の非負性を仮定する．
\[
 \varepsilon_i=\operatorname{sign}(p^i_{i+1,*})
\]
が全て同じであることは不可能である．
従って命題~\ref{cn12:cycle-signs} の二つの型のうち，残る可能性は
$\varepsilon_i=-\varepsilon_{i+1}$ という交互型だけである．
ここで node 数の上限五は，別の仮定ではなく両端の row sums と非負性から従う．
\end{theorem}

\begin{proof}
cycle の向きを逆にすると全 $\varepsilon_i$ の符号が反転するので，
全て $+1$ としてよい．$q_i=q_{i,i+1}>1$ とする．
core $i$ の高分母二 rays 間の entries を
\[
 p^i_{i-1,i+1}=-A_i,\qquad
 p^i_{i+1,i-1}=-B_i,
 \qquad A_iB_i=q_{i-1}q_i-1
\]
と書く．全ての $A_i,B_i$ は正である．

\paragraph{固定向きの低低 entry と reverse divisors．}
共通 puncture $*$ では $\lambda_{i,i+1}=0$ である．
その reverse numerators の絶対値を $c_i$ とすると，補題~\ref{cn12:reverse-bound} より
$1\leq c_i\leq4$ である．
core $i$ の低 ray $i*$ から二本の高 rays へ向かう entries は，順に
\[
 p^i_{*,i-1}=-c_{i-1},\qquad p^i_{*,i+1}=c_i.
\]
もう一つの低 ray $i,i+2$ から高 rays への entries を
\[
 p^i_{i+2,i-1}=\ell_i>0,\qquad
 p^i_{i+2,i+1}=-k_i<0
\]
と書く．固定向きの低低 entries は
\[
 d_i=p^i_{*,i+2}\leq0,\qquad e_i=p^i_{i+2,*}\leq0
\]
であり，ちょうど一方が零，他方が負である．
特に $d_i$ は canonical な非零 entry と同じとは限らない．
低 ray $i*$ の row sum から
\begin{equation}\label{cq12:fixed-row}
 c_i-c_{i-1}+d_i=-1,
 \qquad \sum_i d_i=-4
\end{equation}
を得る．

\paragraph{mixed node による二つの型の伝播．}
$ij=(i,i+1)$ の，第三 label $i+2$ に向かう mixed puncture の両端の outgoing entries は
\[
 -\frac{q_i-1}{k_i},\qquad -A_{i+1}
\]
であり，両者とも負である．従ってその node 数は二以上の整数である．
これを $1+\tau_i$ と書く．共通 puncture の node 数は一であり，もう一つの mixed puncture の node 数も一以上なので
\[
 \tau_i\in\{1,2\}.
\]
全数五を用いて，二つの mixed clutching 式は
\begin{equation}\label{cq12:tau-equations}
 A_{i+1}=\tau_iq_i-\frac{q_i-1}{k_i},\qquad
 B_i=(2-\tau_i)q_i+\frac{q_i-1}{\ell_{i+1}}
\end{equation}
となる．各表示分数は一以上 $q_i-1$ 以下の整数なので，
\[
 \begin{array}{c|cc}
  & A_{i+1} & B_i\\ \hline
 \tau_i=1 & <q_i & >q_i\\
 \tau_i=2 & >q_i & <q_i
 \end{array}
\]
である．一方 $A_iB_i=q_{i-1}q_i-1$ と整数性により，
$A_i<q_{i-1}$ と $B_i>q_i$ は同値であり，逆向きの厳密な不等式も同値である．
従って $\tau_{i-1}=\tau_i$ であり，四つの $\tau_i$ は同じ値となる．

\paragraph{全て $\tau_i=1$ の場合．}
$d_i<0$ の core を考える．このとき $e_i=0$ なので，二高分母の canonical 表示の低 ray $C$ は $i,i+2$，
低 ray $D$ は $i*$ である．$C$ の正の entry は $i,i-1$ に向かうので，
canonical な高 ray $A$ は $i,i-1$，$B$ は $i,i+1$ となる．
canonical 補助変数は
\[
 \alpha=A_i-q_{i-1}<0,\qquad
 \beta=q_i-B_i<0
\]
である．補題~\ref{ml12:dichotomy} より，その非零低低 entry $d_i$ は $-1$ である．
従って全ての $d_i$ は零または $-1$ であり，\eqref{cq12:fixed-row} から四つとも $-1$ となる．
同じ式から $c_i=c$ は一定である．この場合の canonical な低 rays は全て
$C=i+2$，$D=*$ であり，$S=T=c$ となる．
補題~\ref{cc12:proper-negative} の結論 $S\geq2$ により $c\geq2$ である．
以下では分母の下降や $D$ の相殺を使わず，第五の core $*$ の node 条件を用いる．

core $*$ における，ray $*i$ から ray $*j$ への entry を $v_{ij}=p^*_{ij}$ と書く．
これらの ray の分母は全て一なので
\[
 \sum_{j\ne i}v_{ij}=-1,\qquad v_{ij}v_{ji}=0.
\]
低 edge $i*$ の puncture $i+1$ では，core $i$ 側の entry は $c$ である．
従って node 非負性は
\[
 1-c-v_{i,i+1}\geq0,\qquad v_{i,i+1}\leq1-c<0
\]
を与える．添字を一つ戻して pair product を使えば $v_{i,i-1}=0$ であり，row sum から
\[
 v_{i,i+2}=-1-v_{i,i+1}\geq c-2
\]
となる．$c\geq3$ なら $v_{i,i+2}$ と $v_{i+2,i}$ がともに正となり，
その積が零であることに反する．従って $c=2$ である．

補題~\ref{cq13:reverse-two} により，各 core の canonical 高分母は $(a,b)=(9,11)$ である．
ここでは $A=i-1$，$B=i+1$ なので，すべての $i$ に対して
$q_{i-1}=9$，$q_i=11$ を要求する．同じ共有 edge の分母が両方の値を持つことになり矛盾する．

\paragraph{全て $\tau_i=2$ の場合．}
同じ標識の説明により，$d_i<0$ の core では canonical な
$\alpha=A_i-q_{i-1}$，$\beta=q_i-B_i$ がともに正である．
従って $d_i\leq-3$ である．
四つの和が $-4$ なので，ちょうど一つが $-4$，他の三つが零である．
添字を移して $d_0=-4$ とすれば，\eqref{cq12:fixed-row} と $1\leq c_i\leq4$ から
\[
 (c_0,c_1,c_2,c_3)=(4,3,2,1)
\]
となる．
core $3$ では $d_3=0$ なので，canonical な低 ray $C$ は $3*$ であり，
その正の entry は高 ray $3,0$ に向かう．従って canonical な高 ray $A$ は $3,0$ で，
標準形の整数 $r$ は $c_3=1$ である．
ところが $r=1$ の標準形では
\[
 A=a+\frac{a-1}{r}-s=a+(a-1)-s\geq a,
 \qquad s\mid a-1,\quad s>0.
\]
これは $\tau_3=2$ から得る $B_3<q_3$ に反する．
ここでは canonical な $A$ が $B_3$，$a$ が $q_3$ に対応する．
両場合が排除され，結論を得る．
\end{proof}

交互型の四角形成分については，本定理は不存在を結論しない．
同符号型の排除には node 非負性と各 branch の全 node 数五を使っており，
node 条件を外した整数 clutching だけの障害へ読み替えない．
一方で，この改訂された証明では対向端の $D$ の相殺を一切用いていない．

% v12: proper negative low entries and the complete C4 obstruction.
% The local lemma and the no-negative-common-puncture lemma use no node
% nonnegativity.  The C4/C5 global conclusions retain the stated hypotheses.

\subsection{負の低分母 entry の真約数性と四角形成分の完全排除}
\label{cc12:section}

小さい分母の個別分類によらず，二高分母の標準形そのものから
共通 puncture の負の clutching 係数を排除できる．
ここで低分母とは分母が一であることを意味する．

\begin{lemma}[負の低分母 entry は真約数である]
\label{cc12:proper-negative}
正の weights を持つ整数局所模型の分母が $(a,b,1,1)$，$a,b>1$，であるとする．
高分母 $q\in\{a,b\}$ の row から低分母へ向かう負の entry を $p$ と書けば，
\[
 1\leq -p\leq\frac{q-1}{2}.
\]
特に $p=-(q-1)$ は不可能である．
この結論は分母の奇数性，大域 clutching，$D$ の相殺，node 条件のいずれも必要としない．
\end{lemma}

\begin{proof}
正の二高分母標準形により，高分母を $A,B$，低分母を $C,D$ と標識すると
\[
 p_{AC}=x>0,\quad p_{AD}=-s<0,\quad
 p_{BC}=-y<0,\quad p_{BD}=z>0,
\]
\[
 x=\frac{a-1}{r},\quad y=\frac{b-1}{r+1},\quad
 a=sS+1,\quad b=zT+1,
\]
\[
 A=a+x-s>0,\qquad B=b-y+z>0,\qquad AB=ab-1
\]
と書ける．$r,s,z,S,T,x,y$ は正の整数である．
高分母 $B$ の負の entry は $-y$ であり，$r+1\geq2$ から
$y\leq(b-1)/2$ が直ちに従う．

高分母 $A$ の負の entry について $s=a-1$ と仮定する．
このとき $a-1=rx$，$s=rx$，$A=x+1$ である．
$r=1$ なら $A=a$ なので，整数 $B$ に対する $AB=ab-1$ が不可能となる．
以下 $r\geq2$ とする．

$B=(rb+1)/(r+1)+z$ を pair product に代入して整理すると
\[
 A(r+1)z=(1+rs)b-(A+r+1)
\]
を得る．さらに $s=rx$，$A=x+1$，$b=zT+1$ を代入すれば
\begin{equation}\label{cc12:maximal-negative-equation}
 z\{(1+r^2x)T-(x+1)(r+1)\}
 =-(r+1)\{(r-1)x-1\}.
\end{equation}
$r\geq3$ なら $T,x\geq1$ より左辺の括弧内は
\[
 (1+r^2x)T-(x+1)(r+1)
 \geq x(r^2-r-1)-r
 \geq r^2-2r-1>0,
\]
一方で右辺は負なので矛盾する．
$r=2$ かつ $x\geq2$ なら左辺の括弧内は $x-2\geq0$，
右辺は $-3(x-1)<0$ なので，やはり矛盾する．
残る $r=2,x=1$ では式~\eqref{cc12:maximal-negative-equation} は
$z(5T-6)=0$ となり，正の整数 $z,T$ に対して不可能である．

従って $s\ne a-1$ である．$s$ は $a-1$ の正約数なので，
真約数であることから $s\leq(a-1)/2$ を得る．
これで二つの高分母 row の両方が処理された．
\end{proof}

\begin{lemma}[共通の低 puncture は両端で負にならない]
\label{cc12:no-negative-lambda}
高 edge $ij$ の両端の core がともに二高分母型であり，
第三 label $k$ について $q_{ik}=q_{jk}=1$ とする．
正の局所 weights と marked puncture の整数 clutching を仮定すれば，
\[
 p^i_{jk}<0,\qquad p^j_{ik}<0
\]
は同時には成立しない．
従って共通 low-low puncture の outgoing sum を
$\lambda_{ij}q_{ij}$ と書くと $\lambda_{ij}\in\{0,1\}$ である．
この結果には $D$ の相殺も node 非負性も不要である．
\end{lemma}

\begin{proof}
$q=q_{ij}>1$ と置く．両端で負なら補題~\ref{cc12:proper-negative} から
\[
 -(q-1)\leq p^i_{jk}+p^j_{ik}\leq-2<0.
\]
一方，整数 clutching は $q\mid p^i_{jk}+p^j_{ik}$ を要求する．
この開区間 $(-q,0)$ は $q$ の倍数を含まないため矛盾する．
一般には各 outgoing entry は $q-1$ の非零 signed divisor である．
逆符号の場合には和の絶対値が $q$ 未満なので零となり，
両方正の場合には和は正で $2q$ 未満なので $q$ となる．
これで $\lambda_{ij}=0,1$ だけが残る．
\end{proof}

\begin{theorem}[四角形成分の完全排除]
\label{cc12:no-C4}
五成分 plane degeneration の正の局所 charts に対し，
marked puncture の整数 clutching，および全ての node 数の非負性を仮定する．
対向端の $D$ の相殺は仮定しない．
このとき高分母 support の連結成分は $C_4$ ではあり得ない．
各 branch の node 全数五は row sums から従うものであり，追加の上界として仮定しない．
\end{theorem}

\begin{proof}
$C_4$ の vertices を $i\in\mathbb Z/4$，残る label を $*$ とし，
\[
 \varepsilon_i=\operatorname{sign}(p^i_{i+1,*})\in\{1,-1\}
\]
と置く．局所的な正の pairing により
$p^{i+1}_{i,*}$ の符号は $-\varepsilon_{i+1}$ である．
従って補題~\ref{cc12:no-negative-lambda} は
$\varepsilon_i=-1,\varepsilon_{i+1}=1$ という隣接 pair を禁止する．
すなわち全ての $i$ で $\varepsilon_i\geq\varepsilon_{i+1}$ である．
cycle を一周すると全ての $\varepsilon_i$ が同じになる．
これは定理~\ref{cq12:no-constant-pairing} に反する．
ここで符号の一様性そのものは node 条件を使わずに従い，
node 非負性は同符号型を排除する同定理で用いている．
\end{proof}

\begin{corollary}[五角形成分に残る二つの pairing 型]
\label{cc12:C5-two-types}
正の局所 charts，marked puncture の整数 clutching，node 非負性の下で，
高分母 support の連結成分が $C_5$ であるとする．
vertices を $i\in\mathbb Z/5$ とし，
\[
 \varepsilon_i=\operatorname{sign}(p^i_{i+1,i+2})
\]
と置く．このとき $\varepsilon_i=1$ となる vertex は高々一つである．
共通 low-low puncture の係数 $\lambda_{i,i+1}$ の個数は
\[
 \bigl(\#\{\lambda=1\},\#\{\lambda=0\},\#\{\lambda=-1\}\bigr)
 \in\{(5,0,0),(3,2,0)\}.
\]
この二つの型は，後の定理~\ref{cf13:no-C5} により共に排除される．
\end{corollary}

\begin{proof}
高 edge $i,i+1$ の共通 low-low puncture の二つの符号は
$-\varepsilon_i,-\varepsilon_{i+1}$ である．
補題~\ref{cc12:no-negative-lambda} より正の vertices は隣接できない．
また node 非負性により距離二の正の pair も存在しない
（命題~\ref{cn12:cycle-signs}）．
五角形の異なる二 vertices の距離は一または二なので，正の vertex は高々一つとなる．
公式 $\lambda_{i,i+1}=-(\varepsilon_i+\varepsilon_{i+1})/2$ から，
正の vertex が零なら全ての $\lambda$ が一，
正の vertex が一つなら隣接する二つの $\lambda$ が零で他の三つが一となる．
\end{proof}

\begin{lemma}[一方の正の entry が一である標準形の最大分母]
\label{cc12:unit-positive-max}
正の二高分母標準形で $r=1,z=1$ とする．
このとき局所的に最大の分母を持つ高 ray の係数 $D$ は負である．
特に $a=b$ なら $D_A,D_B$ はともに負である．
\end{lemma}

\begin{proof}
標準形の pair product と係数の定義から
\[
 (s+1)(b+3)=6a+2,\qquad S=\frac{a-1}{s},
\]
\[
 D_A=\frac{b+1}{2}-a+S,\qquad
 D_B=2(a-b)+2-s
\]
を得る．$b>a$ なら整数性により $D_B\leq-s<0$ である．
$a\geq b$ の場合，$s=1$ なら $b=3a-2>a$ となり不可能であり，
$s=2$ なら $3\mid6a+2$ を要求するので，やはり不可能である．
従って $s\geq3$ であり，
\[
 D_A\leq\frac{a+1}{2}-a+\frac{a-1}{s}
 =-(a-1)\left(\frac12-\frac1s\right)<0.
\]
さらに $a=b$ なら $D_B=2-s<0$ でもある．
この補題に大域 clutching や node 条件は不要である．
\end{proof}

\begin{theorem}[全ての共通 punctures が正である五角形型の排除]
\label{cc12:no-all-positive-common}
五成分 plane degeneration の正の局所 charts，
marked puncture の整数 clutching，対向端の $D$ の相殺を仮定する．
高分母 support の連結成分が $C_5$ であるとき，
全ての共通 low-low punctures で $\lambda=1$ となることは不可能である．
この排除に node 非負性は必要ない．
\end{theorem}

\begin{proof}
各高 edge の分母は奇数 $q\geq3$ である．
$\lambda=1$ なら両端の共通 puncture への entries はともに正であり，
その和は $q$ である．それぞれ $q-1$ の正約数なので，二つの値は
$\{1,q-1\}$ である．実際，二つとも真約数なら和は $q-1$ 以下となり，
両方とも $q-1$ なら $q\geq3$ と矛盾する．

二高分母標準形の $B$ 側の正の低分母 entry $z$ は $b-1$ になれない．
もし $z=b-1$ なら $T=1$ であり，補題~\ref{td12:low-entry} の $S\leq T$ から $S=1$，
従って $s=a-1$ となるが，これは補題~\ref{cc12:proper-negative} に反する．
一方，$A$ 側の正の低分母 entry $x=(a-1)/r$ が $a-1$ となることは
$r=1$ と同値である．従って各 core には，対応する高分母を $q$ として
値 $q-1$ を持つ正の低分母 entry が高々一つしかない．

全ての $\lambda=1$ のとき，五本の高 edges はそれぞれちょうど一端で値 $q-1$ を持つ．
各高 row の共通 puncture への entry はその row の唯一の正の低分母 entry なので，
この計五個の値は五つの core に高々一個ずつ配られる．
従って各 core にちょうど一個あり，それは canonical な $A$ 側である．
もう一方の $B$ 側の正の entry は一となり，すべての core で $r=1,z=1$ を得る．

五本の高 edges のうち最大の分母を持つものを選ぶ．
その分母は両端の各 core でも局所的に最大なので，
補題~\ref{cc12:unit-positive-max} により対向端の二つの係数 $D$ がともに負となる．
これはそれらの相殺に反する．
\end{proof}


% v13: complete C5 obstruction.  No D-opposition hypothesis is used.
% Uses standard positive two-high charts, integer marked clutching,
% odd high denominators, and nonnegative node counts.

\subsection{五角形成分の完全排除：三本の共通 puncture と低分母 edge}
\label{cf13:section}

前節で残った五角形の符号型は，例外 core の分母を列挙することなく排除できる．
要点は，共通 puncture の三本の連続する正の接着係数が，離れた低分母 edge の
node 数を負にすることである．この議論では対向端の $D$ の相殺を使わない．

\begin{lemma}[正の最大約数と canonical な高 ray]
\label{cf13:full-entry}
正の二高分母標準形で，高 ray の分母を $q$ とする．
その ray から低 ray へ向かう正の entry が $q-1$ なら，その高 ray は
canonical な $A$ であり，$r=1$ である．
各 core にはこのような高 ray が高々一つしかない．
また，二つの高 rays の正の低分母 entries がともに一となることはない．
\end{lemma}

\begin{proof}
canonical な正の entries は $x=(a-1)/r$ と $z=(b-1)/T$ である．
補題~\ref{td12:low-entry} と補題~\ref{cc12:proper-negative} から
$T\geq S\geq2$ なので，$B$ 側の $z$ は $b-1$ ではない．
$A$ 側では $x=a-1$ と $r=1$ が同値である．
さらに $x=z=1$ なら canonical 補助変数が
$\alpha=1-s\leq0$，$\beta=y-1\geq0$ となり，
両者が非零で同符号であることに反する．
\end{proof}

\begin{lemma}[五角形で連続する四つの負の pairing は不可能]
\label{cf13:four-negative}
高分母 support が $C_5$ であり，vertices を $i\in\mathbb Z/5$ と標識する．
正の局所 charts，marked puncture の整数 clutching，node 数の非負性を仮定する．
$q_i=q_{i,i+1}$，$\varepsilon_i=\operatorname{sign}(p^i_{i+1,i+2})$ と置く．
このとき
\[
 \varepsilon_1=\varepsilon_2=\varepsilon_3=\varepsilon_4=-1
\]
は不可能である．$\varepsilon_0$ の値や対向端の $D$ の相殺は仮定しない．
\end{lemma}

\begin{proof}
三価 branch の飽和条件から各 $q_i$ は奇数で，$q_i\geq3$ である．
仮定した符号から，共通 low-low punctures の係数は
$\lambda_1=\lambda_2=\lambda_3=1$ となる．
これらの高 edge の両端で，共通 puncture への正の entries は
それぞれ $q_i-1$ の約数で，その和は $q_i$ である．
従って二つの値は $\{1,q_i-1\}$ である．
実際，両方が真約数なら和は $q_i-1$ 以下となり，両方が $q_i-1$ なら
$q_i\geq3$ に反する．以下，値 $q_i-1$ を持つ端をその edge の最大約数端と呼ぶ．

core $2,3$ では両方の高 rays がこの三辺のいずれかに属する．
その二つの正の低分母 entries は一または最大約数である．
補題~\ref{cf13:full-entry} より，最大約数端は各 core にちょうど一つあり，
それが canonical な $A$ に対応する．
辺 $23$ の最大約数端は一方にしかないため，core $2,3$ の $A$ は
ともに前の高 ray であるか，ともに次の高 ray であるかの二通りとなる．
以下の表では $A,B,C,D$ を，その ray のもう一方の vertex label で表す．
$C$ から $D$ への canonical entry は常に零である．

\paragraph{両方が前の高 ray である場合．}
辺 $12$ の最大約数端は core $2$，辺 $23$ では core $3$ である．
従って辺 $34$ では core $3$ の正の entry は一，core $4$ が最大約数端となる．
$\varepsilon_4=-1$ なので，使用する三つの canonical 標識は
\[
\begin{array}{c|rrrr}
\text{core}&A&B&C&D\\ \hline
2&1&3&4&0\\
3&2&4&0&1\\
4&3&0&1&2
\end{array}
\]
である．特に $p^4_{12}=0$ となる．
一方，辺 $12$ の core $1$ 側の共通 puncture $4$ への entry は一なので，
\[
 p^1_{24}=1,\qquad p^1_{42}=q_1-1
\]
である．後式は pair product $p^1_{24}p^1_{42}=q_{12}q_{14}-1$ から従う．
低分母 edge $14$ の puncture $2$ では $q_{14}=1$ だから，node 式より
\begin{equation}\label{cf13:negative-node-forward}
 n_{14;2}=1-p^1_{42}-p^4_{12}=2-q_1<0.
\end{equation}
これは node 非負性に反する．

\paragraph{両方が次の高 ray である場合．}
辺 $12$ の最大約数端は core $1$，辺 $23$ では core $2$，辺 $34$ では core $3$ となる．
canonical 標識は
\[
\begin{array}{c|rrrr}
\text{core}&A&B&C&D\\ \hline
1&2&0&4&3\\
2&3&1&0&4\\
3&4&2&1&0
\end{array}
\]
である．従って $p^1_{43}=0$ であり，辺 $34$ の core $4$ 側は一の端なので
$p^4_{31}=1$，$p^4_{13}=q_3-1$ となる．
同じ低分母 edge $14$ の puncture $3$ で
\begin{equation}\label{cf13:negative-node-backward}
 n_{14;3}=1-p^1_{43}-p^4_{13}=2-q_3<0
\end{equation}
となり，やはり矛盾する．これで両場合が排除された．
\end{proof}

\begin{theorem}[五角形成分の完全排除]
\label{cf13:no-C5}
五成分 plane degeneration の正の局所 charts に対し，marked puncture の整数 clutching
および全ての node 数の非負性を仮定する．
このとき高分母 support の連結成分は $C_5$ ではあり得ない．
この結果は分母の上界も，対向端の $D$ の相殺も必要としない．
\end{theorem}

\begin{proof}
系~\ref{cc12:C5-two-types} の証明は $D$ の相殺を使わず，
正の pairing の vertex が高々一つであることを示している．
正の vertex があるならそれを $0$ と標識し，ないなら任意の vertex を $0$ とする．
残る四つの $\varepsilon_i$ は全て $-1$ なので，補題~\ref{cf13:four-negative} に反する．
\end{proof}

\begin{corollary}[五次の次数二以下の分母剛性]
\label{cf13:degree-two-rigidity}
正の局所 charts，marked puncture の整数 clutching，対向端の $D$ の相殺，
node 非負性を仮定する．高分母 support $H$ の最大次数が二以下なら $H$ は空であり，
全ての branch denominators は一である．
同値に，高分母 edge が存在すれば，少なくとも一つの core は三本以上の高分母 rays を持つ．
\end{corollary}

\begin{proof}
最大次数二以下の辺を持つ連結成分は道または閉路である．
道成分，三角形成分，四角形成分は既に排除されており，
五頂点の support に残る五角形成分は定理~\ref{cf13:no-C5} により排除される．
\end{proof}

五角形の新しい障害では，三本の高 edges に沿う接着から低 edge $14$ 上の
個別の node 数を計算した．node 全数五や各 node 数の上限五を別に用いてはいない．
また，これは標準 $K_4$ cones と現在の branch graph を保った charts に対する障害であり，
core surgery や追加の branch singularities を許した全ての fibrations の分類ではない．

% v14. A node obstruction for a triangle with two degree-two vertices.
% No opposing-D cancellation and no denominator cutoff are used.

\section{二つの二価頂点を持つ高分母三角形の障害}
\label{tn14:section}

三角形成分の排除は，三つの core が全て二高分母である場合の結果であった．
ここでは，そのうち二つだけが二高分母である場合を扱う．
残る頂点には追加の高分母 branches が接続していてよい．
個別の node 条件を用いる代わりに，対向する $D$ の相殺は必要としない．

二高分母の canonical 表示で，高 rays を $A,B$，低 rays を $C,D$ とする．
高 ray へ向かう二つの低 rays からの reverse numerators は
\[
 \begin{array}{c|cc}
  &C&D\\\hline
  A&r&-S\\
  B&-(r+1)&T
 \end{array}
\]
である．既知の正値性と整数条件により $r\geq1$，$S\geq2$，$T\geq S$ であり，
$r=S$ は起こらない．補題~\ref{ml12:dichotomy} により，
$S>r$ なら $rT>(r+1)S$，$S<r$ なら $T=S$ である．

\begin{lemma}[小さい reverse numerators と負の枝]
\label{tn14:small-negative-branch}
正の二高分母整数標準形で $2\leq S<r\leq4$ とする．
このとき唯一の可能性は
\[
 (r,s,z,S,T)=(4,4,5,2,2),\qquad (a,b)=(9,11)
\]
である．
\end{lemma}

\begin{proof}
$H=r-S$，$\alpha=-u$，$\beta=-v$ と置くと $u,v$ は正整数であり，
式~\eqref{ml12:E-equation} は
\begin{equation}\label{tn14:negative-equation}
 \left(1-\frac{S^2}{H(H+1)}\right)uv=v-u+1
\end{equation}
を与える．候補 $(r,S)$ は $(3,2),(4,2),(4,3)$ の三つである．
$(3,2)$ では $u(1-v)=v+1$，$(4,3)$ では
$u(1-7v/2)=v+1$ となるが，いずれも左辺は非正，右辺は正なので不可能である．

$(4,2)$ では $T=S=2$，$u=s/2$，$v=3z/5$ である．
残る numerator の整数性 $4\mid2s$，$5\mid2z$ から
$s=2u$，$z=5\ell$，$v=3\ell$，$\ell\geq1$ と書ける．
式~\eqref{tn14:negative-equation} は
\[
 u\ell=3\ell-u+1,
 \qquad (3-u)(\ell+1)=2
\]
となる．$\ell+1\geq2$ なので $\ell=1$，$u=2$，従って $s=4,z=5$ である．
分母は $a=sS+1=9$，$b=zT+1=11$ となる．
このデータは元の row sums，pair products を満たし，標準形の正の weights を持つ．
\end{proof}

\begin{lemma}[両 reverse numerators が四以下の高 ray]
\label{tn14:small-reverse-ray}
正の二高分母整数標準形で，一つの高 ray に向かう二つの低 rays からの
reverse numerators の絶対値がともに四以下であるとする．
その高 ray は必ず canonical な $A$ である．
\end{lemma}

\begin{proof}
$B$ だとすると $r+1\leq4$，$T\leq4$ であり，従って $r\leq3$，$S\leq4$ である．
$S>r$ なら $S\geq\max(2,r+1)$ と $rT>(r+1)S$ により，
$r=1,2,3$ のどの場合も $T\geq5$ となって矛盾する．
$S<r$ なら $2\leq S<r\leq3$ なので $(r,S)=(3,2)$ だが，
補題~\ref{tn14:small-negative-branch} によって不可能である．
\end{proof}

\begin{theorem}[二つの二価頂点を持つ三角形の不存在]
\label{tn14:no-degree-two-triangle}
五次の標準 $K_5$ branch graph 上で，正の局所整数 charts，marked puncture の整数 clutching，
および全ての node 数の非負性を仮定する．
高分母 support $H$ は，三角形のうち二つの頂点が $H$ において次数二であるような部分 graph を含まない．
残る三角形頂点の次数や，その先の接続は制限しない．
分母の上限と対向端の $D$ の相殺は，いずれも仮定しない．
\end{theorem}

\begin{proof}
次数二の頂点を $i,j$，両者のもう一つの高隣接頂点を $k$ とし，
残る二つの labels を $\ell,m$ とする．
辺 $ij$ の分母を $q>1$ と書く．両端の core は二高分母であり，
$\ell,m$ は両端で共通の低分母 punctures である．
共有高 ray の二つの低分母への outgoing entries は，各 core で一正一負である．

二つの punctures で同符号同士が対応する場合，一方では負の entries が二つ接着する．
補題~\ref{ls11:low-puncture} はその組を $\{-1,-(q-1)\}$ に限るが，
補題~\ref{cc12:proper-negative} は二高分母の負の高低 entry が
$-(q-1)$ であることを禁じる．従って，両 punctures とも異符号同士が対応する．
それぞれの outgoing entries，および reverse entries の和は零である．

各共通低 label に対して補題~\ref{cn12:reverse-bound} を適用すると，
両端の共有高 ray における二つの reverse numerators の絶対値はともに四以下となる．
この補題の第三 core の二つの rays は低分母なので pair product が零であり，
その少なくとも一方の entry が零である．対応する低 edge の node 数が $1\pm c$ となり，
row sums から従う区間 $[0,5]$ に入るため $|c|\leq4$ となる．
第三 core の他の分母は，この議論に影響しない．

補題~\ref{tn14:small-reverse-ray} により，共有 ray は両 core で canonical な $A$ である．
一方の canonical reverse pair を $(r,-S)$ とすると，他方では符号が反転するため
\[
 (r',S')=(S,r).
\]
従って $2\leq r,S\leq4$ であり，$r\ne S$ なので一方の core は負の枝である．
補題~\ref{tn14:small-negative-branch} により，その core の shared denominator は $q=9$，
reverse pair は $(4,-2)$ となる．他方の core では
\[
 a=q=9,\qquad r=2,\qquad S=4,\qquad s=2
\]
でなければならない．

このデータを二高分母の pair-product 方程式に代入すると
\begin{equation}\label{tn14:terminal-divisor}
 z(5T-33)=9
\end{equation}
を得る．正の枝の不等式は $T>6$ を与えるので $5T-33>0$ であり，
$z$ は九の正約数である．しかし $z=1,3,9$ はそれぞれ
\[
 T=\frac{42}{5},\quad \frac{36}{5},\quad \frac{34}{5}
\]
を与え，いずれも整数ではない．これが最後の矛盾である．
\end{proof}

この定理は，三角形成分に加え，三角形に一枚の葉を付けた形，長さ二の枝を付けた形，
同じ三角形頂点に二枚の葉を付けた形，および二つの三角形が一頂点を共有する形を排除する．
一方，三角形の異なる二頂点に葉を付けた場合や diamond では，
対象となる二価頂点の pair が存在しないので，この定理だけでは排除されない．
ここで排除したのは明示された次数条件を持つ部分 graph であり，
任意の三角形部分 graph の不存在を主張していない．

% v14: C4 with one pendant edge.  The finite coefficient exclusions are
% complete divisor computations, using the proved fixed-denominator theorem.
\section{四角形に一本の葉辺を付けた高分母成分の排除}
\label{cp14:section}

次に，高分母 support が四角形と一本の葉辺からなる場合を扱う．
三つの二価 core を持つこの場合には，中央の二価 core から葉に向かう
puncture の符号を用いて，全ての整数条件を有限計算と一つの無限族へ帰着できる．
以下の有限計算は，分母の上限による探索ではない．
固定分母の完全な約数表示を用い，検査した全約数と棄却理由を
\texttt{verify\_c4\_pendant\_v14.py} の結果に記録する．

\begin{lemma}[小さい負の枝の完全分類と接続障害]
\label{cp14:small-negative}
正の二高分母標準形で $S=T=w\leq4$ とする．
このとき canonical データは次の七つだけである．
\[
\begin{array}{rrrr|rr|rr}
r&s&z&w&a&b&D_A&D_B\\\hline
4&4&5&2&9&11&3&-1\\
5&10&2&3&31&7&-25&23\\
6&6&14&3&19&43&29&-23\\
7&7&2&4&29&9&-22&21\\
7&21&4&4&85&17&-69&63\\
8&6&9&4&25&37&13&-10\\
8&8&27&4&33&109&87&-75
\end{array}
\]
いずれの模型でも，どちらの高 ray についても，その $D$ と反対の係数を持つ
同分母の正の二高分母模型は存在しない．従って，これらの高 ray は
対向 $D$ の相殺を満たしながら別の二高分母 core に接続できない．
\end{lemma}

\begin{proof}
$S\geq2$ である．$H=r-w$，$\alpha=-u$，$\beta=-v$ と置けば
$H,u,v\geq1$ であり，既知の負の枝の式は
\[
 Euv=v-u+1,\qquad E=1-\frac{w^2}{H(H+1)}
\]
となる．$H^2<H(H+1)<(H+1)^2$ より $E\ne0$ である．また $u=1$ は $E=1$ を要求するので不可能である．$E\leq-1$ なら
$u(1+Ev)=v+1$ の左辺が非正となるため不可能である．
$E>0$ のとき $u=1$ は $E=1$ を要求するので $u\geq2$ であり，
$Eu=1-(u-1)/v<1$ より $E<1/2$ を得る．従って全ての場合に
\begin{equation}\label{cp14:H-bound}
 \frac{w^2}{2}<H(H+1)<2w^2
\end{equation}
となる．$w=2,3,4$ に対する $(w,H)$ は
\[
 (2,2),\quad(3,2),(3,3),\quad(4,3),(4,4),(4,5)
\]
だけである．

$E=M/N$ を既約表示し，$N>0$ とする．上の方程式は
\begin{equation}\label{cp14:negative-factor}
 (N-Mu)(N+Mv)=N(N-M)
\end{equation}
と同値である．実際，$N-Mu=N(u-1)/v>0$，$N+Mv=N(v+1)/u>0$ なので，両因子は正の整数である．その全ての正約数の組から
\[
 u=\frac{N-d}{M},\qquad
 v=\frac{N(N-M)/d-N}{M}
\]
を復元し，$u,v$ の正の整数性と
\[
 s=\frac{u(H+w)}H,\qquad
 z=\frac{v(H+w+1)}{H+1}
\]
の整数性を課すと，表の七模型を得る．これは六つの有限な約数計算で尽くされ，
全ての採用模型は元の行和・pair products と正の標準 weights を満たす．

接続障害には，固定分母の完全な factorization を適用する．
表の各 ray に対し，以下の反対係数が存在しないことを得る．
\[
\begin{array}{r|l@{\qquad}r|l}
q&\text{存在しない }D&q&\text{存在しない }D\\\hline
7&-23&9&-3,-21\\
11&1&17&-63\\
19&-29&25&-13\\
29&22&31&25\\
33&-87&37&10\\
43&23&85&69\\
109&75&&
\end{array}
\]
具体的には，固定 ray が $A$ の場合，$r,s\mid q-1$ と $x=(q-1)/r$ に対して
\[
 z\{(r+1)(q+x-s)-(rs+1)T\}=(r+1)(s-x-1)
\]
の非零右辺の全正約数 $z$ を調べる．固定 ray が $B$ の場合には，
$r+1,z\mid q-1$，$y=(q-1)/(r+1)$ に対し
\[
 s\{((r+1)z+1)S-r(q-y+z)\}=r(y-z-1)
\]
の全正約数 $s$ を調べる．既存の固定分母定理により，これらは非共鳴の場合を全て尽くす．
上表の分母のうち共鳴を持つのは七だけであり，その係数 $7k-9$，$k\geq1$，は
$-23$ にはならない．計算では全候補の分母，整数性，行和，pair products，係数を再検査した．
従って上表は相手分母の数値上限なしで成立する．
\end{proof}

\begin{lemma}[中央 core に現れる二つの小さい表示]
\label{cp14:small-middle}
正の二高分母標準形について，次が成立する．
\begin{enumerate}
\item $r=1$ かつ $L=T-S+1\leq5$ なら，
$(a,b)=(3,11),(5,31)$，または
\[
 (r,s,z,S,T)=(1,2t-1,t,2,6),\quad
 (a,b)=(4t-1,6t+1),\quad (D_A,D_B)=(3,-7),\quad t\geq1
\]
である．
\item $z=1$，$S>r$，$S\leq4$ なら，次の五模型，または固定七の共鳴族だけが残る．
\[
\begin{array}{rr|rrrrr}
a&b&r&s&z&S&T\\\hline
3&7&1&1&1&2&6\\
13&13&1&4&1&3&12\\
5&13&1&1&1&4&12\\
13&17&1&3&1&4&16\\
29&19&1&7&1&4&18
\end{array}
\]
この第二のリストでは，少なくとも一方の高 ray が，対向 $D$ の相殺を保って
別の二高分母 core に接続することを妨げられる．
\end{enumerate}
\end{lemma}

\begin{proof}
第一の場合，$S\geq2>r$ なので $T>2S$ である．$T-S+1\leq5$ と合わせると
$(S,T)=(2,5),(2,6),(3,7)$ に限られる．
最初の場合は pair-product 式が $(3-s)(z+2)=8$ となり，
$(s,z)=(1,2),(2,6)$ を得る．第二の場合は $s=2z-1$ となる．
最後の場合は $z(3s-5)+4s+2=0$ であり，$s=1,z=3$ 以外は不可能だが，
この組は $2\mid zT$ を満たさない．各係数は標準式への代入で得る．

第二の場合には $r<S\leq4$ である．$r=1$ なら
\[
 T=6S-4-\frac{6S-8}{s+1},\qquad 2\mid T
\]
となり，$S=2,3,4$ ごとの全正約数を調べると表の五模型を得る．
$r=2,S=3$ なら $T=6$，$s=2k$ となって固定七の共鳴族である．
$r=2,S=4$ では $T=9-3/(2s+1)$ なので，整数になる唯一の値 $s=1,T=8$ は
$3\mid T$ に反する．$r=3,S=4$ では $s=3k$，$k\geq1$ と書け，
\[
 T=6+\frac{2k+2}{9k+1}
\]
は六と七の間にあるので整数ではない．

表の第一・第三模型は分母三・五の係数集合から排除される．
第二模型の $a=13,D_A=-3$ に対して反対係数三は存在せず，
第四模型の $b=17,D_B=-9$ に対して反対係数九は存在せず，
第五模型の $a=29,D_A=-15$ に対して反対係数十五は存在しない．
これらも前補題と同じ完全な約数計算による．十三の共鳴係数 $13k-14$ は三にならない．
固定七の共鳴族の接続は，既存の小分母係数伝播の結果
\ref{sp13:no-resonant-continuation} によって妨げられる．
\end{proof}

\begin{theorem}[四角形と一本の葉辺からなる成分の不存在]
\label{cp14:no-pendant-square}
五次の標準 $K_5$ branch graph 上で，正の局所整数 charts，marked puncture の整数 clutching，
対向 $D$ の相殺，および全ての node 数の非負性を仮定する．
高分母 support は四角形に一本の葉辺を付けた graph ではあり得ない．
分母の数値上限は仮定しない．
\end{theorem}

\begin{proof}
四角形を $0123$，葉の辺を $04$ と標識する．core $1,2,3$ は二高分母，
core $4$ は一高分母，core $0$ は三高分母である．
中央の core $2$ の canonical 高 rays を $A=1,B=3$ とする．
必要なら四角形を反転してこの標識を選べる．低 ray $4$ は $C,D$ のいずれかである．
辺 $12,23$ の共通低 puncture $4$ の接着係数を $\lambda_A,\lambda_B$ と書く．
負の高低 entry の真約数性から，それぞれ零または一である．
core $2$ の二つの高 ray から低 ray $4$ への entries は異符号なので，
両方が一となることはない．以下 $V_{ij}=p^4_{ij}$ と置く．

\paragraph{低 ray $4$ が canonical な $D$ の場合．}
両 $\lambda$ が零なら，共通 puncture の reverse bound により $S,T\leq4$ である．
正の枝では $rT>(r+1)S$ と $S>r$ から $T\geq5$ なので不可能であり，
負の枝 $S=T\leq4$ は補題~\ref{cp14:small-negative} によって接続できない．

一方が一なら，正の entry $z$ を持つ $B$ 側の $\lambda_B=1$ であり，
$\lambda_A=0$ である．canonical $B$ は最大約数 $b-1$ の正 entry を持てないので，
$z=1$ となる．$A$ 側の reverse bound は $S\leq4$ を与える．
$y$ は正整数で $\beta=y-z\ne0$ だから $\beta>0$，従って $S>r$ である．
補題~\ref{cp14:small-middle} の第二の場合がこの接続を排除する．

\paragraph{低 ray $4$ が canonical な $C$ の場合．}
$B$ 側の outgoing entry $-y$ は負なので，$\lambda_B=0$ である．
従って core $3$ の reverse entry は $p^3_{42}=r+1$ となる．
低 edge $34$ の node 非負性は $V_{32}\leq-r<0$ を与え，pair product より $V_{23}=0$ である．
一方，低 edge $24$ の punctures $1,0$ における node 非負性は，
中央 core の $p^2_{41}=r$，$p^2_{40}=0$ を用いて
\[
 V_{21}\leq1-r\leq0,\qquad V_{20}\leq1
\]
を与える．葉の第二 row は $V_{20}+V_{21}=-1$ である．
もし $V_{21}<0$ なら $V_{20}\geq0$，しかも高分母との pair product から $V_{20}\ne0$ である．
しかし一高分母の葉の正の reverse entry は $t(h+1)\geq2$ なので，$V_{20}\leq1$ に反する．
従って
\[
 V_{21}=V_{23}=0,\qquad V_{20}=-1,\qquad r=1
\]
である．この議論は $\lambda_A$ の値を指定しない．

$q_{04}=m+1$ とすれば pair product は $V_{02}=-m$ を与える．
辺 $04$ の puncture $2$ の clutching により
\[
 p^0_{42}\in\{m,-1\},\qquad R:=p^0_{24}\in\{1,-m\}.
\]
中央 core の低低 entry を $p^2_{04}=-L$，$L=T-S+1$ と置く．
$r=1$ では正の枝なので $T>2S$，特に $L\geq4$ である．
低 edge $02$ の node 式は
\begin{equation}\label{cp14:central-low-node}
 n_{02;4}=1-R+L\in[0,5]
\end{equation}
となる．$R=-m$ は $m\geq2$ により不可能なので，$R=1$，$L\leq5$ を得る．

\paragraph{残った無限族の最後の障害．}
以上より，残る場合には常に中央 core で $r=1,L\leq5$ である．
補題~\ref{cp14:small-middle} の第一の場合の二つの有限模型は，
分母三・五の二価接続障害によって排除される．従って
\[
 (a,b)=(4t-1,6t+1),\qquad
 (r,s,z,S,T)=(1,2t-1,t,2,6),\qquad (D_A,D_B)=(3,-7)
\]
だけを調べればよい．$B$ 側の $\lambda_B=0$ なので，core $3$ の共有 ray $2$ は
低 ray $4$ へ正の entry $3t$ と reverse entry 二を持つ．
この共有 ray が canonical $B$ なら $T'=2$ となり，
補題~\ref{cp14:small-negative} に反する．従って canonical $A$ で $r'=2$ である．
そのもう一つの低 ray $1$ への entry を $-s'$ とする．

辺 $23$ の puncture $1$ の outgoing sum は $-(4t+1+s')$ である．
mixed puncture では負の高高 entry と $b-1$ の signed divisor を加えるため，その和は $b$ より小さい．整数 clutching からその和は非正であり，二つの mixed node 数はともに一以上である．従って
\[
 \frac{4t+1+s'}{6t+1}\in\{1,2\}.
\]
値二は $s'=8t+1>b-1$ を与えて整除条件に反するため，
$s'=2t$，$S'=3$ である．対向 $D$ の相殺は $D'_A=7$ を要求し，
\[
 A'=b+3t-2t=7t+1,\qquad B'=b+6=6t+7
\]
となる．core $3$ のもう一つの高分母を $c$ とすれば
\[
 c=\frac{A'B'+1}{b}=7t+8,
 \qquad z'=B'-c+\frac{c-1}{3}=\frac{4(t+1)}3.
\]
従って
\[
 T'=\frac{c-1}{z'}=\frac{21}{4}
\]
となり，整数性に反する．これで全ての場合が排除された．
\end{proof}

この結果は，node 非負性と対向 $D$ の相殺をともに用いる．
三角形に葉を異なる二頂点で付けた形，diamond やその葉付きの形，
その他のより密な support まで排除したとは主張しない．

% v15: a negative two-high branch cannot face a low edge; K_{2,3} follows.
% No denominator cutoff, oddness argument, or opposing-D cancellation is used.
\section{低辺を挟む二高分母の障害と \texorpdfstring{$K_{2,3}$}{K2,3} の排除}
\label{kb15:section}

二高分母 core の二つの高分母近傍が互いに低辺で結ばれるとき，
その低辺の一つの node 条件が，元の core の負の補助変数の枝を排除する．
この局所的な障害では，二つの近傍 core の高分母次数を制限しない．
その結果，三つの二価 core を持つ $K_{2,3}$ support を排除できる．

\begin{lemma}[二つの高分母近傍の間が低辺である場合]
\label{kb15:no-negative-low-wedge}
五次の標準 $K_5$ branch graph 上で，正の局所整数 charts，marked puncture の
整数 clutching，および node 数の非負性を仮定する．
高分母 support の二価頂点 $i$ の高分母近傍を $j,k$ とし，
$q_{jk}=1$ と仮定する．このとき core $i$ の二高分母標準形は正の枝に属する．
すなわち，canonical な記号では
\[
 S>r,\qquad rT>(r+1)S,\qquad T-S+1\geq3
\]
であり，負の枝 $S=T<r$ は不可能である．
この結論に対向 $D$ の相殺と分母の数値上限は不要であり，
core $j,k$ の高分母次数にも仮定を置かない．
\end{lemma}

\begin{proof}
補題~\ref{ml12:dichotomy} の負の枝が存在すると仮定する．
core $i$ の canonical 高 rays を $A=j,B=k$ とし，
その分母を $a=q_{ij}>1,b=q_{ik}>1$ と書く．負の枝の表示を
\[
 S=T=w,\quad H=r-w>0,\quad
 \alpha=x-s=-c<0,\quad \beta=y-z=-e<0
\]
とする．ここで $r,s,z,w,x,y,c,e$ は正整数で，
\[
 a=sw+1,\quad b=zw+1,\quad
 x=\frac{sw}{r},\quad y=\frac{zw}{r+1},\quad
 c=\frac{sH}{r},\quad e=\frac{z(H+1)}{r+1}.
\]
負の枝の pair-product 式~\eqref{ml12:E-equation} は
\begin{equation}\label{kb15:E}
 Ece=e-c+1,\qquad E=1-\frac{w^2}{H(H+1)}
\end{equation}
である．まず，この式から
\begin{equation}\label{kb15:E-bounds}
 c\geq2,\qquad -1<E<\frac12,\qquad E\ne0,\qquad w\geq2
\end{equation}
を得る．実際，$c=1$ なら $E=1$ となって不可能である．
$E\leq-1$ なら $c(1+Ee)=e+1$ の左辺が非正となる．
$E>0$ の場合には $Ec=1-(c-1)/e<1$ なので $E<1/2$ であり，
$E\leq0$ の場合にも同じ上界は自明である．
$H^2<H(H+1)<(H+1)^2$ より $H(H+1)$ は平方数でなく，$E=0$ はない．
最後に $w=1$ なら $E\geq1/2$ となるので，$w\geq2$ である．

高辺 $ij,ik$ を，それぞれ puncture $k,j$ で比較する．
\[
 X=p^j_{ik},\qquad Y=p^k_{ij}
\]
と置くと，core $i$ での高高 entries は
$p^i_{jk}=-(a-c)$，$p^i_{kj}=-(b+e)$ なので，整数 clutching は
\begin{equation}\label{kb15:clutching-residues}
 X\equiv-c\pmod a,\qquad Y\equiv e\pmod b
\end{equation}
を要求する．一方，$q_{jk}=1$ により，core $j,k$ の pair equations は
\[
 Xp^j_{ki}=a-1,\qquad Yp^k_{ji}=b-1
\]
となる．従って $X,Y$ はそれぞれ $a-1,b-1$ の非零 signed divisors であり，
$1\leq|X|\leq a-1$，$1\leq|Y|\leq b-1$ である．
$c=s-x\leq s-1\leq a-2$，$e=z-y\leq z-1\leq b-2$ と合わせると，
合同式~\eqref{kb15:clutching-residues} の全ての候補は
\begin{equation}\label{kb15:two-representatives}
 X\in\{-c,a-c\},\qquad Y\in\{e,e-b\}
\end{equation}
で尽くされる．

正の候補 $a-c$ は
\[
 a-c=s(w-1)+x+1>\frac{a-1}{2},\qquad a-c<a-1
\]
を満たす．後者には $c\geq2$ を用いた．正整数の真約数はその半分以下なので，
$a-c$ は $a-1$ を割らず，$X=-c$ に限られる．
同様に，$e\geq2$ なら
\[
 b-e=z(w-1)+y+1>\frac{b-1}{2},\qquad b-e<b-1
\]
となり，$Y=e-b$ は不可能である．

まず $Y=e$ の場合を考える．$e=1$ でもこの場合は含まれる．
二つの逆向き entries は整数だが，その和は
\begin{align}
 p^j_{ki}+p^k_{ji}
 &= -\frac{a-1}{c}+\frac{b-1}{e}
 =-\frac{rw}{H}+\frac{(r+1)w}{H+1}\notag\\
 &=-\frac{w^2}{H(H+1)}=E-1.
 \label{kb15:reverse-sum}
\end{align}
従って $E$ は整数でなければならないが，式~\eqref{kb15:E-bounds} に反する．

残るのは $e=1$，$Y=1-b$ だけである．このとき $p^k_{ji}=-1$ であり，
低辺 $jk$ の puncture $i$ における node 数は
\begin{equation}\label{kb15:exceptional-node}
 n_{jk;i}=1-p^j_{ki}-p^k_{ji}
 =2+\frac{rw}{H}.
\end{equation}
式~\eqref{kb15:E} は $E=2/c-1$ を与える．$c\geq2$ であり，
$c=2$ は $E=0$ となって不可能なので，$E<0$ である．
従って $w^2>H(H+1)>H^2$，すなわち $w>H$ となる．ゆえに
\[
 \frac{rw}{H}=w+\frac{w^2}{H}>2w\geq4,
 \qquad n_{jk;i}>6.
\]
辺 $jk$ の三つの node 数の和は，その両端の row sums が $-1$ であることから五である．
全て非負なら各 node 数は五以下なので，これは矛盾である．
これで負の枝を全て排除した．正の枝の残りの不等式は
補題~\ref{ml12:dichotomy} の二分から従う．
\end{proof}

\begin{corollary}[負の枝は高分母三角形に属する]
\label{kb15:negative-closes-triangle}
前補題の大域仮定の下で，二価 core の標準形が負の枝に属するなら，
その二つの高分母近傍を結ぶ辺も高分母である．従って，その core は
高分母 support の三角形に属する．
\end{corollary}

\begin{proof}
近傍間の分母が一である場合は前補題で排除されたので，一より大きい．
\end{proof}

この証明で core $j,k$ に使ったのは，整数 pair equations，整数 clutching，
および低辺 $jk$ の node 非負性と row sums である．
高分母次数三の標準形や，特定の一高分母・全低分母 core の分類は用いていない．
従って，分母三に伴う低 core の追加の自由度を捨てる議論にも依存しない．

\begin{theorem}[\texorpdfstring{$K_{2,3}$}{K2,3} 高分母 support の不存在]
\label{kb15:no-K23}
五次の標準 $K_5$ branch graph 上で，正の局所整数 charts，marked puncture の
整数 clutching，および node 数の非負性を仮定する．
高分母 support が $K_{2,3}$ となることはない．
この結論に対向 $D$ の相殺と分母の数値上限は不要である．
\end{theorem}

\begin{proof}
三価頂点を $a,b$，二価頂点を $i\in\mathbb Z/3$ とする．$q_{ab}=1$ なので，
補題~\ref{kb15:no-negative-low-wedge} により三つの二価 core は全て正の枝に属する．
各 core の canonical data を $r_i,S_i,T_i$ と書き，$L_i=T_i-S_i+1$ とすると
\begin{equation}\label{kb15:positive-bounds}
 S_i\geq r_i+1,\qquad T_i\geq r_i+3,\qquad L_i\geq3.
\end{equation}
第二式は $r_iT_i>(r_i+1)S_i\geq(r_i+1)^2$ から従う．

二価 core $i$ の二つの低 rays のうち，高 ray $a$ へ正の逆向き entry を持つ方へ
矢印を向ける．互いに相手へ矢印を向ける二頂点 $i,j$ があれば，
低辺 $ij$ の puncture $a$ で
$n_{ij;a}=1-p^i_{ja}-p^j_{ia}\leq-1$ となるため不可能である．
三頂点の各 outdegree は一なので，矢印は directed three-cycle をなす．
その向きに合わせて $i\to i+1$ と標識する．

canonical 高 ray $A$ が $a$ に対応するときを $\sigma_i=+$，
$b$ に対応するときを $\sigma_i=-$ とする．
低 rays 間の entries は，$\sigma_i=+$ なら
\[
 (p^i_{i+1,i-1},p^i_{i-1,i+1})=(0,-L_i),
\]
$\sigma_i=-$ なら $(-L_i,0)$ である．
もし $\sigma_i$ が一定でなければ，ある cyclic な隣接対 $i,j=i+1$ で
$(\sigma_i,\sigma_j)=(-,+)$ となる．残る二価頂点を $k$ とすると
\[
 n_{ij;k}=1+L_i+L_j\geq7
\]
となり，その低辺の node 全数五と非負性に反する．従って $\sigma_i$ は一定である．

全て $\sigma_i=+$ の場合，低辺 $i,i+1$ の puncture $b$ の node 数は
$2+r_i-T_{i+1}$ である．三辺で足すと
\[
 \sum_i n_{i,i+1;b}=6+\sum_i(r_i-T_i)\leq-3
\]
となり，非負性に反する．全て $\sigma_i=-$ の場合にも，puncture $a$ の node 数を足すと
同じ矛盾を得る．従って全ての場合が排除された．
\end{proof}

三つの二価 core の間の低辺だけを課した必要条件系は，中心間の低辺を含まない．
本節の障害は，中心間の低辺と，そこへ到達する高辺の整数 clutching を用いる．

% v15: diamond exclusion. Entire proof is analytic and has no denominator bound.
\section{Diamond 型高分母 support の排除}
\label{dg15:section}

高分母 support が四頂点上の diamond であり，第五の core は全低分母である場合を考える．
この場合には，四本の低分母辺に沿う対向係数の相殺を加えることで，
二つの二価 core の低低 entry を同期させることができる．
以下では分母の数値上限も，三高分母局所模型の分類も用いない．

\begin{lemma}[全低分母 core における総和恒等式]
\label{dg15:hub-balance}
五つの core を $0,1,2,3,4$ と標識し，$q_{i4}=1$ $(0\leq i\leq3)$ とする．
各 core の整数行列は行和と pair products を満たし，四辺 $i4$ で対向 $D$ が相殺するとする．
$I=\{0,1,2,3\}$ とし，$i\ne j$ に対して
\[
 x_{ij}=p^i_{j4},\qquad r_{ij}=p^i_{4j},\qquad V_{ij}=p^4_{ij}
\]
と置く．このとき
\begin{equation}\label{dg15:balance-total}
 \sum_{i\ne j\in I}x_{ij}=\sum_{i\ne j\in I}r_{ij}=-4.
\end{equation}
さらに，高分母辺 $ij$ の puncture $4$ で整数 clutching が成立すれば，
\begin{equation}\label{dg15:low-defects}
 \sum_{\substack{i<j\in I\\q_{ij}=1}}
 (x_{ij}+x_{ji}-r_{ij}-r_{ji})=0.
\end{equation}
\end{lemma}

\begin{proof}
各 core の第 $4$ 行から $\sum_{j\ne i}r_{ij}=-1$ を得る．一方，
\[
 D^i_4=-1-\sum_{j\ne i}x_{ij},\qquad
 D^4_i=-1-\sum_{j\ne i}V_{ji}
\]
なので，相殺を $i$ について加えると
$\sum x_{ij}+\sum V_{ij}=-8$ となる．$V$ の四行の和は $-4$ だから
\eqref{dg15:balance-total} が従う．

$q=q_{ij}>1$ のとき，$x_{ij},x_{ji}$ は $q-1$ の符号付き約数であり，
$q\mid x_{ij}+x_{ji}$ である．従ってその対は
\begin{equation}\label{dg15:low-puncture-options}
 (t,-t),\qquad \{1,q-1\},\qquad\{-1,-(q-1)\}
\end{equation}
のいずれかである．実際，和の絶対値は $2(q-1)<2q$ であり，
同符号の和が $q$ なら，一方の約数が $q-1$，他方が $1$ でなければならない．
両方が真の正約数なら，その和は $q-1$ 以下になるからである．
各場合に
\[
 r_{ij}+r_{ji}=(q-1)/x_{ij}+(q-1)/x_{ji}=x_{ij}+x_{ji}
\]
が成立する．高分母辺の寄与を総和から除けば \eqref{dg15:low-defects} を得る．
\end{proof}

\begin{theorem}[Diamond の排除]
\label{dg15:no-diamond}
全平面五次の標準 $K_4$ cones と $K_5$ branch graph に対する，正の整数局所 charts，
marked clutching，対向 $D$ の相殺，node 非負性を満たすデータについて，
高分母 support が
\[
 H=\{01,02,03,12,13\}
\]
となることはない．
この排除には分母の上限を要しない．
\end{theorem}

\begin{proof}
二高分母の標準形と補題~\ref{ml12:dichotomy} の二分を用い，
\begin{equation}\label{dg15:canonical}
 P=\begin{pmatrix}
 0&-A&x&-s\\
 -B&0&-y&z\\
 r&-r-1&0&0\\
 -S&T&-L&0
 \end{pmatrix},\qquad
 \begin{gathered}
 a=sS+1,\quad b=zT+1,\\
 x=(a-1)/r,\quad y=(b-1)/(r+1),\\
 A=a+x-s,\quad B=b-y+z,\quad AB=ab-1,\\
 L=T-S+1
 \end{gathered}
\end{equation}
と書く．全ての $r,s,z,S,T,x,y$ は正整数，$S\geq2$ である．
負の枝は $r>S=T$ で $L=1$，正の枝は $S>r$ で
\begin{equation}\label{dg15:positive-range}
 S<r(L-1)
\end{equation}
を満たす．従って $L\in\{1,3,4,\ldots\}$ である．
以下の primed 記号は第二の標準形を表す．

\paragraph{二つの低低 entry の同期．}
\[
 u_2=p^2_{34},\quad v_2=p^2_{43},\qquad
 u_3=p^3_{24},\quad v_3=p^3_{42}
\]
と置く．二高分母の標準形により，各対 $(u_i,v_i)$ は $(0,-L_i)$ または
$(-L_i,0)$ であり，$L_i>0$ である．補題~\ref{dg15:hub-balance} の総和には
低分母辺 $23$ だけが残るので，$u_2+u_3=v_2+v_3$ となる．
従って $L_2=L_3=:L$ であり，向きは反対である．$2,3$ を交換して
\[
 (u_2,v_2)=(-L,0),\qquad (u_3,v_3)=(0,-L)
\]
とする．node $23;4$ の数は $1+L$ だから $L\leq4$，従って
\begin{equation}\label{dg15:three-L}
 L\in\{1,3,4\}.
\end{equation}

core $2$ の第 $3$ 行と core $3$ の第 $2$ 行には，高 ray に向かう正の entry が
それぞれ一つだけある．もしその行先が同じなら，対応する node $23;0$ または $23;1$ の数が負になる．
従って $0,1$ を交換すれば，core $2$ の標準形は $(B,A,D,C)=(0,1,3,4)$，
core $3$ は $(B,A,C,D)=(0,1,2,4)$ とできる．
辺 $23$ の三つの node 数は
\begin{equation}\label{dg15:n23}
 (n_{23;0},n_{23;1},n_{23;4})
 =(r'+2-T,\ 1+S-r',\ 1+L).
\end{equation}
従って
\begin{equation}\label{dg15:rprime-range}
 S+L-3\leq r'\leq S+1.
\end{equation}

\paragraph{共通低 puncture の例外の全列挙．}
中心 core $0,1$ の第 $4$ 行で，共有高 ray に向かう entry を
\[
 R_0=p^0_{41},\qquad R_1=p^1_{40}
\]
と書く．$q_{01}=Q>1$ なので，これらは非零であり，
\eqref{dg15:low-puncture-options} と同じ三択を満たす．特に，異符号なら和は零である．

四辺 $02,12,03,13$ の puncture $4$ で \eqref{dg15:low-puncture-options} の
反対符号の場合を選ぶと，中心側の第 $4$ 行は
\[
 p^0_{42}=r+1,\quad p^1_{42}=-r,\qquad
 p^0_{43}=-T',\quad p^1_{43}=S'.
\]
その行和から
\begin{equation}\label{dg15:central-baseline}
 R_0=T'-r-2,\qquad R_1=r-S'-1,\qquad R_0+R_1=L-4
\end{equation}
を得る．この選択と異なる場合を以下では例外と呼ぶ．全ての例外は次の表で尽くされる．
表の outgoing entry は，二価側の共有高 ray の行から puncture $4$ に向かう成分を指す．
\[
\begin{array}{c|c|c|c}
 \text{辺}&\text{二価側の outgoing entry}&\text{例外の必要条件}&
 \text{中心の第4行の変化}\\\hline
 02&-y&y=1&p^0_{42}:\ r+1\longmapsto-1\\
 12&x&r=1\ \text{または}\ x=1&p^1_{42}:\ -r\longmapsto -r+a\\
 03&z'&z'=1&p^0_{43}:\ -T'\longmapsto1\\
 13&-s'&s'=1&p^1_{43}:\ S'\longmapsto-1
\end{array}
\]
例えば $y=(b-1)/(r+1)<b-1$ だから，負の同符号対は $y=1$ の場合だけである．
他の三行も \eqref{dg15:low-puncture-options} と $S',T'\geq2$ から従う．
例外による $(R_0,R_1)$ の変化は，順に $(b,0),(0,-a),(-b',0),(0,a')$ である．

ここで $x=1$ は常に不可能である．正の枝では $x=sS/r>s\geq1$ である．
負の枝では $H=r-S$，$u=s-x$，$v=z-y$ と置くと，正整数 $u,v$ は
既出の式~\eqref{ml12:E-equation}，すなわち
\begin{equation}\label{dg15:negative-identity}
 Euv=v-u+1,\qquad E=1-\frac{S^2}{H(H+1)}
\end{equation}
を満たす．$u=1$ は $E=1$ を要求するので $u\geq2$ である．
$u(1+Ev)=v+1$ より $E>-1$，また $E>0$ なら
$Eu=1-(u-1)/v<1$ より $E<1/2$ である．従って
\begin{equation}\label{dg15:negative-range}
 \frac{S^2}{2}<H(H+1)<2S^2.
\end{equation}
なお $S=1$ は $H(H+1)\geq2$ と右側の厳密不等式に反するので，負の枝でも $S\geq2$ である．
$x=1$ なら $r=sS$，$s=u+1\geq3$，従って $H\geq2S$ となり，
右側の不等式に反する．これで上表の全例外を吟味する準備ができた．

\paragraph{$L=1$ の排除．}
両方とも負の枝であり，$r>S\geq2$，$r'>S'\geq2$ である．
従って $r=1$ は不可能であり，$s'\geq2$，$z'\geq2$ なので，残る例外は $y=1$ だけである．

例外がなければ $R_0+R_1=-3$ なので，\eqref{dg15:low-puncture-options} から
$Q=3$，$R_0,R_1<0$ となる．特に $r\leq S'$ である．しかし
\[
 S<r\leq S'<r'\leq S+1
\]
は整数について不可能である．

$y=1$ の例外を選ぶ場合は $b=r+2$ なので
\[
 R_0=S'\geq2,\qquad R_0+R_1=r-1>0.
\]
同符号正の対は $\{1,Q-1\}$ だから，$R_1=1$，$S'=r-2$ である．
$r'\geq S'+1=r-1$ と \eqref{dg15:rprime-range} から $S\geq r-2$，
従って $H=r-S\in\{1,2\}$ となる．さらに $y=1$ は
\[
 (z-1)S=H+1
\]
を与える．$S\geq2$ より $(H,S)=(1,2),(2,3)$ の二つだけである．
前者は \eqref{dg15:negative-range} に反する．後者では
$(r,S,S',r')=(5,3,3,4)$ となり，primed 模型の $H'=1,S'=3$ が
同じ不等式に反する．

\paragraph{$L=3$ の排除．}
両方とも正の枝であり，$r<S<2r$，$r'<S'<2r'$ である．
従って $r=1$ は不可能である．また $y=zT/(r+1)>z\geq1$ なので $y=1$ は不可能である．
$s'=1$ なら $x'=S'/r'\in(1,2)$ となる．
$z'=1$ なら
\[
 y'=\frac{S'+2}{r'+1}\in(1,2)
\]
となる．いずれも整数性に反する．従って例外は一つもなく，
$R_0+R_1=-1$ が $Q>1$ による可除性に反する．

\paragraph{$L=4$ の例外の排除．}
この場合 $r'=S+1\geq3$ である．引き続き $y=1$ は不可能である．
$s'=1$ なら $1<S'/r'<3$ と整数性から $S'=2r'$，$T'=2r'+3$ となる．
$y'$ の整数性より $z'=h(r'+1)$，$h\geq1$ と書ける．
$A'B'=a'b'-1$ へ代入すると
\[
 h(r'+1)=2
\]
となり，$r'=1$ を要求するため不可能である．

$z'=1$ なら $1<T'/(r'+1)<3$ なので
$T'=2(r'+1)$，$S'=2r'-1$ となる．$x'$ の整数性より $s'=h r'$ と書け，
同じ pair equation は
\[
 h(r'-2)=0
\]
となる．従って $r'=2$ を要求し，これも不可能である．
これら二つの代入恒等式は任意の正整数に対する恒等式であり，有界探索を用いていない．

残る例外は $r=1$ の場合の辺 $12$ だけである．これを選ぶと
$S=2$，$T=5$，$r'=3$ であり，
\[
 R_0=S'>0,\qquad R_1=-S'-a<0.
\]
異符号でありながら和が $-a\ne0$ となるので，この例外も不可能である．

\paragraph{$L=4$ の反対符号の場合の排除．}
残った場合は \eqref{dg15:central-baseline} そのものであり，
\[
 R_0=R=S'-r+1,\qquad R_1=-R.
\]
$S\geq r+1$，$r'=S+1$，$S'\geq r'+1$ より $R\geq4$ である．
全低 core の行列を引き続き $V$ と書く．node $04;1$ の非負性は
$V_{01}\leq1-R<0$ を与える．$V_{01}V_{10}=0$ より $V_{10}=0$ である．
従って node $14;0$ の数は $1+R$ であり，辺 $14$ の node 総数五より $R\leq4$ となる．
結局
\[
 R=4,\qquad S=r+1,\qquad r'=r+2,\qquad S'=r+3.
\]
node $14;0$ が五を使い尽くすので $n_{14;2}=0$ であり，
$p^1_{42}=-r$ より $V_{12}=r+1>0$ となる．
今度は $V_{12}V_{21}=0$ より $V_{21}=0$ だから，
\[
 n_{24;1}=1-p^2_{41}-V_{21}=1-r\geq0.
\]
従って $r=1$，すなわち primed 模型は $(r',S',T')=(3,4,7)$ でなければならない．

この模型では $s'=3h$，$z'=4k$，$h,k\geq1$ と書け，
\[
 a'=12h+1,\quad b'=28k+1,\quad A'=13h+1,\quad B'=25k+1.
\]
従って pair equation は
\[
 k(11h+3)=h+1
\]
となる．左辺は $11h+3>h+1$ 以上なので不可能である．
以上で \eqref{dg15:three-L} の全場合を排除した．
\end{proof}

\begin{remark}[使用した大域条件の範囲]
\label{dg15:scope}
上の証明で $D$ の相殺を使うのは全低 core に接する四辺 $i4$ の総和だけである．
二つの三高分母中心の正符号 graph の分類や，中心側の正 weights の存在判定は用いていない．
二価 core の局所標準形，低 puncture の整数 clutching，行和，pair products，
指定した低辺の node 非負性だけを組み合わせている．
従って，三高分母の matching 定理や全高 $D$ の符号に依存しない．
\end{remark}

% v15. General one-high normal form, a two-high low-coefficient lemma,
% and an unbounded restriction on the remaining two-leaf triangle support.
% Dependencies: qh9:isolated-local, q9:four-line-criterion,
% td12:one-negative, ml12:dichotomy, ml12:denominator-descent,
% ls11:low-puncture, cc12:proper-negative, cq13:reverse-two.

\section{一高分母の完全標準形と二葉三角形の交差条件}
\label{tl15:section}

一高分母 core の高分母 row だけでなく，全ての低分母 entries を記述する．
その自由度を保持することは，分母三で同じ負の entry が二度現れる場合にも必要である．
以下では，局所係数を
\[
 D_j=-q_j-\sum_{k\ne j}p_{kj}
\]
と書く．共有辺に対しては，同じ marked punctures の clutching と
$D$ の相殺を仮定する．

\begin{proposition}[一高分母 core の全整数標準形]
\label{tl15:one-high-full}
正の weights を持つ分母 pattern $(q,1,1,1)$，$q>1$ の四直線 core は，
高分母 ray を $A$ とし，低分母 rays を適切に標識すると，次の形に書ける：
\begin{equation}\label{tl15:leaf-matrix}
\begin{gathered}
 h,t\in\ZZ_{>0},\quad m=th(h+1),\quad q=m+1,\quad
 u=th,\quad v=t(h+1),\\
 c,d\in\ZZ,\quad cd=0,\qquad
 P=\begin{pmatrix}
 0&-m&-h-1&h\\
 -1&0&0&0\\
 -u&u-1-c&0&c\\
 v&-v-1-d&d&0
 \end{pmatrix}.
\end{gathered}
\end{equation}
逆に，この表示の全ての整数 choices は正の weights を持つ局所模型を与える．
係数は $A,B,C,D$ の順に
\begin{equation}\label{tl15:leaf-D}
 (-m-t,\ m+t+1+c+d,\ h-d,\ -h-1-c)
\end{equation}
である．
\end{proposition}

\begin{proof}
命題~\ref{qh9:isolated-local} により，高分母 row とその reverse entries を
$(-m,-h-1,h)$ と $(-1,-u,v)$ に標準化できる．
$B$ から他の低分母 rays への entries を $(e,-e)$ とする．
$e\ne0$ なら低分母間の pair products が零であるため
$p_{CB}=p_{DB}=0$ である．残る二つの row sums は
\[
 p_{CD}=u-1,\qquad p_{DC}=-v-1
\]
を与える．その積は零であり，$v+1>0$ なので $u=1$，従って $h=t=1$ である．
この場合 $A\to B,A\to C$ はともに $-2$，reverse entries はともに $-1$ であり，
$C$ の低分母 outgoing entries は零である．$B,C$ を交換すれば，
高分母 row を保ったまま $B$ の低分母 outgoing entries を零にできる．
従って全ての場合に $p_{BC}=p_{BD}=0$ としてよい．
$p_{CD}=c,p_{DC}=d$ と置けば，残る row sums と pair equation は
式~\eqref{tl15:leaf-matrix} と $cd=0$ をちょうど与える．
各 column sum が式~\eqref{tl15:leaf-D} を与える．

逆向きには，表示した行列の row sums は $(-q,-1,-1,-1)$ であり，
全ての pair products は $q_iq_j-1$ となる．
正の entry の directed graph を考えると，$B$ から正の edge は出ず，
$C\to A$ は負である．$cd=0$ より $C,D$ 間には正の二閉路がない．
従って唯一あり得る正の directed cycle は $A\leftrightarrow D$ であり，
その normalized gain は
\[
 \frac{p_{AD}p_{DA}}{q_Aq_D}=\frac{hv}{q}=\frac{m}{m+1}<1
\]
である．正の entries に対する不等式
$w_j>(p_{ij}/q_i)w_i$ を対数化した有限差分不等式系は，
全ての directed cycle の gain が一未満なら解を持つ．
実際，有限個の単純 cycles の各平均対数 gain が負なので，
全 edges に十分小さい共通の正数を加えても cycle sums は負のままであり，
通常の差分不等式の potential を取れば元の厳密不等式を得る．
非正の entries に対しては任意の正の weights で basin positivity が成立する．
よって正の weights が存在し，命題~\ref{q9:four-line-criterion} が適用できる．
\end{proof}

\begin{lemma}[二高分母 core の低係数は同時に負でない]
\label{tl15:low-D}
正の二高分母標準形
\[
 P=\begin{pmatrix}
 0&-A&x&-s\\
 -B&0&-y&z\\
 r&-r-1&0&0\\
 -S&T&S-T-1&0
 \end{pmatrix},\quad
 a=sS+1, b=zT+1, x=\frac{sS}{r}, y=\frac{zT}{r+1}
\]
において，低分母係数
\begin{equation}\label{tl15:low-D-formula}
 D_C=y-x+T-S,\qquad D_D=s-z-1
\end{equation}
が同時に負となることはない．
\end{lemma}

\begin{proof}
$D_D<0$ と仮定すると整数性から $z\ge s$ である．
正の枝 $S>r$ では，補題~\ref{ml12:dichotomy} より
$T/(r+1)>S/r$ なので $y>x$，従って $D_C>0$ となる．
負の枝では $S=T=w<r$ と書ける．
定理~\ref{ml12:denominator-descent} の証明の記号
\[
 H=r-w,\quad E=1-\frac{w^2}{H(H+1)},\quad
 s=c+\frac{cw}{H},\quad z=e+\frac{ew}{H+1}
\]
を用いる．$E<0$ の場合は同証明により $s>z$ なので，現在の仮定の下では
$E>0$ である．同証明から $\delta=e-c\ge1$，$c\le H$ を得るため，
\[
 D_C=\frac{w(H\delta-c)}{H(H+1)}\ge0
\]
である．
\end{proof}

\begin{lemma}[葉の負の reverse entry に対する二分]
\label{tl15:leaf-test}
式~\eqref{tl15:leaf-matrix} の葉 core において，低分母 ray $E$ から
唯一の高分母 ray $A$ への entry が負とする．従って $E=B$ または $C$ である．
残る二つの低分母 rays を $J,K$ とし，ある整数 $\kappa\ge1$ について
$p_{JE}\le-\kappa$ とする．このとき次が成立する．
\begin{enumerate}
\item $E=C$ なら $J=D,K=B$，$d\le-\kappa$，$c=0$ であり，
$D_E>0$ および $p_{KE}=0$ である．
\item $E=B$ かつ $D_E\le0$ なら $J=D,K=C$，$d=0$ であり，
\begin{equation}\label{tl15:leaf-test-large}
 p_{KE}=t(h+1)^2-D_E\ge4
\end{equation}
となる．
\end{enumerate}
\end{lemma}

\begin{proof}
$E=C$ の場合，$B\to C=0$ なので $J=D$ であり，結論は標準形から従う．
$E=B$ とする．$J=C$ なら $u-1-c\le-\kappa$ より
$c\ge u+\kappa-1>0$，従って $d=0$ であり $D_B>0$ となる．
$J=D$ で $d\ne0$ なら $c=0$ であり，
$-v-1-d\le-\kappa$ から
\[
 D_B=m+t+1+d\ge m+t+\kappa-v=th^2+\kappa>0
\]
となる．従って $D_B\le0$ の下では $J=D,d=0$ でなければならない．
このとき $c=D_B-m-t-1$ を $p_{CB}=u-1-c$ に代入すると
式~\eqref{tl15:leaf-test-large} を得る．
\end{proof}

\begin{theorem}[二葉三角形に必要な交差 pairing]
\label{tl15:crossed-pairing}
各 core が正の weights を持つ整数局所 chart であり，全共有辺で
同じ marked punctures の整数 clutching，対向 $D$ の相殺，
meridian node 数の非負性を満たすとする．高分母 support が
\[
 H=\{01,02,12,03,14\}
\]
であるとする．support の自己同型 $(0\,1)(3\,4)$ による標識交換を必要なら行い，
三角形 $012$ の二価 core $2$ の正の二高分母標準形で
$A=0,B=1$ とすると，低分母標識 $C=3,D=4$ は不可能である．
従って，この support が残るなら $C=4,D=3$ でなければならない．
すなわち，二価 core の正の高低 pairing は，葉の付く三角形頂点を交差して対応させる．
この結論に分母の上限はない．三価 cores の一般分類も用いない．
\end{theorem}

\begin{proof}
$C=3,D=4$ と仮定する．core $2$ の低 rows は
\[
 (p^2_{30},p^2_{31},p^2_{34})=(r,-r-1,0),\quad
 (p^2_{40},p^2_{41},p^2_{43})=(-S,T,S-T-1).
\]
低辺 $23,24$ の node 非負性から
\[
 p^3_{20}\le1-r\le0,\qquad p^4_{21}\le1-T<0.
\]
前者も pair product が正であるため零ではない．
従って両葉で ray $2$ は一高分母標準形の $B$ または $C$ である．
両方が $B$ なら $r-1\le1,T-1\le1$ なので $r\le2,T\le2$ となる．
$2\le S\le T$ より $S=T=2$ であり，負の枝 $r>S$ と矛盾する．
よって $B$ 型の葉は高々一つである．

core $2$ の負の高低 entries $-y=p^2_{13}$ と $-s=p^2_{04}$ は
補題~\ref{cc12:proper-negative} により各 $q-1$ の真の約数である．
補題~\ref{ls11:low-puncture} を共有高辺 $12,02$ に適用すると，
$y>1,s>1$ の場合には
\[
 p^1_{23}=y,\quad p^1_{32}=r+1,\qquad
 p^0_{24}=s,\quad p^0_{42}=S
\]
しかない．従って低辺 $13,04$ の node 非負性が
\begin{equation}\label{tl15:cross-low-bounds}
 p^3_{12}\le-r,\qquad p^4_{02}\le1-S
\end{equation}
を要求する．
両葉で ray $2$ が $C$ なら，補題~\ref{tl15:leaf-test} によって
両葉の $D_2$ が正となる．低辺 $23,24$ の相殺により
core $2$ の $D_C,D_D$ がともに負となり，補題~\ref{tl15:low-D} に反する．
一方だけが $B$ なら，もう一方の $C$ 型葉に対向する低係数は負である．
補題~\ref{tl15:low-D} により $B$ 型葉に対向する低係数は非負なので，
その葉の $D_2$ は非正である．補題~\ref{tl15:leaf-test} から，
低辺 $34$ の puncture $2$ における二つの outgoing entries は，
一方が少なくとも四，他方が零となる．その和は一以下でなければならず，矛盾する．

例外 $s=1$ を扱う．負の枝では $x=sS/r<s$ が正整数でなければならないため，
現在は正の枝である．従って $T>r+1$，$y>1$ であり，
\[
 D_D=-z<0,\qquad
 D_C-z=z\left(\frac{T}{r+1}-1\right)
          +T-S-\frac Sr>0.
\]
式~\eqref{tl15:cross-low-bounds} の第一の不等式は有効である．
$D_C>0$ と補題~\ref{tl15:leaf-test} により，葉 $3$ では
ray $2$ が $B$，ray $1$ が $D$，ray $4$ が $C$，かつ $d_3=0$ となる．
その parameters を $h,t$ と書くと
\[
 D^3_4=h,\qquad p^3_{42}=D_C+t(h+1)^2.
\]
葉 $4$ の ray $2$ は $C$ である．低辺 $34$ の node 条件から
$p^4_{32}<0$ が必要なので，その ray $3$ は $D$，ray $0$ は $B$，
$d_4<0,c_4=0$ となる．葉 $4$ の parameter を $h_4$ とすると，
低辺 $34,24$ の係数相殺がそれぞれ
\[
 h=h_4+1,\qquad d_4=h_4-z
\]
を与える．従って低辺 $34$ の puncture $2$ は
\[
 D_C+t(h+1)^2+h-1-z\le1
\]
を要求する．しかし $D_C>z$，$h,t\ge1$ なので不可能である．

最後に $y=1$ とする．正の枝では $T>r+1$ により $y>z\ge1$ なので，
現在は負の枝 $S=T$，かつ $s>1$ である．
$D_C=1-x$ であり，$x=1$ なら $sS=r,zS=r+1$，従って
$(z-s)S=1$ となって $S\ge2$ に反する．よって $D_C<0$ であり，
補題~\ref{tl15:low-D} から $D_D\ge0$ となる．
式~\eqref{tl15:cross-low-bounds} の第二の不等式と
補題~\ref{tl15:leaf-test} によって，葉 $4$ の ray $2$ は $B$ でなければならない．
先ほどと同じ node 条件が $T\le2$，従って $S=T=2$ を強制する．
しかし補題~\ref{cq13:reverse-two} の唯一の模型では $y=2$ であり矛盾する．
$s>1,y>1$，$s=1$，$y=1$ の全てを排除した．
\end{proof}

ここまでの議論は低標識を $C=4,D=3$ に限定する．
この交差割当では葉から高 ray への reverse entries が正にもなり得るため，
上の負の reverse entry の二分をそのまま適用できない．
以下では一高分母標準形の自由な整数 $c,d$ を保持して全低辺を復元し，
定理~\ref{tg16:no-two-leaf-triangle} で二葉三角形 support 全体を排除する．

% v15. Additional independent module after one_high_and_crossed_pairing_v15.tex.
% The finite calculation concerns only cores 2,3,4 and has no denominator cutoff.

\subsection{二葉三角形の負の枝の排除と残る小 reverse 条件}
\label{tl15:negative-section}

\begin{lemma}[二高分母標準形の $x$ は一でない]
\label{tl15:x-at-least-two}
正の二高分母標準形では $x=(a-1)/r\ge2$ である．
\end{lemma}

\begin{proof}
正の枝 $S>r$ では $x=sS/r>1$ である．
負の枝 $S=T=w<r$ で $x=1$ と仮定する．
$s w=r$ なので $s\ge2$ である．$y=zw/(r+1)$ の整数性と
$\gcd(w,r+1)=1$ から $z=k(r+1)$，$y=kw$，$k\ge1$ と書ける．
\[
 a=r+1,\quad b=k(r+1)w+1,\quad
 A=r+2-s,\quad B=1+k\{r(w+1)+1\}
\]
を $AB=ab-1$ に代入し，$r=sw$ を用いると
\[
 k\{w(-s^2+3s-1)+2-s\}=s-2
\]
となる．$s=2$ なら左辺は $kw>0$，右辺は零である．
$s\ge3$ なら左辺は負，右辺は正である．いずれも矛盾する．
\end{proof}

\begin{lemma}[交差 pairing の大きい reverse entries の障害]
\label{tl15:large-reverse}
定理~\ref{tl15:crossed-pairing} の全仮定の下で，残る交差標識
$(A,B,C,D)=(0,1,4,3)$ を取る．
$r\ge5,S\ge5,z>1$ が同時に成立することはない．
\end{lemma}

\begin{proof}
core $2$ の低 rows は
\[
 (p^2_{30},p^2_{31},p^2_{34})=(-S,T,S-T-1),\qquad
 (p^2_{40},p^2_{41},p^2_{43})=(r,-r-1,0).
\]
各低辺の node 数は非負で総和五なので，各 puncture で
$-4\le p+p'\le1$ である．従って低辺 $23,24$ は
\[
 p^3_{20}\ge S-4>0,\qquad p^4_{21}\ge r-3>0
\]
を要求する．一高分母標準形~\eqref{tl15:leaf-matrix} により，
両葉の ray $2$ はその標準形の $D$ でなければならない．

補題~\ref{tl15:x-at-least-two} と $r>1$ より
$1<x<a-1$ であり，$z>1,T\ge S>1$ より $1<z<b-1$ である．
共有高辺 $02,12$ の共通低 punctures に
補題~\ref{ls11:low-puncture} を適用すると，反対符号の通常の場合しか残らず，
\[
 p^0_{42}=-r,\qquad p^1_{32}=-T
\]
となる．低辺 $04,13$ の node 上界五から，さらに
\[
 p^4_{02}\ge r-4>0,\qquad p^3_{12}\ge T-4>0
\]
を得る．
各葉では ray $2=D$ であり，$B\to D=0$ なので，
これらの正の entries の出発 ray は $C$ である．
従って，両葉で他の葉へ向かう ray は $B$ であり，$c>0,d=0$ となる．
その係数 $D_B=m+t+1+c$ は正なので，低辺 $34$ の両端の係数は相殺できない．
\end{proof}

\begin{lemma}[負の枝で中間 core を固定したときの葉の有限復元]
\label{tl15:finite-leaves}
交差標識 $(A,B,C,D)=(0,1,4,3)$ を持つ負の枝 $S=T=w<r$ の
core $2$ の整数行列を固定する．低辺 $23,24$ において
marked node 条件と $D$ の相殺を満たす一高分母の葉の行列は，
分母の上限を置かず，有限個の約数計算から全て復元できる．
\end{lemma}

\begin{proof}
葉 $3$ で ray $2=B$ とすると $p^3_{21}=0$ であるが，
$p^2_{31}=w>1$ と合わせると低辺 $23$ の node 非負性に反する．
葉 $4$ で ray $2=B$ とすると $p^4_{20}=0$ であるが，
$p^2_{40}=r>1$ と合わせても同じ矛盾を得る．
従って両葉で ray $2$ は $C$ または $D$ である．

葉を $i\in\{3,4\}$，その高 ray を $H_i$ とし，固定された整数を
\[
 R_i=p^2_{iH_i},\qquad E_i=D^2_i
\]
と書く．ray $2=C$ の場合，$p^i_{2H_i}=-u$，$u=th$ は
\[
 u\in[R_i-1,R_i+4]\cap\ZZ_{>0}
\]
に入る．各 $u$ の全正約数 $h$ を取って $t=u/h$ とすれば，
$h,t$ の全 choices を得る．係数相殺は $d=h+E_i$ を固定する．
$d\ne0$ なら $c=0$ であり，$d=0$ なら残る $c=p^i_{2D}$ は
低辺 $2i$ の puncture $D$ の六整数の node 区間に入る．
残る二つの低 ray の $B,D$ の割当は二通りしかない．

ray $2=D$ の場合，$p^i_{2H_i}=v$，$v=t(h+1)$ は
\[
 v\in[-4-R_i,1-R_i]\cap\ZZ_{\ge2}
\]
に入る．各 $v$ の二以上の全正約数 $h+1$ を取って $t=v/(h+1)$ とする．
係数相殺は $c=E_i-h-1$ を固定する．
$c\ne0$ なら $d=0$ であり，$c=0$ なら $d=p^i_{2C}$ が
同じ六整数の node 区間で固定される．$B,C$ の二つの割当を試せばよい．

最後に元の三つの node 条件と係数相殺を再検査すれば，ちょうど全候補を得る．
どの段階にも葉の分母 $th(h+1)+1$ に対する上限は仮定していない．
\end{proof}

\begin{proposition}[小さい負の枝の低三角形は接続しない]
\label{tl15:small-negative}
交差標識で $S=T\le4$ の場合には，cores $2,3,4$ の低辺
$23,24,34$ における node 条件と係数相殺を同時に満たす模型はない．
\end{proposition}

\begin{proof}
補題~\ref{cp14:small-negative} の七模型がこの場合を全て尽くす．
補題~\ref{tl15:finite-leaves} の全約数復元を実行すると，
低辺 $23,24$ にそれぞれ接続する候補数と，検査する葉の全組数は次のとおりである．
\[
\begin{array}{rrrr|rr|r}
 r&s&z&w&\#\text{leaf }3&\#\text{leaf }4&\#\text{pairs}\\\hline
 4&4&5&2&4&8&32\\
 5&10&2&3&4&8&32\\
 6&6&14&3&4&19&76\\
 7&7&2&4&14&10&140\\
 7&21&4&4&5&10&50\\
 8&6&9&4&5&12&60\\
 8&8&27&4&5&19&95
\end{array}
\]
合計四百八十五組は全て，低辺 $34$ の三つの node 区間
$[-4,1]$ または $D^3_4+D^4_3=0$ の少なくとも一方に反する．
標準 library の整数・有理数演算だけを用いる
\nolinkurl{verify_two_leaves_v15.py} は，七模型を生成する完全な約数 trials，
全葉候補の行列，全組の棄却理由を
\nolinkurl{two_leaves_exact_certificate_v15.json} に出力する．
検算器は復元した各行列の元の row sums，pair products，係数，
および正の cycle inequalities を全て再検査する．

これは三つの cores の低三角形だけで生じる矛盾である．
三価 cores の未知の分母 $q_{01}$ を固定・列挙する必要はなく，
その値によってこの矛盾を回避することもできない．
\end{proof}

\begin{theorem}[二葉三角形に残る正の枝と小 reverse 条件]
\label{tl15:remaining-positive}
定理~\ref{tl15:crossed-pairing} の全仮定の下で，二葉三角形 support
$\{01,02,12,03,14\}$ が残るなら，二価 core $2$ は
\[
 (A,B,C,D)=(0,1,4,3),\qquad S>r,\qquad rT>(r+1)S
\]
を満たす正の枝であり，さらに
\begin{equation}\label{tl15:small-reverse-or-unit}
 r\le4\quad\text{または}\quad z=1
\end{equation}
でなければならない．
\end{theorem}

\begin{proof}
非交差標識は定理~\ref{tl15:crossed-pairing} によって排除される．
交差標識で負の枝 $S=T=w$ を考える．$w\le4$ は
命題~\ref{tl15:small-negative} に反する．$w\ge5$ なら
$r>w\ge5$ であり，$y=zw/(r+1)$ は正整数で $y<z$ なので $z>1$ である．
従って補題~\ref{tl15:large-reverse} に反する．
よって正の枝のみが残り，$rT>(r+1)S$ は既知の二分補題に従う．
その枝で $r\ge5,z>1$ なら $S>r\ge5$ なので，
再び補題~\ref{tl15:large-reverse} に反する．
\end{proof}

式~\eqref{tl15:small-reverse-or-unit} は次の有限化への中間的な条件である．
この段階では $S,T,s,z$ に上限を置いていない．
次の共通低 puncture の補題で $1\le r\le4$ かつ $z=1$ へ強め，
葉の完全復元と中心の pair equations により全体の不存在へ進む．

% v16. A node budget bound at a common low puncture.
% No D opposition or hypothesis on the high degrees of the three cores
% is used in the general lemma.

\section{共通低 puncture における reverse 係数の一様な上界}
\label{cl16:section}

高辺の共通低 puncture では，二つの逆向き係数だけでなく，
第三の core における低低 pair equation を併せて使える．
この条件は，二高分母の標準形や一高分母の葉の分類に依存しない．

\begin{lemma}[共通低 puncture の通常接着における node 上界]
\label{cl16:ordinary-bound}
異なる三つの cores $i,j,k$ について
\[
 q_{ij}=q>1,\qquad q_{ik}=q_{jk}=1
\]
とする．各 core の整数 pair equations と各共有辺の非負 node 条件を仮定する．
各辺の node 総数を五とし，共有高辺の puncture $k$ において
通常の反対符号接着
\[
 p^i_{jk}=-p^j_{ik}
\]
が成立するとする．このとき
\begin{equation}\label{cl16:reverse-bound}
 p^i_{kj}=-p^j_{ki}=u,\qquad 1\le |u|\le4.
\end{equation}
この結論は，三つの cores の高分母次数を制限せず，対向 $D$ の相殺も使わない．
\end{lemma}

\begin{proof}
高低 pair product が $q-1$ であるため，通常接着から reverse entries は
$u,-u$ となり，$u$ は零でない整数である．core $k$ において
\[
 a=p^k_{ij},\qquad b=p^k_{ji}
\]
と置く．この二つの rays は低分母なので $ab=0$ である．
低辺 $ik,jk$ の node 数をそれぞれ $n=n_{ik;j}$，$m=n_{jk;i}$ とすると
\[
 a=1-n-u,\qquad b=1-m+u,\qquad 0\le n,m\le5.
\]
従って
\[
 0=ab=-\{u-(1-n)\}\{u-(m-1)\}.
\]
ゆえに $u=1-n\in[-4,1]$ または $u=m-1\in[-1,4]$ であり，
式~\eqref{cl16:reverse-bound} を得る．
\end{proof}

\begin{remark}[他の node 予算]
同じ証明により，共有辺の node 総数がそれぞれ $N_i,N_j$ なら
$u\in[1-N_i,1]\cup[-1,N_j-1]$ である．特に両辺の総数が
同じ $N\ge2$ なら $|u|\le N-1$ となる．
本稿の五次模型では $N=5$ を用いる．
\end{remark}

\begin{lemma}[正の二高分母枝における $T$ の下界]
\label{cl16:T-five}
二高分母標準形の正の枝 $S>r$，$rT>(r+1)S$ では $T\ge5$ である．
\end{lemma}

\begin{proof}
$S,r$ は正整数なので $S\ge r+1$ である．従って
\[
 rT>(r+1)S\ge(r+1)^2=4r+(r-1)^2\ge4r.
\]
ゆえに $T>4$ である．
\end{proof}

\begin{theorem}[二葉三角形に残る四つの reverse 値と単位係数]
\label{cl16:two-leaf-unit}
定理~\ref{tl15:remaining-positive} の仮定の下で，二葉三角形 support
$H=\{01,02,12,03,14\}$ が存在するとする．同定理の交差標識
$(A,B,C,D)=(0,1,4,3)$ における二高分母標準形は，必ず
\begin{equation}\label{cl16:small-and-unit}
 1\le r\le4,\qquad z=1
\end{equation}
を満たす．さらに
\begin{equation}\label{cl16:forced-exception}
 b=T+1,\qquad
 p^1_{23}=T,\quad p^1_{32}=1,\quad
 p^3_{12}=0,\qquad
 n_{12;3}=n_{13;2}=0
\end{equation}
が強制される．$r>1$ の場合には，併せて
\begin{equation}\label{cl16:forced-ordinary}
 p^0_{24}=-x,\quad p^0_{42}=-r,\quad p^4_{02}=0,\qquad
 n_{02;4}=1,\quad n_{04;2}=r+1
\end{equation}
となる．これらの追加制約を導く際には対向 $D$ の相殺は使わない．
\end{theorem}

\begin{proof}
既知の帰着により $S>r$，$rT>(r+1)S$ の正の枝である．
その交差標識では
\[
 p^2_{13}=z,\quad p^2_{31}=T,\quad b-1=zT;
 \qquad
 p^2_{04}=x,\quad p^2_{40}=r,\quad a-1=rx.
\]
補題~\ref{cl16:T-five} により $T\ge5$ である．

まず $z>1$ と仮定する．$T>1$ より $1<z<b-1$ である．
共通低 puncture の符号付き約数の分類
（補題~\ref{ls11:low-puncture}）により，高辺 $12$ の puncture $3$ では
$p^1_{23}=-z$ の通常接着しか起こらない．従って
補題~\ref{cl16:ordinary-bound} が $T\le4$ を与え，矛盾する．
ゆえに $z=1$ である．このとき $b=T+1$ であり，符号付き約数の二択は
$p^1_{23}=-1$ または $T$ となる．前者は同じ通常接着の矛盾を与えるので，
$p^1_{23}=T$，pair product から $p^1_{32}=1$ である．
また，低辺 $23$ の node 非負性より
$p^3_{21}\le1-T<0$ なので，core $3$ の低低 pair equation が
$p^3_{12}=0$ を与える．従って
\[
 n_{12;3}=1-\frac{1+T}{T+1}=0,\qquad
 n_{13;2}=1-(1+0)=0.
\]

次に $r>1$ とする．正の枝で $x=sS/r>1$ であり，
$x=(a-1)/r<a-1$ なので，高辺 $02$ の共通低 puncture $4$ でも
通常接着しか起こらない．従って $p^0_{24}=-x$，$p^0_{42}=-r$ であり，
補題~\ref{cl16:ordinary-bound} より $r\le4$ である．
$r=1$ の場合にも式~\eqref{cl16:small-and-unit} は成立する．
最後に $r>1$ なら低辺 $24$ の node 条件が
$p^4_{20}\le1-r<0$ を与え，低低 pair equation より $p^4_{02}=0$ となる．
従って
\[
 n_{02;4}=1-\frac{x-x}{a}=1,\qquad
 n_{04;2}=1-(-r+0)=r+1.
\]
これが式~\eqref{cl16:forced-ordinary} である．
\end{proof}

式~\eqref{cl16:small-and-unit} は，従来の
「$r\le4$ または $z=1$」を，同時に成立する二条件へ強める．
その証明の中心は，第三 core の低低 pair equation と，node 数の有限な予算である．
葉の分母や三価中心の未分類の整数模型に探索上限を置く必要はない．

% v16: proof-derived finite reduction for the non-B representation of leaf 4.
% This file is intended to follow the global z=1, r<=4 reduction.

\subsection{単位 entry の枝の恒等式と有限化}
\label{zu16:section}

本節では交差標識 $(A,B,C,D)=(0,1,4,3)$ を固定し，
$z=1$，$1\le r\le4$，$S>r$，$rT>(r+1)S$ を仮定する．
前節までの局所 chart，整数 clutching，node 非負性，対向係数の相殺を保持する．
葉 $4$ の ray $2$ を一高分母標準形の $C$ または $D$ として表せる場合を扱う．
分母三の標識重複によって $B$ としても表せる行列を除外するという意味ではない．

\begin{lemma}[単位 entry の二つの恒等式]
\label{zu16:identities}
$T=(r+1)y$，$x=sS/r$ と置くと，$y,x$ は正整数であり，
\begin{align}
 (r+2)sS-r(rs+1)y&=2r(s-1),\label{zu16:pair}\\
 F_r(s)S&=r(rs+1)E+2r(r+2)(s-1),\label{zu16:E}\\
 F_r(s)&=-rs^2+(4r+3)s-r,\qquad
 E=D^2_4=(r+2)y-\frac{s+r}{r}S.\notag
\end{align}
$1\le r\le4$，$s\in\ZZ_{>0}$ では $F_r(s)\ne0$ である．
さらに $E\ge-2$ なら $s\le9$ である．
\end{lemma}

\begin{proof}
$A=a+x-s$，$B=b-y+1$，$a=sS+1$，$b=(r+1)y+1$ を
$AB=ab-1$ に代入して $r$ を掛けると式~\eqref{zu16:pair} を得る．
それと $E=(r+2)y-(s+r)S/r$ から $y$ を消去すると
式~\eqref{zu16:E} が従う．
$s=1,2,3,4$ での $F_r(s)$ は順に
$2r+3,3r+6,2r+9,12-r$ で正である．
$s=5,6$ では $15-6r,18-13r$ であり，$r=1,2,3,4$ のどれでも零でない．
$s\ge7$ では $F_r(s)=s\{r(4-s)+3\}-r<0$ である．
最後に式~\eqref{zu16:pair} と $s\ge1$ から
\[
 y<\frac{r+2}{r^2}S,\qquad
 E<\left(4+\frac4r-s\right)\frac Sr
\]
となる．$s\ge10$ なら右辺は $-2$ より小さいので，$E\ge-2$ に反する．
\end{proof}

\begin{lemma}[小さい $S$ の完全な処理]
\label{zu16:small-S}
$S\le3$ の場合，接続可能性を残す中央模型は
\[
 (r,s,S,T)=(1,1,2,6),\ (1,4,3,12),\
 (2,2,3,6),\ (2,4,3,6),\ (2,6,3,6)
\]
の五つに限られる．
\end{lemma}

\begin{proof}
$S>r$ と式~\eqref{zu16:pair} から，
$(r,S)=(1,2)$ では $(s,y)=(1,3)$，
$(r,S)=(1,3)$ では $(s,y)=(4,6)$ だけとなる．
$(r,S)=(2,3)$ では $y=2$ が強制され，$s$ は任意の正の偶数である．
この共鳴族では $E=5-3s/2$，$D^2_3=s-2$ である．

後者で $s\ge8$ と仮定する．両葉で ray $2=B$ は node 条件に反する．
ray $2=C$ の場合は $u=th=1$ が必要だが，係数相殺が非零の $d$ を要求し，
$c=0$ となる．すると他方の中心への entry が零になり，
葉 $3$ では中央 row の $T=6$，葉 $4$ では $r=2$ と両立しない．
従って両葉で ray $2=D$ である．node 条件は $2\le v\le4$ を与え，
両葉の parameter $h$ は三以下となる．
葉 $3$ では $c_3=s-h_3-3>0$，
葉 $4$ では $c_4=E-h_4-1<0$ なので，両方で $d=0$ である．
中央の正の entries $T,r$ を打ち消すには，他方の中心の ray が $B$，
他の葉の ray が $C$ でなければならない．
従って $D^3_4=h_3>0$，$D^4_3=h_4>0$ となり，相殺できない．
よって共鳴族では $s=2,4,6$ だけを調べればよい．
\end{proof}

以下，$S\ge4$ とする．葉 $3$ では低辺 $23$ の node 条件により
$p^3_{20}\ge S-4\ge0$ であり，高低 pair product が正なので実際は正である．
従ってその ray $2$ は $D$ である．葉 $3$ の parameters を $h,t,c,d$ と書くと
\begin{equation}\label{zu16:leaf3}
 v=t(h+1),\quad c=s-h-3,\quad
 \alpha=v-S\in[-4,1]\cap\ZZ,\quad
 \beta=p^3_{21}+T\in[-4,1]\cap\ZZ .
\end{equation}
第三の node 和 $\gamma=-2-\alpha-\beta$ も $[-4,1]\cap\ZZ$ に入る．

\begin{lemma}[非 $B$ 葉の有限分岐]
\label{zu16:nonB-split}
葉 $4$ の ray $2$ を $C$ または $D$ として表すと，次のいずれかが成立する．
\begin{enumerate}
\item $-2\le E\le7$．
\item $r=1$，葉 $4$ の ray $3$ が $B$ であり，
\[
 D^4_3=E+\kappa,\qquad \kappa\in\{2,5,8,10\}.
\]
\end{enumerate}
\end{lemma}

\begin{proof}
葉 $4$ の parameters を $H,\tau,c_4,d_4$ とする．
ray $2=C$ なら，node 条件は $1\le u_4=\tau H\le3-r$ を与える．
$c_4\ne0$ なら $d_4=0$，係数相殺より $E=-H$ なので第一の場合となる．
$c_4=0$ なら，残る二低 rays の割当に応じて
$(B,D)=(0,3)$ または $(3,0)$ である．
前者では $r=u_4=H=\tau=1$ であり，$D^4_3=-2$，$p^4_{31}=2$ である．
低辺 $34$ の puncture $1$ より葉 $3$ の ray $4$ は $C$，
$c\ge th>0$ となる．従って $d=0$ であり，
低辺 $23$ は $T\le v+2\le S+3$ を要求する．
一方 $r=1$ の正枝条件は $T>2S$ なので，$S\ge4$ に反する．
後者では $r=1$，$(H,\tau)=(1,1),(1,2),(2,1)$ であり，
\[
 D^4_3=E+\tau(H^2+H+1)+H+1
\]
から $\kappa=5,8,10$ を得る．

ray $2=D$ なら $2\le v_4=\tau(H+1)\le r+2\le6$ である．
$d_4\ne0$ なら $c_4=0$，係数相殺から $E=H+1\le6$ となる．
$d_4=0$，$(B,C)=(0,3)$ の場合は $D^4_3=H$ である．
葉 $3$ の ray $4$ が $C$ なら係数相殺により $d=h+H>0$，従って $c=0$ である．
ray $4=B$ でも $p^3_{42}=0$ なので，いずれも低辺 $34$ の puncture $2$ から
$-4\le c_4=E-H-1\le1$，すなわち $H-3\le E\le H+2$ を得る．
$1\le H\le5$ だから第一の場合である．
最後に $(B,C)=(3,0)$ の場合は $p^4_{20}=0$ なので $r=1$ であり，
$(H,\tau)=(1,1),(2,1)$ に限られる．
\[
 D^4_3=E+\tau(H^2+H+1)-H
\]
から $\kappa=2,5$ を得る．
\end{proof}

\begin{lemma}[自由な係数を残す分岐の有限化]
\label{zu16:free-E}
補題~\ref{zu16:nonB-split} の第二の場合には，次の有限集合だけを調べればよい．
\begin{enumerate}
\item 葉 $3$ の ray $4=C$：
$th\le4$，$-4\le c\le1$，$s=h+3+c$，
$S=t(h+1)-\alpha$，$-4\le\alpha\le1$．
\item 葉 $3$ の ray $4=B$：
$s=h+3$，$1\le h\le12$，$1\le t\le6$，
$S=t(h+1)-\alpha$ であり，node 和 $\alpha,\beta,\gamma$ と $\kappa$ は
\begin{equation}\label{zu16:free-E-identity}
 t(h^2+7h+3)
 =\alpha(h^2+5h+7)+(h+4)(1+\beta+\kappa)-2h-4
\end{equation}
を満たす．
\end{enumerate}
各場合で $y$ は式~\eqref{zu16:pair} により一意に決まり，$T=2y$ である．
\end{lemma}

\begin{proof}
葉 $4$ の ray $3=B$ の row は，高 ray $1$ への entry が $-1$，
他の低 rays への entries が零である．
従って葉 $3$ の ray $4=C$ の場合，低辺 $34$ の punctures $0,2$ が
$th\le4$，$-4\le c\le1$ を与える．

ray $4=B$ の場合，葉 $3$ の ray $1=C$ なので
$d=p^3_{21}=-T+\beta<0$，従って $c=0$，$s=h+3$ である．
低辺 $34$ の係数相殺は
\[
 T=t(h^2+h+1)+1+\beta+E+\kappa
\]
となる．これに $E=3y-(h+4)S$，$T=2y$，
$S=t(h+1)-\alpha$ と式~\eqref{zu16:pair} を代入して
式~\eqref{zu16:free-E-identity} を得る．

$A=h^2+7h+3$ と置く．右辺は
$h^2+15h+51$ 以下であり，$7A-(h^2+15h+51)=6h^2+34h-30>0$
なので $t\le6$ である．
$\alpha\le0$ なら右辺は $10h+44$ 以下なので
$h^2-3h-41\le0$，従って $h\le8$ である．
$\alpha=1,t\ge2$ なら $h^2-h-45\le0$，従って $h\le7$ である．
残る $\alpha=t=1$ では，$K=1+\beta+\kappa$ と置くと
$(K-4)h+4K=0$ である．正整数解は $K=2,h=4$ または $K=3,h=12$
だけなので，すべての場合に $h\le12$ となる．
\end{proof}

\begin{proposition}[非 $B$ 表現の完全な有限候補集合]
\label{zu16:nonB-finite}
葉 $4$ の ray $2$ を $C$ または $D$ として表せる場合には，
中央模型は証明から定まる三十一個の候補のいずれかに属する．
これらの中央模型に低三角形 $234$ の全整数葉を接続すると，
次の十模型，計五十二組だけが残る：
\[
\begin{array}{rrrr|r}
 r&s&S&T&\#\text{低三角形}\\\hline
 1&4&3&12&21\\
 1&7&4&18&3\\
 1&4&8&36&4\\
 1&15&12&64&3\\
 2&4&3&6&4\\
 2&4&12&30&5\\
 2&5&14&36&1\\
 4&4&18&30&1\\
 4&4&35&60&8\\
 4&4&52&90&2
\end{array}
\]
表の復元数には ray $2=B$ と表せる葉も含めており，
同じ標識付き行列の組は一度だけ数えている．
\end{proposition}

\begin{proof}
補題~\ref{zu16:small-S} の五模型を取る．$S\ge4$ の有界係数分岐では，
$r=1,\ldots,4$，$s=1,\ldots,9$，$E=-2,\ldots,7$ を取り，
式~\eqref{zu16:E} から $S$，式~\eqref{zu16:pair} から $y$ を復元する．
それぞれの整数性と元の正枝条件を課す．自由係数分岐は
補題~\ref{zu16:free-E} の有限集合を用いる．
元の row sums，pair products，奇数分母，正の weights を再確認すると，
三十一個の異なる中央模型を得る．その生成 witness は六十六個である．

各中央模型について，一高分母の完全標準形の約数復元と ray $2=B$ の例外復元を
実行し，低辺 $23,24,34$ の三つずつの node 条件と係数相殺を再確認する．
この固定中央模型の復元は，node 区間と正約数に由来する有限手順である．
\nolinkurl{verify_two_leaf_unit_finiteness_v16.py} は全有限 trials，
中央模型と葉の完全行列，生成 witness を
\nolinkurl{two_leaf_unit_finiteness_certificate_v16.json} に保存する．
この計算に中央・葉・中心間の分母の任意の上限はない．
本命題が扱っていない ray $2=B$ のみの表現は，別に処理する必要がある．
\end{proof}

% v16. Complete treatment of the ray-2=B leaf at r=1,z=1.

\subsection{\texorpdfstring{$r=1$}{r=1} における例外的な \texorpdfstring{$B$}{B} 型葉の完全な有限化}
\label{rb16:section}

二葉三角形の交差標識を取り，$r=z=1$ とする．
葉 $4$ で ray $2$ が一高分母標準形の $B$ になる場合は，
この一辺だけを見ても葉の分母を抑えられない．
以下では，他方の葉と低三角形の条件を同時に使ってこの例外を尽くす．

\begin{proposition}[$B$ 型葉を持つ低三角形の完全表]
\label{rb16:complete-list}
交差した正の二高分母 core $2$ で $r=z=1$ とし，
葉 $4$ の ray $2$ が一高分母標準形の $B$ であるとする．
三つの低辺 $23,24,34$ の node 条件と $D$ の相殺を全て満たすなら，
中央 parameters $(s,S,T)$ は
\[
 (4,3,12),\quad(4,8,36),\quad(7,4,18),\quad(15,12,64)
\]
のいずれかである．標識付き行列の組として，低三角形はそれぞれ
三，二，一，一個，合計七個である．
\end{proposition}

\begin{proof}
$y=T/2$ とすると中央の pair equation は
\[
 3sS-(s+1)y=2s-2
\]
となる．従って，ある正整数 $k$ によって
\begin{equation}\label{rb16:middle-param}
 (s+1)k=3S-4,\qquad y=sk+2,\qquad T=2sk+4
\end{equation}
と書ける．また $S\ge2$，$T>2S$ なので，偶数性から $T\ge2S+2$ である．
低辺 $23$ の三つの node 区間は
\begin{equation}\label{rb16:leaf3-row}
 p^3_{20}\in[S-4,S+1],\quad
 p^3_{21}\in[-T-4,1-T],\quad
 p^3_{24}\in[T-S-3,T-S+2]
\end{equation}
を与える．後二者はそれぞれ厳密に負，厳密に正である．

葉 $3$ で ray $2=B$ は不可能である．ray $2=C$ なら
$p^3_{20}=-th$ より $S\le3$ であり，後二者の符号から $c\ne0,d=0$ となる．
係数相殺 $D^3_2=2-s$ は $h=2-s$，従って $s=h=1$ を与える．
式~\eqref{rb16:middle-param} によって唯一の中央模型
$(s,S,T)=(1,2,6)$ が残る．これを後の有限計算に含める．

従って残る場合は葉 $3$ の ray $2=D$ である．
式~\eqref{rb16:leaf3-row} の後二者の一つはその $d$ なので $d\ne0,c=0$ であり，
係数相殺から
\begin{equation}\label{rb16:h-def}
 h=s-3\ge1,\quad u=th,\quad v=t(h+1),\quad
 (h+4)k=3S-4,\quad T=2(h+3)k+4
\end{equation}
を得る．以下の二つの割当を尽くせばよい．

\smallskip\noindent
\textbf{I．葉 $3$ の ray $4=C$，ray $1=B$．}
この場合 $d>0$ である．低辺 $34$ の puncture $1$ において，
葉 $4$ の ray $3=D$ とすると $u-1+v_4>1$ となるので不可能である．
従って葉 $4$ の ray $3=C$，ray $0=D$ であり，その parameters を
$h_4,t_4,c_4,d_4$，$u_4=t_4h_4$ とする．同じ低辺の node 条件から
\[
 c_4-u\in[-4,1],\qquad u-1-u_4\in[-4,1].
\]
$c_4=0$ なら直ちに $u\le4$ である．
$c_4\ne0$ なら $d_4=0$ なので，係数相殺が $h_4=d-h$ を与える．
$d\ge T-S-3$，$h_4\le u_4\le u+3$ より
\[
 T\le S+h(t+1)+6.
\]
一方，$v\le S+1$ より $3t(h+1)\le(h+4)k+7$ である．
これらを式~\eqref{rb16:h-def} に代入すると
\begin{equation}\label{rb16:caseI-bound}
 (4h^2+15h+14)k\le3h^2+20h+10
\end{equation}
となる．$k\ge1$ から $(h-1)(h-4)\le0$，従って $1\le h\le4$ である．
この四値では右辺は左辺の $k$ の係数の二倍より小さいので $k=1$ である．
$3S=h+8$ の整数性から $(h,S)=(1,3)$ または $(4,4)$ に限られ，
$v\le S+1$ からそれぞれ $t\le2$ または $t=1$ である．
従ってこの場合にも $th\le4$ となる．
以上により場合 I の候補は
\begin{equation}\label{rb16:caseI-finite}
 1\le h\le4,\quad 1\le t\le\lfloor4/h\rfloor,\quad
 t(h+1)-1\le S\le t(h+1)+4
\end{equation}
という証明由来の有限集合に入る．

\smallskip\noindent
\textbf{II．葉 $3$ の ray $4=B$，ray $1=C$．}
この場合 $d<0$ であり，葉 $3$ の row $4$ は
$p^3_{40}=-1,p^3_{41}=p^3_{42}=0$ である．
従って低辺 $34$ の puncture $1$ は再び葉 $4$ の ray $3=D$ を排除し，
ray $3=C$，ray $0=D$ を強制する．node 条件は
\[
 1\le u_4=t_4h_4\le4,\qquad u_4-2\le c_4\le2
\]
を与える．$m_4=t_4h_4(h_4+1)$ として
\[
 C_4=m_4+t_4+h_4+c_4+2
\]
と置く．$u_4=1,2,3,4$ の全約数表示を取ると，$C_4$ の最大値は
それぞれ $8,13,20,29$ である．特に $C_4\le29$ となる．

ここで
\[
 \alpha=S-t(h+1)\in[-1,4],\qquad \beta=T+d\in[-4,1]
\]
と置く．低辺 $24,34$ の $D$ 相殺を加減すると
\begin{equation}\label{rb16:caseII-balance}
 th+\alpha(h+1)+k+2=C_4+\beta\le30
\end{equation}
を得る．実際，$E=-D^2_4=(s+1)S-3y$ とすると
$E+T=t(h^2+h+1)+C_4+\beta$ であり，
$E+T=(h+1)S+k+2$ を代入すればよい．
従って $\alpha\ge0$ なら $th\le27$，
$\alpha=-1,t\ge2$ なら $h(t-1)\le28$ である．
残る $\alpha=-1,t=1$ では $S=h$ なので
\[
 k=3-\frac{16}{h+4}.
\]
正整数性により $h+4=8$ または $16$，すなわち $h=4,12$ だけが残る．
これで場合 II も分母上限を仮定せず有限化された．

式~\eqref{rb16:caseI-finite} と場合 II の三つの有限集合に，
式~\eqref{rb16:middle-param}，$S\ge2$，$sS$ の偶数性を課す．
先の例外 $(1,2,6)$ も含めて，中央模型は三十三個となる．
それぞれに対し，葉 $3$ は補題~\ref{tl15:finite-leaves} の
$C/D$ 復元式と同じ node 区間・係数相殺から全て復元する．
固定した葉 $3$ に続く $B$ 型葉 $4$ も有限に復元できる．
実際，$E=-D^2_4,F=-D^3_4,N=p^3_{42}$，$L(H)=H(H+1)+1$ と置く．
葉 $4$ の ray $3=C$ で $d_4=0$ なら $H=F\ge1$ が固定され，
$-4\le t_4(H+1)^2-E+N\le1$ が $t_4$ を有限化する．
同じ ray 割当で $c_4=0$ なら $t_4L(H)+H+1=E+F$ が全候補を有限化する．
ray $3=D$ で $c_4=0$ なら $H=-1-F\ge1$ が固定され，
$-4\le t_4H^2-E+N\le1$ を用いる．
同じ割当で $d_4=0$ なら $t_4L(H)-H=E+F$ と
$t_4L(H)-H\ge H^2+1$ を用いる．
残る自由係数は $D^4_2=E$ によって一意に決まる．
この四場合は $c_4d_4=0$ を尽くし，両係数の符号にも制限を置かない．
元の三つの低辺の条件を再検査した結果は次の完全表である．
\[
\begin{array}{rrr|r|l}
 s&S&T&\#\text{行列組}&(q_{03},q_{14})\\\hline
 4&3&12&3&(3,3),(5,3),(5,3)\\
 4&8&36&2&(9,7),(9,5)\\
 7&4&18&1&(21,3)\\
 15&12&64&1&(157,3)
\end{array}
\]
同じ分母対 $(5,3)$ を持つ二組は，葉 $3$ の標識付き行列が異なる．
\nolinkurl{verify_r1_B_leaf_v16.py} は全生成範囲と棄却理由，
三十三模型，七組の全行列を記録し，row sums と pair products も再検査する．
\end{proof}

\begin{proposition}[$B$ 型葉の中心への延長障害]
\label{rb16:no-B-extension}
命題~\ref{rb16:complete-list} の七組は全て，中心 core $0$ へ延長できない．
この段階の矛盾には，未知分母 $q_{01}$，中心での $D$ 相殺，
または三高分母 core の一般分類は不要である．
\end{proposition}

\begin{proof}
$a=q_{02}=sS+1$，$Q=q_{03}$ とし，葉 $3$ の parameter を $h=s-3$ とする．
中心 $0$ の rays $2,3$ に関する pair equation と node 条件は
\begin{equation}\label{rb16:center-pair}
 \{a(1-n)+s\}\{Q(1-m)-h\}=aQ-1,
 \qquad n,m\in\{0,1,2,3,4,5\}
\end{equation}
を要求する．ここで $n=n_{02;3}$，$m=n_{03;2}$ である．
七組に現れる $(a,Q,h)$ を代入すると，
$(a,Q)=(13,3),(13,5),(29,21),(181,157)$ では全三十六組に解がない．
$(a,Q,h)=(33,9,1)$ の場合だけ $(n,m)=(0,0)$ が残る．

この最後の場合には，二つの葉 $4$ の行列で
$p^4_{02}=18$ または $16$ である．一方，高辺 $02$ の共通低 puncture $4$ では
$p^2_{04}=a-1=32$ なので，符号付き約数の分類から
\[
 (p^0_{24},p^0_{42})=(1,32)\quad\text{または}\quad(-32,-1)
\]
しかない．どちらも $p^0_{42}+p^4_{02}\in[-4,1]$ に反する．
以上で全七組を排除した．
\end{proof}

\subsection{二葉付き三角形の全排除}
\label{tg16:section}

\begin{lemma}[中心の既知三行の有限復元]
\label{tg16:known-rows}
cores $2,3,4$ の整数行列と分母を固定し，中心 $i\in\{0,1\}$ への延長を考える．
$j\in\{2,3,4\}$ に対して $q_j=q_{ij}$ と置く．
未知の中心行列の row $j$ は，三つの node 数を指定すれば
\begin{equation}\label{tg16:row-reconstruction}
 p^i_{jk}=q_j(1-n_{ij;k})-p^j_{ik},
 \qquad k\notin\{i,j\}
\end{equation}
によって一意に定まる．非負 node 数の総和は五なので，各 row の候補は
$\binom{7}{2}=21$ 個である．従って三行の全候補は $21^3$ 個であり，
\begin{equation}\label{tg16:known-pairs}
 p^i_{jk}p^i_{kj}=q_jq_k-1,\qquad
 \{j,k\}\subset\{2,3,4\}
\end{equation}
を全て満たす候補がなければ，中心 $i$ への延長は不可能である．
この必要条件には，中心間の分母 $q_{01}$，未知の第四行，
中心での $D$ の相殺，または正の weights の判定は現れない．
\end{lemma}

\begin{proof}
式~\eqref{tg16:row-reconstruction} は整数 clutching の node 式そのものである．
三つの非負整数の和が五となる組は二十一個である．
既知の反対側の row sum が $-q_j$ であるため，復元された row sum も
$q_j(3-5)+q_j=-q_j$ となる．
局所 pair equations は延長に必要なので，三つの既知 rays の間だけでも
式~\eqref{tg16:known-pairs} が必要である．
\end{proof}

\begin{proposition}[十模型の中心延長障害]
\label{tg16:finite-obstruction}
命題~\ref{zu16:nonB-finite} の十模型・五十二組は，全て中心 $0$ または $1$ へ延長できない．
各模型について，下表の中心では，低三角形の全ての組に対して
式~\eqref{tg16:known-pairs} を満たす三行候補が零である．
\[
\begin{array}{rrrr|r|c}
r&s&S&T&\#\text{低三角形}&\text{矛盾する中心}\\\hline
1&4&3&12&21&0\\
1&7&4&18&3&0\\
1&4&8&36&4&0\\
1&15&12&64&3&0\\
2&4&3&6&4&0\\
2&4&12&30&5&0\\
2&5&14&36&1&0\\
4&4&18&30&1&1\\
4&4&35&60&8&0\\
4&4&52&90&2&0
\end{array}
\]
\end{proposition}

\begin{proof}
五十二組の元の局所行和・pair products・$D$ 係数・正の cycle inequalities，
および低辺 $23,24,34$ の全 node と係数相殺を再確認した上で，
補題~\ref{tg16:known-rows} の二十一通りずつの全復元を行う．
各組で上表の中心における三つの pair products の少なくとも一つが不成立となる．
\nolinkurl{verify_two_leaf_centers_v16.py} は，葉を生成するコードの中心復元関数を使わず，
元の式~\eqref{tg16:row-reconstruction} から三行を構成し，
全 node 配分を整数演算で検査する．
入力の五十二組と検算コードの hash，各組の棄却数と必要条件を満たす候補数を
\nolinkurl{two_leaf_center_obstruction_v16.json} に記録する．
有限集合の完全性は前の二つの有限化の証明によるもので，
分母や未知行列 entries に数値上限を仮定したものではない．
\end{proof}

\begin{theorem}[二葉付き三角形 support の不存在]
\label{tg16:no-two-leaf-triangle}
五次の標準 $K_5$ branch graph において，各 core が正の weights を持つ
整数局所 chart であり，全共有辺が marked punctures の整数 clutching，
対向 $D$ の相殺，非負の node 条件を満たすとする．
高分母 support が
\[
 H=\{01,02,12,03,14\}
\]
となる配置は存在しない．分母の上限は不要であり，
三高分母 core の一般分類は用いない．
\end{theorem}

\begin{proof}
定理~\ref{tl15:remaining-positive} と定理~\ref{cl16:two-leaf-unit} により，
中央 core $2$ は交差標識の正の枝で $z=1$，$1\le r\le4$ を満たす．
一高分母標準形の葉 $4$ における ray $2$ は $B,C,D$ のいずれかである．
$B$ として表せる場合は node 条件から $r=1$ であり，
命題~\ref{rb16:complete-list} と命題~\ref{rb16:no-B-extension} に反する．
$C$ または $D$ として表せる場合は，命題~\ref{zu16:nonB-finite} と
命題~\ref{tg16:finite-obstruction} に反する．
分母三で複数の標準形表示を持つ行列も，この場合分けの少なくとも一方に含まれる．
従って全ての配置が排除される．
\end{proof}

命題~\ref{rb16:complete-list} の七組は，十模型の全葉を復元した五十二組にも含まれる．
これらの個数を足して五十九組と数えてはならない．
二つの有限化は，無限に見える葉の自由度を別々の条件から抑えるために必要である．

% v17 current exposition; original v16 source retained for verification.
\section{\texorpdfstring{$K_4$}{K4} 型 support における低 puncture の流れと例外辺}
\label{kh16:section}

高分母 support が $I=\{0,1,2,3\}$ 上の $K_4$ であり，
第五の core $4$ は全低分母である場合を考える．
三高分母 core の正符号 graph が完全マッチングであるとは仮定しない．
ここでは，hub に接する四辺の node と対向係数だけから，
六本の高辺の接着を二・四・六例外に帰着する．続く各節で三場合を全て排除する．

\subsection{通常辺と例外辺，および node による流れの発散}

\begin{definition}\label{kh16:notation}
$i\ne j\in I$ について
\[
 R_{ij}=p^i_{4j},\qquad X_{ij}=p^i_{j4},\qquad
 V_{ij}=p^4_{ij},\qquad N_{ij}=n_{i4;j}=1-R_{ij}-V_{ij}
\]
と置き，対角成分を零とする．各高辺では $R_{ij}X_{ij}=q_{ij}-1$ である．
共通低 puncture $4$ の整数 clutching により，
補題~\ref{dg15:hub-balance} の三択を reverse entries に書けば，
\begin{equation}\label{kh16:edge-options}
 (R_{ij},R_{ji})=(u,-u),\quad
 \{R_{ij},R_{ji}\}=\{1,q_{ij}-1\},\quad
 \{R_{ij},R_{ji}\}=\{-1,-(q_{ij}-1)\}.
\end{equation}
第一の場合を\emph{通常辺}，第二・第三の場合をそれぞれ正・負の\emph{例外辺}と呼ぶ．
これは $H$ の辺を低 puncture での接着に応じて分類するものであり，
局所三高分母行列の符号型の分類とは異なる．
\end{definition}

\begin{lemma}[通常辺の reverse bound と流れの式]
\label{kh16:flow}
上記の $K_4$ 型において，通常辺の $u$ は $1\leq |u|\leq4$ を満たす．
さらに各 core の行和，四辺 $i4$ の対向 $D$ の相殺，node 非負性から
\begin{equation}\label{kh16:node-budget}
 N_{ij}\geq0,\qquad \sum_{j\ne i}N_{ij}=5
\end{equation}
および
\begin{equation}\label{kh16:divergence}
 \sum_{j\ne i}(X_{ij}-R_{ji})=d_i,
 \qquad d_i:=\sum_{j\ne i}N_{ji}-5
\end{equation}
を得る．例外辺の左辺への寄与は零であり，通常辺の寄与は
\begin{equation}\label{kh16:edge-weight}
 X_{ij}-R_{ji}
 =\operatorname{sgn}(R_{ij})W_{ij},\qquad
 W_{ij}:=\frac{q_{ij}-1}{|R_{ij}|}+|R_{ij}|>0
\end{equation}
である．通常辺を $R_{ij}>0$ の向きに向けると，式~\eqref{kh16:divergence} は
この正の重み付き流れの発散を指定する．
\end{lemma}

\begin{proof}
通常辺で $R_{ij}=u\geq5$ とすれば，$N_{ij}\geq0$ から
$V_{ij}\leq1-u<0$ である．hub の pair product は $V_{ij}V_{ji}=0$ なので
$V_{ji}=0$，従って $N_{ji}=1+u\geq6$ となる．
一方，$R,V$ の行和はいずれも $-1$ だから
$\sum_jN_{ij}=3-(-1)-(-1)=5$ である．これは矛盾であり，負の $u$ は標識交換で同じである．

対向 $D$ の相殺は
\[
 \sum_{j\ne i}X_{ij}+\sum_{j\ne i}V_{ji}=-2
\]
と同値である．$V_{ji}=1-R_{ji}-N_{ji}$ を代入すると
式~\eqref{kh16:divergence} になる．通常辺では
$R_{ji}=-R_{ij}$，$X_{ij}=(q_{ij}-1)/R_{ij}$ であり，
例外辺では $X_{ij}=R_{ji}$ なので，残る主張も従う．
\end{proof}

\begin{proposition}[例外辺数と，通常の橋の分母上界]
\label{kh16:counts-and-bridges}
例外辺数は二，四，六のいずれかである．
通常辺だけからなる graph の橋 $ij$ は
\begin{equation}\label{kh16:bridge-bound}
 W_{ij}\leq10,\qquad q_{ij}\leq25
\end{equation}
を満たす．従って例外辺が四本なら，残る通常辺二本の分母はともに $25$ 以下である．
その通常辺二本が端点を共有しないなら，さらに両方とも $q_{ij}\leq7$ である．
\end{proposition}

\begin{proof}
四行の $R$ の和は $-4$ である．通常辺の両方向の和は零，
正・負の例外辺の和はそれぞれ $q_{ij},-q_{ij}$ だから
\begin{equation}\label{kh16:exception-balance}
 \sum_{\text{正例外辺}}q_{ij}-\sum_{\text{負例外辺}}q_{ij}=-4.
\end{equation}
各高分母は奇数なので，例外辺数は偶数であり，零ではない．

$\varnothing\ne S\subsetneq I$ とする．式~\eqref{kh16:node-budget} と $N_{ii}=0$ より
\[
 \sum_{i\in S}d_i\in
 \begin{cases}
 [-5,10],&|S|=1,\\
 [-10,10],&|S|=2,\\
 [-10,5],&|S|=3.
 \end{cases}
\]
一頂点への流入は他の三行から高々 $15$，二頂点への流入は全四行から高々 $20$ であり，
三頂点の場合は $\sum_i d_i=0$ と補集合を用いた．
通常 graph の橋を一本除き，その一方の連結成分を $S$ とすれば，
流れの式を $S$ 上で足した左辺は $\pm W_{ij}$ だけになる．従って $W_{ij}\leq10$ である．
$u=|R_{ij}|\in\{1,2,3,4\}$ なので
\[
 q_{ij}=1+u(W_{ij}-u)\leq1+u(10-u)\leq25.
\]
通常辺が二本だけなら両方とも橋である．さらに二本が端点を共有しないなら，
各辺の負の発散を持つ端点で $d_i=-W_{ij}\geq-5$ だから $W_{ij}\leq5$ となり，
$q_{ij}\leq1+\max_{1\leq u\leq4}u(5-u)=7$ を得る．
\end{proof}

\subsection{二例外の場合の完全な分母制約}

\begin{proposition}[二本の例外辺の分母]
\label{kh16:two-exceptions}
例外辺が二本なら，一本は正，一本は負であり，その分母の順序付き対は
\begin{equation}\label{kh16:two-exception-pairs}
 (q_+,q_-)=(3,7)\quad\text{または}\quad(5,9)
\end{equation}
に限られる．
\end{proposition}

\begin{proof}
式~\eqref{kh16:exception-balance} と $q_{ij}\geq3$ より，
符号は正負一本ずつで，$q_-=q_++4$ である．
$P=q_+-1$，$M=q_--1=P+4$ と置くと，$P,M$ は正の偶数で $M\geq6$ である．
負例外辺の大きい reverse $-M$ を持つ頂点を $v$ とする．

正例外辺が $v$ に接しないなら，$v$ の他の二つの reverse は通常辺で絶対値四以下だから，
行和から $M\leq9$，偶数性から $M\leq8$ である．
正例外辺が接し，$v$ での reverse が $1$ なら，残る通常 reverse を $u$ として
$-M+1+u=-1$，従って $M=u+2\leq6$ となる．

最後に，正例外辺も $v$ で大きい reverse $P$ を持つ場合を排除する．
第三の辺は通常辺で，行和によりその reverse は $3$ である．
正例外辺では $V$ の $v$ からの成分は $1-P-N\leq-1$ なので，相手からの成分は零である．
負例外辺ではその成分は $1+M-N\geq M-4\geq2$，
通常辺でも $1-3-N\leq-2$ なので，同じく相手からの成分は零である．
従って $V$ の第 $v$ 列は零であり，対向 $D$ の相殺は
$\sum_j X_{vj}=-2$ を要求する．ところが実際の和は
\[
 1-1+\frac{q_{v\ell}-1}{3}>0
\]
である．これは矛盾する．
以上より常に $M\leq8$，従って $P\in\{2,4\}$ となり，主張が従う．
\end{proof}

\subsection{全六辺が例外となる場合の排除}

\begin{theorem}[全例外の $K_4$ は不可能]
\label{kh16:not-all-exceptional}
$K_4$ 型高分母 support の六辺がすべて例外辺となることはない．
この排除には分母上限，三高分母 core の分類，正の weights の条件を用いない．
\end{theorem}

\begin{proof}
各例外辺の reverse は $\pm1$ と $\pm(q-1)$ であり，$q-1$ は正の偶数である．
単位 reverse の端から大きい reverse の端へ辺を向ける．
各行の和が奇数 $-1$ なので，単位 reverse の本数は各頂点で奇数である．
六本の向き付き辺の出次数は，標識交換により $(3,1,1,1)$ となる．
頂点 $0$ を source とすれば，他の三頂点は有向三角形をなす．

source の三つの単位 reverse の符号は一正二負である．
他の各頂点では，大きい reverse が二つと単位 reverse が一つあり，
行和 $-1$ から二つの大きい reverse の符号は反対でなければならない．
正の source 辺を $01$ とし，残る標識を交換すると，
ある正の偶数 $A,C,D\geq2$ を用いて全 reverse 行列は
\begin{equation}\label{kh16:all-exception-R}
 R=\begin{pmatrix}
 0&1&-1&-1\\
 A&0&1&-A-2\\
 -D-2&D&0&1\\
 -C&-1&C&0
 \end{pmatrix}
\end{equation}
となる．例えば，$1\to2\to3\to1$ を単位端から大きい端への三角形とすると，
その辺の符号は順に正，正，負であり，三つの非 source 行の和が
式~\eqref{kh16:all-exception-R} の係数差を与える．

全辺が例外なら，式~\eqref{kh16:divergence} の左辺は零である．
従って $N$ は行和・列和がともに五の非負整数行列である．
三つの正例外辺 $01,12,23$ では，大きい reverse の端の $V$ 成分が厳密に負なので，
hub の pair products により $V_{01}=V_{12}=V_{23}=0$ となる．
$V$ の行和と $N=\boldsymbol{1}_{i\ne j}-R-V$ の列和を解くと，二整数 $u,w$ により
\begin{equation}\label{kh16:all-exception-V}
 V=\begin{pmatrix}
 0&0&u&-1-u\\
 -A-2-u&0&0&A+1+u\\
 D+3-w&w-D-4&0&0\\
 C-1+w+u&2-w&-C-2-u&0
 \end{pmatrix}
\end{equation}
と書ける．対応する node 行列は
\begin{equation}\label{kh16:all-exception-N}
 N=\begin{pmatrix}
 0&0&2-u&3+u\\
 3+u&0&0&2-u\\
 w&5-w&0&0\\
 2-w-u&w&3+u&0
 \end{pmatrix}.
\end{equation}
従って
\begin{equation}\label{kh16:uw-range}
 -3\leq u\leq2,\qquad0\leq w\leq5,\qquad u+w\leq2.
\end{equation}
残る三つの hub pair products は
\begin{equation}\label{kh16:three-products}
 \begin{split}
 u(D+3-w)&=0,\\
 (-1-u)(C-1+w+u)&=0,\\
 (A+1+u)(2-w)&=0.
 \end{split}
\end{equation}

$u\ne0$ なら第一式と $D\geq2,w\leq5$ から $D=2,w=5$ となり，
式~\eqref{kh16:uw-range} から $u=-3$ である．第三式は $A=2$ を要求するが，
第二式は $2(C+1)=0$ となり矛盾する．
$u=0$ なら第三式から $w=2$ となり，第二式は $-(C+1)=0$ となる．
従って全ての場合が排除される．
\end{proof}

\begin{corollary}[$K_4$ の接着に対する中間の二場合]
\label{kh16:remaining-cases}
$K_4$ 型高分母 support が残るなら，例外辺は二本または四本である．
二本の場合，正・負の例外分母は $(3,7)$ または $(5,9)$ であり，
通常四辺の graph は $C_4$ または三角形に葉辺を一本付けた形である．
四本の場合，通常二辺の分母は二十五以下であり，二辺が端点を共有しないなら七以下である．
これらは必要条件であり，いずれの整数接着の存在も主張しない．
\end{corollary}

\begin{remark}[有限化の前に残る自由度]
通常辺上の流れは，発散を指定した有限 graph の流れである．
二例外の場合の通常 graph は連結で第一 Betti 数一を持つので，
固定した node 行列と通常 reverse に対して，重みの自由度は高々一つの cycle circulation である．
これは三高分母 core の残る pair equations と高辺の接着を解いたことを意味しない．
四例外の場合の通常分母上界だけでは，例外分母の上界は得られない．
続く定理~\ref{kt17:no-two-exceptions} は二例外を hub の流れだけで排除し，
定理~\ref{kf17:no-four-exceptions} は例外分母の無限方向を有限型にまとめて四例外を排除する．
\end{remark}

% v17: two exceptional K4 edges are impossible already at the low hub.
\subsection{二本の例外辺の全排除}
\label{kt17:section}

\begin{theorem}[二例外の低 hub 障害]
\label{kt17:no-two-exceptions}
定義~\ref{kh16:notation} の $K_4$ 型において，例外辺がちょうど二本となることはない．
この排除は各 core の整数行和，高辺の奇数分母と共通低 puncture での整数 clutching，
hub の低低 pair products，および hub に接する四辺の node 条件と対向 $D$ 相殺だけを用いる．
正の weights，三高分母 core の一般分類，高辺上の $D$ 相殺，分母の探索上限は不要である．
\end{theorem}

\begin{proof}
式~\eqref{kh16:exception-balance} により，例外辺は正負一本ずつである．
正例外の分母を $P+1$，負例外の分母を $M+1$ とすると，
\[
 P\ge2,\qquad M=P+4\ge6,
 \qquad P,M\in2\ZZ.
\]
負例外辺の大きい reverse $-M$ を持つ端を $v$ とする．
各 node は五以下なので，この辺の $v$ からの hub entry は
\[
 V_{vj}=1+M-N_{vj}\ge M-4\ge2.
\]
従って $V_{jv}=0$ である．また，通常辺の $R_{v\ell}>1$ に対しても
$V_{v\ell}=1-R_{v\ell}-N_{v\ell}<0$ なので，$V_{\ell v}=0$ となる．
以下では hub への対向 $D$ 相殺を
\begin{equation}\label{kt17:hub-D}
 \sum_{j\ne v}X_{vj}+\sum_{j\ne v}V_{jv}=-2
\end{equation}
の形で用いる．

\emph{正例外辺が $v$ に接しない場合．}
残る二辺は通常辺で，その reverse を $u,w$ とすると
\[
 u+w=M-1\ge5,\qquad 1\le |u|,|w|\le4.
\]
従って $u,w>0$ であり，一に等しいものは高々一つである．
$u,w\ge2$ なら，$v$ の hub 列はすべて零であり，
$X$ の和は $-1$ と二つの正整数の和なので一以上となる．
これは式~\eqref{kt17:hub-D} に反する．

一方が一なら，それを $u=1$ とする．その通常辺の $X$ は $q-1$ であり，
奇数 $q>1$ により二以上である．もう一方の $X$ は一以上であり，
負例外の $X$ は $-1$ だから，$\sum_jX_{vj}\ge2$ となる．
零とは限らない hub 列成分は $u=1$ の辺からのものだけで，
\[
 V_{\ell v}=1-(-1)-N_{\ell v}\ge-3.
\]
従って式~\eqref{kt17:hub-D} の左辺は $-1$ 以上であり，やはり矛盾する．

\emph{正例外辺も $v$ で大きい reverse $P$ を持つ場合．}
第三の辺は通常辺で，行和はその reverse を $3$ と定める．
正例外辺，負例外辺，通常辺のすべてで $v$ の hub 列成分が零となる．
しかし $X$ の和は
\[
 1-1+\frac{q-1}{3}>0
\]
であり，式~\eqref{kt17:hub-D} に反する．

\emph{正例外辺が $v$ で単位 reverse $1$ を持つ場合．}
第三の通常辺の reverse は $M-2$ である．通常辺の上界四から
$M=6,P=2$ となり，その reverse は $4$ である．
標識を $v=0$，正例外 $01$，負例外 $02$ と選ぶと
\[
 R_{01}=1,\quad R_{10}=2,\quad
 R_{02}=-6,\quad R_{20}=-1,\quad
 R_{03}=4,\quad R_{30}=-4.
\]
$V_{03}<0$ と $V_{02}>0$ から，それぞれ
\begin{equation}\label{kt17:forced-nodes}
 V_{30}=V_{20}=0,\qquad N_{30}=5,\qquad N_{20}=2
\end{equation}
を得る．一方，$R$ の第 $1$ 行の和は $-1$ なので
$R_{12}+R_{13}=-3$ である．少なくとも一方，例えば $R_{1j}$ は $-2$ 以下である．
辺 $1j$ は通常辺なので $R_{j1}\ge2$，従って $V_{j1}<0$ から $V_{1j}=0$ であり，
\[
 N_{1j}=1-R_{1j}\ge3,\qquad N_{10}\le2.
\]
式~\eqref{kt17:forced-nodes} と合わせると
\[
 d_0=N_{10}+N_{20}+N_{30}-5=N_{10}+2\le4.
\]
ところが頂点 $0$ に接する通常辺は $03$ だけなので，
式~\eqref{kh16:divergence} は
\[
 d_0=W_{03}=4+\frac{q_{03}-1}{4}\ge5
\]
を要求する．最後の不等式は $(q_{03}-1)/4$ が正整数であることによる．
以上で正例外辺の全ての位置と向きを尽くした．
\end{proof}

この証明は命題~\ref{kh16:two-exceptions} の二つの分母対を先に列挙する必要もない．
全低 hub に接する条件だけが矛盾するので，残る高高 entries や
通常閉路の circulation をどのように選んでも延長できない．
独立の有限点検として，正例外を $R_{01}=1,R_{10}=P$ に固定した全 reverse 行列と
全 node 配分を \nolinkurl{verify_k4_two_exceptions_v17.py} で復元する．
既知の $P=2,4$ に対して行和を満たす reverse 行列は五十八個，
hub の低低 pair products まで通過する組は九百五十四個である．
その各組には通常辺だけを外向きに持つ頂点 $i$ があり，
\[
 d_i<\sum_{i\to j}\ell(|R_{ij}|),\qquad
 \ell(u)=\begin{cases}u+2,&u\text{ が奇数},\\u+1,&u\text{ が偶数}\end{cases}
\]
となる．右辺は奇数分母の通常 flow の下界なので，この有限点検も全組を棄却する．
一般の排除は上の三場合の証明によっており，有限点検はその補助である．

\subsection{四例外の場合の有限型帰着と全排除}
\label{kf17:section}

\begin{lemma}[例外 reverse の行ごとの有限型]
\label{kf17:finite-row-types}
$K_4$ の六辺のうち四辺が例外辺であるとする．通常二辺を固定し，
それぞれの reverse $u$ を $\{-4,-3,-2,-1,1,2,3,4\}$ から選ぶ．
各例外辺の符号と単位 reverse の端を指定すれば，残る reverse の絶対値は
各行独立の線形方程式で決まる．
その全解は有限個の確定値と，ある正の偶数 $h$ に対する
\begin{equation}\label{kf17:free-row}
 m_1=h+\max(\delta,0),\qquad
 m_2=h+\max(-\delta,0),\qquad \delta\in2\ZZ
\end{equation}
の組で表せる．この自由な組は異なる二本の例外辺の大きい reverse が
同じ頂点で反対符号になる場合に限られる．

hub の row/pair 条件，四つの低辺の node 条件と $D$ 相殺に関しては，
各 $h$ を $2$，$4$，または全偶数 $h\ge6$ の三型に分ければ十分である．
第三型は $h=6$ で代表させてよく，これは分母の探索上限ではない．
\end{lemma}

\begin{proof}
各例外辺には絶対値一の reverse と，正の偶数 $m=q-1$ の reverse が一つずつある．
$m$ が現れるのは，その大きい端に対応する $R$ の一行だけである．
頂点 $i$ での大きい端の集合を $B_i$ と書くと，行和は
\begin{equation}\label{kf17:row-equation}
 \sum_{j\in B_i}\varepsilon_{ij}m_{ij}
 =c_i:=-1-\sum_{j\notin B_i}R_{ij},
 \qquad m_{ij}\in2\ZZ_{>0}
\end{equation}
となる．右辺には，既に選んだ通常 reverse と単位 reverse しか含まれない．

$|B_i|=3$ なら，その頂点は通常辺にも単位端にも接しないため $c_i=-1$ であり，
左辺の偶数性に反する．$|B_i|=0$ なら $c_i=0$ が必要である．
$|B_i|=1$ なら $m=\varepsilon c_i$ が正の偶数であることが必要十分である．
$|B_i|=2$ で符号が同じなら，二つの正偶数の和が固定されるので全解は有限である．
符号が異なるなら，$m_1-m_2=\delta$ を解いたものが
式~\eqref{kf17:free-row} である．各例外辺の大きい端は一つしかないため，
異なる行の自由変数の間に未処理の同一視はない．

残る無限部分 $h\ge6$ を考える．どちらの大きい reverse も絶対値六以上である．
対応する hub entry は
\[
 V_{ij}=1-R_{ij}-N_{ij},\qquad0\le N_{ij}\le5
\]
だから，正の大きい reverse なら厳密に負，負の大きい reverse なら厳密に正である．
従って $V_{ij}V_{ji}=0$ は，単位端の $V$ が零であることと同値であり，
その node は正例外なら零，負例外なら二に固定される．
この条件は $h\ge6$ の値に依存しない．

$R$ の行和は全 $h$ で $-1$ であるから，node の行和五を課した後の
$V$ の行和も全 $h$ で $-1$ となる．また例外辺は
式~\eqref{kh16:divergence} の左辺に寄与しない．
通常辺の流れは $N$ だけから決まるため，四つの $D$ 相殺も $h$ に依存しない．
よって無限部分は各行で $h=6$ とした代表と，ここで課す hub 条件について完全に同値である．
\end{proof}

\begin{proposition}[四例外で残る hub の全候補]
\label{kf17:complete-hub-list}
四例外の場合，通常二辺が端点を共有しない配置は存在しない．
通常二辺が端点を共有する場合，標識を取り直して $01,02$ を通常辺とすると，
hub の row/pair 条件，四つの低辺の node 条件と $D$ 相殺を満たす候補は六個である．
これらは $1,2$ の交換の下で，次の三型に尽くされる．
全ての自由な組は有限型 $h=2$ に属し，無限型 $h\ge6$ は残らない．
\[
\begin{split}
 R^{(1)}&=\begin{pmatrix}
 0&-4&1&2\\4&0&-4&-1\\-1&-1&0&1\\1&-4&2&0
 \end{pmatrix},\qquad
 R^{(2)}=\begin{pmatrix}
 0&-4&2&1\\4&0&-4&-1\\-2&-1&0&2\\2&-4&1&0
 \end{pmatrix},\\[.4em]
 R^{(3)}&=\begin{pmatrix}
 0&-1&4&-4\\1&0&2&-4\\-4&1&0&2\\-1&-1&1&0
 \end{pmatrix}.
\end{split}
\]
対応する node 行列と分母は
\[
 N^{(1)}=N^{(2)}=\begin{pmatrix}
 0&5&0&0\\0&0&0&5\\3&2&0&0\\0&5&0&0
 \end{pmatrix},\qquad
 N^{(3)}=\begin{pmatrix}
 0&5&0&0\\0&0&0&5\\5&0&0&0\\2&3&0&0
 \end{pmatrix},
\]
\[
\begin{array}{c|rrrrrr}
 \text{型}&q_{01}&q_{02}&q_{03}&q_{12}&q_{13}&q_{23}\\\hline
 1&13&5&3&5&5&3\\
 2&13&7&3&5&5&3\\
 3&3&5&5&3&5&3
\end{array}
\]
である．$V$ は $V_{ij}=1-R_{ij}-N_{ij}$ で復元される．
ここで数えるのは hub の必要条件を満たすデータであり，
高 core の局所 chart が存在するという主張は含まない．
\end{proposition}

\begin{proof}
二本の通常辺は，標識交換後，$\{01,23\}$ または $\{01,02\}$ である．
通常 reverse の全 $8^2$ choices，四例外の大きい端の全 $2^4$ choices，
符号の全 $2^4$ choices を取り，補題~\ref{kf17:finite-row-types} の行方程式を解く．
同符号の正偶数分割は全て取り，異符号の自由変数は $h=2,4,6$ の三型で代表させる．
行和を満たす $R$ の有限型は，端点を共有しない場合に百十二，
共有する場合に二百四となる．そのうち無限部分を代表するものはそれぞれ三十二，八十二である．

各 $R$ に対し，三つの非負整数の和が五となる二十一通りの全 node 行を用い，
四行の全組に $V_{ij}V_{ji}=0$ を課す．通常辺は forest なので，
流れは $d_i=\sum_jN_{ji}-5$ から一意に決まる．
通常辺を小さい標識から大きい標識への向きで記すと，
$\{01,23\}$ の場合は
\[
 F_{01}=d_0,\quad F_{23}=d_2,\qquad d_0+d_1=d_2+d_3=0,
\]
$\{01,02\}$ の場合は
\[
 F_{01}=-d_1,\quad F_{02}=-d_2,\qquad d_0+d_1+d_2=d_3=0
\]
が必要十分である．各通常辺では
\begin{equation}\label{kf17:recover-q}
 X_{ij}=F_{ij}-R_{ij},\qquad
 q_{ij}=1+R_{ij}(F_{ij}-R_{ij})
\end{equation}
が奇数かつ三以上でなければならない．例外辺の分母は
$q_{ij}=1+R_{ij}R_{ji}$ である．全検査の個数は次のとおりである．
\[
\begin{array}{c|rr}
 &\text{端点非共有}&\text{端点共有}\\\hline
 R\text{ の有限型}&112&204\\
 \text{うち無限部分の代表}&32&82\\
 \text{hub pair を満たす }(R,N)&3552&4922\\
 \text{forest 成分の流れが釣り合う }(R,N)&320&558\\
 \text{正の奇分母を持つ候補}&0&6\\
 \text{うち無限部分の代表}&0&0
\end{array}
\]
六個は上に表示した三型とその $1,2$ 交換である．
\nolinkurl{verify_k4_four_exceptions_v17.py} は有限型を全て生成し，
元の行和，hub pair，node，四辺の $D$ 相殺を整数演算で再検査する．
全 $R$ 型と各検査段階の候補数，六候補の全行列を
\nolinkurl{k4_four_exceptions_verification_v17.json} に保存する．
無限型の置換の完全性は前補題によって証明されており，
分母十三という上界は残った候補の結果であって探索の仮定ではない．
\end{proof}

\begin{lemma}[残存候補に共通する実数局所障害]
\label{kf17:negative-discriminant}
低 ray を一つ持つ三高分母 core で，高 rays を $a,b,c$ とし，
\[
 (q_a,q_b,q_c)=(5,3,3),\qquad
 (p_{4a},p_{4b},p_{4c})=(-4,1,2),\qquad
 (p_{a4},p_{b4},p_{c4})=(-1,2,1)
\]
を指定する．これを満たす局所行列は，整数の条件を外しても存在しない．
\end{lemma}

\begin{proof}
三つの高 rows で高高 entries の和はそれぞれ $-4,-5,-4$ である．
$A=p_{ab}$ と置くと pair products と行和から
\[
 p_{ba}=\frac{14}{A},\quad
 p_{ac}=-4-A,\quad
 p_{ca}=\frac{14}{-4-A},\quad
 p_{bc}=-5-\frac{14}{A},\quad
 p_{cb}=-4-\frac{14}{-4-A}
\]
となる．ここで $A\ne0,-4$ である．残る pair product
$p_{bc}p_{cb}=3\cdot3-1=8$ は，分母を払うと
\[
 (-5A-14)(4A+2)=8A(-4-A),
 \qquad 6A^2+17A+14=0
\]
を与える．判別式は $17^2-4\cdot6\cdot14=-47<0$ なので実数解も存在しない．
\end{proof}

\begin{theorem}[四例外の $K_4$ は不可能]
\label{kf17:no-four-exceptions}
五次の標準 $K_5$ branch graph で，高分母 support が $K_4$ となる配置において，
四本の高辺が例外接着となることはない．
この排除には，分母の上限，三高分母 chart の一般分類，正の weights，
または高辺における $D$ 相殺を用いない．
\end{theorem}

\begin{proof}
命題~\ref{kf17:complete-hub-list} の第一・第二型では core $3$，
第三型では core $1$ において，ray 標識を取り直すと
補題~\ref{kf17:negative-discriminant} の低 row と column が現れる．
$1,2$ を交換した三候補も同じである．従って全ての hub 候補が排除される．
\end{proof}

% v17: combine the three complete exceptional-edge obstructions.
\subsection{四頂点支持の全排除と全低分母 core からの伝播}
\label{kg17:section}

\begin{theorem}[$K_4$ 型高分母 support の全排除]
\label{kg17:no-K4}
五次の標準 $K_5$ branch graph において，高分母 support
$H=\{ij:q_{ij}>1\}$ が四頂点上の $K_4$ と孤立頂点一つからなることはない．
この結論は，整数の局所行和と pair products，高辺の奇数分母，
共通低 puncture における整数 clutching，および全低 hub に接する
四辺の node 非負性と対向 $D$ 相殺から従う．
正の weights，高辺の $D$ 相殺，分母の上限，三高分母 chart の一般分類は必要ない．
\end{theorem}

\begin{proof}
命題~\ref{kh16:counts-and-bridges} により，例外辺数は二，四，六のいずれかである．
二本の場合は定理~\ref{kt17:no-two-exceptions}，四本の場合は
定理~\ref{kf17:no-four-exceptions}，六本の場合は
定理~\ref{kh16:not-all-exceptional} により，それぞれ排除される．
三つの証明は上記の条件だけを用いるため，これで全場合を尽くす．
\end{proof}

\begin{corollary}[一つの全低分母 core からの大域伝播]
\label{kg17:low-core-propagation}
正の局所整数 charts，marked puncture の整数 clutching，対向 $D$ の相殺，
および node 非負性を満たす五次の標準 $K_5$ branch graph を考える．
ある core $i$ において $q_{ij}=1$ が全ての $j\ne i$ で成り立つなら，
全ての辺について $q_{jk}=1$ が成り立つ．
同値に，非空の高分母 support は全五頂点を使う連結 graph である．
\end{corollary}

\begin{proof}
系~\ref{fo13:remaining-support} により，高分母の辺を持つ連結成分は一つであり，
少なくとも四頂点を使い，木ではない．全低分母 core があれば，この成分は
ちょうど残る四頂点を使う．四頂点の連結 graph で四辺のものは
$C_4$ または三角形に葉辺を付けた形であり，前者は
系~\ref{cf13:degree-two-rigidity}，後者は定理~\ref{tn14:no-degree-two-triangle}
に反する．五辺の場合は diamond なので定理~\ref{dg15:no-diamond} に反し，
六辺の場合は定理~\ref{kg17:no-K4} に反する．従って高分母の辺は存在しない．
\end{proof}

この伝播の仮定は固定した全平面五次の branch graph と作用適合整数 charts に関する．
任意の Lagrangian torus fibration の分母剛性を述べるものではない．
また，定理~\ref{kg17:no-K4} は support 全体が四頂点の $K_4$ となる場合を排除する．
五頂点の support が $K_4$ を部分 graph として含む場合には，
全低 hub が存在しないので，この定理をそのまま適用できない．

% v18. Exact exclusion of G_{5,5}; no denominator cutoff.
% Uses standard positive two-high form, no-negative-low-wedge,
% common-low reverse bound, and low-edge D opposition only at 23 and 24.

\section{五頂点六辺の house 型 support の排除}
\label{hs18:section}

三角形 $012$ と四角形 $0134$ が辺 $01$ を共有する graph
\[
 H=\{01,02,03,12,14,34\}=G_{5,5}
\]
を扱う．二価 cores は $2,3,4$ であり，従来の四頂点障害とは異なり，
全五頂点を使用する．外側の二価 cores $3,4$ の高隣接頂点の間は低辺なので，
補題~\ref{kb15:no-negative-low-wedge} により，この二つは正の枝に限られる．
中央の core $2$ は三角形上にあり，その負の枝をあらかじめ捨てない．

\subsection{共有高辺の共通低 ray の三つの型}

二高分母の canonical 標識を $(A,B,C,D)$ とし，
\[
 P=\begin{pmatrix}
 0&-A_0&x&-s\\
 -B_0&0&-y&z\\
 r&-r-1&0&0\\
 -S&T&-L&0
 \end{pmatrix},\qquad
 L=T-S+1
\]
と書く．高分母は $a=sS+1,b=zT+1$，
$x=(a-1)/r,y=(b-1)/(r+1)$ である．
正の枝では
\begin{equation}\label{hs18:positive-bounds}
 S>r\ge1,\quad S\ge2,\quad T\ge S+2,\quad T\ge5,\quad L\ge3.
\end{equation}
中央の負の枝を含む全正局所模型についても $S\ge2,T\ge S$ は成り立つ．
負の枝の $S\ge2$ は式~\eqref{kb15:E-bounds} の $w\ge2$，
$T\ge S$ は補題~\ref{td12:low-entry} による．
以下の $C,D$ は canonical な低 rays の名称であり，出発 core の番号とは別である．

core $3,4$ における ray $2$ は低 ray であり，それぞれ $C$ または $D$ である．
共有高 ray の分母を $q=q_{34}$ とすると，その ray から puncture $2$ へ向かう
entry と reverse entry は次表で尽くされる．
\begin{center}
\begin{tabular}{cccc}
ray $2$ & 共有高 ray & outgoing entry & reverse entry\\\hline
$C$ & $A$ & $x>0$ & $r>0$\\
$C$ & $B$ & $-y<0$ & $-r-1<0$\\
$D$ & $A$ & $-s<0$ & $-S<0$\\
$D$ & $B$ & $z>0$ & $T>0$
\end{tabular}
\end{center}
負の entries $-y,-s$ は $q-1$ の真約数の負であり，$z$ も真約数である．
また $x\ge2$ であり，$x=q-1$ と $r=1$ は同値である．
これらは式~\eqref{hs18:positive-bounds} と pair products から従う．

\begin{lemma}[共有辺 $34$ に残る低 ray の型]
\label{hs18:shared-types}
正局所 charts，整数 clutching，および非負 nodes を仮定する．
共有高辺 $34$ の puncture $2$ においては，次のいずれかだけが残る．
\begin{enumerate}
\item 両 core で ray $2=C$ であり，通常の異符号接着である．
一方の共有高 ray は $A$，他方は $B$ である．
\item 一方の core $i$ で ray $2=D$，他方 $j$ で ray $2=C$ であり，
通常接着である．共有高 ray は両方で $A$ であり，
\[
 r_j=S_i,\qquad 2\le S_i\le4.
\]
\item 一方の core $i$ で ray $2=D$，他方 $j$ で ray $2=C$ であり，
正の例外接着である．共有高 ray は $i$ 側で $B$，$j$ 側で $A$ であり，
\[
 z_i=1,\quad T_i=q-1,\quad r_j=1,\quad x_j=q-1.
\]
\end{enumerate}
両 core で ray $2=D$ となる場合はない．
\end{lemma}

\begin{proof}
共通低 puncture の二つの outgoing entries は $q-1$ の signed divisors である．
整数 clutching により，異符号なら和が零となり，同符号なら絶対値の組は
$\{1,q-1\}$ となる（補題~\ref{ls11:low-puncture}）．
負の entries はいずれも真約数なので，負の例外は存在しない．

両 ray $2=C$ なら，正の entries はともに $x\ge2$ なので正例外も存在せず，第一の場合となる．
両 ray $2=D$ なら，正の entries はともに真約数 $z$ なので正例外はない．
異符号接着では一方の reverse が $T\ge5$ となり，
共通低 puncture の通常 reverse 上界四（補題~\ref{cl16:ordinary-bound}）に反する．

$C,D$ が一つずつの場合，通常接着で $C$ 側が $B$，$D$ 側が $B$ だと，
やはり $T\ge5$ が上界四に反する．従って通常接着では両共有 rays が $A$ となり，
reverse の相殺から $r_j=S_i\le4$ を得る．
正例外では $C$ 側が $A$，$D$ 側が $B$ である．
$z_i$ は真約数なので，最大約数は $x_j$ に限られ，第三の場合を得る．
\end{proof}

\subsection{node 条件による唯一の枝への帰着}

graph の自己同型 $(0\ 1)(3\ 4)$ により，中央 core $2$ の canonical 高 rays を
$A=0,B=1$ としてよい．低 rays は $(C,D)=(3,4)$ または $(4,3)$ である．
その parameters を添字 $2$ で表す．低辺に沿う node 式は
\begin{equation}\label{hs18:low-node}
 n_{2i;k}=1-p^2_{ik}-p^i_{2k}\in\{0,1,2,3,4,5\}.
\end{equation}
特に canonical $D$ row の低低 entry は $-L$ であり，零と置いてはいけない．

\begin{lemma}[node 条件を通過する唯一の中央族]
\label{hs18:node-family}
前補題の仮定の下で，$G_{5,5}$ が接着可能なら，標識交換後に次が必要である：
\[
 \begin{array}{c|c|ccccc}
 \text{core}&(A,B,C,D)&r&s&z&S&T\\\hline
 2&(0,1,4,3)&1&2t-1&t&2&6\\
 3&(0,4,1,2)&1&1&1&2&6
 \end{array}\qquad (t\ge1).
\]
従って
\[
 (q_{02},q_{12},q_{03},q_{34})=(4t-1,6t+1,3,7).
\]
core $4$ の canonical 標識は $(A,B,C,D)=(3,1,2,0)$ であり，
その $a=q_{34}=7,r=1$ が固定される．
\end{lemma}

\begin{proof}
補題~\ref{hs18:shared-types} の三場合を順に調べる．

\paragraph{両外側で ray $2=C$ の場合．}
このとき $p^3_{21}=p^4_{20}=0$ である．
中央が $(C,D)=(4,3)$ なら $n_{23;1}=1-T_2<0$ となる．
中央が $(C,D)=(3,4)$ の場合，core $3$ の共有 ray が $A$ なら，
通常接着から $r_3=r_4+1\ge2$ なので
$n_{23;4}=1-r_3<0$ となる．共有 ray が $B$ なら，
core $3$ のもう一つの高 ray $0$ は $A$ なので
$n_{23;0}=1-r_2-r_3<0$ となる．従って第一の場合は存在しない．

\paragraph{$C,D$ 間の通常接着．}
$D$ 側を $i$，$C$ 側を $j$ とすると，reverse はそれぞれ $-S_i,S_i$ である．
$S_i\ge2$ なので，中央の低 $C$ が $j$ であると
$p^2_{ji}=0$ から $n_{2j;i}=1-S_i<0$ となる．
従って中央の低 $C$ は $i$，低 $D$ は $j$ である．

$i=3,j=4$ なら，core $3$ の共有 ray $4$ が $A$，もう一つの高 ray $0$ が $B$ なので
$n_{23;0}=1-r_2-T_3<0$ となる．
$i=4,j=3$ なら，core $3$ の共有 ray $4$ が $A$，もう一つの高 ray $0$ が $B$ であり，
\[
 p^2_{30}=-S_2,\qquad p^3_{20}=-(S_4+1).
\]
従って
$n_{23;0}=S_2+S_4+2\ge6$ となって上界五に反する．
これで第二の場合も排除された．

\paragraph{$C,D$ 間の正例外接着．}
$D$ 側を $i$，$C$ 側を $j$ とすると，reverse はそれぞれ $T_i\ge5,1$ である．
中央の低 $C$ が $i$ だと $n_{2i;j}=1-T_i<0$ なので，
中央の低 $D$ が $i$，低 $C$ が $j$ である．

$i=4,j=3$ なら，中央の $p^2_{40}=-S_2$ と
core $4$ の $p^4_{20}=-L_4$ により
$n_{24;0}=1+S_2+L_4\ge6$ となる．
従って $i=3,j=4$ であり，canonical 標識は
\[
 (A,B,C,D)_2=(0,1,4,3),\qquad
 (A,B,C,D)_3=(0,4,1,2),\qquad
 (A,B,C,D)_4=(3,1,2,0)
\]
に限られる．ここでは
\[
 n_{23;0}=1+S_2+S_3\le5
\]
なので $S_2=S_3=2$ である．残る二つの node 数は
\[
 n_{23;1}=T_3-T_2,\qquad n_{23;4}=T_2-T_3
\]
となるから $T_2=T_3$ である．

外側の正の枝から $r_3=1$，例外接着から $z_3=1$ である．
$S_3=2$ を高高 pair equation に代入すると
\[
 T_3=8-\frac4{s_3+1},\qquad s_3\ge1,\qquad 2\mid T_3.
\]
$s_3+1$ は四の正約数で二以上だから，唯一の可能性は
$s_3=1,T_3=6$ である．中央も $S_2=2,T_2=6$ なので負の枝ではなく，
$r_2=1$ である．中央の高高 pair equation は $s_2=2z_2-1$ となる．
$t=z_2$ と置くと表示した族を得る．共有分母は $q_{34}=z_3T_3+1=7$ であり，
core $4$ は正例外の最大約数端なので $r_4=1$ である．
\end{proof}

\subsection{二本の低辺の係数相殺による終結}

\begin{lemma}[固定した $a=7,r=1$ の完全な七模型]
\label{hs18:fixed-seven}
二高分母の正の canonical 整数模型で $a=7,r=1$ とする．
全ての可能性は次の七つであり，いずれも $D_C>0$ である．
\begin{center}
\begin{tabular}{rrrrr|r}
$b$&$s$&$z$&$S$&$T$&$D_C$\\\hline
13&3&2&2&6&4\\
19&1&1&6&18&15\\
19&2&2&3&9&9\\
31&1&2&6&15&18\\
41&2&5&3&8&19\\
43&1&3&6&14&23\\
79&1&6&6&13&40
\end{tabular}
\end{center}
これは相手分母 $b$ に上界を置く分類ではない．
\end{lemma}

\begin{proof}
$s\mid6$ なので $s=1,2,3,6$ であり，$x=6,S=6/s$ である．
pair-product の因数分解は
\[
 z\{2(13-s)-(s+1)T\}=2(s-7),\qquad b=zT+1
\]
である．$s=1,2,3,6$ の場合はそれぞれ
\[
 z(T-12)=6,\qquad z(3T-22)=10,\qquad
 z(T-5)=2,\qquad 7z(T-2)=2.
\]
最後は不可能である．前の三つについて，右辺の全正約数 $z$ を取り，
$T$ の整数性と $2\mid zT$ を課すと，表示した七模型をちょうど得る．
それぞれ元の row sums，pair products，正値性を満たす．
$D_C=T-S-x+y$，$y=zT/2$ から最後の列を計算する．
\end{proof}

\begin{theorem}[$G_{5,5}$ support の不存在]
\label{hs18:no-house}
五次の標準 $K_5$ branch graph 上で，正の局所整数 charts，marked puncture の
整数 clutching，非負 nodes，および対向 $D$ の相殺を仮定する．
高分母 support は $G_{5,5}=\{01,02,03,12,14,34\}$ ではあり得ない．
証明は分母の上界も三高分母の一般分類も使わず，$D$ の相殺は低辺 $23,24$ だけに用いる．
\end{theorem}

\begin{proof}
補題~\ref{hs18:node-family} の唯一の族について，低 rays の係数は
\[
 D^2_3=s_2-z_2-1=t-2,\qquad D^2_4=6-t,\qquad
 D^3_2=s_3-z_3-1=-1
\]
である．低辺 $23$ の相殺は $t-2-1=0$，従って $t=3$ を要求する．
低辺 $24$ の相殺は，core $4$ の ray $2=C$ に対し
\[
 D^4_2=-(6-t)=-3
\]
を要求する．一方，core $4$ は $a=7,r=1$ なので，
補題~\ref{hs18:fixed-seven} によって $D^4_2=D_C>0$ である．
これが最後の矛盾である．
\end{proof}

唯一残る族は低三角形 $234$ の node 条件から得られた必要な部分データであり，
未知の三価 cores $0,1$ の整数行列や $q_{01}$ の存在を主張しない．
実際，二本の低辺で既に矛盾するので，それらの未知データを復元する必要はない．
本定理は現在の標準 cones と branch graph を保つ算術的接着の障害である．
一般の Lagrangian torus fibrations に対する不存在定理には拡張していない．

% v18. G_{5,10}; no denominator cutoff. Depends on mh18:matching.
\section{\texorpdfstring{$G_{5,10}$}{G5,10} 型の完全排除}
\label{gten18:section}

\begin{theorem}[$G_{5,10}$ の障害]
\label{gten18:no-support}
五次の算術的接着データの高分母 support は
\[
 G_{5,10}=\{01,02,03,12,13,24,34\}
\]
と同型ではあり得ない．分母の上限は仮定しない．
証明には正の局所整数データ，marked clutching，非負 nodes，
対向 $D$ 相殺を用いる．
\end{theorem}

\subsection{二価 core と二つの低 row}

頂点 $4$ の高分母近傍 $2,3$ の間は低辺である．従って
補題~\ref{kb15:no-negative-low-wedge} により，core $4$ は正の二高分母枝に属する．
graph の自己同型 $0\leftrightarrow1$，$2\leftrightarrow3$ を使って
canonical な標識を $(A,B,C,D)=(2,3,0,1)$ にそろえる．その行列は
\begin{equation}\label{gten18:central-matrix}
 P^4=
 \begin{pmatrix}
 0&-A&x&-s\\
 -B&0&-y&z\\
 r&-r-1&0&0\\
 -S&T&S-T-1&0
 \end{pmatrix},
\end{equation}
ただし行・列の順は $(2,3,0,1)$，
\[
 \begin{gathered}
 a=q_{24}=sS+1,\quad b=q_{34}=zT+1,\quad
 x=\frac{sS}{r},\quad y=\frac{zT}{r+1},\\
 A=a+x-s,\quad B=b-y+z,\quad AB=ab-1.
 \end{gathered}
\]
全ての $r,s,z,S,T,x,y$ は正整数で，
\[
 S>r,\qquad rT>(r+1)S,\qquad S\ge2.
\]
$\delta=T-S$ と置くと $r\delta>S>r$ なので $\delta\ge2$ である．

低辺 $04,14$ の node 数を
\[
 (n_1,n_2,n_3)=(n_{04;1},n_{04;2},n_{04;3}),\quad
 (m_0,m_2,m_3)=(n_{14;0},n_{14;2},n_{14;3})
\]
と書く．いずれも非負整数で，各組の和は五である．
core $0,1$ は三高分母であり，それらの唯一の低 ray $4$ の rows は
\begin{align}
 R^0&=(p^0_{41},p^0_{42},p^0_{43})
       =(1-n_1, 1-r-n_2, r+2-n_3),\label{gten18:R0}\\
 R^1&=(p^1_{40},p^1_{42},p^1_{43})
       =(\delta+2-m_0, S+1-m_2, 1-T-m_3).
       \label{gten18:R1}
\end{align}
定理~\ref{mh18:matching} により，これらの各 row はちょうど一つの正成分を持つ．
三高分母なので全成分は非零であり，特に
$R^0_2<0$，$R^1_3<0$ である．

$q=q_{01}>1$ とする．共通低 puncture $(01;4)$ の signed-divisor 分類により，
$(R^0_1,R^1_0)$ は通常型 $(u,-u)$，正例外型 $\{1,q-1\}$，
負例外型 $\{-1,-(q-1)\}$ のいずれかである．
以下で三場合を排除する．

\subsection{通常接着の場合}

まず $R^0_1>0$ とする．式~\eqref{gten18:R0} から $R^0_1=1,n_1=0$ であり，
$R^1_0=-1$ から $m_0=\delta+3\le5$ を得る．従って
$\delta=2,m_0=5,m_2=m_3=0$ である．
matching 性により $R^0_3<0$ なので $n_3\ge r+3$，従って $r\le2$ である．
$r<S<2r$ より，唯一の可能性は
\begin{equation}\label{gten18:ordinary-central}
 (r,S,T)=(2,3,5)
\end{equation}
である．このとき $R^0=(1,-1,-1)$，$R^1=(-1,4,-4)$ となる．

次に $R^0_1<0$ とすると $R^1_0>0$ である．matching 性から
$S+1-m_2<0$，すなわち $m_2\ge S+2$ なので $S\le3$ である．
また $R^1_0=-R^0_1\le4$ と $m_0\le3-S$ より
\[
 \delta\le m_0+2\le5-S.
\]
$S=2$ なら $r=1$，正の枝から $\delta\ge3$ であるため
$\delta=3,m_0=1,m_2=4,m_3=0,R^1_0=4$ となる．
すると $n_1=5,n_2=n_3=0$ なので $R^0_2=0$ となり不可能である．
$S=3$ なら $\delta=2,m_0=m_3=0,m_2=5$ であり，
再び式~\eqref{gten18:ordinary-central} を得る．
この場合の rows は $R^0=(-4,-1,4)$，$R^1=(4,-1,-4)$ である．

式~\eqref{gten18:ordinary-central} の整数条件から $s=2u,z=3v$ と書ける．
pair equation は
\[
 (7u+1)(13v+1)=(6u+1)(15v+1)-1,
 \qquad (u-2)(v+1)=-3
\]
となる．$u,v\ge1$ より $u=1,v=2$，従って
\[
 (r,s,z,S,T)=(2,2,6,3,5),\qquad (a,b)=(7,31),\qquad B=27.
\]
しかし $q_{23}=1$ なので，core $3$ の $p^3_{42}$ は $30$ の非零 signed divisor である．
辺 $34$ の puncture $2$ の clutching は
$p^3_{42}-27\equiv0\pmod{31}$ を要求する．
$|p^3_{42}|\le30$ の下でその候補は $27,-4$ だけで，どちらも $30$ を割らない．
これは矛盾である．この通常型の排除には $D$ 相殺を用いていない．

\subsection{正の例外接着の場合}

$R^0_1\le1$ なので，正例外の向きは
\[
 R^0_1=1,\quad R^1_0=q-1,\quad n_1=0,\quad
 q=\delta+3-m_0
\]
である．matching 性と row sum により $R^0=(1,-1,-1)$ となる．
行列~\eqref{gten18:central-matrix} の低 rays における係数は
\begin{equation}\label{gten18:D-low}
 D^4_0=\delta+y-x,\qquad D^4_1=s-z-1.
\end{equation}
辺 $04$ の $D$ 相殺を，core $0$ の低 column を用いて書くと
\[
 q_{02}+q_{03}=5-m_0+x-y.
\]
左辺は六以上なので $x-y>0$ である．
辺 $14$ の $D$ 相殺は
\[
 p^1_{04}+p^1_{24}+p^1_{34}=s-z-2
\]
を与える．ここで $p^1_{04}=1$，他の二成分は負整数なので $s-z\le1$ である．
一方，正の枝の不等式から
\[
 x-y=\frac{sS}{r}-\frac{zT}{r+1}
       <\frac{(s-z)S}{r}
\]
であり，$x-y>0$ は $s-z\ge1$ を要求する．従って $s=z+1$ である．
上の低 column の和は $-1$ となり，二つの負成分はともに $-1$ でなければならない．
pair equations と低 row の和から
\[
 q_{12}+q_{13}=2-(p^1_{42}+p^1_{43})=q+2
\]
を得るが，左辺は偶数，右辺は奇数である．これが矛盾である．

\subsection{負の例外接着と差二の共鳴}

負例外の場合，式~\eqref{gten18:R1} と $\delta\ge2$ から
$R^1_0\ge-1$ であるため，必ず
\begin{equation}\label{gten18:negative-exception}
 R^1_0=-1,\quad R^0_1=-(q-1),\quad
 \delta=2,\quad m_0=5,\quad m_2=m_3=0,\quad
 n_1=q,\quad q\in\{3,5\}
\end{equation}
となる．$h=S-r$ と置けば $1\le h\le r-1$，特に $r\ge2$ である．
正整数
\[
 \alpha=x-s=\frac{sh}{r},\qquad
 \beta=y-z=\frac{z(h+1)}{r+1}
\]
を導入する．二高分母の pair equation は
\begin{equation}\label{gten18:E}
 E\alpha\beta=\alpha-\beta+1,\qquad
 E=\frac{r^2}{h}-\frac{(r+1)^2}{h+1}
\end{equation}
を与える．

辺 $24$ の puncture $3$ では，$p^2_{43}$ は $a-1$ の signed divisor で，
$p^2_{43}\equiv\alpha\pmod a$ である．
$0<\alpha<(a-1)/2$ なので候補は $\alpha$ と $-(a-\alpha)$ だけである．
後者の絶対値は $(a-1)/2$ より大きく，$a-1$ を割るなら $a-1$ 自身である．
従って後者は $\alpha=1$ の場合に限られる．
同様に，辺 $34$ の puncture $2$ では，$p^3_{42}$ は $b-1$ の signed divisor で，
通常の値 $-\beta$ 以外が可能なのは $\beta=1$ のときだけである．

$\alpha,\beta>1$ なら，これらの整除性から
$r^2/h$，$(r+1)^2/(h+1)$ はともに整数となる．
式~\eqref{gten18:E} の右辺を $\alpha\beta$ で割った値は絶対値一未満なので，
$E=0$ である．これは $h=r^2/(2r+1)$ を要求するが，
$\gcd(r,2r+1)=1$ より不可能である．
$\alpha=1,\beta>1$ の場合も $r=sh$ により $r^2/h$ は整数であり，
他方の通常整除性と合わせて $E$ は整数となる．
$E=2/\beta-1$ なので $\beta=2,E=0$ となり，やはり不可能である．
従って必ず $\beta=1$ である．

このとき $E=1$，$r+1=z(h+1)$，$z\ge2$ である．これらを代入すると
\[
 (h+1)z(z-2)=h-1
\]
を得る．$z\ge3$ は不可能であり，$z=2,h=1$ に限られる．
以上により全ての残存模型は，探索上限なしに
\begin{equation}\label{gten18:resonance}
 r=3,\quad S=4,\quad T=6,\quad z=2,\quad s=3\alpha,\quad
 a=12\alpha+1,\quad b=13,\qquad \alpha\in\ZZ_{>0}
\end{equation}
へ帰着した．これは固定分母十三の既知の共鳴族と一致する．

\subsection{共鳴族の \texorpdfstring{$D$}{D} と node による排除}

式~\eqref{gten18:negative-exception}--\eqref{gten18:resonance} により
\[
 R^1=(-1,5,-5).
\]
その低 column を $(-(q-1),U,-V)$ と書けば
$q_{12}=5U+1$，$q_{13}=5V+1$ であり，$U,V$ は正の偶整数である．
式~\eqref{gten18:D-low} と辺 $14$ の $D$ 相殺は
\begin{equation}\label{gten18:UV}
 U-V=3\alpha+q-5
\end{equation}
を与える．従って $\alpha$ は偶数で，$\alpha\ge2$，
$U\ge3\alpha$，$q_{12}\ge15\alpha+1$ である．

core $0$ の低 row は
$R^0=(-(q-1),-(2+n_2),5-n_3)$ である．
初めの二成分は負なので最後は正であり，その低 column は
$(-1,-u,v)$，$u,v>0$ と書ける．辺 $04$ の $D$ 相殺から
\begin{equation}\label{gten18:uv}
 u=v+4\alpha-5\ge4\alpha-4,\qquad
 q_{02}=(2+n_2)u+1\ge8\alpha-7
\end{equation}
を得る．

$\alpha\ge2$ なので，上の signed-divisor 分類は
$p^2_{43}=\alpha$ を強制する．従って $p^2_{34}=12$ である．
core $2$ は三高分母であり，matching 性によって低 ray $3$ は高 ray $4$ と対になる．
そのため $p^2_{40},p^2_{41}$ はともに負である．
辺 $24$ の puncture $3$ の node は二なので，残る二つの node を $\nu_0,\nu_1$ とすると
$\nu_0+\nu_1=3$ である．clutching と node の式は
\[
 p^2_{40}=(8-12\nu_0)\alpha+1-\nu_0,\qquad
 p^2_{41}=(12\nu_0-21)\alpha+\nu_0-2
\]
を与える．両方が負になるのは $\nu_0=1,\nu_1=2$ の場合だけであり，
\[
 p^2_{40}=-4\alpha,\qquad p^2_{41}=-(9\alpha+1)
\]
である．正整数
\[
 M=-p^2_{04}=\frac{a q_{02}-1}{4\alpha}>3q_{02},\qquad
 N=-p^2_{14}=\frac{a q_{12}-1}{9\alpha+1}>q_{12}
\]
を考える．core $4$ の高 ray $2$ の係数は
$D^4_2=12-12\alpha$ なので，辺 $24$ の $D$ 相殺は
\[
 M+N=24\alpha+1
\]
を要求する．一方，式~\eqref{gten18:UV},\eqref{gten18:uv} から
\[
 M+N>3q_{02}+q_{12}
       \ge39\alpha-20
       \ge24\alpha+10
\]
となり矛盾する．これで三つの接着型を全て排除し，
定理~\ref{gten18:no-support} が従う．

この証明で有限に数えたものは node の三成分配分および符号の場合分けである．
分母を切る探索ではなく，残り得る無限族~\eqref{gten18:resonance} も不等式で排除した．
$D$ 相殺を実際に用いた辺は，canonical 標識で $04,14,24$ の三本だけである．

% v18. K4 plus one high leaf: an unbounded obstruction.
% Dependencies: mh18:matching, tl15:one-high-full, ls11:low-puncture.
\section{\texorpdfstring{$K_4$}{K4} に一葉を付けた support の排除}
\label{lh18:section}

本節では，高分母 support
\[
 H=\{01,02,03,12,13,23,04\}
\]
を排除する．四頂点 $K_4$ の障害を部分 graph に適用するのではなく，
一高分母 core $4$ の完全標準形を使う．
三つの三高分母 cores $1,2,3$ に対しては
定理~\ref{mh18:matching} を用いる．
以下では各低 row に正の entry がちょうど一つある，という同定理の帰結が必要である．

\subsection{一高 hub の補正された node と係数の式}

命題~\ref{tl15:one-high-full} に従い，core $4$ の標識を
$(A,B,C,D)=(0,1,2,3)$ とする．$K_4+\mathrm{leaf}$ は $1,2,3$ の任意の置換を許すので，
この正規化で一般性を失わない．
\begin{equation}\label{lh18:hub}
 \begin{gathered}
 h,t\ge1,\quad u=th,\quad v=t(h+1),\quad Q=q_{04}=th(h+1)+1=hv+1,\\
 c,d\in\ZZ,\quad cd=0,\qquad
 V=P^4=\begin{pmatrix}
 0&1-Q&-h-1&h\\
 -1&0&0&0\\
 -u&u-1-c&0&c\\
 v&-v-1-d&d&0
 \end{pmatrix}.
 \end{gathered}
\end{equation}
$I=\{0,1,2,3\}$，$\ell_0=Q$，$\ell_1=\ell_2=\ell_3=1$ とし，
$i\ne j\in I$ について
\[
 R_{ij}=p^i_{4j},\qquad X_{ij}=p^i_{j4},\qquad N_{ij}=n_{i4;j}
\]
と置く．対角成分は零とする．hub に接する四辺の node 条件は
\begin{equation}\label{lh18:node-reconstruction}
 N_{ij}\in\ZZ_{\ge0},\qquad \sum_{j\ne i}N_{ij}=5,\qquad
 R_{ij}=\ell_i(1-N_{ij})-V_{ij}
\end{equation}
と同値である．特に各 node 行は二十一通りである．局所の pair products と
四辺 $i4$ の対向 $D$ 相殺はそれぞれ
\begin{align}
 R_{ij}X_{ij}&=\ell_iq_{ij}-1,\label{lh18:pair}\\
 \sum_{j\ne i}X_{ij}+\sum_{j\ne i}V_{ji}&=-2\ell_i.\label{lh18:D}
\end{align}
従って，補正された発散の式は
\begin{equation}\label{lh18:weighted-flow}
 \sum_{j\ne i}(X_{ij}-R_{ji})
 =\sum_{j\ne i}\ell_jN_{ji}-(Q+3)-\ell_i
\end{equation}
である．これは全低 hub の発散公式を，重み $\ell_j$ 付きで置き換える．
また $V_{0j}V_{j0}=Q-1$ であり，この三 pair に零積条件を課してはならない．

$i,j\in\{1,2,3\}$ の共通低 puncture $4$ では，補題~\ref{ls11:low-puncture} の符号付き約数の分類により
\begin{equation}\label{lh18:reverse-partners}
 R_{ji}=-R_{ij},\quad\text{または}\quad
 \{R_{ij},R_{ji}\}=\{s,s(q_{ij}-1)\},\quad s\in\{1,-1\}
\end{equation}
である．通常の場合に限って $q_{ij}$ が自由に残り，例外の場合には
$q_{ij}=1+\max(|R_{ij}|,|R_{ji}|)$ となる．

\subsection{特別な低 ray における四つの接続型}

\begin{lemma}[標準形の $B$ ray を通る接続]\label{lh18:B-types}
五次の算術的接着データが上記 support を持つとする．
$\kappa=1-n_{01;4}\in\{-4,-3,-2,-1,0,1\}$ と置くと，
$(N_{01},N_{10},q_{01},\kappa)$ は次の四型のいずれかである：
\begin{equation}\label{lh18:B-table}
\begin{array}{c|cccl}
 \text{型}&N_{01}&N_{10}&q_{01}&\kappa\\\hline
 \mathrm I&4&0&2Q-1&0\\
 \mathrm {II}&3&1&Q&0\\
 \mathrm {III}&0&0&2Q+1&1\\
 \mathrm {IV}&2&1&\text{任意の奇数 }q_{01}\ge3&1-Q,\quad Q\in\{3,5\}
\end{array}
\end{equation}
型 IV も残しており，$q_{01}$ に上限を置いていない．
\end{lemma}

\begin{proof}
$r=R_{10}=2-N_{10}$ は $[-3,2]$ の非零整数である．
$R_{12}=1-N_{12}$，$R_{13}=1-N_{13}$ は一以下である．
core $1$ の低 row は正 entry をちょうど一つ持ち，行和は $-1$ だから，
$r<0$ なら $r=-1$ でなければならない．

$r=-1$ を排除する．$q=q_{01}$，$a=2-N_{01}\in[-3,2]$ と置くと
$R_{01}=aQ-1$，$X_{10}=1-q$，$X_{01}=(\kappa+1)q-1$ である．
$b=\kappa+1\in[-3,2]$ と書けば pair equation は
\[
 (aQ-1)(bq-1)=Qq-1
\]
となる．二因子は同符号なので，$a,b\ge1$ または $a,b\le0$ である．
前者では $(a,b)=(1,1)$ の差は $2-Q-q<0$，$(1,2),(2,1)$ の差は
それぞれ $(Q-2)(q-1)>0$，$(Q-1)(q-2)>0$，$(2,2)$ でも差は正である．
後者で両者が負なら左辺は $Qq-1$ より大きい．一方だけが零なら
$Q$ または $q$ が二を割る必要があり，両者が零でも不可能である．
従って $r=1,2$ に限られる．

$r=2$ では $x=(q-1)/2\ge1$ と置く．
$X_{01}=(2\kappa-1)x+\kappa$ である．
$\kappa=1$ なら $a\in\{1,2\}$ であり，pair equation は
$a=2,x=Q$ だけを許す．
$\kappa=0$ なら $a=-A$，$A\in\{0,1,2,3\}$ と書けて，
\[
 \{(A-2)Q+1\}x=Q-1
\]
となる．正整数解は $A=2,x=Q-1$ だけである．
$\kappa=-k\le-1$ の場合，$A\ge1$ なら
$(AQ+1)((2k+1)x+k)>Q(2x+1)-1$ である．
$A=0$ なら $e=Q-k$ として $(2e-1)x=1-e$ となり，
$e\ge1$，$e\le0$ のいずれも符号が矛盾する．
これで型 I，III を得る．

$r=1$ では $X_{10}=q-1$，$X_{01}=(\kappa-1)q+1$ である．
$\kappa=1$ は $a=q\ge3$ を要求して不可能である．
$\kappa=0$ では $a=-A$ として
$\{(A-1)Q+1\}(q-1)=Q-1$ となり，$A=1,q=Q$ だけが残る．
$\kappa=-k\le-1$ では，$A\ge1$ の pair product は再び過大であり，
$A=0$ だけが残る．このとき $Q=k+1$ であり，$k\le4$ と $Q$ の奇数性から
$Q=3,5$ である．これは型 II，IV を与える．
\end{proof}

\begin{lemma}[型 I は葉分母三に限られる]\label{lh18:typeI-small}
補題~\ref{lh18:B-types} の型 I では $Q=3$ である．
\end{lemma}

\begin{proof}
型 I では $X_{01}=1-Q$ であり，式~\eqref{lh18:D} の $i=0$ は
\[
 X_{02}+X_{03}=-Q-t
\]
を与える．$N_{01}=4$ なので $N_{02}+N_{03}=1$ である．
$N_{03}=0$ なら $R_{02},R_{03}>0$ となって矛盾するから，
$(N_{02},N_{03})=(0,1)$ であり，$(R_{02},R_{03})=(Q+h+1,-h)$ となる．
$(h+1)u\equiv-1\pmod Q$ と $0<u<Q$ より，ある整数 $z\ge1$ を用いて
\[
 X_{02}=u+(z-1)Q,\qquad X_{03}=-v-zQ,\qquad q_{03}=1+hz
\]
と書ける．$R_{30}=-(v-1+N_{30})$ なので，
\[
 X_{30}=-w,\qquad
 w=\frac{hz}{v-1+N_{30}}\in\ZZ_{>0},\qquad w\le z.
\]
辺 $03$ の node 条件は，ある $1\le a\le4$ に対して
$v+zQ+w=a(1+hz)$ を要求する．$e=a-v$ と置き，$Q=hv+1$ を代入すると
\[
 e(hz+1)=z+w\le2z.
\]
従って $e>0$，$eh<2$ なので $e=h=1$ である．
さらに $v=2t$，$a=v+1\le4$ より $t=1$，従って $Q=3$ となる．
\end{proof}

\begin{proposition}[葉分母の一様な上界]\label{lh18:leaf-bound}
上記 support を持つ五次の算術的接着データでは $Q\le21$ である．
\end{proposition}

\begin{proof}
$Q>21$ と仮定する．$u=th$，$h\le u$ なので，$u\le4$ なら
$Q=u(h+1)+1\le21$ となる．従って $u\ge5$，$v=u+t\ge6$ である．
補題~\ref{lh18:B-types} と補題~\ref{lh18:typeI-small} により型 II，III だけが残る．
いずれも
\[
 R_{12},R_{13}\in\{-1,-2\},\qquad
 R_{20}=u+1-N_{20}\ge u-4>0,\qquad R_{30}<0
\]
を満たす．core $2$ の低 row の唯一の正 entry は $R_{20}$ である．

$d=0$ なら $R_{31}=v+2-N_{31}\ge v-3\ge3$ である．
$R_{13}<0$ なので，式~\eqref{lh18:reverse-partners} は
$R_{31}=-R_{13}\le2$ を要求し，矛盾する．
従って $c=0$ である．このとき
$R_{21}=2-u-N_{21}\le-3$ だから，式~\eqref{lh18:reverse-partners} は
$R_{12}=-1$ を要求する．また $R_{23}=1-N_{23}<0$ なので $|R_{23}|\le4$ である．

core $3$ では $R_{30}<0$ であり，$R_{31},R_{32}$ のちょうど一方が正である．
もし $R_{32}<0$ なら $R_{31}=-R_{13}\le2$ である．しかし row sum から
\[
 R_{31}+R_{32}=v-2+N_{30}
\]
となるため，$R_{32}\ge v-4>0$ となり矛盾する．
従って $R_{32}>0$，$R_{31}<0$ であり，辺 $23$ は通常接着なので $R_{32}\le4$ となる．
一方，同じ row sum は $R_{32}\ge v-1\ge5$ を与える．これも矛盾である．
\end{proof}

\subsection{残る自由整数を切らない完全な有限型検査}

$Q\le21$ と式~\eqref{lh18:hub} から，$(h,t)$ は
\[
 (1,1),\ldots,(1,10),\qquad (2,1),(2,2),(2,3),\qquad (3,1),(4,1)
\]
の十五個だけである．ただし，この上界だけから $c,d$ や他の高分母が有限になるとは言えない．
次の手順では，それらを厳密な無限型のまま残す．

\begin{lemma}[固定 node 行からの完全な型の復元]\label{lh18:finite-types}
上の十五個の $(h,t)$ を固定し，式~\eqref{lh18:node-reconstruction} の四行を取る．
$c=0,d=z$ または $d=0,c=z$ として，$z$ は全整数を動くものとする．
次の有限な判定は，hub に関する row/pair，四辺 $i4$ の node 条件，
六辺 $ij\subset I$ の puncture $4$ における clutching と非負 node 条件，
および三高分母 cores の低 row の一正二負条件を完全に尽くす．
\begin{enumerate}
\item $i=1,2,3$ について $R=R_{0i}$，$r=R_{i0}$ と置く．
これらは $z$ に依存しない．$Rr\ne0$ を要求し，各 $\kappa\in[-4,1]\cap\ZZ$ に対して
\begin{equation}\label{lh18:center-linear}
 \{Qr+R-\kappa Rr\}q_{0i}=R+r
\end{equation}
を調べる．左係数が非零なら得られた有理数を取り，
\[
 q_{0i}\in2\ZZ+1,\quad q_{0i}\ge3,\quad
 r\mid q_{0i}-1,\quad R\mid Qq_{0i}-1
\]
を検査する．左係数が零で右辺が非零なら棄却する．両方が零なら
\begin{equation}\label{lh18:free-center-q}
 q_{0i}=1+\operatorname{lcm}(2,|r|)a,\qquad a\in\ZZ_{>0}
\end{equation}
という全無限族を残す．
\item $z$ を含まない二つの低 rows の間の pair には
式~\eqref{lh18:reverse-partners} をそのまま適用する．
$z$ を含む row の残り二 entries は，それぞれ $s=z+b$，$s=-z+b'$ の形である．
その相手側の固定 entry を $r\ne0$ とすると，$s$ の許される集合はちょうど
\begin{equation}\label{lh18:partner-set}
 \{-r\}\ \cup\
 \begin{cases}\{\operatorname{sgn}(r)\},&|r|\in2\ZZ_{>0},\\ \varnothing,&\text{それ以外}
 \end{cases}
 \ \cup\
 \begin{cases}\{2a\operatorname{sgn}(r):a\in\ZZ_{>0}\},&|r|=1,\\
 \varnothing,&|r|\ne1.
 \end{cases}
\end{equation}
これは有限個の点と偶奇を指定した半直線の和集合である．
二つの entries に対する条件を $z$ 上の集合に直して交差させ，
その row が正 entry をちょうど一つ持つ型だけを残す．
\end{enumerate}
\end{lemma}

\begin{proof}
式~\eqref{lh18:center-linear} は
$X_{0i}=(Qq_{0i}-1)/R$，$X_{i0}=(q_{0i}-1)/r$ と
$X_{0i}+X_{i0}=\kappa q_{0i}$ から分母を払ったものである．
両係数が零の場合は $R=-r$ かつ $r\mid Q-1$ なので，
二つの pair 整除条件と奇数性は式~\eqref{lh18:free-center-q} にちょうど同値である．
$\kappa\in[-4,1]$ は $0\le n_{0i;4}\le5$ を保持する．

後半は式~\eqref{lh18:reverse-partners} を，固定した非零整数 $r$ に対して解いたものである．
通常の $s=-r$ では $q=1+\operatorname{lcm}(2,|r|)a$ の全正整数族が残る．
例外では $q=1+\max(|r|,|s|)$ である．半直線を有限個の値で置き換えていない．
$cd=0$ は指定した二つの枝で尽くされ，その交わり $c=d=0$ が二重に現れることは
不存在判定に影響しない．
\end{proof}

\begin{theorem}[$K_4$ に高分母の一葉を付けた support は不可能]\label{lh18:no-K4-leaf}
五次の算術的接着データの高分母 support が
$\{01,02,03,12,13,23,04\}$ となることはない．
特に従来の候補 $G_{5,12}$ は排除される．
\end{theorem}

\begin{proof}
命題~\ref{lh18:leaf-bound} により十五個の $(h,t)$ だけを扱えばよい．
各 node 行の二十一通りを全て取り，補題~\ref{lh18:finite-types} の条件を検査する．
固定された二つの低 rows の条件までを通過する node 配置数を $A$，
残る二つの partner 型の組のうち一正二負条件を満たす数を $B$ とすると，結果は次表となる．
\begin{center}
\small
\begin{tabular}{rrr|rr|r}
 $h$&$t$&$Q$&$A$&$B$&\text{非空の交差}\\\hline
 1&1&3&76&135&0\\
 1&2&5&9&18&0\\
 1&3&7&23&37&0\\
 1&4&9&18&27&0\\
 1&5&11&23&37&0\\
 1&6&13&18&27&0\\
 1&7&15&23&37&0\\
 1&8&17&18&27&0\\
 1&9&19&23&37&0\\
 1&10&21&18&27&0\\
 2&1&7&8&16&0\\
 2&2&13&20&35&0\\
 2&3&19&22&40&0\\
 3&1&13&24&39&0\\
 4&1&21&22&40&0\\\hline
 &&\text{合計}&345&579&0
\end{tabular}
\end{center}
型の組は一正二負条件を課す前には $1192$ 個であり，そのうち $613$ 個が符号条件で落ちる．
残る $579$ 個は全て，点または偶奇付き半直線の交差が空である．
例えば交差を $[L,U]\cap(2\ZZ+\epsilon)$ と書ける場合，端点を指定した偶奇へ丸めて
$L>U$ となることを判定する．無限端点は無限のまま保持する．
\nolinkurl{verify_k4_leaf_obstruction_v18.py} は式~\eqref{lh18:center-linear} と
式~\eqref{lh18:partner-set} だけを用い，この完全表と空交差の証明書を再生成する．
未知の分母や $c,d$ への探索上限はない．従って必要条件だけで全ての型が排除される．
\end{proof}

\begin{remark}[量化と標識]
上界 $Q\le21$ の証明では一高分母の完全標準形，三高分母の正値性，
指定した clutching と node，および辺 $04$ の $D$ 相殺を使う．
その後の有限型検査は追加の $D$ 条件や三高分母行列の高高 entries を用いない．
表の $A,B$ は手順内の型の数であり，互いに異なる局所 charts や幾何的同型類の数ではない．
$Q=3$ の標準形における同じ負 entry の重複は低標識を区別したまま扱う．
$K_4$ の自己同型で標準形を選んだ後，全ての node 行を走査するため，
$B,C$ の標識交換でしか現れない表示を捨てていない．
\end{remark}

% v19 G5_7 obstruction. Dependencies: mh18:matching, tl15:one-high-full,
% ml12:dichotomy, ls11:low-puncture, td12:low-entry.
% Exact arithmetic certificate: verify_g5_7_obstruction_v19.py.
\section{三角形を共有する菱形に一葉を付けた支持の排除}
\label{g719:section}

高支持を
\[
 H=\{01,02,12,13,23,04\}
\]
とする．これは標準代表 $G5\_7=\{01,02,03,12,13,24\}$ を
旧頂点 $(2,0,1,3,4)$ の順に新頂点 $(0,1,2,3,4)$ へ写したものである．
以下では葉 $4$ の標準形の低 rays を $B,C,D$ と呼び，$A=0$ とする．
三つの低 rays のうち $0$ との辺が低分母であるものを $m$ と書く．
従って $m$ は物理的には頂点 $3$ であるが，標準形内での位置は
$B,C,D$ の三通りを全て保つ．残る二頂点は互換である．

\begin{lemma}[二高 core の二つの小さい正 reverse は共存しない]
\label{g719:small-positive}
正の整数二高局所模型の二つの低行を取り，各低行中の唯一の正 entry を
$e_1,e_2$ とする．$e_1,e_2\le2$ は不可能である．
\end{lemma}
\begin{proof}
二高標準形ではこれらは $r,T$ であり，低行は
$(r,-r-1,0)$ と $(-S,T,S-T-1)$ である．
正の枝では $S>r$ および $rT>(r+1)S$ なので $T>r+1\ge2$ となる．
負の枝では $S=T<r$ なので，$r,T\le2$ のもとでは
$r=2,S=T=1$ に限られる．この場合 $a=s+1,b=z+1$，
$x=s/2,y=z/3$，$A=1+s/2,B=1+5z/3$ であり，
$AB=ab-1$ は
\[
 (s-4)(z+3)=-6
\]
と同値である．整数性から $s\ge2$ は偶数，$z\ge3$ は三の倍数である．
$s=2$ なら左辺は $-12$ 以下，$s\ge4$ なら非負であり矛盾する．
負の枝を捨てず，その唯一の小さい場合をここで明示的に排除した．
\end{proof}

葉の完全標準形を
\[
 V=P^4=\begin{pmatrix}
 0&1-Q&-h-1&h\\
 -1&0&0&0\\
 -u&u-1-c&0&c\\
 v&-v-1-d&d&0
 \end{pmatrix},
\qquad Q=th(h+1)+1,
\]
\[
 u=th,\quad v=t(h+1),\quad cd=0
\]
と書く．$h,t\ge1$ であり，行列の順序は $(0,B,C,D)$ である．
$N_{ij}=n_{i4;j}$ とし，各行の三成分は非負で総和五である．
\[
 \ell_0=Q,\quad \ell_B=\ell_C=\ell_D=1,\qquad
 R_{ij}=p^i_{4j}=\ell_i(1-N_{ij})-V_{ij},\quad
 X_{ij}=p^i_{j4}
\]
と置けば
\begin{equation}\label{g719:basic}
 R_{ij}X_{ij}=\ell_iq_{ij}-1,
 \qquad
 \sum_{j\ne i}X_{ij}+\sum_{j\ne i}V_{ji}=-2\ell_i.
\end{equation}
第二式は辺 $i4$ の $D$ 相殺である．特に
\begin{equation}\label{g719:hub-D}
 \sum_{j=B,C,D}X_{0j}=1-2Q-t.
\end{equation}
中央 core $0$ は三高なので，matching 定理から行 $R_0$ はちょうど一つの
正 entry を持つ．core $m$ は二高で，その低 rays は $0,4$ である．
従って
\begin{equation}\label{g719:low-pair}
 R_{m0}\le0,\quad X_{m0}\le0,\quad R_{m0}X_{m0}=0,\quad
 R_{0m}X_{0m}=Q-1,\quad X_{0m}+X_{m0}\in\{-4,-3,-2,-1,0,1\}.
\end{equation}

\begin{lemma}[正の最大 divisor の二枝の排除]
\label{g719:positive-branches}
次の二つの枝はいずれも不可能である：
\[
 (m=B,\ N_{0B}=1),\qquad
 (m=D,\ h=1,\ N_{0D}=0).
\]
\end{lemma}
\begin{proof}
両方の場合で $R_{0m}=Q-1>0$，$X_{0m}=1$ である．
中央の matching と低行 $m$ の和 $-1$ から，残る二高 rays を $a,b$ と書くと
\[
 p^0_{ma}=p^0_{mb}=-1.
\]
従って低辺 $0m$ の clutch により
$p^m_{0a}=2-n_{0m;a}\le2$，$p^m_{0b}=2-n_{0m;b}\le2$ である．

$m=B$ の場合，葉の $B$ 行は $(-1,0,0)$ である．
従って core $m$ の低行 $4$ の二つの高 entries は
$1-N_{ma},1-N_{mb}$ であり，唯一の正 entry は一である．
低行 $0$ の正 entry は二以下なので，補題~\ref{g719:small-positive} に反する．

$m=D,h=1$ の場合には残る高 rays を葉の $B,C$ の順に置き，
$A=q_{0B},B'=q_{0C}$ とする．中央の正高間 entries を
$H=p^0_{BC}>0,K=p^0_{CB}>0$ と書けば，row sums と pair equations は
\[
 p^0_{B4}=-H-1,\quad p^0_{C4}=-K-1,\quad HK=AB'-1.
\]
式~\eqref{g719:hub-D} と $t=(Q-1)/2$ から
\[
 H+K=\frac{5(Q-1)}2.
\]
$R_{0B},R_{0C}<0$ なので $N_{0B},N_{0C}\ge2$ であり，
$N_{0B}+N_{0C}=5$ から分配は $(2,3)$ または $(3,2)$ である．
前者では $R_{0B}=-1$，従って $H=QA-2\ge3Q-2>5(Q-1)/2$ となり矛盾する．
ここで高分母は奇数なので $A\ge3$ を用いた．
後者では
\[
 R_{0B}=-(Q+1),\quad R_{0C}=-(Q-2),\quad
 H+1=\frac{QA-1}{Q+1}.
\]
従って $A=k(Q+1)-1$，$H=Qk-2$ であり，$H,K>0$ により $k=1,2$ に限る．
それぞれについて
\[
\begin{array}{c|cccc|c}
 k&A&B'&H&K&HK-(AB'-1)\\\hline
 1&Q&(3Q-5)/2&Q-2&(3Q-1)/2&2-Q\\
 2&2Q+1&(Q-1)/2&2Q-2&(Q-1)/2&(5-3Q)/2
\end{array}
\]
となり，$Q\ge3$ では最後の列が零にならない．
\end{proof}

\begin{proposition}[葉分母と配置の完全な有限帰着]
\label{g719:leaf-bound}
上記支持で正の整数局所模型，marked clutch，非負 node 配分（各辺総和五），
および全ての対向 $D$ 相殺が成立するなら $Q\le21$ であり，
さらに次の表のいずれかに属する：
\[
\begin{array}{c|c|c|c}
 m&N_{0m}&R_{0m}&\text{追加条件}\\\hline
 B&2&-1&Q\le5\\
 C&1&h+1&th\le4\\
 C&2&-1&h=t=1\\
 D&1&-h&t(h+1)\le4
\end{array}
\]
この帰着では他の高分母，および $c,d$ に上限を置かない．
\end{proposition}
\begin{proof}
式~\eqref{g719:low-pair} により $R_{0m}$ は $Q-1$ の非零 signed divisor である．
$m=B$ なら $R_{0B}=Q(2-N_{0B})-1$ なので，その絶対値から
$N_{0B}=1,2$ に限る．前者は補題~\ref{g719:positive-branches} で排除済み．
後者では $X_{0B}=-(Q-1)$ であり，$X_{B0}\le0$ と clutch から
$-(Q-1)\ge-4$，従って $Q\le5$ である．

$m=C$ なら $R_{0C}=Q(1-N_{0C})+h+1$ である．
$N_{0C}=0$ または $N_{0C}\ge3$ では絶対値が $Q-1$ より大きい．
$N_{0C}=2$ では $Q-h-1$ が $Q-1$ を割るので $Q-h-1\mid h$，
従って $th(h+1)\le2h$ であり，$h=t=1$ に限る．
$N_{0C}=1$ なら $R_{C0}=u+1-N_{C0}\le0$ と $N_{C0}\le5$ から
$u=th\le4$ である．$h\le u$ より
$Q=u(h+1)+1\le u(u+1)+1\le21$ となる．

$m=D$ なら $R_{0D}=Q(1-N_{0D})-h$ である．$N_{0D}\ge2$ は大きすぎる．
$N_{0D}=0$ なら $Q-h\mid Q-1$，従って $Q-h\mid h-1$ である．
$h\ge2$ では $Q-h>h-1$ なので不可能であり，$h=1$ は
補題~\ref{g719:positive-branches} で排除済みである．
$N_{0D}=1$ なら $X_{0D}=-v$，かつ
$R_{D0}=1-N_{D0}-v<0$ なので $X_{D0}=0$ である．
従って $-v\in\{-4,-3,-2,-1,0,1\}$ より $v\le4$ であり，
$Q=hv+1\le(v-1)v+1\le13$ である．
\end{proof}

\begin{lemma}[全整数の自由係数を保つ有限型検査]
\label{g719:finite-types}
命題~\ref{g719:leaf-bound} の下で，必要な node，matching，共通低 puncture 条件，
および辺 $04$ の $D$ 相殺を全て検査すると，候補は残らない．
自由係数 $c,d$ および $q_{0j}>1$ に人為的な上限を置く必要はない．
\end{lemma}
\begin{proof}
標準形内の $B,C,D$ を以下では $1,2,3$ と書く．
$Q\le21$ より $h(h+1)\le20$ なので $1\le h\le4$，
$1\le t\le\lfloor20/(h(h+1))\rfloor$ である．
$(h,t)$ は十五通りであり，$m=B,C,D$ と合わせて四十五通りとなる．
このうち命題~\ref{g719:leaf-bound} の配置表だけを保つ．
各 $N_i$ は非負三成分で和五なので二十一通りである．
従って式~\eqref{g719:basic} の $R$ は，$c=0$ または $d=0$ の二枝それぞれにおいて，
一つの行の二成分だけが自由整数 $z$ の一次式となる．

$i\ne m$ の中央高辺に対して $R=R_{0i},r=R_{i0}$ と置くと，
同じ marked puncture $4$ における clutch は
\begin{equation}\label{g719:central-equation}
 (Qr+R-\kappa Rr)q_{0i}=R+r,\qquad -4\le\kappa\le1.
\end{equation}
左辺係数が非零なら $q_{0i}$ はただ一つの有理数として決まり，
奇数整数 $\ge3$ と二つの pair 整除を満たすものだけを保つ．
左辺係数が零なら，右辺が零である場合のみ可能であり，
\[
 r=-R,\quad Q-1=-\kappa r,\quad
 q_{0i}=1+\operatorname{lcm}(2,|r|)a,\qquad a\in\ZZ_{>0}
\]
という完全な無限等差数列を保つ．
$i=m$ の低辺では $q_{0m}=1$ と式~\eqref{g719:low-pair} をそのまま用いる．

$1,2,3$ 間の各辺は高分母であり，共通 puncture $4$ は両端で低分母である．
従って補題~\ref{ls11:low-puncture} による reverse partner の必要条件は
\begin{equation}\label{g719:partner}
 r=-s\ne0\quad\text{または}\quad
 rs>0,\quad\min(|r|,|s|)=1,\quad\max(|r|,|s|)\in2\ZZ_{>0}.
\end{equation}
固定 $r\ne0$ に対する $s$ の集合は，一点 $-r$ に加えて，
$|r|$ が偶数なら一点 $\operatorname{sgn}(r)$，
$|r|=1$ なら全半直線 $2a\operatorname{sgn}(r)$，$a\ge1$ である．
従って $s=\pm z+b$ に対しても，一点または偶奇を指定した全半直線に厳密に移される．
変数行の二つの条件の共通部分は，端点と偶奇の比較だけで空・非空を決定できる．
各行の正 entry がちょうど一つという条件も，これらの型の符号で決定できる．
二高 core $m$ では $R_{m0}=0$ も保ち，他の二高 entries は非零で逆符号とする．

この有限整数演算は，同梱の標準 Python ライブラリのみの検算
\nolinkurl{verify_g5_7_obstruction_v19.py} に完全に記載した．
固定行間の条件を通る node 配置は $529$，自由行の partner 型対は $1988$ であり，
そのうち符号適合は $910$，空交差は $750$，非空の型は $160$ となる．
中央辺の $\kappa$ の選択まで展開すると $212$ 型であり，
$D$ 相殺に必要な情報でまとめると次の十三行だけが残る．
表中 $z$ は高分母なので奇数整数 $\ge3$ でなければならず，
$\Delta=\sum_jX_{0j}-(1-2Q-t)$ とする．
\begin{center}
\small
\begin{tabular}{c|c|c|c|c|c|r}
$m$&$(h,t)$&$Q$&$(R_{0B},R_{0C},R_{0D})$&$(q_{0B},q_{0C},q_{0D})$&$D$ 条件の障害&型数\\\hline
$B$&$(1,1)$&3&$(-1,2,-4)$&$(1,z,7)$&$z=1$&12\\
$B$&$(1,1)$&3&$(-1,2,-4)$&$(1,z,3)$&$z=-1$&2\\
$B$&$(1,1)$&3&$(-1,-4,2)$&$(1,7,z)$&$z=1$&6\\
$B$&$(1,1)$&3&$(-1,-4,2)$&$(1,3,z)$&$z=-1$&2\\
$B$&$(1,2)$&5&$(-1,-8,4)$&$(1,5,z)$&$z=-3$&1\\
$C$&$(1,1)$&3&$(-1,2,-4)$&$(z,1,7)$&$z=1$&24\\
$C$&$(1,2)$&5&$(-1,2,-6)$&$(z,1,11)$&$z=1$&24\\
$D$&$(1,1)$&3&$(5,-7,-1)$&$(7,5,1)$&$\Delta=6$&42\\
$D$&$(1,1)$&3&$(-7,5,-1)$&$(5,7,1)$&$\Delta=6$&42\\
$D$&$(1,2)$&5&$(-11,7,-1)$&$(9,3,1)$&$\Delta=5$&12\\
$D$&$(2,1)$&7&$(-15,10,-2)$&$(13,13,1)$&$\Delta=14$&24\\
$D$&$(2,1)$&7&$(-15,10,-2)$&$(13,3,1)$&$\Delta=7$&9\\
$D$&$(3,1)$&13&$(-27,17,-3)$&$(25,21,1)$&$\Delta=26$&12
\end{tabular}
\end{center}
最後の計算では $X_{0j}=(Qq_{0j}-1)/R_{0j}$ を用いるだけである．
全ての無限等差数列では $D$ 方程式の一意的な解が三未満であり，
分母が全て決まった行では $\Delta\ne0$ である．
これで全てが矛盾する．なお型数は branch と node 配分を保持した個数であり，
$c=d=0$ は二枝に現れるので同型類の個数ではない．
\end{proof}

\begin{theorem}[$G5\_7$ の無上界排除]
\label{g719:no-G5-7}
五次 all-plane SNC，標準四直線 cones，固定された marked punctures の枠内で，
局所正 weights，非負 node 配分（各辺総和五），marked clutch，
および対向 $D$ 相殺を全て要求するとき，高分母支持 $G5\_7$ は存在しない．
この結論は分母または numerator の探索 cutoff に依存しない．
\end{theorem}
\begin{proof}
葉を頂点 $4$，その高隣接を $0$ とする．中央 $0$ は三高，
残る非隣接 $m$ は二高，残る二つの cores は三高である．
一高完全標準形の全標識を $m=B,C,D$ に分ければ漏れはない．
命題~\ref{g719:leaf-bound} と補題~\ref{g719:finite-types} が全場合を排除する．
\end{proof}

% v19 complete obstruction. Canonical labels A=0,B=1,C=2,D=3.
% The physical three-high vertex is denoted J in {1,2,3}.
\section{Diamond の四価頂点に一葉を付けた support}
\label{dl19:section}
本節の support は $G_{5,9}=\{01,02,03,04,12,13\}$ である．
葉 $4$ の一高標準形に合わせて残る標識を交換するため，以下では四価頂点を $0$，
三価頂点を $J\in\{1,2,3\}$ と書く．標準形の $(A,B,C,D)$ は $(0,1,2,3)$ であり，
$J$ 以外の二つの標識を $E,F$ とする．従って $I=\{0,1,2,3\}$ の高辺は
$01,02,03,JE,JF$，唯一の低辺は $EF$ である．$J=1,2,3$ の全配置を残す．

一高標準形（命題~\ref{tl15:one-high-full}）は
\[
 Q=q_{04}=th(h+1)+1,\quad u=th,\quad v=t(h+1),\quad h,t\ge1,\quad cd=0,
\]
\[
 V=P^4=\begin{pmatrix}
0&1-Q&-h-1&h\\-1&0&0&0\\-u&u-1-c&0&c\\v&-v-1-d&d&0
\end{pmatrix}.
\]
ここでは $c,d$ は任意の整数であり，$Q=3$ の標準標識の重複も除かない．
$R_{ij}=p^i_{4j}$，$X_{ij}=p^i_{j4}$，$N_{ij}=n_{i4;j}$，
$\ell_0=Q$，$\ell_1=\ell_2=\ell_3=1$ とする．四つの hub 辺の node 条件と $D$ 相殺は
\begin{align}
 R_{ij}&=\ell_i(1-N_{ij})-V_{ij},\quad N_{ij}\ge0,\quad\sum_{j\ne i}N_{ij}=5,
 \label{dl19:reconstruct}\\
 R_{ij}X_{ij}&=\ell_iq_{ij}-1,\qquad
 \sum_{j\ne i}X_{ij}+\sum_{j\ne i}V_{ji}=-2\ell_i.
 \label{dl19:pairD}
\end{align}
core $J$ の低 row は一正二負である（定理~\ref{mh18:matching}）．
cores $E,F$ では，定理~\ref{ls11:two-high-signs} と補題~\ref{ml12:dichotomy} により，
低 row から二つの高 rays へは一正一負，低低 entry は非正である．
低低 entry は一方が零，他方が $-1$ または $-3$ 以下であり，$-2$ は起こらない．
非零の低低 entry を持つ row の高方向 entries は $-S,T$ で，$S\ge2$，$T\ge S$ である．
これらは二高 core の高隣接頂点間の辺が高である場合にも成立する．
従って本節では負の枝 $S=T$ を除かない．

\begin{lemma}[標準 $B$ の接続]
\label{dl19:B-types}
$r=R_{10}$ とすると，$J=1$ では $r=1,2$，$J\ne1$ では $r=1,2,-2$ である．
$r>0$ の場合の $(N_{01},N_{10},q_{01},\kappa)$ は補題~\ref{lh18:B-types} の
四型 I--IV に限られる．ここで $\kappa=1-n_{01;4}\in[-4,1]$ である．
型 I では $Q=3$，型 IV では $Q=3,5$ である．
$r=-2$ の場合にも $Q=3$ である．
\end{lemma}
\begin{proof}
$r=2-N_{10}\in[-3,2]$ は非零である．$J=1$ なら他の二 entries は
$1-N_{12},1-N_{13}\le1$ であり，一正二負と row sum $-1$ より，$r<0$ は $r=-1$ のみである．
$J\ne1$ なら，もう一つの高方向 entry を $s=1-N_{1J}$，低低 entry を $z$ とすると，
$r<0$ の場合には $s=1$，$z=-r-2\le0$ なので $r=-1,-2$ である．
$r=-1$ を除く factorization は補題~\ref{lh18:B-types} の証明と同一であり，
三高 core の仮定を用いない．$r=1,2$ の四型の計算も同様である．
型 I で $Q=3$ となる証明（補題~\ref{lh18:typeI-small}）が用いるのは，
$0B,0C,0D$ が高辺であること，式~\eqref{dl19:pairD} の $i=0$ と辺 $0D$ の node 条件だけである．
従って本 support にも適用できる．

$r=-2$ を調べる．$q=q_{01}=2x+1$，$x\ge1$，$a=2-N_{01}\in[-3,2]$ とすると
\[
 R_{01}=aQ-1,\quad X_{10}=-x,\quad X_{01}=(2\kappa+1)x+\kappa,
\]
\begin{equation}\label{dl19:negative-two-factor}
 (aQ-1)\{(2\kappa+1)x+\kappa\}=Q(2x+1)-1.
\end{equation}
$\kappa=1$ なら $a=1,2$ である．$a=1$ は $(Q-3)x=0$ を与え，
$a=2$ では左辺が右辺を $(4Q-3)x+Q>0$ だけ超える．
$\kappa=0$ でも $a=1,2$ であり，それぞれ
$(Q+1)x+Q-1=0$，$x=1-Q$ を要求して不可能である．
$\kappa=-k<0$ では $1\le k\le4$，$a=-A$，$0\le A\le3$ である．
$A\ge1,k\ge2$ または $A\ge2,k=1$ では左辺が右辺より大きい．
$A=k=1$ では $(Q-1)x=2$ なので $Q=3,x=1$ である．
$A=0$ では $e=Q-k$ と置くと $(2e+1)x=1-e$ であり，
$e\ge1$ と $e\le-1$ は符号が矛盾する．$e=0$ は $Q=k=3,x=1$ を与える．
従って常に $Q=3$ である．
\end{proof}

\begin{proposition}[全ての標準配置における葉分母の上界]
\label{dl19:leaf-bound}
$G_{5,9}$ を support とする五次の算術的接着データが存在するなら $Q\le13$ である．
\end{proposition}
\begin{proof}
$Q>13$ と仮定する．$h\le u$ と $Q=u(h+1)+1$ より，$u\le3$ なら $Q\le3(3+1)+1=13$ となる．従って $u\ge4$，$v\ge5$ である．
前補題により $B$ の接続は型 II または III，$r=R_{10}>0$ に限られる．
常に
\[
 \begin{gathered}
 R_{30}=1-v-N_{30}<0,\\
 R_{21}=2-u+c-N_{21},\qquad R_{23}=1-c-N_{23},\\
 R_{31}=v+2+d-N_{31},\qquad R_{32}=1-d-N_{32}.
 \end{gathered}
\]

$J=1$ の場合，型 II/III と一正二負により $R_{12},R_{13}\in\{-1,-2\}$ である．
core $3$ は二高なので $R_{31}>0$，$R_{32}\le0$ である．
高辺 $13$ の共通低 puncture $4$ における reverse partner の分類
（補題~\ref{ls11:low-puncture}）により $R_{31}=-R_{13}\le2$ である．
しかし row sum は
\[
 v-2+N_{30}=R_{31}+R_{32}\le2
\]
を与え，$v\ge5$ に反する．

以下 $J\ne1$ とする．core $1$ は二高であり，$r=1$ の場合には
$(R_{1J},R_{1E})=(-2,0)$ だけが可能である．
$r=2$ の場合には $(-3,0)$ または $(-2,-1)$ だけが可能である．
実際，低低 entry が $-2$ になる可能性を除き，低低 entry が $-1$ の場合には
二つの高方向 entries の絶対値は同じで二以上でなければならない．
従って常に $R_{1J}\in\{-2,-3\}$ である．

$J=2$ の場合，二高 core $3$ の低低 entry は $R_{31}\le0$ なので
$d\le3-v<0$，従って $c=0$ となる．すると $R_{21}=2-u-N_{21}\le-2$ である．
高辺 $12$ の共通低 puncture において $R_{12},R_{21}\le-2$ となるが，
同符号の負の reverse partners には必ず $-1$ が含まれるので矛盾する．

$J=3$ の場合，高辺 $13$ の reverse partner 分類から
\[
 R_{31}\in\{-1,2,3\}.
\]
$d=0$ なら $R_{31}\ge v-3\ge2$ であり，三高 core $3$ の一正二負条件は
$R_{32}\le-1$ を与える．従って
$v-2+N_{30}=R_{31}+R_{32}\le2$ となり矛盾する．よって $c=0$ である．
$R_{31}>0$ なら同じ row-sum 矛盾になるため，$R_{31}=-1$，
$R_{32}=v-1+N_{30}\ge4$ である．core $2$ は二高で，
$R_{20}=u+1-N_{20}\ge u-4\ge0$ は高方向 entry なので非零，従って正である．
よって $R_{23}=1-N_{23}<0$ であり，辺 $23$ の partner 分類から
$R_{23}=-R_{32}\ge-4$ である．
等号しか残らず，$v=5$，$N_{30}=0$，$N_{23}=5$ となる．
$v=t(h+1)=5$ から $(h,t,u,Q)=(4,1,4,21)$ であり，$N_{20}=N_{21}=0$ なので
$R_{21}=2-u=-2$ となる．これは core $2$ の低低 entry であり，二分補題に反する．
全配置で矛盾したので $Q\le13$ である．
\end{proof}

\subsection{無限の分母族を保持する完全な型検査}

上界 $Q\le13$ によって $(h,t)$ は
\[
 (1,1),\ldots,(1,6),\qquad (2,1),(2,2),\qquad (3,1)
\]
の九組になる．ここから他の分母や $c,d$ が有限であるとは結論しない．
実際，以下の検査では無限の型が現れ，その全体を処理する必要がある．

\begin{lemma}[hub の必要条件の厳密な型表示]
\label{dl19:finite-types}
上の九組，各 $J\in\{1,2,3\}$，各 $cd=0$ の枝，および四つの node 行を固定する．
$cd=0$ の枝を $(c,d)=(z,0)$ または $(0,z)$ と書き，$z$ は全整数を動くものとする．
以下の手順は，hub 行の必要な符号条件，hub の node 条件，および
$I$ 内の高辺の puncture $4$ における pair/clutching 条件を満たす可能性を全て含む．
\begin{enumerate}
\item 各 $0i$ について $R=R_{0i}$，$r=R_{i0}$ とし，$Rr\ne0$ を要求する．
各 $\kappa\in\{-4,-3,-2,-1,0,1\}$ に対し
\[
 \{Qr+R-\kappa Rr\}q_{0i}=R+r
\]
を解く．左係数が非零なら唯一の有理値について，奇数 $q_{0i}\ge3$，
$r\mid q_{0i}-1$，$R\mid Qq_{0i}-1$ を検査する．両辺の係数がともに零なら
\[
 q_{0i}=1+\operatorname{lcm}(2,|r|)a_{0i},\qquad a_{0i}\in\ZZ_{>0}
\]
の全無限族を保持する．
\item $JE,JF$ では，逆方向 entries $r,s$ の組を
\[
 s=-r,\qquad\text{または}\qquad
 \{r,s\}=\{\epsilon,\epsilon(q-1)\},\quad\epsilon\in\{1,-1\}
\]
に限る．通常の $s=-r$ では $q=1+\operatorname{lcm}(2,|r|)a$，$a\ge1$ の全族を残す．
例外の場合には $q=1+\max(|r|,|s|)$ である．
\item core $J$ の低 row に一正二負を課す．core $E,F$ では，高方向 entries が
異符号非零，低低 entry が非正であることを課す．低低 entry が非零の場合，
その row の高方向 entries は $-S,T$ と書け，$S\ge2$，$T\ge S$，$T\ne S+1$ を課す．
これらは必要条件であり，他の row の全局所実現を主張するものではない．
\end{enumerate}
固定された node 行に対し，以上の $z$ の集合は，有限個の点と，偶奇を指定した
完全な半直線の和として厳密に計算できる．全型のうち残るのは，点九百八型と半直線十八型である．
\end{lemma}

\begin{proof}
最初の式は $X_{0i}=(Qq_{0i}-1)/R$，$X_{i0}=(q_{0i}-1)/r$ と
$X_{0i}+X_{i0}=\kappa q_{0i}$ の分母を払ったものである．
零係数の場合は $R=-r$，$r\mid Q-1$ であり，指定した等差数列が全ての整除条件を尽くす．
二番目は共通低 puncture の補題である．固定した $r\ne0$ の許される相手 $s$ は
\[
 \{-r\}\ \cup\
 \begin{cases}\{\operatorname{sgn}(r)\},&|r|\text{ が正の偶数},\\
 \varnothing,&\text{それ以外}\end{cases}
 \ \cup\
 \begin{cases}\{2a\operatorname{sgn}(r):a\ge1\},&|r|=1,\\
 \varnothing,&|r|\ne1\end{cases}
\]
である．可変 row の二 entries は $z+b,-z+b'$ なので，これらの集合の交差は
点と偶奇付き半直線の演算で計算できる．
二高 row の条件も同様である．特に $R_{i0}=r>0$ なら他の高 entry $s$ は
$-r-1\le s\le-1$ の有限区間にあり，$r=-S<0$ なら $S\ge2$ とともに
\[
 s\in\{S-1,S\}\cup\{S+2,S+3,\ldots\}
\]
に限られる．これは低低 entry が零か負かの全ての場合を含む．
従って任意の打切りを導入せずに，有限な node 行と集合演算だけで全型を尽くせる．

付属の \texttt{verify\_g5\_9\_obstruction\_v19.py} はこの手順を整数演算で直接実装する．
四行には各二十一通りを全て使い，九組の $(h,t)$ と三配置を全て検査する．
$z=0$ は $cd=0$ の二つの枝に重複して現れるが，不存在判定に影響しない．
表示した型数はこの標識付き計数であり，同型類の個数ではない．
\end{proof}

\begin{lemma}[残る十八の半直線型は不可能]
\label{dl19:halflines}
補題~\ref{dl19:finite-types} の十八の半直線型は，全て正の二高局所模型の必要条件に反する．
\end{lemma}
\begin{proof}
十八型は全て $J=2$，$(c,d)=(0,z)$ であり，
\[
 (R_{10},R_{12},R_{13})=(1,-2,0),\qquad R_{21}=-1
\]
を満たす．共通低 puncture の例外接着から $q_{12}=3$ となる．
core $1$ の ray $4$ は二高標準形の零低低 row $C$ であり，その高方向 entries が
$1,-2$ なので，canonical な高 rays は $A=0,B=2$ である．
従って canonical な第二高分母 $b$ が三でなければならない．

しかし正の二高標準形で $b=3$ は不可能である．実際，$r+1\mid b-1$ より $r=1$，
$x=a-1$，$y=1$ であり，$u=-s<0$，$s\mid a-1$ から
\[
 a\le A=2a-1-s\le2a-2,\qquad B=(3a-1)/A<3.
\]
$B=1$ は $A=3a-1$ と矛盾し，$B=2$ は $v=B-b+y=0$ を与える．
いずれも標準形の条件に反する．従って十八の半直線の全点が排除される．
\end{proof}

\begin{theorem}[四価頂点に葉を持つ diamond support の不存在]
\label{dl19:no-G5-9}
五次の算術的接着データの高分母 support が $G_{5,9}=\{01,02,03,04,12,13\}$ となることはない．
葉以外の分母にも任意の上限を仮定しない．
\end{theorem}
\begin{proof}
命題~\ref{dl19:leaf-bound}，補題~\ref{dl19:finite-types}，補題~\ref{dl19:halflines} により，
九百八の点型だけを残せばよい．ここで点型とは $z$ の値が定まったことを意味し，
他の分母の自由な等差数列はそのまま残っている．

唯一の低辺 $EF$ については，二高模型の低低 pair の一方が零，他方が負である．
従って $R_{EF}<0$ なら $X_{EF}=0$，$R_{EF}=0$ なら $X_{EF}<0$ であり，
後者は $-2$ にはならない．同じことが $F$ 側にも成り立つ．
低辺 $EF$ の puncture $4$ における node 条件は
\[
 -4\le X_{EF}+X_{FE}\le1
\]
である．両 entries は非正なので，$R$ が零である側の $X$ は
$\{-1,-3,-4\}$ のいずれかに限られる．以上から $X_{EF},X_{FE}$ を有限に選べる．

他の五辺について，自由な分母をそれぞれ $q_{ij}=1+L_{ij}a_{ij}$，$a_{ij}\ge1$ と書く．
全 $X_{ij}$ はこの自由変数の整数係数の affine 式として復元される．
式~\eqref{dl19:pairD} の四本の $D$ 相殺は，有理数係数の線形系になる．
以下は三配置と二つの $cd=0$ 枝ごとの全計数である．可変 row 二は $c$，三は $d$ を表す．
\[
\begin{array}{cc|rr|rr}
 J&\text{可変 row}&z\text{ 点型}&z\text{ 半直線型}&D\text{ 不整合}&D\text{ 正値性矛盾}\\\hline
 1&2&69&0&179&4\\
 1&3&84&0&227&5\\
 2&2&108&0&321&3\\
 2&3&121&18&338&5\\
 3&2&142&0&398&4\\
 3&3&384&0&1100&12\\\hline
 \multicolumn{2}{c|}{\text{合計}}&908&18&2563&33
\end{array}
\]
低辺の $X$ の複数の選択により，点型九百八個から線形系二千五百九十六個が生じる．
二千五百六十三個では有理係数の行基本変形が $0=\gamma\ne0$ を与える．
残る三十三個は全て $Q=3$，$h=t=1$ であり，row $R_{0*}$ は順序を除き $(-7,2,2)$，
負 entry に対応する分母は五，残りの二分母は $1+2a,1+2b$ である．
$\sum_jV_{j0}=-1-u+v=0$ なので，$i=0$ の $D$ 相殺は
\[
 -2+(1+3a)+(1+3b)=-6,
 \qquad\text{すなわち}\qquad 3a+3b=-6
\]
を要求する．$a,b\ge1$ に反する．

付属検算は全ての線形系を \texttt{fractions.Fraction} によって処理し，
各矛盾について元の四式の有理線形結合を証明書に記録する．
さらに，記録する線形結合を元の係数行列へ再乗算して検査する．
自由変数には数値上限を課していないので，これは有限探索範囲内の不存在ではなく，
残る全分母族の排除である．
\end{proof}

本定理が排除するのは，固定された標準 $K_5$ branch graph と marked punctures を持つ
全平面五次 SNC 構成の当該 support である．一般型曲面上の任意の Lagrangian torus fibration
に対する非存在定理ではない．一高の任意整数 $c,d$，分母三の標準形の重複，二高の負の枝，
および他の分母の全自由族をいずれも保持している．

% v20. Complete G5_11 obstruction.  No denominator or scale cutoff.
\section{\texorpdfstring{$G_{5,11}$}{G5,11} 型の無上界排除}
\label{g51120:section}

\begin{theorem}[二つの三価頂点を持つ七高辺型の障害]
\label{g51120:no-support}
五次の算術的接着データの高分母 support は
\[
 H=G_{5,11}=\{01,02,03,04,12,13,24\}
\]
と同型ではあり得ない．分母の上限は仮定しない．
正の整数局所 charts，marked clutching，非負 node，対向 $D$ の相殺を用いる．
\end{theorem}

この graph では $3,4$ が二価，$1,2$ が三価，$0$ が四価である．
低辺は $14,34,23$ であり，$0$ は二つの二価頂点の共通高近傍である．
証明では最初に二価 cores の reverse data を数学的に有限化する．
その後，分母と scale の全無限族を保持した厳密な整数演算で排除する．

\subsection{四つの向きと単位 entry の禁止}

二価 core の canonical order $(A,B,C,D)$ では
\begin{equation}\label{g51120:canonical}
 P=\begin{pmatrix}
 0&-U&x&-s\\
 -V&0&-y&z\\
 r&-r-1&0&0\\
 -S&T&-L&0
 \end{pmatrix},
\end{equation}
\[
 \begin{gathered}
 a=sS+1,\quad b=zT+1,\quad x=sS/r,\quad y=zT/(r+1),
 \\
 U=a+x-s,\quad V=b-y+z,\quad L=T-S+1.
 \end{gathered}
\]
全ての $r,s,z,S,T,x,y$ は正整数で，$a,b$ は奇数，$UV=ab-1$ である．
既知の二分補題により，正の枝は
\begin{equation}\label{g51120:positive}
 S>r,\qquad rT>(r+1)S,\qquad L\ge3,
\end{equation}
負の枝は
\begin{equation}\label{g51120:negative}
 S=T=w\ge2,\quad r=w+H,\quad
 \frac{w^2}{2}<H(H+1)<2w^2
\end{equation}
を満たす．特に $S\ge2$ である．正の枝では $T\ge5$
（補題~\ref{cl16:T-five}）である．

\begin{lemma}[正の単位 entry の一つの禁止]\label{g51120:no-x-one}
式~\eqref{g51120:canonical} では $x=1$ は起こらない．
\end{lemma}
\begin{proof}
正の枝では $x=sS/r>1$ である．負の枝で $x=1$ と仮定すると
$r=sS$，$s\ge2$，$T=S$ である．高高 pair equation を代入して整理すると
\[
 z\{(-s^2+3s-1)S+2-s\}+(2-s)(sS+1)=0
\]
を得る．$s=2$ では左辺は $zS>0$，$s\ge3$ では両項が負であり，
どちらも不可能である．
\end{proof}

二価 core を $i\in\{3,4\}$，もう一方を $j$，
$i$ 固有の高近傍を $h$，残る三価頂点を $k$ と書く．
すなわち $(i,j,h,k)=(3,4,1,2)$ または $(4,3,2,1)$ である．
canonical 標識の四通りを次表で固定する．最後の欄は ray $j$ の行を
列 $(0,h,k)$ の順で記したものである．
\begin{center}\small
\begin{tabular}{ccccccl}
型&$A$&$B$&$C$&$D$& &$(p^i_{j0},p^i_{jh},p^i_{jk})$\\\hline
I&$0$&$h$&$j$&$k$&&$(r,-r-1,0)$\\
II&$h$&$0$&$j$&$k$&&$(-r-1,r,0)$\\
III&$0$&$h$&$k$&$j$&&$(-S,T,-L)$\\
IV&$h$&$0$&$k$&$j$&&$(T,-S,-L)$
\end{tabular}
\end{center}

共有高辺 $ih$ の共通低 puncture $j$ で
\[
 u=p^i_{jh},\qquad v=p^h_{ji}
\]
と置く．補題~\ref{ls11:low-puncture} を reverse entries に書き換えると
\begin{equation}\label{g51120:partner}
 v=-u,\qquad \{u,v\}=\{1,q_{ih}-1\},\qquad
 \{u,v\}=\{-1,-(q_{ih}-1)\}
\end{equation}
の三択である．通常型 $v=-u$ では $|u|\le4$
（補題~\ref{cl16:ordinary-bound}）である．また $u<0$ なら，
相手の二価 core $j$ における $p^j_{ih}$ が低低 entry で非正なので
\[
 1-u-p^j_{ih}=n_{ij;h}\le5
\]
から例外型も含めて $|u|\le4$ となる．

\subsection{正の枝の III 型を先に排除する}

\begin{lemma}\label{g51120:no-positive-III}
二つの二価 cores のいずれも，正の枝の III 型には属さない．
\end{lemma}
\begin{proof}
対称性により core $3$ が正の III 型であるとする．このとき $u=T\ge5$ である．
低辺 $34$ の puncture $1$ の node 非負性から
$p^4_{31}\le1-T<0$ であり，core $4$ では $D=3,C=1$ となる．
従って $p^4_{13}=0$ であり，低辺 $14$ の node を用いると
\[
 v=p^1_{43}=1-n_{14;3}\le1.
\]
通常型は $T>4$ に反するので，式~\eqref{g51120:partner} は
$v=1$，$q_{13}=T+1$ を強制する．従って $z=1$，$n_{14;3}=0$ である．

三高 matching 定理~\ref{mh18:matching} より，core $1$ の低 row の残り二成分は負である．
core $4$ の $C=1$ の行は，高 rays に $r',-r'-1$ を持つ．
後者に対応する node は少なくとも $r'+3$ なので，総数五から $r'\le2$ を得る．
$S'\ge2$ より core $4$ は正の枝である．

もし core $4$ が IV 型なら，低辺 $34$ の puncture $2$ の node は
$1+L+S'\ge6$ となり不可能である．従って core $4$ も III 型である．
puncture $0$ の node は $1+S+S'\le5$ なので $S=S'=2$，$r=r'=1$ となる．
残る二つの node は $T'-T,T-T'$ であり，$T=T'$ である．
同じ議論を反対側に適用すると $z'=1$ である．

$r=1,S=2,z=1$ の pair equation は
\[
 s(8-T)=T-4
\]
である．$T>4$ と $y=T/2\in\ZZ$ より $T$ は六以上の偶数なので，
$T=6,s=1$ に限られる．両 core は
\[
 (r,s,z,S,T)=(1,1,1,2,6),\qquad (a,b)=(3,7)
\]
となる．core $4$ の ray $1=C$ の係数は $D^4_1=5$ である．
core $1$ の低 row は，matching と node 総数五から
\[
 (p^1_{40},p^1_{42},p^1_{43})=(-1,-1,1)
\]
に強制される．pair equations による低 column の和から
$D^1_4=q_{01}+q_{12}-9$ を得るので，辺 $14$ の $D$ 相殺は
$q_{01}+q_{12}=4$ を要求する．両分母は三以上なので矛盾する．
\end{proof}

\subsection{全 reverse data に対する証明された有限上界}

\begin{proposition}\label{g51120:reverse-bounds}
残る二価 cores の canonical reverse data は全て
\begin{equation}\label{g51120:box}
 1\le r\le9,\qquad 2\le S\le7,\qquad S\le T\le13
\end{equation}
を満たす．さらに I 型では $r\le3$，II 型では $r\le4$，
IV 型では $S\le4$ であり，III 型は負の枝に限られる．
この上界は分母 $a,b$ または scales $s,z$ の上限を意味しない．
\end{proposition}

\begin{proof}
I 型は $u=-r-1<0$ なので $r\le3$ である．
補題~\ref{tn14:small-negative-branch} により，この範囲に負の枝はない．
正の枝では $y=zT/(r+1)>1$ なので，$u$ の負例外型は起こらず，
必ず $v=r+1$ の通常型である．

II 型で $r>4$ なら通常型は不可能であり，正例外型は
$v=1,q_{ih}=r+1$ を要求する．これは $x=1$ に反する．従って $r\le4$ である．
III 型は負の枝だけが残る．その $u=T=w$ が四より大きい場合も
正例外型が $q_{ih}=T+1$，すなわち $z=1$ を要求するが，
$y=T/(r+1)<1$ となって不可能である．従って $w\le4$ である．
IV 型では $u=-S$ なので $S\le4$ である．
以上から，どの型でも負の枝なら $w\le4$ であり，
式~\eqref{g51120:negative} は $H\le5$，$r=w+H\le9$ を与える．
以下では正の枝だけを調べる．

\paragraph{IV 型．}
相手が I または IV 型なら，両者の ray $34$ の行が高 ray $0$ に正成分を持ち，
$n_{34;0}<0$ となる．相手が II 型なら
$n_{34;0}=r'+2-T\ge0$ なので $T\le6$ である．
相手が III 型なら，既に負の枝で $S'=T'\le4$ であるから，
$n_{34;0}=1+S'-T\ge0$ により $T\le5$ である．
従って IV 型は $r<S\le4,T\le6$ を満たす．

\paragraph{I 型の相手と上界．}
相手が I または IV 型なら，やはり puncture $0$ の node が負になる．
相手が III 型なら，I 型の低低 entry が零であるため，
相手固有の高 puncture の node が $1-T'<0$ となる．
従って相手は II 型であり，その同じ node が $1-r'\ge0$ を与えるので $r'=1$ である．

core $3$ を I 型，core $4$ を II 型とする．core $1$ の低 row の成分
$v=p^1_{43}=r+1$ は正なので，残る二成分は負である．
core $4$ の ray $1=D$ における逆向き成分は $T',-S',-L'$ であり，
必要な node は
\[
 n_{14;3}=L'-r,\qquad n_{14;2}\ge S'+2.
\]
従って $S'+L'\le r+3$，すなわち $T'\le r+2\le5$ である．
正の枝の $T'\ge5$ と $r'=1$ を合わせると
\[
 r=3,\qquad (r',S',T')=(1,2,5)
\]
を得る．この相手の pair equation と整数性は
\[
 (s'-3)(z'+2)=-8,\qquad z'\in2\ZZ_{>0}
\]
であり，$(s',z')=(1,2),(2,6)$ に限られる．従って $q_{24}=3$ または $5$ である．

辺 $24$ の共通低 puncture $3$ では，相手の reverse が $r'=1$ なので，
$R=p^2_{34}$ は $-1$ または $q_{24}-1$ である．後者なら matching により
$p^2_{30}<0$ なので，辺 $23$ の node を $m$ と書くと
\[
 m_0\ge S+2,\qquad m_4=L+2-q_{24}.
\]
その総数五から $T\le q_{24}\le5$ を得るが，$r=3,S>3$ では $T\ge6$ であり矛盾する．
従って $R=-1$，$m_4=L+2\le5$ である．正の枝の $L\ge3$ より
$L=3$，$3<S<6$ となる．従って $(S,T)=(4,6),(5,7)$ だけが残る．

\paragraph{II 型と II 型．}
両者の共通高 ray $0$ への成分は $-r-1,-r'-1$ なので
$n_{34;0}=r+r'+3\le5$，従って $r=r'=1$ である．両者は正の枝に属する．
辺 $13$ の通常型は $p^1_{43}=-1$ を要求するが，
このとき $n_{14;3}=L'+2\le5$ なので $L'=3$ となる．
これは $r'=1$，$S'>1$，$r'(T'-S')>S'$ と両立しない．
従って両方の共通 punctures は正例外型である．matching と node 総数から
\[
 T'\le q_{13}\le T'-S'+3,\qquad
 T\le q_{24}\le T-S+3
\]
となり，$S,S'\le3$ が従う．
$r=1$ の pair equation は
\[
 z(s+1)T=(4S-2)sz+2(S-1)s+2z+2
\]
であり，$z\ge1$ から
\[
 T\le\frac{(6S-4)s+4}{s+1}<6S-4\le14
\]
を得る．従って $T\le13$ である．

\paragraph{II 型と IV 型．}
II 型の相手が III 型なら，その固有高 puncture の node は $1-T'<0$ である．
I 型との組は既に扱ったので，残るのは IV 型との組だけである．
core $3$ を正の II 型，core $4$ を IV 型とする．$S'\le4$ は既知である．
辺 $24$ の共通低 puncture $3$ では，core $4$ の reverse は $-S'<-1$ なので
$R=p^2_{34}$ は $S'$ または $-1$ に限られる．
通常の $R=S'>0$ なら，matching は $p^2_{31}<0$ を要求し，
\[
 m_1\ge S+2,\qquad m_4=L+1-S'
\]
から $T\le S'+1\le5$ を得る．負例外の $R=-1$ なら
$m_4=L+2\le5$ であり，正の枝より $L=3$ となる．
従って $r<S<2r\le8$，$S\le7,T\le9$ である．
これで四つの向きの全てと全ての相手の組を尽くし，主張が従う．
\end{proof}

\subsection{全 scale の点・半直線表示}

$h=S-r$，$\delta=T-S$ と置くと，高高 pair equation は
\begin{equation}\label{g51120:bilinear}
 Fsz+As+Bz+C=0
\end{equation}
に等しい．ただし
\[
 F=(2r+1)h-r^2(\delta-1),\quad A=h(r+1),\quad
 B=-r(h+\delta-1),\quad C=r(r+1).
\]
本節では $L\le4$ を仮定できないため，$F=0$ の場合も保持する．
$F\ne0$ なら
\begin{equation}\label{g51120:factor}
 (Fs+B)(Fz+A)=K,\qquad K=AB-FC
\end{equation}
である．$K\ne0$ ならその全 signed divisors を列挙すればよい．
$K=0$ なら $Fs+B=0$ または $Fz+A=0$ の全直線を残す．
$F=0$ では式~\eqref{g51120:bilinear} 自体が整数一次式になる．
負の枝の $F=0$ は $2r+1\mid r^2$ を要求するので起こらず，
この場合は $A>0,B<0$ であり，正の整数解を全等差半直線として表示できる．

各直線上で
\[
 r\mid sS,\qquad r+1\mid zT,\qquad
 sS\equiv zT\equiv0\pmod2
\]
を課す．例えば parameter を $2r(r+1)$ を法として分ければ全整数解を尽くす．
正値性で始点をずらすと，各族は
\begin{equation}\label{g51120:halfline}
 (s,z)=(s_0,z_0)+t(d_s,d_z),\qquad t\in\ZZ_{\ge0}
\end{equation}
となる．どの半直線も途中で打ち切らない．

命題~\ref{g51120:reverse-bounds} の範囲で，全ての点と半直線を作ると，
一つの二価 core あたり次の表になる．
\begin{center}\small
\begin{tabular}{crrr}
向き&点型&全半直線型&合計\\\hline
I&27&30&57\\
II&31&30&61\\
III&8&0&8\\
IV&24&18&42\\\hline
合計&90&78&168
\end{tabular}
\end{center}
これは整除・偶奇の剰余類による分割も含む表示数であり，幾何的模型の個数ではない．
無限族を持つ reverse triples は，次の四つだけである：
\begin{center}\small
\begin{tabular}{ccccl}
$r$&$S$&$T$&場合&全正整数 scale の表示\\\hline
1&2&6&$F=0$&$s=2z-1$, $z\ge1$\\
2&6&12&$F=0$&$s=3a-2$, $z=2a-1$, $a\ge1$\\
2&3&6&$K=0$&$s=2a$, $z=1$, $a\ge1$\\
3&4&6&$K=0$&$s=3a$, $z=2$, $a\ge1$
\end{tabular}
\end{center}
後二つは固定分母七・十三の共鳴族である．前二つの affine 族も残すことが必要である．

\subsection{低行の整数接着と最後の parity 障害}

各二価 core の族を一つずつ選び，まず低辺 $34$ の三つの node
\[
 n_{34;k}=1-p^3_{4k}-p^4_{3k},\qquad k=0,1,2
\]
が非負であるものを残す．これらは reverse data だけで定まり，
scale の自由 parameter には依存しない．この段階の型対は $3592$ 個である．

低辺 $14,23$ の node 分配をそれぞれ
\[
 (n_0,n_2,n_3),\qquad(m_0,m_1,m_4)
\]
とし，非負整数で各和が五である全二十一通りを使う．
三高 cores $1,2$ の唯一の低 row は
\begin{equation}\label{g51120:low-rows}
 R^1_k=1-n_k-p^4_{1k}\quad(k=0,2,3),\qquad
 R^2_k=1-m_k-p^3_{2k}\quad(k=0,1,4)
\end{equation}
である．全成分が非零で，正成分がちょうど一つであるものだけを残す．
さらに $(p^3_{41},R^1_3)$ と $(p^4_{32},R^2_4)$ に
式~\eqref{g51120:partner} の三択を課すと，$1395$ 配置が残る．

通常型は分母に新しい式を課さない．例外型では，二つの reverse entries の和の絶対値が
共有高分母 $q_{13}$ または $q_{24}$ に等しい．各分母は式~\eqref{g51120:halfline} の
parameter に関して一次式であり，この条件は正確な整数の等式として処理する．
定数の不一致で $759$ 配置，parameter が整数にならないことで $550$ 配置が消え，
合計 $1309$ 配置が排除される．残る $86$ 配置では，それぞれの parameter は
一つに固定されるか，全ての非負整数を動くままである．

最後に，低辺 $34$ の対向係数を足す．canonical な低 rays の係数は
\[
 D_C=T-S+y-x,\qquad D_D=s-z-1
\]
なので，その残差 $D^3_4+D^4_3$ は残る自由 parameters に関する一次式である．
全配置の結果は次表のとおりであり，零になるものはない．
\begin{center}\small
\begin{tabular}{crl}
向きの組&配置数&$D^3_4+D^4_3$ の障害\\\hline
I--II または II--I&32&全自由 parameters に対して奇数\\
II--II&5&四配置で $8$，一配置で $28$\\
II--IV または IV--II&40&全自由 parameters に対して奇数\\
III--III&9&常に $-4$\\\hline
合計&86&全配置で相殺不可能
\end{tabular}
\end{center}

最後の表は，node 配分と scale の剰余類表示の重複を除くと，
$3\leftrightarrow4$ の交換を除いて十組の canonical models へ圧縮できる．
五つ組を $(r,s,z,S,T)$ の順とし，
\[
 \begin{gathered}
 F_a=(3,3a,2,4,6)\quad(a\ge1),\\
 A=(1,1,2,2,5),\quad B=(1,2,6,2,5),\quad C=(1,3,2,2,6),\\
 E=(2,2,6,3,5),\quad N_0=(4,4,5,2,2),\\
 N_1=(7,7,2,4,4),\quad N_2=(7,21,4,4,4),\\
 N_3=(8,6,9,4,4),\quad N_4=(8,8,27,4,4)
 \end{gathered}
\]
と書くと，残差は次のようになる．
\begin{center}\small
\begin{tabular}{ccl}
一方の向き・模型&他方の向き・模型&$D^3_4+D^4_3$\\\hline
I, $F_a$&II, $A$&$11-4a$\\
I, $F_a$&II, $B$&$19-4a$\\
II, $F_a$&IV, $A$&$3-4a$\\
II, $B$&II, $B$&$28$\\
II, $C$&II, $C$&$8$\\
II, $E$&IV, $N_1$&$13$\\
II, $E$&IV, $N_2$&$25$\\
II, $E$&IV, $N_3$&$5$\\
II, $E$&IV, $N_4$&$-11$\\
III, $N_0$&III, $N_0$&$-4$
\end{tabular}
\end{center}
従って，残る唯一の無限族 $F_a$ についても，全正整数 $a$ に対して矛盾が成立する．

この有限表は，証明から得た reverse 上界と signed-divisor の完全因数分解によるものである．
分母や半直線 parameter の数値上限は使用していない．
\nolinkurl{verify_g5_11_obstruction_v20.py} は，全 $168$ 型の affine 行列と
全 $1395$ 配置を生成し，元の行和と pair products を parameter の多項式恒等式として確認する．
結果 JSON には，各 node 行，二つの低 rows，例外型の要求分母，parameter の全範囲，
および最後の奇数残差または非零定数を保存する．
別実装は主生成器を import せず，canonical データからこれらを再構成できる．

補題~\ref{g51120:no-positive-III} と最後の完全表により，四つの向きの全てが排除される．
従って定理~\ref{g51120:no-support} が従う．

% v21. G5_15: complete reduction of the common-low-puncture branches.
% This is not a G5_15 exclusion theorem. No denominator cutoff is used.
\section{\texorpdfstring{$G_{5,15}$}{G5,15} 型の有限中心と小分母への還元}
\label{g51521:section}

本節では低辺が $23,24$ の二本だけである支持
\[
 H=\{01,02,03,04,12,13,14,34\}=G_{5,15}
\]
を扱う．core $2$ は二高，cores $3,4$ は三高，cores $0,1$ は四高である．
標識 $0,1$ と $3,4$ の交換により，core $2$ を次の canonical order $(0,1,3,4)$ に置ける：
\begin{equation}\label{g51521:central}
 P^2=\begin{pmatrix}
 0&-U&x&-s\\
 -V&0&-y&z\\
 r&-r-1&0&0\\
 -S&T&-L&0
 \end{pmatrix},\qquad
 \begin{gathered}
 a=sS+1,\quad b=zT+1,\\
 x=sS/r,\quad y=zT/(r+1),\\
 U=a+x-s,\quad V=b-y+z,\quad L=T-S+1.
 \end{gathered}
\end{equation}
全ての $r,s,z,S,T,x,y$ は正整数，$a,b$ は奇数，$UV=ab-1$ である．
二高分母の二分補題によって，正の枝では
\begin{equation}\label{g51521:positive-central}
 S>r,\qquad rT>(r+1)S,\qquad L\ge3,\qquad T\ge5,
\end{equation}
負の枝では
\begin{equation}\label{g51521:negative-central}
 S=T=w\ge2,\qquad r=w+H,\qquad
 w^2<2H(H+1)<4w^2
\end{equation}
となる．特に負の枝では $H\ge2$ である．低 rays の係数は
\begin{equation}\label{g51521:low-D}
 D_C=D^2_3=T-S+y-x,\qquad D_D=D^2_4=s-z-1
\end{equation}
である．

\begin{theorem}[有限中心または分母三・五への完全還元]
\label{g51521:reduction}
正の整数局所 charts，奇数分母，marked clutching，非負 node と node 総数五，
および対向 $D$ の相殺を持つ上記の $G_{5,15}$ 支持を考える．
共有高辺 $34$ と共通低 puncture $2$ における reverse entries を
\[
 u=p^3_{24},\qquad v=p^4_{23},\qquad Q=q_{34}>1
\]
と置くと，次のいずれかが必要である．
\begin{enumerate}
\item 通常型 $v=-u$ であり，core $2$ は表~\ref{g51521:center-table} の四模型の一つである．
三高 cores の低 rows と低 $D$ は表~\ref{g51521:ordinary-row-table} の十四配置に限られる．
\item 負例外型であり，$Q\in\{3,5\}$ である．全ての可能性は
命題~\ref{g51521:negative-templates} の二つの template に含まれる．
\end{enumerate}
正例外型 $\{u,v\}=\{1,Q-1\}$ は，分母の上界を仮定せず排除される．
\end{theorem}

したがって $Q\ge7$ なら二価 core の二高分母は四つの組
\[
 (a,b)=(7,31),\ (9,11),\ (31,7),\ (19,43)
\]
のいずれかであり，両方とも四十三以下である．
これは $Q$ 自身の上界でも，支持 $G_{5,15}$ の排除でもない．
低辺の補 graph が $P_3$ である支持全体についての主張であり，
大きい支持に同じ部分 graph が含まれるという理由だけで適用してはならない．

\subsection{低 rows と reverse の三分岐}

低辺の node 分配を
\[
 (n_{23;0},n_{23;1},n_{23;4})=(n_0,n_1,n_4),\qquad
 (n_{24;0},n_{24;1},n_{24;3})=(m_0,m_1,m_3)
\]
とする．各成分は非負整数で，各 triple の和は五である．
三高 cores の唯一の低 row は
\begin{equation}\label{g51521:low-rows}
 \begin{split}
 R^3&=(1-r-n_0,\ r+2-n_1,\ 1-n_4),\\
 R^4&=(S+1-m_0,\ 1-T-m_1,\ L+1-m_3),
 \end{split}
\end{equation}
である．列順はそれぞれ $(0,1,4)$，$(0,1,3)$ である．
三高 matching 定理~\ref{mh18:matching} により，全成分は非零で，
正成分はちょうど一つである．特に
\[
 u=1-n_4\le1,\qquad v=L+1-m_3.
\]
共通低 puncture の補題~\ref{ls11:low-puncture} により
\begin{equation}\label{g51521:three-branches}
 v=-u,\qquad \{u,v\}=\{1,Q-1\},\qquad
 \{u,v\}=\{-1,-(Q-1)\}
\end{equation}
の三択である．通常型では補題~\ref{cl16:ordinary-bound} から $|u|\le4$ である．

\begin{lemma}\label{g51521:positive-reverse}
$u>0$ なら
\[
 u=1,\qquad R^3=(-1,-1,1),\qquad r\in\{1,2\}.
\]
このとき $(n_0,n_1,n_4)=(1,4,0)$ または $(0,5,0)$ である．
また $v>0$ なら $m_0\ge S+2$，従って $S\le3$ である．
\end{lemma}
\begin{proof}
$u>0$ は $n_4=0,u=1$ を与える．残りの二成分が負なので
$r+n_0\ge2$，$n_1\ge r+3$ であり，$n_0+n_1=5$ から
$r+n_0=2$ を得る．これが最初の主張を与える．
core $4$ の第二成分 $1-T-m_1$ は常に負なので，$v>0$ なら
第一成分も負でなければならない．従って $m_0\ge S+2$ である．
\end{proof}

\subsection{通常型の四つの中心模型}

\begin{proposition}\label{g51521:ordinary-reverse}
通常型 $v=-u$ では，core $2$ の reverse triple $(r,S,T)$ は
\[
 (2,3,5),\quad (4,2,2),\quad (5,3,3),\quad (6,3,3)
\]
の四つに限られる．
\end{proposition}
\begin{proof}
まず $u=1,v=-1$ とする．前補題から $r\le2$ なので，負の枝は不可能である．
$m_3=L+2\le5$ と $L\ge3$ により $L=3$ である．
従って $T=S+2$，$r<S<2r$ であり，$(r,S,T)=(2,3,5)$ に限られる．

次に $u=-h,v=h$，$1\le h\le4$ とする．
$m_0\ge S+2$，$m_3=L+1-h=T-S+2-h$ により
\[
 T+4-h\le m_0+m_3\le5,
\]
従って $T\le h+1\le5$ である．正の枝では $T\ge5$ なので
$T=5,h=4$ となる．このとき $n_4=5,n_0=n_1=0$ であり，
$R^3=(1-r,r+2,-4)$ の matching 性から $r\ge2$ を得る．
$S\le3$ と $S>r$ を合わせると，やはり $(r,S,T)=(2,3,5)$ となる．

負の枝では $S=T=w\le3$，$L=1$ であり，
$m_3=2-h\ge0$ から $h\le2$ である．
式~\eqref{g51521:negative-central} は
$(w,H)=(2,2),(3,2),(3,3)$ だけを許す．これが残りの三つの reverse triples を与える．
\end{proof}

\paragraph{全 scale の復元．}
$h=S-r$，$\delta=T-S$ と置く．高高 pair equation は
\begin{equation}\label{g51521:bilinear}
 Fsz+As+Bz+C=0,
\end{equation}
\[
 F=(2r+1)h-r^2(\delta-1),\quad
 A=h(r+1),\quad B=-r(h+\delta-1),\quad C=r(r+1)
\]
となる．$F\ne0$ では
\begin{equation}\label{g51521:factorization}
 (Fs+B)(Fz+A)=AB-FC
\end{equation}
である．上記の四 triples では右辺は順に $-18,-80,-270,-378$ で，零ではない．
全 signed divisors を取り，$s,z>0$，$r\mid sS$，$r+1\mid zT$，
$sS,zT\in2\ZZ$ を課すと，表~\ref{g51521:center-table} の四点を得る．
分母にも scale にも探索上限を課していない．

\begin{table}[htbp]
\centering\small
\begin{tabular}{c|ccccc|cc|rr}
模型&$r$&$s$&$z$&$S$&$T$&$a$&$b$&$D_C$&$D_D$\\\hline
$M_1$&2&2&6&3&5&7&31&9&$-5$\\
$M_2$&4&4&5&2&2&9&11&0&$-2$\\
$M_3$&5&10&2&3&3&31&7&$-5$&7\\
$M_4$&6&6&14&3&3&19&43&3&$-9$
\end{tabular}
\caption{通常型に残る全中心模型．}\label{g51521:center-table}
\end{table}

式~\eqref{g51521:low-rows} で各 node triple の全二十一通りを取り，
matching と $u+v=0$ を課すと十八配置を得る．
低 row の三成分が全て奇数なら，奇数分母と pair products により低 column の三成分は全て偶数なので，
低 $D$ は奇数である．$M_2$ における
\[
 R^3=(-j,j,-1),\quad 3\le j\le6,\qquad R^4=(-1,-1,1)
\]
の四配置は，対向端相殺が $D^4_2=2$ を要求するため排除される．
残る十四配置を表~\ref{g51521:ordinary-row-table} に示す．
各 $j$ と各表示された row の選択を別の配置として数える．

\begin{table}[htbp]
\centering\small
\begin{tabular}{c|c|c|c|c|r}
中心&$R^3$&$R^4$&$j$ の範囲&$(D^3_2,D^4_2)$&個数\\\hline
$M_1$&$(-1,4,-4)$&$(-1,-4,4)$&---&$(-9,5)$&1\\
$M_1$&$(-1,-1,1)$&$(4,-4,-1)$&---&$(-9,5)$&1\\
$M_2$&$(-j,j+1,-2)$&$(-1,-2,2)$ または $(-2,-1,2)$&$3,4,5$&$(0,2)$&6\\
$M_3$&$(-j,j+1,-2)$&$(-1,-2,2)$&$4,5,6$&$(5,-7)$&3\\
$M_4$&$(-j,j+1,-2)$&$(-1,-2,2)$&$5,6,7$&$(-3,9)$&3
\end{tabular}
\caption{通常型に残る十四の低 row 配置．}\label{g51521:ordinary-row-table}
\end{table}

これらの配置では $Q$ は未固定であり，
$m=|u|=|v|$ に対して $m\mid Q-1$ と $Q$ の奇数性が必要である．
表は必要条件の完全還元であり，十四配置のいずれについても大域接着の存在を主張しない．

\subsection{正例外型の無上界排除}

\begin{proposition}\label{g51521:no-positive-exception}
正例外型 $\{u,v\}=\{1,Q-1\}$ は起こらない．
\end{proposition}
\begin{proof}
$u\le1$ なので $u=1,v=Q-1\ge2$ である．
補題~\ref{g51521:positive-reverse} により
\[
 R^3=(-1,-1,1),\quad r\le2,\quad S\le3.
\]
中心は正の枝に属し，$(r,S)=(1,2),(1,3),(2,3)$ に限られる．
高高 pair equation を $T$ について解くと
\[
 T=\frac{r+1}{r}\,
 \frac{\{(r+1)S-r\}s+r+\{(S-r)s+r\}/z}{rs+1}.
\]
$z\ge1$ を用いると，$r=1$ では
\[
 T\le\frac{(6S-4)s+4}{s+1}<6S-4,
\]
$(r,S)=(2,3)$ では $T\le6$ である．
従って reverse triples は次の十二個の範囲で全て尽くされる：
\[
 (r,S,T)=(1,2,5\ldots7),\quad (1,3,7\ldots13),\quad (2,3,5\ldots6).
\]

core $4$ の低 row を $(-c,-d,Q-1)$ と書くと $c+d=Q$ であり，
低 column は $(-A,-B,1)$ と書ける．ここで $A,B>0$ である．
$Q$ は奇数なので $c,d$ の一方は奇数であり，その ray の奇数分母は対応する scale の偶数性を要求する．
従って $A+B\ge3$ であり，
\[
 D^4_2=A+B-2\ge1.
\]
対向端相殺 $D^4_2=-D_D=z+1-s$ は
\begin{equation}\label{g51521:scale-inequality}
 s\le z
\end{equation}
を与える．

全十二 reverse triples に式~\eqref{g51521:bilinear} を適用する．
$F(AB-FC)\ne0$ なら式~\eqref{g51521:factorization} の全 signed divisors を取り，
$F=0$ または $AB-FC=0$ の全直線も保持する．
この範囲で無限族を持つ reverse triples は
\[
 (r,S,T)=(1,2,6):\quad s=2z-1,\ z\ge1,
 \qquad
 (r,S,T)=(2,3,6):\quad z=1,\ s\in2\ZZ_{>0}
\]
だけである．式~\eqref{g51521:scale-inequality} は前者を $s=z=1$ に限定し，後者を排除する．
全有限点にも同じ不等式を課すと，中心は次の八個だけになる：
\[
 \begin{gathered}
 (r,s,z,S,T)=(1,1,2,2,5),\ (1,2,6,2,5),\ (1,1,1,2,6),\\
 (1,2,5,3,8),\ (1,2,2,3,9),\ (1,4,6,3,9),\\
 (1,6,26,3,9),\ (2,2,6,3,5).
 \end{gathered}
\]

core $3$ の低 row $(-1,-1,1)$ に対して，低 column は $(-A',-B',Q-1)$ である．
二つの高分母が奇数で三以上なので $A',B'\ge2$ であり，
$D^3_2=-D_C$ は
\begin{equation}\label{g51521:positive-low-D-obstruction}
 A'+B'=Q-D_C\ge4
\end{equation}
を要求する．残る node 分配を全て入れると，次の八配置となる．
$(1,2,5,3,8)$ は $Q=8$ を要求するため，奇数性で既に消えている．

\begin{center}\small
\begin{tabular}{c|c|c|r|r}
$(r,s,z,S,T)$&$(n_0,n_1,n_4)$&$(m_0,m_1,m_3)$&$Q$&$Q-D_C$\\\hline
$(1,1,2,2,5)$&$(1,4,0)$&$(4,0,1)$&5&$-1$\\
$(1,2,6,2,5)$&$(1,4,0)$&$(4,0,1)$&5&$-9$\\
$(1,1,1,2,6)$&$(1,4,0)$&$(4,1,0)$&7&2\\
$(1,1,1,2,6)$&$(1,4,0)$&$(5,0,0)$&7&2\\
$(1,2,2,3,9)$&$(1,4,0)$&$(5,0,0)$&9&0\\
$(1,4,6,3,9)$&$(1,4,0)$&$(5,0,0)$&9&$-12$\\
$(1,6,26,3,9)$&$(1,4,0)$&$(5,0,0)$&9&$-96$\\
$(2,2,6,3,5)$&$(0,5,0)$&$(5,0,0)$&5&$-4$
\end{tabular}
\end{center}
全て式~\eqref{g51521:positive-low-D-obstruction} に反する．
従って正例外型は排除される．
\end{proof}

\subsection{残る負例外型の正確な形}

\begin{proposition}\label{g51521:negative-templates}
負例外型では，次の二つの template が全ての可能性を覆う．
いずれの template でも中心の scale 条件，三高局所方程式，および未使用の接着条件は引き続き必要である．
\begin{enumerate}
\item $u=-1,v=-(Q-1)$ の場合は $Q=3,L=1$ であり，中心は負の枝に属する．
$S=T=w$ とし，$0\le a_0\le3$，$0\le b_0\le1$ に対し
\[
 (n_0,n_1,n_4)=(a_0,3-a_0,2),\qquad
 (m_0,m_1,m_3)=(b_0,1-b_0,4),
\]
\[
 R^3=(-j,j,-1),\quad j=r+a_0-1,\qquad
 R^4=(k+1,-k,-2),\quad k=w-b_0
\]
となる．
\item $u=-(Q-1),v=-1$ の場合は $Q=3$ または五，$L=1$ または三である．
$0\le a_0\le5-Q$，$0\le b_0\le3-L$ に対し
\[
 (n_0,n_1,n_4)=(a_0,5-Q-a_0,Q),\qquad
 (m_0,m_1,m_3)=(b_0,3-L-b_0,L+2),
\]
\[
 R^3=(-j,j+Q-2,-(Q-1)),\quad j=r+a_0-1\ge1,
\]
\[
 R^4=(k,-k,-1),\quad k=S+1-b_0\ge1
\]
となる．$L=1$ は負の枝であり，$L=3$ は正の枝
$T=S+2$，$r<S<2r$ である．
\end{enumerate}
\end{proposition}
\begin{proof}
$u=-1$ なら $n_4=2$ であり，$v=-(Q-1)$ より $m_3=L+Q\le5$ である．
$Q\ge3$ と二分補題を合わせると $L=1,Q=3$ だけが残る．
node 総数と式~\eqref{g51521:low-rows} が第一の template を与える．

$u=-(Q-1)$ なら $n_4=Q\le5$ であり，奇数性から $Q=3,5$ となる．
$v=-1$ は $m_3=L+2\le5$ を与えるので $L=1,3$ である．
再び node 総数から二番目の表示を得る．
$L=3$ では $T=S+2$ を式~\eqref{g51521:positive-central} に入れると $r<S<2r$ となる．
\end{proof}

通常型の完全還元，命題~\ref{g51521:no-positive-exception}，および
命題~\ref{g51521:negative-templates} を式~\eqref{g51521:three-branches} に適用すると，
定理~\ref{g51521:reduction} が従う．

\subsection{通常型の固定分母に対する完全な有限 lifting}

\begin{proposition}\label{g51521:fixed-Q-finite}
表~\ref{g51521:ordinary-row-table} の一配置と奇数 $Q\ge3$ を固定する．
三 cores $2,3,4$ の局所方程式と，相互の三辺 $23,24,34$ の全 node 条件・対向 $D$ を満たす
整数データは有限個であり，他の分母や scale の上限を仮定せず全て復元できる．
\end{proposition}
\begin{proof}
共有高 ray への低 row 成分が正である側を先に復元する．
その低 row を $(-c,-d,m)$，低 column を $(-A,-B,U)$ と書くと
\[
 U=(Q-1)/m,\qquad A+B=U+1+D_+,
\]
である．$D_+$ は表で指定された低 $D$ であり，$A,B>0$ なので全ての低 columns は有限個である．
固定低 row・column の二次式による復元は各 column に高々二模型を与える
（\ref{fc17:section} 節）．

その模型を固定する．通常型では辺 $34$ の puncture $2$ の node は一なので，
残りの node 分配は $(n_0,n_1)=(0,4),(1,3),\ldots,(4,0)$ の五つである．
clutching により反対側の共有高 row が定まる．外部二標識を必要なら入れ替えて，
反対側の低 row と共有高 row を
\[
 (-a,b,-m),\quad b=a+m-1,\qquad (X,-W,-U),\quad W=X+Q-U
\]
と書く．matching のため $X,W>0$ が必要である．
反対側の低 column は $(-A_-,B_-,-U)$ であり，
表で指定された低 $D_-$ によって
\[
 B_-=A_-+U-1-D_-
\]
となる．残る二分母は $q_-=1+aA_-$，$q_+=1+bB_-$ である．
共有高 ray の $D$ 相殺は，$A_-$ に関する一次式であり，その係数は
\[
 Q\left(-\frac aX+\frac bW\right)
\]
である．これが非零なら候補は高々一つである．

係数が零なら $(m-1)X=a(Q-U)$ であるから，$m=1$ は不可能である．
従って $m\ge2$ とする．共有高 ray との pair equations と残る二つの row sums を使うと，
外部二 rays の相互 entries は $A_-$ の一次式に決まる．
その pair equation の残差の二次係数は
\[
 \frac{(m-1)^2U+m-2}{\{(m-1)U+1\}^2}>0
\]
である．実際，$t=1-Qa/X=1-Qb/W$ とするとこの係数は
$-t(t+m-1)$ であり，
$t=-(m-1)+1/\{(m-1)U+1\}$ を代入すれば表示式を得る．
従って $D$ の一次式が恒等的に零になる場合も，残る pair equation は genuine な二次式で，
候補は高々二つである．整数性，奇数性，正値性，元の全 row/pair equations を確認して完了する．
\end{proof}

\nolinkurl{ordinary_fixed_q_lifting_v21.py} はこの手順の独立した標準ライブラリ実装である．
ここで完成するのは三 cores とその相互の三辺であり，四高 cores $0,1$ を含む大域接着の完成ではない．
固定した $Q$ ごとの有限性は，全 $Q$ にわたる有限性を意味しない．

\paragraph{検証範囲．}
\nolinkurl{verify_g5_15_reduction_v21.py} は標準ライブラリの整数演算だけを用い，
全 scale の signed-divisor 分解と全 affine 半直線，十八の通常配置，四つの parity 排除，
および正例外型の八つの矛盾証明を JSON に保存する．
上記の全結論は分母の数値 cutoff を用いない．
別に，通常型の固定した $Q$ に対する有限 lifting 手順を研究ノートに収録する．
その手順を有限範囲の $Q$ で実行した結果を，$Q$ 全体の排除に読み替えてはならない．

% v19: a finite, exact completion reduction for the three bivalent cores
% of G5_13.  This is NOT an exclusion of G5_13.
\section{Book support の低三角形と固定 core の有限完成}
\label{bk19:section}

本節では高分母 support
\[
 G_{5,13}=\{01,02,03,04,12,13,14\}
\]
の二価 cores $2,3,4$ のみを扱う．その高 rays は共通して $0,1$ であり，
二価 cores 間の三辺は低辺である．正の局所整数 charts，奇数分母，
marked puncture の整数 clutching，非負 node，対向 $D$ の相殺を仮定する．
core $0,1$ の完成は本節の結論に含めない．

各二価 core の canonical order $(A,B,C,D)$ を用い，
\begin{equation}\label{bk19:canonical}
 P=\begin{pmatrix}
 0&-(a+x-s)&x&-s\\
 -(b-y+z)&0&-y&z\\
 r&-r-1&0&0\\
 -S&T&-L&0
 \end{pmatrix},
\end{equation}
\[
 a=sS+1,\quad b=zT+1,\qquad
 x=sS/r,\quad y=zT/(r+1),\quad L=T-S+1
\]
と書く．命題~\ref{kl10:two-high-parametrization} と
補題~\ref{ml12:dichotomy} により，正の枝では
$S>r$，$rT>(r+1)S$，$L\ge3$ であり，負の枝では
$S=T=w<r$，$L=1$ である．後者では $H=r-w$ として
\begin{equation}\label{bk19:H-bound}
 w\ge2,\qquad \frac{w^2}{2}<H(H+1)<2w^2
\end{equation}
が成り立つ（式~\eqref{cp14:H-bound}，または
式~\eqref{kb15:E-bounds}）．特に $H\ge2$ である．ここで負の枝は削除していない．
book では $01$ が高辺なので，低 wedge の負枝排除補題は適用できない．

\begin{lemma}[低三角形の node 制限]\label{bk19:triangle}
各二価 core の $L$ は $1,3,4$ のいずれかである．
各 core から，高 ray $0$ へ正の逆向き成分を持つ低 ray の先へ矢印を向けると，
その矢印は有向三角形をなす．三つの core が全て正の枝であることはない．
三つが全て負の枝である場合，canonical 高 ray $A$ が $0$ か $1$ かの選択は
三つで一定にはならない．これらの結論には対向 $D$ の相殺は不要である．
\end{lemma}

\begin{proof}
各 core の低低 entries は $0,-L$ である．$-L$ を含む低 ray の先との辺で，
残る低 core における node は $1+L$ 以上である．三つの node は非負で総数五なので
$L\le4$ であり，二分補題と合わせて $L\in\{1,3,4\}$ を得る．
二つの矢印が互いに向き合えば，その低辺の puncture $0$ における node は
$1$ から二つの正整数を引いた値となり負である．各 outdegree は一なので，
矢印は有向三角形である．その方向に cores を $i\in\mathbb Z/3$ と標識する．

$A=0$ のときを $\sigma_i=+$，$A=1$ のときを $\sigma_i=-$ とする．
低辺 $i,i+1$ の残る低 core における node は
\begin{equation}\label{bk19:low-node}
 n_{i,i+1;i-1}=1+[\sigma_i=-]L_i+[\sigma_{i+1}=+]L_{i+1}.
\end{equation}
全て正の枝なら，全て $L_i\ge3$ である．$(-,+)$ の隣接対は
式~\eqref{bk19:low-node} を七以上にするので，$\sigma$ は一定である．
全て $+$ の場合，puncture $1$ の三つの node の和は
$6+\sum_i(r_i-T_i)$ である．正の枝では $T_i\ge r_i+3$ なので，この和は負となる．
全て $-$ の場合も puncture $0$ で同じ矛盾を得る．

全て負の枝で，全て $+$ なら，puncture $0$ の node の和は
$3+\sum_i(S_i-r_i)=3-\sum_iH_i\le-3$ である．全て $-$ の場合は
高 rays の名前を交換すればよい．
\end{proof}

\begin{proposition}[有限 offset による低 node の表示]\label{bk19:offsets}
有向三角形の順序で，core $i$ の低 ray $i+1,i-1$ の行を
\[
 (p^i_{i+1,0},p^i_{i+1,1},p^i_{i+1,i-1})=(a_i,-b_i,-c_i),
\]
\[
 (p^i_{i-1,0},p^i_{i-1,1},p^i_{i-1,i+1})=(-d_i,e_i,-f_i)
\]
と書くと，$a_i,b_i,d_i,e_i$ は正整数であり，
\[
 b_i=a_i+1-c_i,\quad d_i=e_i+1-f_i,\quad
 c_if_i=0,\quad c_i+f_i\in\{1,3,4\}.
\]
低辺 $i,i+1$ の node 非負性は，次の有限 offset 条件と同値である：
\begin{equation}\label{bk19:finite-offset}
 c_i+f_{i+1}\le4,\qquad
 e_{i+1}=a_i+t_{i+1},\qquad
 f_{i+1}-2\le t_{i+1}\le2-c_i.
\end{equation}
\end{proposition}

\begin{proof}
行和と式~\eqref{bk19:canonical} から最初の表示を得る．三つの node は
\[
 (n_{i,i+1;0},n_{i,i+1;1},n_{i,i+1;i-1})
 =(1-a_i+d_{i+1},\ 1+b_i-e_{i+1},\ 1+c_i+f_{i+1}).
\]
その和は五である．初めの二つの非負性を代入すると
$c_i-2\le a_i-e_{i+1}\le2-f_{i+1}$ となる．この区間が空でない条件が
$c_i+f_{i+1}\le4$ であり，式~\eqref{bk19:finite-offset} を得る．
\end{proof}

\begin{theorem}[固定二価 core に対する低三角形の有限完全完成]
\label{bk19:finite-completion}
二価 core $i$ の分母と局所整数行列を一つ固定する．すると，
低辺の node 非負性と対向 $D$ の相殺を満たす残り二価 cores の組は有限個である．
その全てを整数演算と平方判定だけで完全列挙できる．
この列挙は分母・分子の数値 cutoff を導入しない．
固定分母 $7,13$ に属する共鳴族も含む．
ただし，これは固定した一つの core に対する有限性であり，
その core 自体を動かす無限集合の有限性も，book support の不存在も主張しない．
\end{theorem}

\begin{proof}
まず固定 core $i$ の低 ray $j$ の相手 core $j$ を考える．
残る三つの標識 $k$ について，非負整数の node 分配 $n_k$ は
$\sum_kn_k=5$ を満たすので二十一通りである．各分配は
\begin{equation}\label{bk19:fixed-row}
 p^j_{ik}=1-n_k-p^i_{jk}
\end{equation}
によって，相手の低 ray $i$ の行全体を固定する．
対向 $D$ は同時に $D^j_i=-D^i_j$ を固定する．

この行が canonical $C$ なら，その高 ray への entries $r,-r-1$ によって
$r$ と高 rays の順序が決まる．負の枝では $2\le w<r$ であり有限である．
正の枝では $\delta=T-S=L-1\in\{2,3\}$ で，
$rT>(r+1)S$ は $S<r\delta$ と同値なので
\[
 r<S<r\delta
\]
により $S$ と $T$ は有限個である．

固定行が canonical $D$ なら，その entries $-S,T,-L$ によって
$S,T$ と高 rays の順序が決まる．正の枝では
$S/\delta<r<S$，負の枝では $r=w+H$ かつ
式~\eqref{bk19:H-bound} である．従ってこの場合も $r$ は有限個である．
canonical $C,D$ のどちらにも合わない行は棄却する．

こうして固定された $r,S,T$ に対し，$h=S-r$，$\delta=T-S$ と置く．
高高 pair equation は
\begin{equation}\label{bk19:bilinear}
 Fsz+As+Bz+C=0,
\end{equation}
\[
 F=(2r+1)h-r^2(\delta-1),\quad
 A=h(r+1),\quad B=-r(h+\delta-1),\quad C=r(r+1)
\]
と同値である．これは式~\eqref{bk19:canonical} の高高 entries の積を
$(sS+1)(zT+1)-1$ と等置し，$r(r+1)$ 倍して展開した式である．
ここで $\gcd(2r+1,r^2)=1$ である．もし $F=0$ なら
$2r+1$ は $\delta-1\in\{-1,1,2\}$ を割るが，$2r+1\ge3$ なので不可能である．
従って $F\ne0$ である．

固定された低 ray の $D$ を $d$ とする．
\[
 D_C=\delta-\frac Sr s+\frac T{r+1}z,\qquad
 D_D=s-z-1
\]
より，その $D$ の条件は
\begin{equation}\label{bk19:linear}
 \ell s+mz+n=0
\end{equation}
という一次式であり，係数は
\[
 (\ell,m,n)=
 \begin{cases}
 (-S(r+1),\ rT,\ (\delta-d)r(r+1)),&\text{固定 ray が }C,\\
 (1,-1,-d-1),&\text{固定 ray が }D.
 \end{cases}
\]
両場合とも $\ell m\ne0$ である．式~\eqref{bk19:linear} から $z$ を消去すると
\begin{equation}\label{bk19:quadratic}
 -F\ell s^2+(-Fn+Am-B\ell)s+(Cm-Bn)=0.
\end{equation}
最高次係数 $-F\ell$ は非零なので，この式には高々二つの整数解しかない．
各 $s$ から $z$ を一意に復元し，正整数性，$x,y$ の整数性，奇数分母，
元の局所方程式を検査する．これは退化した双一次式にも適用できる：
$(Fs+B)(Fz+A)=AB-FC$ の右辺が零でも，
式~\eqref{bk19:quadratic} の最高次係数は非零のままである．
従って共鳴直線を捨てたり，その自由変数を cutoff したりする必要はない．

もう一つの低辺の相手についても同じ有限列挙を行う．
二つの有限集合の直積に対し，残る低辺の三つの node と対向 $D$ を検査すれば，
必要条件を満たす低三角形を過不足なく得る．
\end{proof}

この証明の算法は
\texttt{book\_low\_triangle\_completion\_v19.py} に標準ライブラリのみで実装した．
この定理は高中心 cores $0,1$ の局所方程式，その高辺の clutching，または
幾何学的な atlas の存在を判定したものではない．
固定 core 自体が動く問題は本節の有限性だけでは解決しないが，後の定理~\ref{bc22:book-obstruction} が追加の整数剛性により book 全体を排除する．

% v20: complete all-negative branch obstruction for the book low triangle.
% Depends on bk19:canonical, bk19:triangle, bk19:finite-offset, bk19:H-bound.
\section{Book の全負枝低三角形の排除}
\label{bk20:negative-section}

前節の固定 core 有限完成を用いるだけでは，固定 core 自身が動く無限集合を
排除したことにはならない．本節では，全ての二価 cores が負の枝に属する場合に，
低 node の比較から reverse parameters の一様上界を導く．
scale と高分母には探索上限を置かず，非零の因子積から全て復元する．

\begin{theorem}[全負枝の book 低三角形は存在しない]
\label{bk20:no-all-negative}
Book support $G_{5,13}$ の三つの二価 cores が，全て負の枝
$S_i=T_i=w_i<r_i$ に属すると仮定する．
正の整数局所 charts，奇数高分母，三本の低辺の非負 node，
および対向低 $D$ の相殺を同時に満たす低三角形は存在しない．
高中心 cores $0,1$ の完成に関する条件は不要である．
\end{theorem}

\begin{proof}
有向低三角形の順に二価 cores を $i=1,2,3$ と呼び，添字を三を法として読む．
これは高中心 $0,1$ の実際の標識とは別の巡回添字である．
$H_i=r_i-w_i$ と置くと，式~\eqref{bk19:H-bound} により
\begin{equation}\label{bk20:negative-range}
 w_i\ge2,\quad H_i\ge2,\qquad
 \frac{w_i^2}{2}<H_i(H_i+1)<2w_i^2
\end{equation}
である．補題~\ref{bk19:triangle} により canonical 高 ray の向き
$\sigma_i$ は一定ではない．高中心の交換は有向低三角形の向きを反転し，
全ての $\sigma_i$ を交換する．従って，必要なら高中心を交換し巡回添字を選び直すと
\[
 (\sigma_1,\sigma_2,\sigma_3)=(+,+,-)
\]
としてよい．

全て $L_i=1$ なので，命題~\ref{bk19:offsets} の有限 offset は
\begin{equation}\label{bk20:negative-offsets}
 w_2=r_1+t_2,\qquad r_3=r_2+t_3,\qquad w_1=w_3+t_1,
\end{equation}
\[
 -1\le t_2\le2,\qquad -2\le t_3\le2,\qquad -1\le t_1\le1
\]
となる．ここで node 条件は三つの高分母そのものを使わず，reverse entries のみを使う．

\paragraph{第一段階：証明による一様上界．}
式~\eqref{bk20:negative-range} より
\[
 H+\frac12>\frac{w}{\sqrt2}>\frac7{10}w,
 \qquad H<\sqrt2w<\frac{10}{7}w
\]
なので
\begin{equation}\label{bk20:r-bounds}
 \frac{17}{10}w-\frac12<r<\frac{17}{7}w
\end{equation}
を得る．式~\eqref{bk20:negative-offsets} と合わせると
\[
 w_2>\frac{17}{10}w_1-\frac32,
\]
\[
 r_3>\frac{17}{10}w_2-\frac52
      >\frac{289}{100}w_1-\frac{101}{20},
 \qquad
 r_3<\frac{17}{7}w_3\le\frac{17}{7}(w_1+1).
\]
従って
\begin{equation}\label{bk20:proved-bound}
 \frac{323}{700}w_1<\frac{1047}{140},\qquad
 w_1<\frac{5235}{323}<17.
\end{equation}
整数性から $w_1\le16$，$w_3\le17$ であり，
$w_2=r_1+t_2<\frac{17}{7}\cdot16+2<41$ から $w_2\le40$ である．
これらは全て仮定から導いた上界である．

\paragraph{第二段階：固定 reverse data の全 scale．}
負枝の canonical 表示で
\[
 c=\frac{sH}{r},\qquad e=\frac{z(H+1)}{r+1},\qquad
 E=1-\frac{w^2}{H(H+1)}=\frac MN
\]
と置き，$M,N$ を互いに素，$N>0$ とする．
式~\eqref{ml12:E-equation} と式~\eqref{kb15:E-bounds} は
\[
 Ece=e-c+1,\qquad c\ge2,\ e\ge1,\quad -1<E<\frac12,\quad E\ne0
\]
を与える．従って
\begin{equation}\label{bk20:complete-factorization}
 (N-Mc)(N+Me)=N(N-M)>0.
\end{equation}
両因子は正である．実際，$Ec=1-(c-1)/e<1$ と
$c(1+Ee)=e+1>0$ を用いればよい．
よって $d$ を $N(N-M)$ の全ての正の約数として
\[
 c=\frac{N-d}{M},\qquad
 e=\frac{N(N-M)/d-N}{M},\qquad
 s=\frac{rc}{H},\qquad z=\frac{(r+1)e}{H+1}
\]
を復元し，正整数性と奇数分母 $sw+1,zw+1$ を確認すれば，全ての scale が得られる．
式~\eqref{bk20:complete-factorization} の右辺は非零なので，共鳴直線の見落としはない．
$s,z$ または高分母に上限を加えていない．

\paragraph{第三段階：残る有限データの厳密検査．}
$2\le w_1\le16$，式~\eqref{bk20:negative-range}，および
式~\eqref{bk20:negative-offsets} を満たす reverse data は九十一型である．
第二段階の全 scale 復元が三つの cores 全てで空でないものは二十五型であり，
scale を展開した三模型の組は六十四個となる．
その全てで，最初の低辺の $D$ 相殺が既に破れる．

次の表は，第一 core の canonical $(r,s,z,S,T)$ ごとに，
可能な第二 core との残差 $D_C^{(1)}+D_D^{(2)}$ をまとめたものである．
異なる第一・第二 core の組は二十六組であり，表のいずれの集合にも零はない．
\begin{center}
\small
\begin{tabular}{ll}
第一 core の canonical parameters & 第一低辺の $D$ 残差の全値\\
\hline
$(4,4,5,2,2)$ & $-20,-9,-4,4,7,16$\\
$(5,10,2,3,3)$ & $-40,-25,-9,-1,11,24$\\
$(6,6,14,3,3)$ & $-32,32$\\
$(7,7,2,4,4)$ & $-57,-20,-9,3,12,43$\\
$(7,21,4,4,4)$ & $-64,-27,-16,-4,5,36$\\
\end{tabular}
\end{center}

この有限検査は \nolinkurl{verify_book_all_negative_v20.py} に実装し，
九十一型の reverse data と全 scale，六十四個の node と $D$ 証明書を
\nolinkurl{book_all_negative_verification_v20.json} に保存した．
行和，pair products，明示的な正 weights，三つの node の非負性と総数五は
元の行列から再計算している．
数学的に導いた有限集合上の整数演算であり，有限範囲の探索結果から
無上界の結論へ外挿する段階は含まれない．
以上で全負枝を排除した．
\end{proof}

\begin{corollary}[Book には正負両枝が必要]
\label{bk20:mixed-necessary}
仮定を満たす book 低三角形が存在するなら，三つの二価 cores の中に
正の枝と負の枝がともに現れる．従って負の枝は一個または二個である．
\end{corollary}

\begin{proof}
全正枝は補題~\ref{bk19:triangle}，全負枝は
定理~\ref{bk20:no-all-negative} により排除される．
\end{proof}

この系は全負枝を除く中間結果である．混合枝は後の定理~\ref{bc22:book-obstruction} までの議論によって全て排除する．
Book を排除した後も，残る密な支持の一般分母剛性は別の問題として残る．

% v20: no zero low-D edge in a book low triangle; hence a nonzero D cycle.
% Depends on the canonical two-high form and the v19 finite-completion theorem.
\section{Book 低三角形の非零 \texorpdfstring{$D$}{D} 周期}
\label{bd20:section}

低辺の対向 $D$ 相殺は，node の向きとは別の符号の向きを定める．
零の $D$ を捨てて議論する前に，その場合を完全に排除する．

\begin{lemma}[低係数の零と負の隣接]
\label{bd20:local-zero}
正の整数二高分母標準形において，二つの低係数 $D_C,D_D$ が
同時に零となることはない．一方が負でもう一方が零となるのは，
canonical parameters
\begin{equation}\label{bd20:unique-zero}
 (r,s,z,S,T)=(4,4,5,2,2),\qquad (D_C,D_D)=(0,-2)
\end{equation}
の場合に限られる．この局所主張には分母の奇数性は不要である．
\end{lemma}

\begin{proof}
$\delta=T-S$ と置くと
\[
 D_C=\delta-\frac Sr s+\frac T{r+1}z,\qquad
 D_D=s-z-1.
\]
正の枝では $\delta\ge2$，$r\delta>S$ である．
$D_D<0$ なら整数性から $z\ge s$ であり，
$T/(r+1)>S/r$ によって $D_C>0$ となる．
$D_D=0$ のときは $s=z+1$ を直接代入して
\begin{equation}\label{bd20:positive-zero-identity}
 D_C=\frac{(r+z+1)(r\delta-S)}{r(r+1)}>0
\end{equation}
を得る．従って正の枝には負と零の組も，二つの零もない．

負の枝では $S=T=w$，$H=r-w$ とし，
\[
 c=\frac{sH}{r},\quad e=\frac{z(H+1)}{r+1},\qquad
 E=1-\frac{w^2}{H(H+1)}
\]
を用いる．既知の恒等式 $Ece=e-c+1$ と
$c\ge2,e\ge1,-1<E<1/2,E\ne0$ が成り立つ．
$E<0$ なら整数性から $c\ge e+2$ であり，
\[
 D_C=w\left(\frac e{H+1}-\frac cH\right)<0,
\]
\[
 s-z=(c-e)+w\left(\frac cH-\frac e{H+1}\right)>2.
\]
従って $D_D>1$ であり，零は生じない．

$E>0$ なら，定理~\ref{ml12:denominator-descent} の証明と同じく，
$j=e-c\ge1$，$c\le H$ となる．これらを再確認すると，
$e\le c-1$ は $Ece=e-c+1>0$ に反し，
$e=c$ は pair equation の $\gcd(c,e)=1$ と $c\ge2$ に反する．
また $E>0$ は整数性から $w\le H$ を与えるので
$E\ge1/(H+1)$，一方 $Ec=1-(c-1)/e<1$ だから $c\le H$ である．
従って
\[
 s-z=-j+\frac{w(c-Hj)}{H(H+1)}<0,\qquad D_D<0,
\]
\begin{equation}\label{bd20:negative-low-D}
 D_C=\frac{w(Hj-c)}{H(H+1)}\ge0.
\end{equation}
等号には $j=1,c=H$ が必要十分である．そのとき
$E=2/[H(H+1)]$ なので
\[
 w^2=H(H+1)-2.
\]
$H\ge3$ なら右辺は $H^2$ と $(H+1)^2$ の間に厳密にあるため，
$H=w=2$ に限られる．$c=2,e=3$ から $r=4,s=4,z=5$ を得る．
以上で零の全場合が尽くされた．
\end{proof}

\begin{lemma}[唯一の負・零 core は低三角形へ完成しない]
\label{bd20:zero-core-obstruction}
Book の仮定の下で，式~\eqref{bd20:unique-zero} の core を含む
低三角形は存在しない．
\end{lemma}

\begin{proof}
その core を $2$ とし，canonical order を $(0,1,3,4)$ とする．
行列と低係数は
\[
 P^2=
 \begin{pmatrix}
 0&-7&2&-4\\
 -14&0&-2&5\\
 4&-5&0&0\\
 -2&2&-1&0
 \end{pmatrix},\qquad
 (q_{20},q_{21},q_{23},q_{24})=(9,11,1,1),
\]
\[
 D^2_3=0,\qquad D^2_4=-2.
\]
低辺 $23$ の三つの node を $(n_0,n_1,n_4)$ とすると，
core $3$ の ray $2$ の行は
\[
 (-3-n_0,\ 6-n_1,\ 1-n_4),\qquad n_0+n_1+n_4=5
\]
に固定され，その $D$ は零である．
$n_4=1$ ならこの行は canonical $C$ で，高 ray の順序は $(1,0)$，
$2\le r\le6$ となる．負枝は $2\le w<r$ と $H$ の既知の範囲を，
正枝は $\delta=2,3$ と $r<S<r\delta$ を満たす．
$n_4=L+1$，$L=1,3,4$ なら canonical $D$ で，高順序は $(0,1)$ である．
$L=1$ では $S=T\in\{3,4,5,6\}$，
$L=3,4$ では可能な $(r,S,T)$ は
\[
 (2,3,5),\quad(3,4,6),\quad(2,3,6)
\]
に限られる．それ以外の行は二高標準形に合わない．

同様に低辺 $24$ の node を $(n_0,n_1,n_3)$ とすると，
core $4$ の ray $2$ の行は
\[
 (3-n_0,\ -1-n_1,\ 2-n_3),\qquad n_0+n_1+n_3=5
\]
で，その $D$ は二である．
canonical $C$ なら $n_3=2$，高順序 $(0,1)$，
$1\le r\le3$ である．canonical $D$ なら，可能なのは
高順序 $(1,0)$ の負枝 $w=2,3$ だけである．

以上の有限 reverse data に対し，指定した低 $D$ と高高 pair equation を連立する．
完全復元は定理~\ref{bk19:finite-completion} の二次式で実行できるが，
ここでは次の直接表示でも確認できる．
\[
 u=S+\frac Sr-1,\qquad v=T-\frac T{r+1}+1
\]
と置くと，高高 pair equation は
\[
 (uv-ST)sz+(u-S)s+(v-T)z+1=0.
\]
指定された低係数を $d$ とすると
\[
 z=
 \begin{cases}
 \displaystyle\frac{(r+1)S}{rT}s+\frac{(r+1)(d-\delta)}T,
       &\text{固定 ray が }C,\\[4pt]
 s-d-1,&\text{固定 ray が }D.
 \end{cases}
\]
この代入による二次式の整数根を取り，元の行列の整数性と奇数分母を検査する．
core $3$ の reverse data は六十八個，core $4$ では十五個であり，
全 scale の復元後に残る相手は次表の四個と一個だけである．
scale や高分母の上限は使っていない．

\begin{center}
\small
\begin{tabular}{cll}
相手 core & canonical $(r,s,z,S,T)$ & canonical order\\
\hline
$3$ & $(2,2,1,3,6)$ & $(0,1,4,2)$\\
$3$ & $(3,3,2,4,6)$ & $(0,1,4,2)$\\
$3$ & $(4,4,5,2,2)$ & $(1,0,2,4)$\\
$3$ & $(4,4,2,7,10)$ & $(1,0,2,4)$\\
\hline
$4$ & $(2,2,1,3,6)$ & $(0,1,2,3)$\\
\end{tabular}
\end{center}

残る低辺 $34$ を比べると，core $3$ の上表の順で次の値となる．
node の順序は punctures $(0,1,2)$ である．
\begin{center}
\small
\begin{tabular}{lcc}
core $3$ の canonical parameters & 最後の node 三つ & $D^3_4+D^4_3$\\
\hline
$(2,2,1,3,6)$ & $(2,-2,5)$ & $2$\\
$(3,3,2,4,6)$ & $(1,-1,5)$ & $1$\\
$(4,4,5,2,2)$ & $(2,-3,6)$ & $-2$\\
$(4,4,2,7,10)$ & $(-6,2,9)$ & $1$\\
\end{tabular}
\end{center}
いずれも node 非負性と $D$ 相殺に反する．
全二十一 node 行からの復元とこの表は
\nolinkurl{verify_book_low_D_cycle_v20.py} で独立に再計算し，
\nolinkurl{book_low_D_cycle_verification_v20.json} に保存した．
\end{proof}

\begin{theorem}[Book 低三角形の非零 $D$ 周期]
\label{bd20:nonzero-cycle}
正の整数局所 charts，奇数高分母，低辺の非負 node，
対向 $D$ 相殺を満たす book 低三角形が存在すると仮定する．
三本の低辺の全てで両端の $D$ は非零である．
各二価 core は一つの正の低 $D$ と一つの負の低 $D$ を持ち，
各低辺を負の端から正の端へ向けると有向三角形をなす．
\end{theorem}

\begin{proof}
二本以上の低辺が零の $D$ を持てば，その交点 core では二つの低 $D$ が
零になり，補題~\ref{bd20:local-zero} に反する．
零の辺がちょうど一本とする．その両端 core の残る低 $D$ は非零である．
両方が正なら第三 core の二つの低 $D$ が負になり，
補題~\ref{tl15:low-D} に反する．
従って少なくとも一方の端 core は負と零の組を持つ．
補題~\ref{bd20:local-zero} により式~\eqref{bd20:unique-zero} の模型に限られるが，
補題~\ref{bd20:zero-core-obstruction} がその完成を排除する．

これで全低 $D$ は非零である．三本の辺の相殺により，
六つの低端のうち正が三つ，負が三つとなる．
補題~\ref{tl15:low-D} により一つの core に負が二つ現れることはないので，
三つの cores はそれぞれ一つずつ負を持ち，残りは正である．
従って各頂点の入次数と出次数はともに一であり，有向三角形になる．
\end{proof}

この $D$ の向きは，高 ray $0$ への正 entry から定めた
補題~\ref{bk19:triangle} の向きとは別であり，
一致または反対のいずれも現時点では許している．
符号の周期性だけから book 全体の不存在は結論しない．

% v21: exact sector control for positive two-high branches with L=3 or 4.
\section{小さい低辺係数を持つ正枝の符号と分母順序}
\label{ps21:section}

Book の低三角形では，正枝の非零低 entry は $-3$ または $-4$ である．
この二つの値に限ると，pair equation の係数の符号は高分母の順序と結び付き，
低 $D$ が負となる側には reverse parameter の強い上界がある．
以下の局所補題には分母の奇数性，node 条件，clutching は不要である．

\begin{lemma}[正枝の二つの sector]\label{ps21:positive-sector}
正の整数二高分母標準形において
\[
 \delta=T-S\in\{2,3\},\quad d=\delta-1,\quad
 h=S-r,\quad k=T-r-1=h+d,
\]
\[
 \alpha=\frac{sh}{r},\quad\beta=\frac{zk}{r+1},\quad
 F=(2r+1)h-r^2d
\]
と置く．このとき $F\ne0$ であり，次が成り立つ．
\begin{enumerate}
\item $F>0$ なら $s\le z$，$D_D<0<D_C$，$a<b$ である．
\item $F<0$ なら $D_D\ge0$ であり，
\begin{equation}\label{ps21:h-bound}
 h\le
 \begin{cases}
 \lfloor(r-1)/2\rfloor,&\delta=2,\\
 r-1,&\delta=3.
 \end{cases}
\end{equation}
また $a>b$ であるか，次の二つの例外のいずれかである：
\begin{equation}\label{ps21:equal-base-points}
 (r,s,z,S,T)=(2,2,1,3,6),\qquad (3,3,2,4,6).
\end{equation}
この二点ではそれぞれ $a=b=7,13$，$D_D=0$，$D_C=2,1$ である．
\end{enumerate}
特に $D_C\le0$ なら $F<0$，$a>b$ であり，
\begin{equation}\label{ps21:T-bound}
 T\le
 \begin{cases}
 \lfloor(3r+3)/2\rfloor,&\delta=2,\\
 2r+2,&\delta=3.
 \end{cases}
\end{equation}
\end{lemma}

\begin{proof}
正枝の既知の不等式 $rT>(r+1)S$ により $0<h<rd$ である．
補題~\ref{ml12:dichotomy} と式~\eqref{td12:bezout} より，
$\alpha,\beta$ は正の整数であり，
\begin{equation}\label{ps21:two-identities}
 F\alpha\beta=(\beta-\alpha-1)hk,\qquad
 M\beta-N\alpha=1,
 \quad M=rs+1,
 \quad N=(r+1)z+1.
\end{equation}
特に $\gcd(\alpha,\beta)=1$ である．$F=0$ なら
$2r+1\mid r^2d$ となるが，$\gcd(2r+1,r^2)=1$ と $d\in\{1,2\}$ に反する．

まず $F>0$ とする．式~\eqref{ps21:two-identities} により
$j=\beta-\alpha\ge2$ である．$s\ge z+1$ と仮定すると，
第二式から
\[
 M\{r\beta-(r+1)\alpha\}
 =r+\{r(r+1)(z-s)-1\}\alpha<0
\]
を得る．従って $\alpha>rj$，$\beta>(r+1)j$ となる．
一方，$h<rd$，$k<d(r+1)$，$F\ge1$ なので
\[
 j-1=\frac{F\alpha\beta}{hk}
 >\frac{j^2}{d^2}\ge\frac{j^2}{4}\ge j-1,
\]
ここで最後の不等式は $(j-2)^2\ge0$ である．これは矛盾であり，$s\le z$ である．
従って $D_D=s-z-1<0$，また
$T/(r+1)>S/r$ から $D_C=\delta-sS/r+zT/(r+1)>0$ となる．
さらに
\[
 a-b=rs-(r+1)z+r\alpha-(r+1)\beta<0
\]
である．

次に $F<0$ とする．整数 $h$ に対する
$h<r^2d/(2r+1)$ は式~\eqref{ps21:h-bound} を与える．
$j=\alpha-\beta\ge0$ と置くと，
\[
 E=-\frac{F}{hk}>0,\qquad E\alpha\beta=j+1.
\]
また
\[
 -F-k=(r+1)\{d(r-1)-2h\}\ge0
\]
なので $E\ge1/h$ であり，
\[
 E\beta=\frac{j+1}{\beta+j}\le1
 \quad\Longrightarrow\quad \beta\le h<r.
\]
$rk-(r+1)h=rd-h>0$ と $\alpha\ge\beta$ により
\[
 \frac sz=\frac\alpha\beta\frac{rk}{(r+1)h}>1,
\]
従って $D_D\ge0$ である．$j\ge1$ なら
$M\beta-N\alpha=1$ より $M>N$，すなわち
$rs-(r+1)z\ge1$ である．ゆえに
\[
 a-b=rs-(r+1)z+rj-\beta>0.
\]

残る $j=0$ では $\alpha=\beta=1$ である．
$rs-(r+1)z=1$ なので，$m=s-z\ge1$ と置くと
$z=rm-1$，$s=(r+1)m-1$ である．
$\alpha=sh/r=1$ から $s\le r$ となり，$m=1$，$s=r$，$z=r-1$ を得る．
従って $h=1$，$k=(r+1)/(r-1)$ であり，
$r-1\mid2$ から式~\eqref{ps21:equal-base-points} の二点に限られる．
直接代入によって表示した分母と低 $D$ を得る．
最後に $D_C\le0$ では $F>0$ とこの二点が排除され，
式~\eqref{ps21:h-bound} に $T=r+h+\delta$ を代入すると
式~\eqref{ps21:T-bound} を得る．
\end{proof}

\begin{corollary}[Book の $D$ 周期に沿う高分母の順序]
\label{ps21:book-denominator-order}
定理~\ref{bd20:nonzero-cycle} の book 低三角形を考える．
二つの高頂点を $0,1$ とし，低 row の正 entry が高 ray $0$ を向く側を
next とする有向 node 三角形を用いる．この向きに沿う各低 $D$ が正なら
\[
 q_{i0}<q_{i1}\qquad(i=2,3,4),
\]
全て負なら三つの不等号は全て逆になる．
\end{corollary}

\begin{proof}
正枝では，非零低 $D$ が一正一負であることと補題~\ref{ps21:positive-sector} から
\[
 D_C<0\Longleftrightarrow F<0\Longleftrightarrow a>b,
 \qquad
 D_C>0\Longleftrightarrow F>0\Longleftrightarrow a<b
\]
となる．負枝でも式~\eqref{bd20:negative-low-D} と
定理~\ref{ml12:denominator-descent} の証明により，同じ $D_C$ と $a-b$ の対応がある．
canonical $A=0$ なら next は $C$，$a=q_{i0}$，$b=q_{i1}$ である．
canonical $A=1$ なら next は $D$，$a=q_{i1}$，$b=q_{i0}$ であるため，
$D$ の符号と分母順序が同時に反転する．
対向 $D$ の相殺と各 core の一正一負性から，node 三角形上の next の $D$ は
三つ全て同符号である．これで両方の場合が従う．
\end{proof}

この系は中心二 core に接続する三つずつの分母を比較する必要条件である．
この順序条件だけでは book 全体の不存在は従わない．後の定理~\ref{bc22:book-obstruction} は正負枝の追加の整数剛性を用いて排除を完成する．

% v21: the two-negative-branch book obstruction, with proved finite bounds.
% Requires ps21:positive-sector, bd20:nonzero-cycle, bk19:finite-completion.
\section{Book の二負枝の排除と一負枝への還元}
\label{bn21:section}

全正枝と全負枝の排除だけでは，book の低三角形には
負枝を一つまたは二つ持つ可能性が残る．
本節では正枝の sector 補題を node 周期と結び付け，
負枝を二つ持つ全ての場合に解析的な有限上界を与える．
有限上界は低三角形が存在するという仮定から導くものであり，
探索範囲の仮定ではない．

\subsection{周期に沿う reverse parameter と sector の不等式}

低頂点を node の有向三角形に沿って巡回添字 $i$ で表す．
高 ray を $0,1$ とし，各 core で次の低 ray の row を
$(u_i,-(u_i+1-\gamma_i),-\gamma_i)$，
前の低 ray の row を $(-(v_i+1-\phi_i),v_i,-\phi_i)$ と書く．
最初の二成分は高 ray $0,1$，最後の成分は他方の低 ray に対応する．
canonical $A=0$ を $\sigma_i=+$，$A=1$ を $\sigma_i=-$ と記すと
\begin{equation}\label{bn21:node-orientation}
 (u_i,v_i,\gamma_i,\phi_i)=
 \begin{cases}
 (r_i,T_i,0,L_i),&\sigma_i=+,\\
 (T_i,r_i,L_i,0),&\sigma_i=-.
 \end{cases}
\end{equation}
辺 $i\to i+1$ の node 非負性は
\begin{equation}\label{bn21:offset}
 v_{i+1}=u_i+t_{i+1},\qquad
 \phi_{i+1}-2\le t_{i+1}\le2-\gamma_i
\end{equation}
を与える．特に $\gamma_i+\phi_{i+1}\le4$ である．
定理~\ref{bd20:nonzero-cycle} により，next 側の低係数は三つ全て
同じ非零符号 $\tau\in\{+,-\}$ を持つ．

\begin{lemma}[二つの枝の定量的不等式]\label{bn21:bounds}
負枝を $S=T=w$，$H=r-w$ と書くと
\begin{equation}\label{bn21:negative-linear}
 \frac{17}{10}w-\frac12<r<\frac{17}{7}w,
 \qquad w,H\ge2.
\end{equation}
また，非零低 $D$ を一正一負で持つ負枝では
\begin{equation}\label{bn21:negative-sector}
 D_C>0\Longrightarrow r\ge2w,\qquad
 D_C<0\Longrightarrow r\le2w-1.
\end{equation}

正枝では $\delta=T-S\in\{2,3\}$，$h=S-r$，$d=\delta-1$，
$k=h+d$ と置く．その低 $D$ が一正一負であるとき
\begin{equation}\label{bn21:positive-lower}
 D_C>0\Longrightarrow
 T\ge
 \begin{cases}
 (3r+4)/2,&\delta=2,\\
 2r+3,&\delta=3.
 \end{cases}
\end{equation}
また $D_C<0$ なら
\begin{equation}\label{bn21:small-h}
 d r^2\le2hr+h(h+\delta).
\end{equation}
従って $D_C<0$ かつ $T-r\le5$ なら
$\delta=2$ では $r\le7$，$\delta=3$ では $r\le3$ である．
\end{lemma}

\begin{proof}
負枝の既知の範囲 $w^2/2<H(H+1)<2w^2$ から
$H+1/2>w/\sqrt2>7w/10$，$H<\sqrt2w<10w/7$ を得る．
これが式~\eqref{bn21:negative-linear} である．
$E=1-w^2/[H(H+1)]$ とすると，負枝の低 $D$ の符号計算より
$D_C>0$ は $E>0$，$D_C<0$ は $E<0$ に対応する．
$E>0$ なら整数性から $w\le H$ であり，$E<0$ なら $w\ge H+1$ である．
これで式~\eqref{bn21:negative-sector} が従う．

正枝の $F=(2r+1)h-r^2d$ は，補題~\ref{ps21:positive-sector} と
一正一負性により，$D_C$ と同じ符号である．
$F>0$，$d=1$ なら
$h>r^2/(2r+1)>(r-1)/2$ なので $h\ge\lceil r/2\rceil$ である．
$d=2$ なら $h>r-1/2$ なので $h\ge r$ である．
$T=r+h+\delta$ に代入して式~\eqref{bn21:positive-lower} を得る．

$F<0$ の場合，式~\eqref{ps21:two-identities} の正の整数
$\alpha,\beta$ に対して $\alpha\ge\beta\ge1$ であり，
\[
 \frac{-F}{hk}\alpha\beta=\alpha-\beta+1\le\alpha.
\]
従って $-F\le hk$，すなわち式~\eqref{bn21:small-h} が成り立つ．
$T-r\le5$ は $h\le3$（$\delta=2$）または $h\le2$（$\delta=3$）を意味する．
前者では $r^2\le6r+15$ より $r\le7$，後者では
$2r^2\le4r+10$ より $r\le3$ である．
\end{proof}

\begin{proposition}[二負枝に対する解析的な有限上界]
\label{bn21:analytic-finite}
Book の低三角形が負枝をちょうど二つ持つと仮定する．
その二つのうち少なくとも一つは $2\le w\le50$ を満たす．
この結論には scale または高分母の上界を仮定しない．
\end{proposition}

\begin{proof}
唯一の正枝を $P$ とする．高 ray の交換は node 周期を反転し
全ての $\sigma$ を交換するので，$\sigma_P=+$ としてよい．
負枝二つの $\sigma$ による四場合を尽くす．

\paragraph{負枝二つがともに $\sigma=-$ の場合．}
式~\eqref{bn21:node-orientation} では，正枝で $v-u=T-r\ge L\ge3$，
負枝で $v-u=H\ge2$ となる．三本の低辺の高 ray $1$ 側の node 総和は
\[
 \sum n_1=6+\sum(u_i-v_i)-\sum\gamma_i
 \le6-(3+2+2)-(0+1+1)=-3,
\]
となり，node 非負性に反する．

\paragraph{三枝とも $\sigma=+$ の場合．}
node 周期を $N_1\to N_2\to P\to N_1$ と書き，
正枝の parameters を $r,S,T,L$ とする．
式~\eqref{bn21:offset} は
\begin{equation}\label{bn21:same-offsets}
 w_2=r_1+t_2,\quad T=r_2+t_3,\quad w_1=r+t_1,
 \qquad
 -1\le t_1,t_2\le2,\quad L-2\le t_3\le2
\end{equation}
を与える．$\tau=+$ なら二つの負枝で $r_i\ge2w_i$ である．
$L=3$ では
\[
 4w_1-1\le T\le2r+1\le2w_1+3,
\]
より $w_1\le2$ であり，$L=4$ では
\[
 4w_1\le T\le3r+2\le3w_1+5
\]
より $w_1\le5$ である．上界は正枝の $r<S<r\delta$ から従う．

$\tau=-$ なら $D_C(P)<0$ である．式~\eqref{bn21:negative-linear} を
式~\eqref{bn21:same-offsets} に二回適用すると
\[
 r_2>\frac{289}{100}w_1-\frac{61}{20}.
\]
$L=3$ では
\[
 \frac{289}{100}w_1-\frac{41}{20}<T\le2w_1+3,
\]
$L=4$ では補題~\ref{ps21:positive-sector} の $T\le2r+2$ によって
\[
 \frac{289}{100}w_1-\frac{21}{20}<T\le2w_1+4.
\]
いずれも $w_1<505/89<6$ であり，$w_1\le5$ となる．

\paragraph{周期が $P^+\to N_1^-\to N_2^+\to P^+$ の場合．}
node 条件は
\begin{equation}\label{bn21:A-offsets}
 r_1=r+t_1,\quad w_2=w_1+t_2,\quad T=r_2+t_0,
 \qquad
 -2\le t_1\le2,\quad-1\le t_2\le1,\quad L-2\le t_0\le2
\end{equation}
である．$\tau=-$ なら $N_1$ の next は $D$ なので
$D_C(N_1)>0$，一方 $D_C(N_2)<0$ である．従って
\[
 r_1\ge2w_1,\qquad r_2\le2w_2-1,
\]
から
\[
 T\le r_2+2\le2w_2+1\le2w_1+3\le r_1+3\le r+5
\]
を得る．補題~\ref{bn21:bounds} により $r\le7$ であり，
$w_1\le r_1-2\le r\le7$ となる．

$\tau=+$ なら $D_C(P)>0$ である．
式~\eqref{bn21:A-offsets} と式~\eqref{bn21:negative-linear} より
\[
 r>\frac{17}{10}w_1-\frac52,\qquad
 T<\frac{17}{7}w_1+\frac{31}{7}.
\]
一方，両方の $L$ に対して式~\eqref{bn21:positive-lower} は
\[
 T\ge\frac32r+2>\frac{51}{20}w_1-\frac74
\]
を与える．従って
\begin{equation}\label{bn21:A-bound}
 \frac{17}{140}w_1<\frac{173}{28},\qquad
 w_1<\frac{865}{17}<51.
\end{equation}
特に $w_1\le50$ である．なお $L=4$ なら強い下界
$T\ge2r+3$ を使って $w_1<225/34<7$ と改善できる．

\paragraph{周期が $P^+\to N_1^+\to N_2^-\to P^+$ の場合．}
$N_2^-\to P^+$ では $\gamma_{N_2}+\phi_P=1+L\le4$ が必要なので
$L=3$ である．offset は
\begin{equation}\label{bn21:B-offsets}
 w_1=r+t_1,\quad r_2=r_1+t_2,\quad T=w_2+1,
 \qquad-1\le t_1\le2,\quad-2\le t_2\le2.
\end{equation}
$\tau=-$ なら $r_1\le2w_1-1$，$r_2\ge2w_2$ なので
\[
 2(T-1)=2w_2\le r_2\le r_1+2\le2w_1+1\le2r+5.
\]
整数性から $T-r\le3$ となる．$h=T-r-2\ge1$ より $h=1$ であり，
式~\eqref{bn21:small-h} は $r^2\le2r+3$，従って $r\le3$ を与える．
よって $w_1\le r+2\le5$ である．

$\tau=+$ なら $r\ge w_1-2$ と式~\eqref{bn21:positive-lower} により
\[
 w_2=T-1\ge\frac32w_1-2.
\]
二つの負枝の線形評価と式~\eqref{bn21:B-offsets} から
\[
 \frac{51}{20}w_1-\frac{39}{10}<r_2
 \le r_1+2<\frac{17}{7}w_1+2.
\]
従って $w_1<826/17<49$，すなわち $w_1\le48$ となる．
以上で全ての $\sigma$ と $\tau$ の場合を尽くした．
\end{proof}

\subsection{導出された有限範囲の全 scale 復元}

\begin{lemma}[小さい負 core の有限 certificate]
\label{bn21:finite-certificate}
正の整数局所 charts，奇数高分母，低辺の非負 node，対向低 $D$ の相殺を
満たす book 低三角形は，$2\le w\le50$ の負 core を含まない．
\end{lemma}

\begin{proof}
負 core を $2$，その canonical order を $(0,1,3,4)$ としてよい．
$2\le w\le50$ と $w^2/2<H(H+1)<2w^2$ を満たす
reverse parameters $(w,H)$ は八百九十八個である．
各々について
\[
 E=1-\frac{w^2}{H(H+1)}=\frac MN\ne0,\qquad
 c=\frac{sH}{w+H},\quad e=\frac{z(H+1)}{w+H+1}
\]
と置く．負枝の完全 factorization
\begin{equation}\label{bn21:factorization}
 (N-Mc)(N+Me)=N(N-M)>0
\end{equation}
では，$Ece=e-c+1$ と $c\ge2,e\ge1$ により
$1-Ec=(c-1)/e>0$，$1+Ee=(e+1)/c>0$ である．
従って両因子は正であり，$N(N-M)$ の正の約数を全て取ることにより
正の整数 $c,e$，従って全ての scale $s,z$ を復元できる．
行列の整数性，pair equation，奇数分母を課すと三百十三個の canonical cores が残る．
ここでは $s,z$ または高分母に探索上限を設定していない．

各固定 core の二本の低辺に対し，総和五の二十一個の非負 node triples を取る．
相手の低 row は $p^j_{ik}=1-n_k-p^i_{jk}$ に固定される．
定理~\ref{bk19:finite-completion} により，相手の reverse data は有限であり，
指定された低 $D$ と高高 pair equation の連立は最高二次の非零方程式に帰着する．
従って共鳴する reverse data を含めて全ての相手の scale を復元できる．

計算の独立性のため，verifier では元の pair equation を直接用いる．
\[
 u=S+S/r-1,\qquad v=T-T/(r+1)+1
\]
とすると，その方程式は
\[
 (uv-ST)sz+(u-S)s+(v-T)z+1=0.
\]
指定された低係数を $D_*$ としたとき
\[
 z=\begin{cases}
 \displaystyle\frac{(r+1)S}{rT}s+
 \frac{(r+1)(D_*-(T-S))}{T},&\text{固定 ray が }C,\\[4pt]
 s-D_*-1,&\text{固定 ray が }D
 \end{cases}
\]
を代入し，整数根を正確に求める．二次係数が非零であることは
$uv-ST=F/[r(r+1)]\ne0$ と正の線形係数から従う．
再構成した行列では全ての row sum と pair equation を直接再検査する．
さらに正の weights
\[
 \left(a,b,\frac{2a-1}{2r},\frac{2b-1}{2T}\right)
\]
により，全ての局所 strict positivity inequalities を直接確認できる．

全三百十三 core のうち，二本の低辺の双方に相手が存在するのは次の六個だけである．
$\mathcal N_3,\mathcal N_4$ は二本の低辺の全相手集合である．
最後の低辺 $34$ における残差は $D^3_4+D^4_3$ を意味する．
\begin{center}
\small
\begin{tabular}{lrrl}
固定 $(r,s,z,w,w)$ & $|\mathcal N_3|$ & $|\mathcal N_4|$ & 最後の全 $D$ 残差\\
\hline
$(4,4,5,2,2)$ & 4 & 1 & $2,1,-2,1$\\
$(6,6,14,3,3)$ & 1 & 1 & $-1$\\
$(7,7,2,4,4)$ & 1 & 2 & $-8,-5$\\
$(7,21,4,4,4)$ & 2 & 1 & $8,6$\\
$(8,8,27,4,4)$ & 1 & 1 & $-9$\\
$(35,35,16,18,18)$ & 1 & 1 & $-6$\\
\end{tabular}
\end{center}
十一組全てで残差は非零である．従って最後の辺の対向 $D$ 相殺が成立しない．
Python 標準ライブラリのみを用い，旧 generator を import しない
\nolinkurl{verify_book_two_negative_obstruction_v21.py} が
この有限全数復元を実装する．その全固定 core，個数，十一組の最後の node と
$D$ 残差は \nolinkurl{book_two_negative_obstruction_verification_v21.json}
に記録している．
\end{proof}

\begin{theorem}[Book 低三角形の二負枝不存在]
\label{bn21:two-negative-obstruction}
正の整数局所 charts，奇数高分母，低辺の非負 node，対向低 $D$ の相殺を
満たす book 低三角形が，負枝をちょうど二つ持つことはない．
従って既知の全正枝・全負枝の排除と合わせると，
book 低三角形が存在するなら負枝はちょうど一つ，正枝はちょうど二つである．
\end{theorem}

\begin{proof}
二負枝の低三角形が存在すれば，命題~\ref{bn21:analytic-finite} によって
その一つは $w\le50$ を満たす．補題~\ref{bn21:finite-certificate} がこれを排除する．
\end{proof}

\begin{corollary}[残る一負枝の二つの node 型]
\label{bn21:one-negative-normal-forms}
Book 低三角形が存在すると仮定する．唯一の負枝を $N$ とし，
高 ray の交換により $\sigma_N=+$ と正規化する．
node 周期は次の二型のいずれかであり，$N$ は $w\ge51$ を満たす：
\[
 \mathrm{I}:\quad N^+\longrightarrow P_1^+\longrightarrow P_2^+
 \longrightarrow N^+,
\]
\[
 \mathrm{II}:\quad N^+\longrightarrow P_1^+\longrightarrow P_2^-
 \longrightarrow N^+.
\]
型 II では $L_{P_2}=3$，$S_{P_2}=w-1$，$T_{P_2}=w+1$ であり，
辺 $P_2\to N$ の node は高 ray $0,1$ と残る低 ray の順で $(0,0,5)$ となる．
\end{corollary}

\begin{proof}
正枝二つがともに $\sigma=-$ なら，三本の低辺の高 ray $0$ 側の node 総和は
\[
 \sum n_0=6+\sum(v_i-u_i)-\sum\phi_i
 \le6-(2+3+3)-(1+0+0)=-3
\]
となり，不可能である．周期 $N^+\to P_1^-\to P_2^+\to N^+$ も，
辺 $P_1^-\to P_2^+$ が
$\gamma_{P_1}+\phi_{P_2}=L_{P_1}+L_{P_2}\ge6$ を与え，
式~\eqref{bn21:offset} に反する．残るのが型 I と II である．
型 II の辺 $P_2^-\to N^+$ では $L_{P_2}+1\le4$ により $L_{P_2}=3$ である．
この辺の offset は $t_N=-1$ に固定されるので $w=T_{P_2}-1$ となる．
残る低 ray の node は $1+L_{P_2}+1=5$ であり，
node 総和五と非負性から他の二つは零である．
最後に $w\ge51$ は補題~\ref{bn21:finite-certificate} から従う．
\end{proof}

この結論は全ての scale と高分母について成り立つ．
一方，負枝を一つ持つ型の全排除は本節では主張しない．

% v21 supplement: sharpen the residual one-negative book sectors.
% The earlier two-negative proof and its certificate remain unchanged.
\section{一負枝 book の成長不等式と二つの残余 sector}
\label{bs21:section}

系~\ref{bn21:one-negative-normal-forms} は，仮想的な book 低三角形を
唯一の負枝 $N$ と二つの正枝 $P_1,P_2$ に帰着し，
node 周期の二つの型 I, II を与えた．
本節ではこの二型をさらに低 $D$ の符号と $L$ によって絞る．
前節の二負枝排除を変更するものではなく，残余一負枝に対する追加の必要条件である．

高 ray の交換によって $\sigma_N=+$ とし，
node 周期に沿う next の低 $D$ の共通符号を $\tau$ とする．
前節により負枝の $w$ は少なくとも五十一である．
正枝には引き続き
\[
 \begin{gathered}
 \delta=T-S\in\{2,3\},\quad d=\delta-1,\quad h=S-r,\quad k=h+d,\\
 \alpha=sh/r,\quad\beta=zk/(r+1),\quad F=(2r+1)h-r^2d .
 \end{gathered}
\]
を用いる．

\begin{lemma}[共鳴尺度の分類と残余三角形での排除]
\label{bs21:beta}
正枝において $\beta=1$ となる reverse data は
\begin{equation}\label{bs21:resonances}
 (r,S,T,z)=(3,4,6,2),\qquad(2,3,6,1)
\end{equation}
の二つに限られる．型 I, II の残余 book 低三角形には，
これらの reverse data を持つ正枝は現れない．
\end{lemma}

\begin{proof}
$\beta=1$ なら $r+1=zk$ である．
式~\eqref{ps21:two-identities} の Bezout 恒等式に $\alpha=sh/r$ を代入すると
\[
 r^2=((r+1)z+1)h
\]
を得る．$r=zk-1$，$k=h+d$ を用いて整理すれば
\[
 k(1+2z-dz^2)=\delta.
\]
$k=h+d\ge d+1=\delta$ であり，括弧内は正の整数なので，
$k=\delta$，$1+2z-dz^2=1$ である．従って $h=1$，$dz=2$ となり，
式~\eqref{bs21:resonances} が従う．自由な $s$ はそれぞれ $3m,2m$ であり，
この議論は既知の分母十三・七の共鳴族の全ての scale を含む．

両方の node 型で，$P_1$ がこの枝なら $T_1=6$ である．
辺 $N^+\to P_1^+$ の offset は $T_1=r_N+t_1$，$t_1\ge1$ なので
$r_N\le5$，従って $w\le r_N-2\le3$ となる．
型 I の $P_2$ がこの枝なら $r_2\le3$ であり，辺 $P_2^+\to N^+$ から
$w\le r_2+2\le5$ となる．
型 II の $P_2$ がこの枝なら $T_2=w+1=6$ より $w=5$ となる．
いずれも $w\ge51$ に反する．
\end{proof}

\begin{lemma}[負の $F$ sector の成長下界]
\label{bs21:growth}
残余 book の正枝が $F<0$ を満たすなら
\begin{equation}\label{bs21:half-sector}
 0<\frac{-F}{h(h+\delta-1)}<\frac12,
 \qquad
 2(\delta-1)r^2<h^2+4rh+(\delta+1)h.
\end{equation}
特に
\begin{equation}\label{bs21:linear-growth}
 h>\frac37r-1,
 \qquad T>\frac{10}{7}r+1.
\end{equation}
\end{lemma}

\begin{proof}
$F<0$ なら $\alpha\ge\beta\ge1$ であり，
$E_+=-F/(hk)>0$ は
\[
 E_+\alpha\beta=\alpha-\beta+1
\]
を満たす．補題~\ref{bs21:beta} により $\beta\ge2$ なので
\[
 E_+\beta=1-\frac{\beta-1}{\alpha}<1.
\]
従って $E_+<1/\beta\le1/2$ であり，分母を払えば
式~\eqref{bs21:half-sector} の第二式を得る．

$h\le3r/7-1$ と仮定する．右辺が負なら $h\ge1$ に直ちに反する．
そうでなければ，$H\mapsto H^2+4rH+(\delta+1)H$ は $H\ge0$ で増加するので，
式~\eqref{bs21:half-sector} の右辺は $H=3r/7-1$ を代入した値以下である．
その値を $2(\delta-1)r^2$ から引くと
\[
 \begin{cases}
 \displaystyle\frac5{49}r^2+\frac{25}{7}r+2,&\delta=2,\\[3pt]
 \displaystyle\frac{103}{49}r^2+\frac{22}{7}r+3,&\delta=3
 \end{cases}
\]
となり，常に正である．これは式~\eqref{bs21:half-sector} に反する．
従って $h>3r/7-1$ であり，$T=r+h+\delta$，$\delta\ge2$ から
式~\eqref{bs21:linear-growth} を得る．
\end{proof}

\begin{proposition}[追加で排除する sector の解析的上界]
\label{bs21:bounded-sectors}
型 I, II の book 低三角形が次表のいずれかの sector に属すると仮定する．
その唯一の負枝は，対応する有限上界を満たす．
\begin{center}
\begin{tabular}{llc}
node 型 & 低 $D$ の符号と正枝の $L$ & 負枝の上界\\
\hline
I & $\tau=-$ & $w\le134$\\
I & $\tau=+$，少なくとも一方の正枝で $L=4$ & $w\le8$\\
II & $\tau=-$ & $w\le11$\\
II & $\tau=+$，$L_{P_1}=3$ & $w\le6$\\
\end{tabular}
\end{center}
これらの上界には scale または高分母の上界の仮定はない．
\end{proposition}

\begin{proof}
型 I の node 周期は $N^+\to P_1^+\to P_2^+\to N^+$ である．
offset を書けば
\begin{equation}\label{bs21:I-offsets}
 T_1=r_N+t_1,\quad T_2=r_1+t_2,\quad w=r_2+t_N,
 \qquad
 L_i-2\le t_i\le2,\quad-1\le t_N\le2.
\end{equation}
$\tau=-$ なら $D_C(N)<0$ より $r_N\le2w-1$ であり，
両正枝は $F<0$ を満たす．従って
\[
 T_1\le r_N+2\le2w+1,
\]
\[
 r_1\ge T_2-2>\frac{10}{7}r_2-1,
 \qquad r_2\ge w-2.
\]
さらに補題~\ref{bs21:growth} を $P_1$ に適用すると
\[
 T_1>\frac{10}{7}r_1+1
 >\frac{100}{49}r_2-\frac37
 \ge\frac{100w-221}{49}.
\]
これと $T_1\le2w+1$ から $2w<270$，すなわち $w<135$ を得る．

型 I の $\tau=+$ では両正枝で $F>0$ である．
補題~\ref{bn21:bounds} により，$L=3$ なら $T\ge3r/2+2$，
$L=4$ なら $T\ge2r+3$ となる．
少なくとも一方が $L=4$ の場合，式~\eqref{bs21:I-offsets} の
$r_1\ge T_2-2$ と合わせると，いずれの順序でも
\[
 T_1\ge3r_2+3\ge3w-3
\]
を得る．一方，負枝の $r_N<17w/7$ より
$T_1\le r_N+2<17w/7+2$ である．
従って $w<35/4<9$，すなわち $w\le8$ である．

次に型 II を考える．この型では $L_{P_2}=3$，$T_2=w+1$ であり，
残る offset は
\begin{equation}\label{bs21:II-offsets}
 T_1=r_N+t_1,\quad r_2=r_1+t_2,
 \qquad L_{P_1}-2\le t_1\le2,\quad -2\le t_2\le2.
\end{equation}
$\tau=-$ なら $P_2$ の next は $D$ なので $D_C(P_2)>0$，従って $F_2>0$ である．
その成長下界は
\[
 w+1=T_2\ge\frac32r_2+2,
 \qquad r_2\le\frac{2(w-1)}3
\]
を与える．一方 $P_1$ は $F_1<0$ なので，両方の $L_{P_1}$ に対して
$T_1\le2r_1+2$ である．従って
\[
 T_1\le2(r_2+2)+2\le\frac{4w+14}{3}.
\]
負枝の下界と $t_1\ge1$ から
\[
 T_1\ge r_N+1>\frac{17}{10}w+\frac12.
\]
両者を比較して $w<125/11<12$，すなわち $w\le11$ を得る．

最後に型 II，$\tau=+$，$L_{P_1}=3$ とする．
$P_2$ は $D_C(P_2)<0$，従って $F_2<0$ なので
\[
 w+1=T_2>\frac{10}{7}r_2+1,
 \qquad r_2<\frac7{10}w.
\]
正枝 $P_1$ の一般不等式 $r_1<S_1<2r_1$ と式~\eqref{bs21:II-offsets} から
\[
 T_1\le2r_1+1\le2(r_2+2)+1<\frac75w+5.
\]
一方，$D_C(N)>0$ より $r_N\ge2w$ であり，$T_1\ge r_N+1\ge2w+1$ である．
従って $w<20/3<7$，すなわち $w\le6$ となる．
\end{proof}

\begin{lemma}[追加の有限 certificate]
\label{bs21:finite-certificate}
Book 低三角形は $2\le w\le134$ の負 core を含まない．
\end{lemma}

\begin{proof}
補題~\ref{bn21:finite-certificate} と同じ完全復元を，今回の
命題~\ref{bs21:bounded-sectors} で導いた範囲 $2\le w\le134$ に適用する．
初期負 core の全 scale は正の factorization
式~\eqref{bn21:factorization} の全約数から得られる．
各低辺の二十一 node triples と指定された低 $D$ から，
相手の reverse data と全 scale を完全に復元する．
最後に二つの相手の間の低辺の node と対向 $D$ を直接確認する．
scale，高分母，行列 entry の上限は設定しない．
初期 reverse parameters は六千三百九十二個，
全 scale の復元後に残る負 canonical cores は千八十九個であり，
左右の全相手集合に含まれる相手の総数は百六十九個である．

この全数復元でも，両隣の相手集合が双方非空の初期 core は
補題~\ref{bn21:finite-certificate} の表と同じ六個だけであり，
最後の比較も同じ十一組だけである．従って最後の低辺の $D$ 残差は全て非零となる．
新しい検証器 \nolinkurl{verify_book_one_negative_sector_reduction_v21.py} は，
旧 v19 generator を使用しない第21版の独立復元実装をそのまま用い，
今回の全固定 core と左右の全相手集合を
\nolinkurl{book_one_negative_sector_reduction_verification_v21.json} に保存する．
これは旧来の $w\le150$ の探索結果を外挿したものではなく，
ここで証明した解析的上界に対する全 scale の再計算である．
\end{proof}

\begin{theorem}[一負枝 book の二つの残余 sector]
\label{bs21:two-sectors}
正の整数局所 charts，奇数高分母，低辺の非負 node，対向低 $D$ 相殺を満たす
book 低三角形が存在すると仮定する．
唯一の負枝 $N$ を $\sigma_N=+$ と正規化すると，次のいずれかに限られる：
\begin{enumerate}
 \item 型 I：$N^+\to P_1^+\to P_2^+\to N^+$，
 $\tau=+$，$(L_{P_1},L_{P_2})=(3,3)$；
 \item 型 II：$N^+\to P_1^+\to P_2^-\to N^+$，
 $\tau=+$，$(L_{P_1},L_{P_2})=(4,3)$．
\end{enumerate}
いずれも負枝は $w\ge135$，$H\ge w$ を満たす．
型 II では引き続き $S_{P_2}=w-1$，$T_{P_2}=w+1$，
辺 $P_2\to N$ の node は $(0,0,5)$ である．
\end{theorem}

\begin{proof}
命題~\ref{bs21:bounded-sectors} の四 sector は
補題~\ref{bs21:finite-certificate} により全て排除される．
残るのが表示した二 sector である．
負 core に対する同じ有限 certificate は $w\ge135$ を与える．
$\sigma_N=+$，$\tau=+$ より $D_C(N)>0$ なので，
式~\eqref{bn21:negative-sector} によって $r_N\ge2w$，従って $H=r_N-w\ge w$ である．
\end{proof}

この定理は二つの残余 sector の存在を示すものでも，それらの不存在を示すものでもない．
次の課題はこの二 sector に限って，node offset，低 $D$ の大きさ，
および中心二 core に接続する高分母の順序を同時に制御することである．

% v22: exact small-L rigidity closes the remaining book sectors.
% This section uses the v21 two-sector reduction and its existing certificate.
\section{小さい \texorpdfstring{$L$}{L} の正枝の剛性と book 支持の完全排除}
\label{bc22:section}

定理~\ref{bs21:two-sectors} の二つの残余 sector は，いずれも存在しない．
鍵は，正枝の $F>0$ 側における reverse data の成長下界が，
整数尺度の条件を使うと等式に強まることである．
この追加の議論には探索範囲の拡大を要しない．

\begin{lemma}[小さい $L$ の正枝の厳密な reverse data]
\label{bc22:positive-rigidity}
正の整数二高局所模型で $L\in\{3,4\}$，$F>0$ とする．
正枝の記号
\[
 \delta=L-1,\quad d=\delta-1,\quad h=S-r,\quad
 \alpha=\frac{sh}{r},\quad \beta=\frac{z(h+d)}{r+1},\quad
 F=(2r+1)h-r^2d
\]
を用いると，次が成り立つ：
\begin{equation}\label{bc22:rigidity}
 \begin{array}{c|ccc}
 L & r&S&T\\\hline
 3&2R&3R&3R+2\quad(R\in\mathbb Z_{>0})\\
 4&r&2r&2r+3
 \end{array}
\end{equation}
この結論は book に組み込まれることや低 $D$ の相殺を仮定しない．
\end{lemma}

\begin{proof}
整数行列から $\alpha,\beta$ は正の整数であり，
$0<h<rd$ である．式~\eqref{ps21:two-identities} により
\[
 F\alpha\beta=(\beta-\alpha-1)h(h+d).
\]
従って $E_+=F/[h(h+d)]>0$ と置くと
\begin{equation}\label{bc22:scale-bound}
 E_+\alpha=1-\frac{\alpha+1}{\beta}<1,
 \qquad
 E_+\alpha=\frac{r(2\alpha-ds)+\alpha}{h+d}.
\end{equation}
第二式には $r\alpha=sh$ を使った．

$L=3$，すなわち $d=1$ とする．
もし $h\le(r-1)/2$ なら $F\le-(r+1)/2<0$ なので，
$h$ の整数性から $h\ge r/2$ である．
従って $n=2\alpha-s=s(2h-r)/r$ は非負整数である．
$n\ge1$ なら，$h\le r-1$ と式~\eqref{bc22:scale-bound} により
\[
 E_+\alpha=\frac{rn+\alpha}{h+1}
 \ge\frac{r+1}{r}>1,
\]
矛盾する．従って $n=0$，$h=r/2$ であり，
$r=2R$，$S=r+h=3R$，$T=S+2=3R+2$ を得る．

$L=4$，すなわち $d=2$ とする．
$h\le r-1$ なら $F\le-r-1<0$ なので $h\ge r$ である．
従って $n=\alpha-s=s(h-r)/r$ は非負整数である．
$n\ge1$ なら $h\le2r-1$ から
\[
 E_+\alpha=\frac{2rn+\alpha}{h+2}
 \ge\frac{2r+1}{2r+1}=1,
\]
再び矛盾する．従って $h=r$，$S=2r$，$T=2r+3$ である．
\end{proof}

\begin{lemma}[負枝の厳密な reverse data]
\label{bc22:negative-rigidity}
負枝を $S=T=w$，$r=w+H$ と書き，
\[
 E=1-\frac{w^2}{H(H+1)}
\]
と置く．正の整数局所模型では $E\ne0$ であり，
\begin{equation}\label{bc22:negative-rigidity-data}
 \begin{array}{c|ccc}
  &H&r&E\\\hline
  E>0&w&2w&1/(w+1)\\
  E<0&w-1&2w-1&-1/(w-1)
 \end{array}
\end{equation}
に限られる．この結論も低辺の接合を仮定しない．
\end{lemma}

\begin{proof}
連続する平方数の間に $H^2<H(H+1)<(H+1)^2$ があるので，
$H(H+1)=w^2$ は不可能であり，$E\ne0$ である．
式~\eqref{bn21:factorization} で用いた正整数
\[
 c=\frac{sH}{w+H}\ge2,\qquad
 e=\frac{z(H+1)}{w+H+1}\ge1
\]
は $Ece=e-c+1$ を満たす．

$E>0$ なら $H\ge w$ であり，$Ec=1-(c-1)/e<1$ である．
$H>w$ と仮定すると
\[
 u=s-c=\frac{cw}{H}\in\mathbb Z_{>0},\qquad u\le c-1.
\]
従って
\[
 Ec=c-\frac{u^2}{c}\frac{H}{H+1}
 >c-\frac{u^2}{c}
 \ge2-\frac1c>1
\]
となり，矛盾する．従って $H=w$ である．

$E<0$ なら $H\le w-1$ であり，$(-E)e=(c-e-1)/c<1$ である．
$H\le w-2$ と仮定すると
\[
 v=z-e=\frac{ew}{H+1}\in\mathbb Z_{>0},\qquad v\ge e+1.
\]
従って
\[
 (-E)e=\frac{v^2}{e}\frac{H+1}{H}-e
 >\frac{v^2}{e}-e
 \ge2+\frac1e>1,
\]
再び矛盾する．従って $H=w-1$ である．
それぞれ $r=w+H$ と $E$ を計算すれば表示式を得る．
\end{proof}

\begin{proposition}[型 I の解析的排除]
\label{bc22:type-I}
定理~\ref{bs21:two-sectors} の型 I は存在しない．
より具体的には，型 I の reverse data と正の整数尺度だけで
$w\le18$ が必要になる．
\end{proposition}

\begin{proof}
この型では二つの正枝で $L=3$，$F>0$ なので，
補題~\ref{bc22:positive-rigidity} により
$T_i=(3r_i+4)/2$ である．
式~\eqref{bs21:I-offsets} は
\[
 T_1=r_N+t_1,\qquad T_2=r_1+t_2,\qquad w=r_2+t_N,
 \qquad t_1,t_2\in\{1,2\},\quad t_N\in\{-1,0,1,2\}
\]
となる．$r_1,r_2$ を消去すると
\begin{equation}\label{bc22:type-I-line}
 4r_N=9w+\kappa,\qquad
 \kappa=20-4t_1-6t_2-9t_N,\qquad -18\le\kappa\le19.
\end{equation}
この型の負枝は $H\ge w$ を満たすので $E>0$ であり，
補題~\ref{bc22:negative-rigidity} により $H=w$，$r_N=2w$ である．
式~\eqref{bc22:type-I-line} に代入すると $w=-\kappa\le18$ が必要となるが，
定理~\ref{bs21:two-sectors} は $w\ge135$ を与えている．
\end{proof}

\begin{proposition}[型 II の解析的排除]
\label{bc22:type-II}
定理~\ref{bs21:two-sectors} の型 II は存在しない．
この型の最初の低辺は，負枝の $r_N$ が同時に偶数と奇数であることを強制する．
\end{proposition}

\begin{proof}
$P_1$ は $L=4$，$F_1>0$ なので
補題~\ref{bc22:positive-rigidity} から $T_1=2r_1+3$ である．
辺 $N^+\to P_1^+$ の offset は
式~\eqref{bs21:II-offsets} により $t_1=2$ に固定されるので
\begin{equation}\label{bc22:II-first}
 r_N=2r_1+1.
\end{equation}
一方，負枝は $H\ge w$，従って $E>0$ を満たすので，
補題~\ref{bc22:negative-rigidity} が $r_N=2w$ を与える．
式~\eqref{bc22:II-first} と合わせると $2w=2r_1+1$ となり，
整数性に反する．この矛盾自体は $w\ge135$ を必要としない．
\end{proof}

\begin{theorem}[Book 支持の完全排除]
\label{bc22:book-obstruction}
標準 $K_5$ 分岐配置における book 支持
\[
 G_{5,13}=\{01,02,03,04,12,13,14\}
\]
では，二価 cores $2,3,4$ の低三角形が，正の整数局所 charts，
奇数高分母，低辺の非負 node および総和五，対向低 $D$ の相殺を
同時に満たすことはない．
従ってこの支持は，同じ局所条件を要求する標準構成から排除される．
\end{theorem}

\begin{proof}
定理~\ref{bs21:two-sectors} が全ての仮想的な低三角形を
型 I または型 II に帰着する．
命題~\ref{bc22:type-I} と命題~\ref{bc22:type-II} が両者を排除する．
\end{proof}

\paragraph{証明書の範囲．}
第21版の有限復元は，そこで解析的に導かれた $w\le134$ の範囲を
全尺度について検証している．本節はその既存の定理を用い，
$w\ge135$ の残余二 sector を整数性と恒等式だけで排除する．
新しい \nolinkurl{verify_book_complete_obstruction_v22.py} は，
正負両枝の剛性に用いた恒等式，十六個の node offset，
局所全尺度表示の row 和と pair 方程式，
型 II の強制式を標準ライブラリの厳密有理数演算で確認する．
これを新たな大きな $w$ の探索に置き換えてはいない．
また結論はここで指定した標準配置・整数局所条件に対するものであり，
一般型曲面の任意の Lagrangian torus fibration の不存在を主張するものではない．

\subsection{\texorpdfstring{$L\in\{1,3,4\}$}{L=1,3,4} の二高枝の完全な尺度分類}
\label{rs22:section}

前節の局所剛性を正枝の $F<0$ 側にも拡張すると，
$L\in\{1,3,4\}$ の全二高 canonical 模型が六族に分類される．
各族は全整数尺度を含み，高分母の奇数性も自動的に成り立つ．
これは book の排除後にも，他の支持の低辺を接続するために使える．
以下では式~\eqref{g51521:central} の canonical 行列を用い，
$a=sS+1$，$b=zT+1$，$x=sS/r$，$y=zT/(r+1)$ とする．
本節の尺度 $s,z$ と低係数 $D_C,D_D$ の記号は同式および
式~\eqref{g51521:low-D} に従う．

\begin{lemma}[正枝の $F<0$ 側の剛性]
\label{rs22:negative-F-rigidity}
正枝で $L\in\{3,4\}$ とする．$F=0$ は起こらず，$F<0$ の場合には
\[
 \begin{array}{c|ccc}
 L&r&S&T\\\hline
 3&2R-1&3R-2&3R\quad(R\ge2)\\
 4&r&2r-1&2r+2\quad(r\ge2)
 \end{array}
\]
に限られる．
\end{lemma}

\begin{proof}
引き続き $d=L-2\in\{1,2\}$，$h=S-r$，$k=h+d$，
$\alpha=sh/r$，$\beta=zk/(r+1)$ を用いる．
$F=(2r+1)h-r^2d=0$ なら，$\gcd(r^2,2r+1)=1$ より
$2r+1\mid d$ となるが，$2r+1\ge3>d$ なので不可能である．

$F<0$ とし，$\rho=r+1$ と置くと
$F=(2\rho-1)k-\rho^2d$ である．
$d=1$ で $k>\rho/2$ なら，$k$ の整数性から
$k\ge(\rho+1)/2$，従って $F\ge(\rho-1)/2>0$ となる．
ゆえに $k\le\rho/2$ である．
$d=2$ で $k\ge\rho+1$ なら $F\ge\rho-1>0$ なので，$k\le\rho$ である．

尺度方程式を $E_-=(-F)/(hk)>0$ によって書くと
\[
 E_-\beta=1-\frac{\beta-1}{\alpha}\le1,
 \qquad
 E_-\beta=\frac{\rho(dz-2\beta)+\beta}{h}.
\]
$d=1$ では $n=z-2\beta=z(1-2k/\rho)$ が非負整数であり，
$n\ge1$ なら分子は $\rho+1$ 以上，分母は
$h=k-1\le\rho/2-1$ なので矛盾する．
従って $n=0$，$k=\rho/2$ であり，第一行を得る．
$d=2$ では $n=z-\beta=z(1-k/\rho)$ が非負整数であり，
$n\ge1$ なら分子は $2\rho+1$ 以上，分母は
$h=k-2\le\rho-2$ なので同じ矛盾になる．
従って $k=\rho$ であり，第二行を得る．
\end{proof}

\begin{proposition}[正枝の全尺度]
\label{rs22:positive-scales}
補題~\ref{bc22:positive-rigidity} の $F>0$ の模型は，次の表示で尽くされる．
\begin{enumerate}
\item $L=3$ では $R\ge1$，$\ell\mid R+2$，$1\le\ell\le R$ とし，
\[
 t=(R+2)/\ell-1,\quad
 (r,S,T)=(2R,3R,3R+2),\quad
 s=2(R+1-\ell),\quad z=(2R+1)t.
\]
高分母はともに奇数である．
\item $L=4$ では $r\ge1$，$\ell\mid r+3$，$1\le\ell\le r+1$ とし，
\[
 t=(r+3)/\ell-1,\quad (S,T)=(2r,2r+3),\quad
 s=r+2-\ell,\quad z=(r+1)t.
\]
高分母はともに奇数である．
\end{enumerate}
これらの表示は正の整数局所 charts を与える．低辺の node や対向 $D$ の相殺は仮定しない．
\end{proposition}

\begin{proof}
$L=3$ では $\alpha=s/2$，$E_+=1/(R+1)$，
$\beta=z(R+1)/(2R+1)$ である．$\beta$ の整数性と
$\gcd(R+1,2R+1)=1$ から，$z=(2R+1)t$，$\beta=(R+1)t$ と書ける．
尺度方程式は $(t+1)(R+1-\alpha)=R+2$ なので，
$\ell=R+1-\alpha$ と置くと必要な全表示を得る．
$\alpha\ge1$ と $\ell>0$ は $1\le\ell\le R$ に一致する．
$L=4$ では $\alpha=s$，$E_+=1/(r+2)$，
$\beta=z(r+2)/(r+1)$ を使う．同様に
$(t+1)(r+2-s)=r+3$ となり，第二表示を得る．

逆に表示式を canonical 行列へ代入すれば，$x,y$ は正整数であり，
全 row sums と六つの pair products が成り立つ．
正の weights は両方で
\[
 (w_A,w_B,w_C,w_D)=
 \left(a,b,x+\frac1{2r},z+\frac1{2T}\right)
\]
と取れる．正の四 entry に対応する basin gaps は
$1/(2r),1/2,1/(2T),1/2$ であり，他の entry は非正である．
最後に，第一族では $Rt=\ell t(t+1)-2t$，
第二族では $(r+1)t=\ell t(t+1)-2t$ が偶数なので，
$sS,zT$ は両方とも偶数であり，高分母は奇数となる．
\end{proof}

\begin{proposition}[正枝の $F<0$ 側の全尺度と共鳴族]
\label{rs22:negative-F-scales}
補題~\ref{rs22:negative-F-rigidity} の全模型は，
$\ell\in\mathbb Z_{\ge0}$，$t\in\mathbb Z_{>0}$ により次のように表示される：
\begin{enumerate}
 \item $L=3$ では
 \[
 R=\ell(t+1)+2,\quad (r,S,T)=(2R-1,3R-2,3R),\quad
 s=(2R-1)t,\quad z=2(\ell t+1).
 \]
 \item $L=4$ では
 \[
 r=\ell(t+1)+2,\quad (S,T)=(2r-1,2r+2),\quad
 s=rt,\quad z=\ell t+1.
 \]
\end{enumerate}
両族は正の整数局所 charts を与え，高分母は自動的に奇数となる．
$\ell=0$ は，それぞれ任意の $t\ge1$ を持つ共鳴族
\[
 (r,s,z,S,T)=(3,3t,2,4,6),\qquad(2,2t,1,3,6)
\]
である．$\ell>0$ では，固定した $R$ または $r$ に対する尺度は
$(t+1)\ell=R-2$ または $(t+1)\ell=r-2$ の約数から全て得られる．
\end{proposition}

\begin{proof}
$L=3$ では $\alpha=s(R-1)/(2R-1)$，$\beta=z/2$，
$E_-=1/(R-1)$ である．$\gcd(R-1,2R-1)=1$ なので
$s=(2R-1)t$，$\alpha=(R-1)t$ と書ける．
尺度方程式から
\[
 (t+1)(R-1-\beta)=R-2
\]
を得る．$R\ge2$ なので $\ell=R-1-\beta\ge0$ であり，
表示式に一致する．$R=2$ のときも $\ell=0$ として全ての $t\ge1$ を含む．

$L=4$ では $\alpha=s(r-1)/r$，$\beta=z$，$E_-=1/(r-1)$ なので，
$s=rt$ と書けば $(t+1)(r-1-z)=r-2$ となる．
$\ell=r-1-z\ge0$ と置いて第二族を得る．
逆向きの row/pair 方程式と正の weights は
命題~\ref{rs22:positive-scales} と同じ直接代入で確認できる．
第一族の $Rt$ と第二族の $rt$ は，どちらも
$\ell t(t+1)+2t$ に等しく偶数なので，高分母は奇数である．
\end{proof}

\begin{proposition}[負枝の全尺度]
\label{rs22:negative-scales}
補題~\ref{bc22:negative-rigidity} の負枝は $S=T=w$ とし，
正の約数 $\ell\mid w$ および $t=w/\ell-1\ge1$ によって次のように尽くされる：
\[
 \begin{array}{c|c|c|c|c}
 &r&\ell\text{ の範囲}&s&z\\\hline
 E>0&2w&1\le\ell\le w-1&2(w+1-\ell)&(2w+1)t\\
 E<0&2w-1&1\le\ell\le w-2&(2w-1)t&2(w-1-\ell)
 \end{array}
\]
両方とも高分母は自動的に奇数となる．
これらは全て正の整数局所 charts を与え，低係数は
\[
 \begin{array}{c|cc}
 &D_C&D_D\\\hline
 E>0&\ell t^2-1&-2\ell t^2-t+1\\
 E<0&-\ell t^2-1&2\ell t^2-t+1
 \end{array}
\]
となる．特に $-D_D-2D_C=t+1$ が成り立つ．
\end{proposition}

\begin{proof}
$E>0$ の場合，$H=w$ なので $c=s/2$ であり，
$e=z(w+1)/(2w+1)$ の整数性から $z=(2w+1)t$ と書ける．
$Ece=e-c+1$ は $(t+1)(w+1-c)=w$ となる．
$\ell=w+1-c$ と置くと，$c\ge2$ より第一行の範囲と尺度を得る．
$E<0$ の場合は $H=w-1$，$e=z/2$ であり，
$c=s(w-1)/(2w-1)$ の整数性から $s=(2w-1)t$ と書ける．
尺度方程式は $(t+1)(w-1-e)=w$ なので，$\ell=w-1-e$ と置けば第二行を得る．
$e\ge1$ は $\ell\le w-2$ に対応する．
逆向きの row/pair equations と正の weights は前命題と同じ直接代入で確認できる．
$wt=\ell t(t+1)$ は偶数なので高分母は奇数である．
低係数の表も $w=\ell(t+1)$ を用いた計算である．
\end{proof}

$F>0$ の正枝の低係数にも短い転送式がある．$p=-D_D>0$，$q=D_C>0$ と書けば
\[
 \begin{array}{c|c|c|c}
 L&p&q&\text{関係}\\\hline
 3&2\ell t^2-3t+3&3\ell t^2-4t+5&2q-3p=t+1\\
 4&\ell t^2-2t+2&2\ell t^2-3t+5&q-2p=t+1
 \end{array}
\]
となる．$F<0$ の二族では
\[
 \begin{array}{c|cc|c}
 L&D_C&D_D&\text{関係}\\\hline
 3&5-3\ell t^2-4t&2\ell t^2+3t-3&2D_C+3D_D=t+1\\
 4&5-2\ell t^2-3t&\ell t^2+2t-2&D_C+2D_D=t+1
 \end{array}
\]
である．$\ell=0,t=1$ は，それぞれ低 $D$ が $(1,0)$，$(2,0)$ となる
既知の零低係数模型を含む．さらに $L=4$，$(\ell,t)=(1,1)$ では
$(r,s,z,S,T)=(4,4,2,7,10)$，$(D_C,D_D)=(0,1)$ となる．
この三つの零低係数端点を除くと $D_C<0<D_D$ である．
これら全六族の row/pair 方程式と低係数の恒等式は
\nolinkurl{verify_book_complete_obstruction_v22.py} が係数ごとに検査する．
固定した reverse data では，明示した二つの共鳴族を除いて尺度の集合は有限である．
共鳴族には無限個の尺度があり，また $R,r,w$ が動く全模型の一様有限性を
主張するものでもない．

\section{未完成の外側 core による尺度の線形復元}
\label{oc23:section}

一つの低辺の両端だけを先に完成しようとすると，その二つの局所模型には
無限個の尺度が残り得る．しかし第三の core について，四本の rays の全行列を
完成する前でも，一つの pair-product 方程式を課すことができる．
本節はこの必要条件を線形の形に書き，共鳴する全ての場合を明示する．
以下の $a,j,p,S$ は前節の canonical 変数とは独立に定める．

\begin{lemma}[第三 core の単一 pair による復元]
\label{oc23:linear-transfer}
三つの異なる cores を $i,j_0,k$ とし，$q_{ij_0}=1$，$q_{ik}=a>1$ とする．
core $i$ の entry を $p=p^i_{k j_0}$ と置く．
core $j_0$ の低 row と反対 column の対応する二成分を
\[
 p^{j_0}_{ik}=\varepsilon\,j,\qquad
 p^{j_0}_{ki}=\varepsilon\,S,
 \qquad \varepsilon\in\{1,-1\},\quad j,S\in\mathbb Z_{>0}
\]
と書くと，$q_{j_0k}=jS+1$ である．
二つの node に対して
\[
 u=1-n_{ik;j_0},\qquad v=1-n_{j_0k;i},
 \qquad -4\le u,v\le1
\]
と置き，$E=au-p$ とする．
core $k$ の rays $i,j_0$ の pair-product 方程式は，厳密に
\begin{equation}\label{oc23:linear-equation}
 \bigl(E(jv-\varepsilon)-aj\bigr)S=a-1-Ev
\end{equation}
と同値である．特に左辺の係数が非零なら，$S$ の候補は高々一つである．
左辺の係数を $A$ とし，$q=jS+1$ と置くと，さらに
\begin{equation}\label{oc23:denominator-divisor}
 qA=-\varepsilon E-j
\end{equation}
が成り立つ．従って非共鳴の解では $q\mid-\varepsilon E-j$ であり，
$q\le|\varepsilon E+j|$ である．
\end{lemma}

\begin{proof}
二つの clutching 方程式は
\[
 p^k_{ij_0}=au-p=E,\qquad
 p^k_{j_0i}=(jS+1)v-\varepsilon S
\]
を与える．これらの積を $a(jS+1)-1$ に等しく置き，$S$ で整理すれば
式~\eqref{oc23:linear-equation} を得る．各 node は非負であり，各辺の総数は五なので，
$u,v\in\{-4,-3,-2,-1,0,1\}$ である．
他の rays に関する core $k$ の entries や，その正の weights を先に決める必要はない．
式~\eqref{oc23:linear-equation} の右辺を $B$ とすると，
$A+jB=-\varepsilon E-j$ である．$AS=B$ を用いれば
式~\eqref{oc23:denominator-divisor} も得られる．
\end{proof}

\begin{proposition}[全共鳴の分類]
\label{oc23:resonances}
前補題において，$a$ は三以上の奇数，$p$ は $a-1$ の非零の signed divisor とする．
式~\eqref{oc23:linear-equation} の係数と定数がともに零となるのは，
次の表のいずれかの場合に限り，逆に表の全ての場合には両者が零となる．
\begin{center}\small
\begin{tabular}{c|c|c|c|c|l}
 $\varepsilon$&$j$&$p$&$u$&$v$&条件\\\hline
 $-1$&$a-1$&$1$&$1$&$1$&---\\
 $-1$&$a-1$&$-(a-1)$&$0$&$1$&---\\
 $+1$&$(a-1)/h$&$(a-1)/h$&$0$&$-h$&$h\mid a-1,\ 1\le h\le4$\\
 $+1$&$a-1$&$-1$&$-1$&$-1$&---\\
 $+1$&$1$&$-(a-1)$&$-1$&$-(a-1)$&$a\in\{3,5\}$
\end{tabular}
\end{center}
これは単一 pair 条件の共鳴であり，局所模型や接着全体の存在を意味しない．
\end{proposition}

\begin{proof}
定数が零である条件は $Ev=a-1$ である．これを係数が零である条件へ代入すると
\[
 -\varepsilon E=j,\qquad
 E=-\varepsilon j,\qquad
 jv=-\varepsilon(a-1)
\]
を得る．逆も直接成り立つ．また $p\mid a-1$ から
$1\le |p|\le a-1$ である．

$\varepsilon=-1$ では $v>0$ なので $v=1$，$j=a-1$ である．
$p=au-j$ と上記の絶対値上限から，$u=1,p=1$ または
$u=0,p=-(a-1)$ だけが残る．

$\varepsilon=+1$ では $v=-h$，$1\le h\le4$，$jh=a-1$ と書ける．
$p=au+j$ と $1\le j\le a-1$，$|p|\le a-1$ より，$u=0$ または $u=-1$ である．
$u=0$ は表の第三行を与える．$u=-1$ では $p=j-a$ なので
\[
 a-j\mid a-1\quad\Longleftrightarrow\quad a-j\mid j-1.
\]
$h=1$ なら $a-j=1$ であり，第四行を得る．
$h\ge2$ なら $a-j=(h-1)j+1>j-1\ge0$ であるから，$j=1$ が必要である．
このとき $a=h+1$ は三以上の奇数で，$h\le4$ なので $a=3,5$ に限られる．
これが第五行である．全表示を代入すれば逆向きも確認できる．
\end{proof}

\begin{corollary}[非共鳴時の有限尺度候補]
\label{oc23:finite-scale}
固定した $(a,p,j,\varepsilon)$ が前命題のどの共鳴にも当たらないなら，
式~\eqref{oc23:linear-equation} と node 条件を満たす正整数尺度 $S$ の候補は
高々三十六個であり，分母の上限を仮定せずに全て求められる．
\end{corollary}

\begin{proof}
$(u,v)$ の三十六通りを取り，係数が零で定数が非零なら除く．
係数が非零なら，定数を係数で割った値が正整数である場合だけを残す．
共鳴は仮定によってない．異なる node の組が同じ尺度を与える場合は重複を除いてよい．
\end{proof}

実際の接着では，さらに $jS+1$ の奇数性，他の pair equations，row sums，
全 node 配分，対向 $D$，正の weights を課す必要がある．
この補題は，その前段階で尺度の無限自由度を切るための必要条件である．
特に，一つの低辺の両側がともに完成しても，第三 core の pair を満たさなければ
大域の atlas には延長できない．

\begin{corollary}[低辺を横切る分母比較と残る共鳴]
\label{oc23:denominator-comparison}
前命題の設定で，さらに低辺 $ij_0$ 自身の puncture $k$ における
非負 node 条件を課す．非共鳴の解では
\[
 q_{j_0k}\le6q_{ik}+2
\]
である．$q_{ik}=a\ge7$ の場合，残る共鳴は次の二族だけである：
\[
\begin{array}{c|c|c|c|c|c}
 \varepsilon&p&j&u&v&n_{ij_0;k}\\\hline
 -1&1&a-1&1&1&1\\
 +1&-1&a-1&-1&-1&1
\end{array}
\]
いずれも $q_{j_0k}=1+(a-1)S$ であり，単一 pair と三つの表示された
node 条件は $S>0$ 全体で成り立つ．他の局所・大域条件は引き続き必要である．
$a=3,5$ には，この二族に加えて
\[
\begin{array}{c|c|c|c|c|c|c}
 a&\varepsilon&p&j&u&v&n_{ij_0;k}\\\hline
 3&-1&-2&2&0&1&4\\
 3&+1&-2&1&-1&-2&1\\
 5&+1&-4&1&-1&-4&1
\end{array}
\]
の三場合だけが追加される．
\end{corollary}

\begin{proof}
core $i$ の反対 entry は $r=(a-1)/p$ である．
低辺の clutching は，$w=1-n_{ij_0;k}\in[-4,1]$ により
\[
 r+\varepsilon j=w
\]
と書ける．従って $j\le|r|+4\le a+3$ である．
非共鳴の解では式~\eqref{oc23:denominator-divisor} と $|u|\le4$ を使い，
\[
 q\le|\varepsilon(au-p)+j|
 \le4a+(a-1)+(a+3)=6a+2
\]
を得る．

共鳴表の第三行では $r=h$，$\varepsilon j=j$ だから $w=h+j\ge2$ となり除かれる．
第二行では $w=-a$ なので $a=3$ だけが残る．
第一行と第四行では $w=0$，第五行でも $w=0$ であり，表示した全表を得る．
各場合を元の三つの clutching と単一 pair へ代入すれば，逆向きも成り立つ．
\end{proof}

\section{三高 core 同士を結ぶ単位枝の三角形分類}
\label{ut24:section}

前節の単一 pair による線形復元は，二つの端点がともに三高 core である場合に
さらに強くなる．本節では，単位枝の node 配分を二種類に限定し，必ず現れる
小さい隣接分母を決定する．この結果は $G_{5,14}$ の二つの単位枝と
$G_{5,16}$ の唯一の単位枝のいずれにも適用できる．
第三 core の残りの行列を完成したという主張ではない．

\subsection{二つの低行と共通の負位置}

単位枝を $ij$，残る三つの cores を $a,b,c$ とする．両端 $i,j$ は，
この枝だけを低 ray とする正の三高整数 core と仮定する．
定理~\ref{mh18:matching} により，各低行は負二成分と正一成分を持つ．
二つの正成分が同じ位置にあると，そこで
$p^i_{jk}+p^j_{ik}\ge2$ となり $n_{ij;k}\ge0$ に反する．
従って標識を選び，両低行を
\begin{equation}\label{ut24:low-rows}
 \begin{array}{c|ccc}
       &a&b&c\\\hline
 i\text{ の低行}&u&-t&-s\\
 j\text{ の低行}&-j_0&v&-s'
 \end{array}
 \qquad
 u=s+t-1,\quad v=s'+j_0-1
\end{equation}
と書ける．全ての表示した大きさは正整数である．
共通の負位置 $c$ では
\begin{equation}\label{ut24:small-negative-row}
 n_{ij;c}=1+s+s',\qquad s+s'\le4.
\end{equation}
位置 $a$ の node を $n$ と書くと
\begin{equation}\label{ut24:opposite-row}
 j_0=u+n-1,\qquad
 0\le n\le2,\qquad n_{ij;b}=4-s-s'-n.
\end{equation}
ここで $j_0$ は正整数の大きさであり，core の標識 $j$ とは異なる．

\subsection{共通の負位置での完全な pair 分類}

位置 $c$ における低 column を $-S,-T$ と書く．
\[
 q_{ic}=sS+1,\qquad q_{jc}=s'T+1,\qquad S,T>0.
\]
さらに $\alpha=1-n_{ic;j}$，$\beta=1-n_{jc;i}$ と置く．
clutching により core $c$ の該当する二 entries は
\[
 (sS+1)\alpha+S,\qquad(s'T+1)\beta+T
\]
となる．従って pair-product 条件は
\begin{equation}\label{ut24:negative-pair}
 \bigl((s\alpha+1)S+\alpha\bigr)
 \bigl((s'\beta+1)T+\beta\bigr)
 =(sS+1)(s'T+1)-1
\end{equation}
である．ここでも $-4\le\alpha,\beta\le1$ である．

\begin{lemma}[共通負位置の全解]
\label{ut24:negative-classification}
$s,s',S,T>0$，$s+s'\le4$ とし，$sS+1,s'T+1$ は奇数とする．
式~\eqref{ut24:negative-pair} の全解は次表で与えられる．
自由な $R$ は任意の正整数である．
\begin{center}\small
\begin{tabular}{cc|rr|cc|cc}
 $s$&$s'$&$\alpha$&$\beta$&$S$&$T$&$q_{ic}$&$q_{jc}$\\\hline
 2&1&1&0&$R$&2&$2R+1$&3\\
 1&2&0&1&2&$R$&3&$2R+1$\\\hline
 1&2&-1&-3&2&1&3&3\\
 1&2&-2&-1&2&1&3&3\\
 2&1&-1&-2&1&2&3&3\\
 2&1&-3&-1&1&2&3&3\\
 1&3&-2&-1&2&2&3&7\\
 3&1&-1&-2&2&2&7&3\\
 2&2&-1&-2&3&1&7&3\\
 2&2&-1&-2&2&2&5&5\\
 2&2&-2&-1&2&2&5&5\\
 2&2&-2&-1&1&3&3&7
\end{tabular}
\end{center}
最初の二行では core $c$ の pair は正であり，残りの十行では負である．
特に $s=s'=1$ には解がない．表は三角形の必要条件の全解であって，
各行が両端の完全な局所模型や大域接着に延長することまでは述べない．
\end{lemma}

\begin{proof}
左辺第一因子は $\alpha\ge0$ のとき正，$\alpha\le-1$ のとき負であり，
第二因子も同様である．右辺は正なので，二つの因子の符号は等しい．

両方が正なら $\alpha,\beta\in\{0,1\}$ である．
$(0,0)$ では左辺が小さ過ぎ，$(1,1)$ では大き過ぎる．
$(1,0)$ の場合の差は
\[
 [s+1-ss']ST-sS+(1-s')T.
\]
$s'\ge2$ なら全て非正で少なくとも一項は負になる．
$s'=1$ なら $S(T-s)=0$ なので $T=s$ である．
奇数分母は $s$ を偶数にし，$s+s'\le4$ より $s=2$ となる．
$(0,1)$ は両端を交換した場合である．

両方が負なら $\alpha=-A$，$\beta=-B$，$1\le A,B\le4$ とする．
差は
\begin{equation}\label{ut24:bilinear-coefficients}
 C ST+D S+E T+AB,
 \quad
 \begin{cases}
 C=ss'(AB-1)-sA-s'B+1,\\
 D=s(AB-1)-B,\\
 E=s'(AB-1)-A.
 \end{cases}
\end{equation}
$A,B\ge2$ なら
$(sA-1)(s'B-1)\ge ss'$ より $C\ge0$，また $D,E>0$ なので解はない．
$A=B=1$ なら差は $-(s+s'-1)ST-S-T+1<0$ である．
両端交換により，残りは $A=1$，$2\le B\le4$ だけ調べればよい．

$s=1$ の場合，式は
\[
 S=B-1-\frac{T-1}{s'T+1}
\]
と書ける．$T>1$ では最後の分数が零と一の間にあるため，$T=1$，$S=B-1$ である．
奇数分母と $s'\le3$ から $s'=2$，$B=3$ を得る．

$s=2$ の場合，$D=B-2$，$C=s'(B-2)-1$，$E=s'(B-1)-1$ である．
$B\ge3$ では全係数が非負で定数項が正なので解はない．
$B=2$ では
\[
 T(S-s'+1)=2.
\]
$s'=1$ では奇数性から $(S,T)=(1,2)$，
$s'=2$ では $(S,T)=(3,1),(2,2)$ となる．

$s=3$ では $s'=1$ である．$B\ge3$ は再び係数の符号で除かれ，
$B=2$ では $S(T-1)=2$ となる．奇数性により $(S,T)=(2,2)$ だけが残る．
これらと両端交換が表の十点である．表への直接代入で逆向きも確認できる．
\end{proof}

\subsection{符号が反対の位置での全尺度族}

位置 $a$ の低 row を $u,-j_0$，低 column を $U,-T$ とする．
従って $q_{ia}=uU+1$，$q_{ja}=j_0T+1$ である．
$n=n_{ij;a}$，$j_0=u+n-1$ とし，
$\alpha=1-n_{ia;j}$，$\beta=1-n_{ja;i}$ と置くと，
第三 core の pair は
\begin{equation}\label{ut24:opposite-pair}
 \bigl((u\alpha-1)U+\alpha\bigr)
 \bigl((j_0\beta+1)T+\beta\bigr)
 =(uU+1)(j_0T+1)-1
\end{equation}
となる．

\begin{proposition}[反対符号位置の全族]
\label{ut24:opposite-classification}
$n\in\{0,1,2\}$，$u,j_0,U,T>0$，$j_0=u+n-1$ とし，
$uU+1,j_0T+1$ は奇数とする．
式~\eqref{ut24:opposite-pair} の全解は，次の四通常族と四種類の小例外である．
\begin{center}\small
\begin{tabular}{c|rr|l|l}
 $n$&$\alpha$&$\beta$&尺度条件&整数・奇偶条件\\\hline
 0&1&1&$T=U-1$&$u,U\ge2$，$u,U$ は異なる偶奇\\
 1&1&1&$U=1$&$u\ge2$ 偶数，$T\ge1$\\
 1&-1&-1&$T=1$&$u\ge2$ 偶数，$U\ge1$\\
 2&-1&-1&$T=U+1$&$u,U\ge1$，$u,U$ は異なる偶奇
\end{tabular}
\end{center}
\begin{center}\small
\begin{tabular}{ccc|rr|cc|l}
 $n$&$u$&$j_0$&$\alpha$&$\beta$&$U$&$T$&条件\\\hline
 0&2&1&0&-3&$2r$&$2r$&$r\ge1$\\
 1&1&1&0&-2&$2r$&$2r$&$r\ge1$\\
 1&2&2&0&-3&1&1&---\\
 1&1&1&$-k$&-1&$2r$&$k$&$r\ge1,\ k\in\{2,4\}$
\end{tabular}
\end{center}
第三 core の pair が正なのは最初の表の最初の二族だけである．
\end{proposition}

\begin{proof}
第一因子が正であることは $\alpha=1$ と同値であり，第二因子が正であることは
$\beta\in\{0,1\}$ と同値である．両方が正の場合，$(\alpha,\beta)=(1,0)$ では
差の $UT$ 係数 $u(1-j_0)-1$ が負，$U$ 係数が $-u$，$T$ 係数が $1-j_0$ なので不可能である．
$(1,1)$ の場合は
\[
 T(1-nU)=U-1
\]
となり，正整数性から最初の二族を得る．

両方が負の場合は $\alpha=-A$，$\beta=-B$ とし，
$0\le A\le4$，$1\le B\le4$ とする．差は
\begin{equation}\label{ut24:opposite-coefficients}
 \bigl[uj_0(AB-1)+j_0B-uA-1\bigr]UT
 +\bigl[u(AB-1)+B\bigr]U
 +\bigl[j_0(AB-1)-A\bigr]T+AB.
\end{equation}
$A\ge1,B\ge2$ なら全係数が非負で $U$ 係数が正であるため不可能である．
$A=B=1$ なら
\[
 (n-2)UT+U-T+1=0
\]
となり，第三・第四の通常族を得る．
$A>1,B=1$ とする．$T$ 係数は $(j_0-1)(A-1)-1$ であり，
これが非負なら他の全係数も非負で定数項が正である．
従って $j_0=1$ である．$n=0,u=2$ では
$T=A-1+(U+1)/(2U+1)$ が整数でなく，$n=1,u=1$ では $T=A$ となる．
奇数性が $A\in\{2,4\}$，$U$ 偶数を強制し，最後の小例外を与える．

最後に $A=0$ とする．$d=B-u$ と書くと
\[
 T\bigl((j_0d-1)U-j_0\bigr)=-dU.
\]
$d\le0$ では正整数解がない．$d\ge1$ の場合，右辺が負なので
$(j_0d-1)U<j_0$ が必要である．$j_0d=1$ なら $j_0=d=1$，$T=U$ であり，
$n=0,u=2,B=3$ または $n=1,u=1,B=2$ となる．これが最初の二小例外である．
$j_0d>1$ なら $j_0(d-1)<1$ より $d=1$，さらに $U=1$，$T=1$ である．
$B=u+1\le4$ と奇数性により $u=j_0=2,n=1,B=3$ が残る．
全ての族で奇数性を課すと表示した偶奇条件を得る．直接代入で逆向きも成立する．
\end{proof}

\subsection{単位枝に必ず現れる小分母と五node境界}

\begin{theorem}[二種類の node 配分と小分母]
\label{ut24:unit-edge-boundary}
標準五次接着において，単位枝 $ij$ の両端が正の三高 core であるとする．
三つの node の多重集合は
\[
 \{n_{ij;a},n_{ij;b},n_{ij;c}\}
   =\{5,0,0\}\quad\text{または}\quad\{4,1,0\}
\]
に限られる．最大 node の位置 $c$ は，両端の低行がともに負である唯一の位置である．
\begin{enumerate}
 \item $n_{ij;c}=5$ なら，$(q_{ic},q_{jc})$ は
 $(3,7),(5,5),(7,3)$ のいずれかで，core $c$ の pair $ij$ は負である．
 \item $n_{ij;c}=4$ なら，$\min(q_{ic},q_{jc})=3$ である．
 core $c$ の pair が負なら両分母は三，正なら分母対は
 $(3,2R+1)$ または $(2R+1,3)$，$R\ge1$ である．
\end{enumerate}
\end{theorem}

\begin{proof}
式~\eqref{ut24:small-negative-row} と補題~\ref{ut24:negative-classification} により
$s+s'$ は三または四である．総 node 数が五なので配分は表示した二種類となる．
分母対と第三 core の符号は同補題の表から直接従う．
\end{proof}

\begin{proposition}[五node境界での二つの正 pair]
\label{ut24:five-node-positive}
前定理の仮定に対向低係数 $D^i_j+D^j_i=0$ を加え，$n_{ij;c}=5$ とする．
他の二位置 $a,b$ では core $a,b$ の pair $ij$ がともに正であり，
それぞれの隣接高辺の該当 node は零である．
\end{proposition}

\begin{proof}
$n_{ij;a}=n_{ij;b}=0$ である．各位置について，正の低行を持つ側の低 column を
$U_a,U_b$，反対側の負 column の絶対値を $T_a,T_b$ とする．
命題~\ref{ut24:opposite-classification} の $n=0$ から，
\[
 U_a-T_a\in\{0,1\},\qquad U_b-T_b\in\{0,1\},
\]
である．値一は通常族 $\alpha=\beta=1$，値零は小例外を意味する．
位置 $c$ の負 column の絶対値は補題~\ref{ut24:negative-classification} により
$S+T=4$ を満たす．一方，対向 $D$ 相殺は全低 column の和を $-2$ とするので
\[
 (U_a-T_a)+(U_b-T_b)=S+T-2=2.
\]
両差は一となり，両位置で $\alpha=\beta=1$，すなわち隣接高辺の node が零である．
\end{proof}

\begin{proposition}[少なくとも一つの正の外側 pair]
\label{ut24:positive-witness}
前定理の仮定に対向低 $D$ の相殺を加えると，三つの外側 cores の少なくとも一つで
単位枝の両端を結ぶ pair が正になる．
\end{proposition}

\begin{proof}
五nodeの場合は前命題で示したので，四nodeの場合に三つの pair が全て負と仮定する．
標識を選び，$n_{ij;a}=0$，$n_{ij;b}=1$，$n_{ij;c}=4$ とする．
位置 $a$ に正の低行を持つ側を $i$ とすれば，式~\eqref{ut24:low-rows} において
$s+s'=3$，$j_0=u-1$，$v=t$ である．
位置 $a$ は負 pair なので命題~\ref{ut24:opposite-classification} の小例外により
\[
 u=2,\quad j_0=1,\quad U_a=T_a=2r,\qquad r\ge1.
\]
位置 $c$ も負 pair なので $S+T=3$ である．対向低 $D$ 相殺から，
位置 $b$ の低 column の正側と負側の絶対値の差は一でなければならない．

$u=s+t-1=2$ より $(s,t)=(1,2)$ または $(2,1)$ である．
$t=1$ の場合，位置 $b$ では $n=1,u=1$ の小例外しか残らず，
column の差は零または偶数であって一にはならない．
従って $s=1,t=2,s'=2$ であり，位置 $b$ は通常負族の
正 column 二，負 column の絶対値一に限られる．
位置 $c$ の column の絶対値は $(S,T)=(2,1)$ である．

core $i$ の高 rays を $b,c,a$ の順に並べると，低 row と column は
\[
 (-2,-1,2),\qquad(-1,-2,2r),
\]
高分母は $(3,3,4r+1)$ となる．$K=6r+1$，高 ray $a$ の $b$ 成分を $A$ とする．
その残りの高成分は行和から $C=-K-A$ であり，pair products から $A,C\ne0$ である．
残る二つの高行の和は $-2,-1$，共通 pair product は八であるから
\[
 \left(-2-\frac{2K}{A}\right)
 \left(-1-\frac{2K}{C}\right)=8.
\]
$C=-K-A$ を代入して非零の $AC$ を払うと
\[
 (3A+K)(A+K)=0
\]
となる．$A=-K$ は $C=0$ に反し，$A=-K/3$ は $K\equiv1\pmod3$ に反する．
従ってこの低 row/column は整数局所模型へ完成できず，矛盾する．
\end{proof}

\begin{corollary}[$G_{5,14}$ の五node位置]
\label{ut24:g514-five-node}
$G_{5,14}$ の単位枝を $ij,kl$，残る四高 core を $o$ とする．
$ij$ に五nodeの位置があるなら，その位置は $k$ または $l$ であり，$o$ ではない．
両端の低行の正位置の一つは $o$ である．
\end{corollary}

\begin{proof}
仮に五node位置が $o$ なら，命題~\ref{ut24:five-node-positive} により core $k,l$ の pair $ij$ は両方とも正である．
三高 matching 性により，core $k$ では低 ray $l$ と ray $o$ が正に結ばれ，
core $l$ でも低 ray $k$ と ray $o$ が正に結ばれる．
従って単位枝 $kl$ の位置 $o$ における両低行 entry が正となり，node 非負性に反する．
残る二位置が低行の正位置なので，その一つは $o$ である．
\end{proof}

これらの結果は $G_{5,14}$ や $G_{5,16}$ 全体の排除ではない．
大域接着には，残りの第三 core の pair，行和，正 weights，全高辺の node 配分，
対向高 $D$ の相殺をさらに課す必要がある．
\nolinkurl{verify_unit_three_high_triangles_v24.py} は恒等式を厳密係数で確認し，
共通負位置の全九十六負符号場合を整数因数分解から完全に列挙する．
反対符号の無限族は上の手計算による分類であり，有界回帰を無上界証明として用いない．

\section{\texorpdfstring{$G_{5,14}$}{G5,14} の五node単位枝の無上界排除}
\label{ug25:section}

本節では二つの単位枝を同時に用いる．単位枝一つの三角形条件だけでは残った
$\{5,0,0\}$ 配分が，もう一つの単位枝と正の三高 core の行和を合わせると
排除されることを示す．分母や行列成分に探索上限を置かない．

\subsection{仮定と，利用する局所事実}

五つの cores を $0,1,2,3,4$ とし，
$q_{01}=q_{23}=1$，他の八分母は一より大きい奇数とする．
core $i$ の ray $j$ から $k$ への整数成分を $p^i_{jk}$ と書き，
元の条件を
\begin{align}
 \sum_{k\ne i,j}p^i_{jk}&=-q_{ij},&
 p^i_{jk}p^i_{kj}&=q_{ij}q_{ik}-1,\label{ug25:local}\\
 w^i_k-\frac{p^i_{jk}}{q_{ij}}w^i_j&>0\quad(w^i_j>0),&
 n_{ij;k}&=1-\frac{p^i_{jk}+p^j_{ik}}{q_{ij}}\in\mathbb Z_{\ge0},
 \label{ug25:positive-node}\\
 \sum_{k\ne i,j}n_{ij;k}&=5,&
 D^i_j+D^j_i&=0,\qquad
 D^i_j=-q_{ij}-\sum_{k\ne i,j}p^i_{kj}
 \label{ug25:budget-D}
\end{align}
とする．本節の結論はこの標準五次接着の必要条件に対するものであり，
任意の Lagrangian torus fibration の分類を主張しない．

正の三高 core の正 pair graph は完全マッチングである
（三高 matching 定理~\ref{mh18:matching}）．従って唯一の低行は
正一成分・負二成分を持ち，その二つの負位置同士も正 pair を作る．
以下ではこの帰結を明示的に使う．四高 core の正 graph の分類は使わない．

単位枝の両低行の正位置は異なる．同じ位置ならその node が負になるからである．
正位置を $a,b$，共通の負位置を $c$ として
\[
 \begin{array}{c|ccc}
 &a&b&c\\\hline
 i&u&-t&-s\\
 j&-j_0&v&-s'
 \end{array},\qquad
 u=s+t-1,\quad v=s'+j_0-1
\]
と書くと，$n_{ij;c}=1+s+s'$，従って $s+s'\le4$ である．
共通負位置の低 column を $-S,-T$ と書けば
$q_{ic}=sS+1$，$q_{jc}=s'T+1$ である．
高枝の node から得る $\alpha,\beta\in\{-4,-3,-2,-1,0,1\}$ により，
外側 core $c$ の pair equation は
\begin{equation}\label{ug25:negative-pair}
 ((s\alpha+1)S+\alpha)((s'\beta+1)T+\beta)
 =(sS+1)(s'T+1)-1
\end{equation}
となる．ここから得る必要条件を再掲する：
\begin{center}\small
\begin{tabular}{c|c|c|c}
外側 pair&$s+s'$&$(q_{ic},q_{jc})$&必要な低成分\\\hline
正&3&$(2R+1,3)$&$(s,s',S,T)=(2,1,R,2)$\\
正&3&$(3,2R+1)$&$(s,s',S,T)=(1,2,2,R)$\\
負&3&$(3,3)$&---\\
負&4&$(3,7)$&$(1,3,2,2)$ または $(2,2,1,3)$\\
負&4&$(5,5)$&$(2,2,2,2)$\\
負&4&$(7,3)$&$(3,1,2,2)$ または $(2,2,3,1)$
\end{tabular}
\end{center}
$R$ は任意の正整数である．特に全配分は $\{4,1,0\}$ または
$\{5,0,0\}$ であり，五nodeの場合には $S+T=4$ である．
この表は補題~\ref{ut24:negative-classification} の全解から従うが，
必要な無上界性は次の因数分解でも確認できる．
負 pair に対して $\alpha=-A,\beta=-B$ とすると，
\eqref{ug25:negative-pair} の左辺引く右辺は
\[
 CST+DS+ET+AB,
 \quad C=ss'(AB-1)-sA-s'B+1,
 \quad D=s(AB-1)-B,
 \quad E=s'(AB-1)-A.
\]
$A,B\ge2$ は係数の非負性で，$A=B=1$ は狭義負で除かれる．
交換により $A=1,2\le B\le4$ としてよい．
$s=1$ なら $S=B-1-(T-1)/(s'T+1)$ より $T=1$，
$s=2,B=2$ なら $T(S-s'+1)=2$，
$s=3,s'=1,B=2$ なら $S(T-1)=2$ である．
$s\ge2,B\ge3$ は係数の符号で除かれる．奇数性が上表の全点を与える．
正 pair では $\alpha,\beta\in\{0,1\}$ となり，
$(1,0)$ が $s'=1,T=s$，$(0,1)$ がその交換を与える．
奇数性と $s+s'\le4$ が二つの正族を与える．

次に反対符号位置の低行を $v,-j$，低 column を $V,-J$ とする．
$n$ を単位枝上のその位置の node とすると $j=v+n-1$ である．
外側 pair equation は
\begin{equation}\label{ug25:opposite-pair}
 ((v\alpha-1)V+\alpha)((j\beta+1)J+\beta)
 =(vV+1)(jJ+1)-1.
\end{equation}
本節で必要な負 pair，$n\in\{0,1\}$ の全族は
\begin{center}\small
\begin{tabular}{c|c|c|c|c|c}
$n$&$v$&$j$&$V$&$J$&$(\alpha,\beta)$\\\hline
0&2&1&$2r$&$2r$&$(0,-3)$\\
1&正の偶数&$v$&任意の正整数&1&$(-1,-1)$\\
1&1&1&$2r$&$2r$&$(0,-2)$\\
1&2&2&1&1&$(0,-3)$\\
1&1&1&$2r$&$k\in\{2,4\}$&$(-k,-1)$
\end{tabular}
\end{center}
であり，$r\ge1$ である．正 pair，$n=0$ の全族は
\begin{equation}\label{ug25:positive-zero}
 \alpha=\beta=1,\qquad j=v-1,\qquad J=V-1,
 \qquad v,V\ge2,\quad v\not\equiv V\pmod2
\end{equation}
である．これらも命題~\ref{ut24:opposite-classification} の一部だが，
本節で使う範囲の完全性を短く確認しておく．負 pair では
$\alpha=-A,\beta=-B$，$0\le A\le4,1\le B\le4$ と置く．差は
\[
 [vj(AB-1)+jB-vA-1]VJ
 +[v(AB-1)+B]V+[j(AB-1)-A]J+AB.
\]
$A\ge1,B\ge2$ は係数の非負性で除かれる．
$A=B=1$ は $(n-2)VJ+V-J+1=0$ なので，$n=1,J=1$ となる．
$A>1,B=1$ では $J$ 係数が負でなければならず $j=1$ であり，
$n=0,v=2$ は整数解なし，$n=1,v=1$ は $J=A$ である．
$A=0$ では $d=B-v$ と置くと
\[
 J((jd-1)V-j)=-dV.
\]
$d\le0$ は不可能，$jd=1$ は $j=d=1,J=V$，
$jd>1$ は $d=V=J=1$ となる．これに奇数性を課すと上表を得る．
正 pair では $(\alpha,\beta)=(1,1)$ 以外が排除され，
$J(1-nV)=V-1$ から \eqref{ug25:positive-zero} を得る．

\subsection{五node仮定の正規形}

\begin{lemma}\label{ug25:normal-form}
単位枝 $01$ が五nodeを持つなら，標識を取り直してその位置を $2$ とし，
低 row と低 column を，列順 $(3,4,2)$ で
\begin{equation}\label{ug25:normal-rows}
 \begin{array}{c|c|c}
 &\text{低 row}&\text{低 column}\\\hline
 0&(u,-t,-s)&(A,1-B,-S)\\
 1&(1-u,t+1,-s')&(1-A,B,-T)
 \end{array}
\end{equation}
と書ける．ここで
\[
 u=s+t-1,\quad s+s'=S+T=4,\quad u,A,B\ge2,
 \quad u\not\equiv A\pmod2,
\]
かつ $(s,s',S,T)$ は上の五node表に属する．従って
\begin{align}
 q_{02}&=sS+1\in\{3,5,7\},&q_{12}&=s'T+1,\label{ug25:qnormal1}\\
 q_{03}&=uA+1\ge7,&q_{13}&=(u-1)(A-1)+1,\label{ug25:qnormal2}\\
 q_{04}&=t(B-1)+1,&q_{14}&=(t+1)B+1.\label{ug25:qnormal3}
\end{align}
core $0$ の pair $23$ は負，core $1$ の pair $23$ は正である．
core $3$ の低行の正位置は $4$ である．
\end{lemma}

\begin{proof}
五nodeの他の二位置は node 零である．
\eqref{ug25:opposite-pair} の $n=0$ の二族により，
正側低 column と負側低 column の絶対値の差は零または一である．
五node位置では $S+T=4$ であり，対向低 $D$ の相殺は全低 column の和を
$-2$ とするので，残る二差の和は二である．従って両方の差は一で，
両外側 pair は正となる．

五node位置が $4$ なら，core $2,3$ の pair $01$ がともに正となる．
matching によりそれぞれの低行の正位置が $4$ となり，単位枝 $23$ の
位置 $4$ の node が負になる．従って五node位置は $2$ または $3$ である．
交換して $2$ とし，$0,1$ を交換して core $0$ の低正位置を $3$，
core $1$ の低正位置を $4$ とする．
\eqref{ug25:positive-zero} を二つの位置に適用すると
\eqref{ug25:normal-rows} と分母式を得る．
$u,A\ge2$ は異なる偶奇なので $uA\ge6$ である．
matching が残りの符号を決め，core $3$ で pair $01$ が正であることから
その低正位置は $4$ となる．
\end{proof}

\subsection{もう一方の単位枝による二つの還元}

core $2$ で pair $01$ は負であるため，matching によりその低正位置は
$0$ または $1$ である．$4$ なら pair $01$ が正になるので不可能である．

\begin{lemma}\label{ug25:case-A}
core $2$ の低正位置が $0$ なら
\[
 (q_{02},q_{03},q_{12},q_{13})=(7,7,3,3),
\]
かつ $(s,s',S,T,u,A,t)$ は
\begin{equation}\label{ug25:A-cases}
 (3,1,2,2,3,2,1),\quad
 (2,2,3,1,2,3,1),\quad
 (2,2,3,1,3,2,2)
\end{equation}
の三つのいずれかである．
\end{lemma}

\begin{proof}
単位枝 $23$ の両低正位置は $0,4$ なので，共通負位置は $1$ である．
core $1$ の pair $23$ が正であるから，その位置の node は四であり，
位置 $0$ の node $n$ は零または一である．
core $0$ の pair $23$ は負である．その反対符号位置の表を用い，
core $2$ 側を正の低行 $v$，低 column $V$，core $3$ 側を
負の低行 $-j$，低 column $-J$ とする．
\[
 q_{02}=vV+1\le7,\qquad q_{03}=jJ+1\ge7.
\]
$n=0$ の例外は $v=2,j=1,V=J=2r$ であり，第一不等式から $r=1$，
従って $q_{03}=3$ となり矛盾する．
$n=1$ の通常負族は $j=v,J=1$ であり，
$v+1\le vV+1\le7\le v+1$ から $v=6,V=1$ となる．
例外 $v=j=1,V=J=2r$ も両分母を七に強制する．
他の二例外では $q_{03}\le5$ なので不可能である．
従って常に $q_{02}=q_{03}=7$ である．
最初の単位枝の五node表から $q_{12}=3$，
$uA=6$ から $(u,A)=(2,3),(3,2)$ かつ $q_{13}=3$ を得る．
五node表と $t=u-s+1\ge1$ を組み合わせると
\eqref{ug25:A-cases} が尽くされる．
\end{proof}

\begin{lemma}\label{ug25:case-B}
core $2$ の低正位置が $1$ なら
\[
 (q_{02},q_{03},q_{12},q_{13})=(3,7,7,3),
\]
かつ $(s,s',S,T,u,A,t)$ は
\begin{equation}\label{ug25:B-cases}
 \begin{split}
 &(1,3,2,2,2,3,2),\quad (1,3,2,2,3,2,3),\\
 &(2,2,1,3,2,3,1),\quad (2,2,1,3,3,2,2)
 \end{split}
\end{equation}
の四つのいずれかである．
\end{lemma}

\begin{proof}
単位枝 $23$ の共通負位置は $0$ で，そこでの pair $23$ は負である．
共通負位置の表により $(q_{02},q_{03})$ は
$(3,3),(3,7),(5,5),(7,3)$ のいずれかである．
$q_{03}\ge7$ により $(3,7)$ だけが残る．
最初の単位枝の五node表から $q_{12}=7$，また $uA=6$ から
$q_{13}=3$ を得る．同じ表の $(3,7)$ の二つの低成分と
$(u,A)=(2,3),(3,2)$ を組み合わせると
\eqref{ug25:B-cases} を得る．
\end{proof}

\subsection{二つの固定小分母を持つ core の有限因子障害}

ここまで残った core $1$ の低 row と低 column を，列順 $(3,2,4)$ で
\begin{equation}\label{ug25:fixed-core}
 (-(u-1),-s',v),\qquad (-(A-1),-T,B),\qquad v=t+1
\end{equation}
と書く．最初の二分母を $m,n$ とすると，七つのケースは次表になる．
$a=m-(A-1)$，$b=n-T$ とした．
\begin{center}\small
\begin{tabular}{c|cc|ccc|c}
ケース&$m$&$n$&$v$&$a$&$b$&$B$ の当初の条件\\\hline
$A_1$&3&3&2&2&1&$B\ge2$\\
$A_2$&3&3&2&1&2&$B\ge2$\\
$A_3$&3&3&3&2&2&$B\ge2$\\\hline
$B_1$&3&7&3&1&5&$B\ge2$\\
$B_2$&3&7&4&2&5&$B\ge2$\\
$B_3$&3&7&2&1&4&$B\ge2$\\
$B_4$&3&7&3&2&4&$B\ge2$
\end{tabular}
\end{center}

\begin{lemma}[固定二分母の因子判定]\label{ug25:factor-test}
\eqref{ug25:fixed-core} を持つ正の三高整数 core が存在するなら，
ある正整数 $L,M$ に対して
\begin{equation}\label{ug25:factor-equations}
 LM=mn-1,\qquad
 \frac{mvB+m-1}{L+a},\quad
 \frac{nvB+n-1}{M+b}\quad\text{はいずれも正整数},
\end{equation}
かつ
\begin{equation}\label{ug25:linear-equation}
 CB+E=0,
 \quad
 \begin{cases}
 C=mv(M+b)+nv(L+a)-(v+1)(L+a)(M+b),\\
 E=(m-1)(M+b)+(n-1)(L+a)-(L+a)(M+b)
 \end{cases}
\end{equation}
が成り立つ．これは任意に大きい $B$ に対する必要条件である．
\end{lemma}

\begin{proof}
matching により，二つの負低行位置の間の pair は正である．
その二 entries を $L,M>0$ とすれば $LM=mn-1$ である．
当該二行の残りの高 entry は行和から $-(L+a),-(M+b)$ である．
分母 $vB+1$ の ray からの逆 entries は pair equations により
\eqref{ug25:factor-equations} の二商の負数となる．
それらの絶対値の和は最後の行和から $(v+1)B+1$ である．
分母を払うと \eqref{ug25:linear-equation} を得る．
整除条件を落としていない点が重要である．
\end{proof}

$A$ の三ケースでは $LM=8$，$B$ の四ケースでは $LM=20$ である．
全正因子を代入して得る $B=-E/C$ は次の通りである．
\begin{center}\small
\begin{tabular}{c|rrrr}
&$L=1$&$L=2$&$L=4$&$L=8$\\\hline
$A_1$&$-1/3$&$-1/3$&$\ast$&$-1/3$\\
$A_2$&$-1/3$&$\ast$&$-1/3$&$-1/3$\\
$A_3$&$-4/3$&$-2/3$&$-2/3$&$-4/3$
\end{tabular}
\end{center}
$\ast$ では $C=E=0$ だが，二商の一方が
$(6B+2)/6=B+1/3$ となり，整数性に反する．
\begin{center}\small
\begin{tabular}{c|rrrrrr}
&$L=1$&$L=2$&$L=4$&$L=5$&$L=10$&$L=20$\\\hline
$B_1$&$-12/67$&$-1/6$&0&0&$3/14$&$4/3$\\
$B_2$&$7/9$&$-3/4$&$-1/3$&$-3/11$&$\dagger$&$-3/7$\\
$B_3$&$-3/7$&$\dagger$&$3/11$&$1/3$&$3/2$&$-31/9$\\
$B_4$&$-2/3$&$-2/7$&0&$2/5$&$-2/3$&$-32/67$
\end{tabular}
\end{center}
$\dagger$ では $C=0,E\ne0$ である．従ってどの列にも正整数 $B$ はない．
この三十六場合は固定された八と二十の全正因子であり，有限範囲の数値探索ではない．

\begin{theorem}[$G_{5,14}$ の五node排除]\label{ug25:no-five}
\eqref{ug25:local}--\eqref{ug25:budget-D} を満たす標準五次接着で，
高分母 support が $G_{5,14}=K_5\setminus\{01,23\}$ ならば，
二つの単位枝の node 配分はともに $\{4,1,0\}$ である．
特にどの単位枝も $\{5,0,0\}$ を持たない．
\end{theorem}

\begin{proof}
単位枝 $01$ が五nodeを持つと仮定する．
補題~\ref{ug25:normal-form} で正規化し，core $2$ の低正位置により
補題~\ref{ug25:case-A} または補題~\ref{ug25:case-B} を適用する．
残る七つの局所型は全て補題~\ref{ug25:factor-test} と全因子表に反する．
単位枝 $23$ については二つの枝を交換すれば同じである．
一般の単位枝配分は既に二種類に限定されているので結論を得る．
\end{proof}

\begin{corollary}[残る配置の小分母]\label{ug25:small-three}
$G_{5,14}$ の各単位枝について，共通負位置の二つの隣接高分母の少なくとも
一つは三である．その外側 pair が負なら両方とも三であり，正なら
$(3,2R+1)$ または $(2R+1,3)$，$R\ge1$ である．
\end{corollary}
\begin{proof}
五nodeが排除された後の共通負位置の表を読むだけでよい．
\end{proof}

\paragraph{検証と残る問題．}
\nolinkurl{verify_g514_five_node_v25.py} は旧 verifier を import せず，
必要な共通負位置の有限係数分類，七局所型の再構成，三十六正因子の
整数係数線形式と整除障害を Python 標準ライブラリの整数・有理数で確認する．
無限尺度を除くのは式と整除障害であり，有界 scan の不在証明への外挿ではない．
本定理は $G_{5,14}$ support 自体を排除していない．
二つの $\{4,1,0\}$ 配分の同時接着，$G_{5,16}$ の五node sector，
四高 core の完全 matching 型における整数模型の分類と大域整合性は未解決である．

\section{\texorpdfstring{$G_{5,14}$}{G5,14} の中央 core の matching の固定}
\label{cm25:section}

本節では定理~\ref{ug25:no-five} の後に残る二つの $\{4,1,0\}$ 配分を
符号レベルで同時に接着する．単位枝を $01,23$，唯一の四高 core を $4$ とする．
core $i\in\{0,1,2,3\}$ の低行の唯一の正位置を $f_i$ と書く．
たとえば $f_0=2$ は $p^0_{12}>0$ を意味する．
元の条件は \eqref{ug25:local}--\eqref{ug25:budget-D} のままである．

\begin{lemma}\label{cm25:negative-common}
単位枝の共通負位置の外側 pair が負なら，node 零の位置の外側 pair は正である．
その位置の正低行側分母は少なくとも七である．
\end{lemma}
\begin{proof}
共通負位置の node は四で，両隣接高分母は三，低 column の絶対値の和は三である．
各反対符号位置で「正側 column 引く負側 column の絶対値」を $\delta_0,\delta_1$ と書く．
対向低 $D$ は $\delta_0+\delta_1=1$ を与える．
node 零で外側 pair が負ならその小例外は $\delta_0=0$ である．
node 一で外側 pair が正の場合には正通常族により
$\delta_1=1-J\le0$ となり矛盾する．
従って node 一の pair も負となるが，これは全三外側 pair が負になるため，
命題~\ref{ut24:positive-witness} に反する．
node 零で外側 pair が正なら \eqref{ug25:positive-zero} により分母は
$vV+1$，$v,V\ge2$ は異なる偶奇なので少なくとも七である．
\end{proof}

\begin{lemma}[外側正 pair が共通負位置だけの場合]\label{cm25:one-positive}
ある単位枝 $ij$ の外側 pair が共通負位置 $c$ だけで正とする．
node 零の位置を $a$，node 一の位置を $b$ とし，
node 零位置の低行が正である端点を $i$ とする．
列順 $(a,b,c)$ において，局所必要条件は
\begin{equation}\label{cm25:one-positive-normal}
 \begin{array}{c|c|c}
 &\text{低 row}&\text{低 column}\\\hline
 i&(2,-1,-2)&(2r,-k,-R)\\
 j&(-1,1,-1)&(-2r,2h,-2)
 \end{array}
 \quad
 R=2h-k>0,\quad r,h\ge1,\quad k\in\{2,4\}
\end{equation}
に限られる．従って
\begin{equation}\label{cm25:one-positive-q}
 (q_{ia},q_{ja},q_{ib},q_{jb},q_{ic},q_{jc})
 =(4r+1,2r+1,k+1,2h+1,2R+1,3).
\end{equation}
外側 cores での clutching entries は
\begin{align}
 (p^a_{ij},p^a_{ji})&=(-2r,-4r-3),\label{cm25:E0}\\
 (p^b_{ij},p^b_{ji})&=(-1,-2h(k+1)-k).\label{cm25:E3}
\end{align}
\end{lemma}

\begin{proof}
node 零の pair は負なので，反対符号の完全表から両低行は $2,-1$，
両低 column は $2r,-2r$ である．共通負位置の低行を $-s,-s'$ とすれば
$s+s'=3$ であり，行和から残りの低行は $-(3-s),3-s$ となる．
共通負 pair が正なので，低 column の絶対値は
$(s,s',S,T)=(1,2,2,R)$ または $(2,1,R,2)$ である．
対向低 $D$ により node 一位置の column 差は $R>0$ となる．

前者では node 一位置の正低行は二で，負通常族から負 column の絶対値は一である．
core $i$ の低 row/column は $(2,-2,-1),(2r,-1,-2)$ となる．
これは補題~\ref{ug25:factor-test} の $A_1$ 型であり，全正因子表と
$(6B+2)/6$ の整除障害により整数 core に完成しない．
後者では node 一位置の正低行は一である．
負例外表のうち正の column 差を持てるものは
$(U,J)=(2h,k)$，$k\in\{2,4\}$ の族だけであり，差は $R=2h-k$ である．
これが \eqref{cm25:one-positive-normal} を与える．
node 零族の $(\alpha,\beta)=(0,-3)$ を元の clutching に代入すると
\eqref{cm25:E0} を得る．node 一族では正低行側が $j$ で，
その係数は $(-k,-1)$ なので，添字を戻すと \eqref{cm25:E3} となる．
\end{proof}

\begin{theorem}[中央 core の正 pair の強制]\label{cm25:central}
$G_{5,14}$ の標準五次接着では，core $4$ の二つの pair $01,23$ が
ともに負になることはない．従って四 ray matching 定理の下で，
core $4$ の正 matching は必ず $\{01,23\}$ である．
交差する二つの matching $\{02,13\},\{03,12\}$ は起こらない．
\end{theorem}

\begin{proof}
core $4$ の pair $01,23$ がともに負と仮定する．
各単位枝には外側正 pair が少なくとも一つある．
三高 matching により，単位枝 $01$ の外側 core $2,3$ における正 pair は
それぞれ $f_2=4,f_3=4$ と同値であり，$f_2,f_3$ が同時に四になることは
単位枝 $23$ の node 非負性に反する．従ってちょうど一方が四である．
同様に $f_0,f_1$ のちょうど一方が四である．
各単位枝の端点を交換して
\begin{equation}\label{cm25:four-sign-cases}
 f_1=f_3=4,\qquad (f_0,f_2)\in\{(3,0),(3,1),(2,1),(2,0)\}
\end{equation}
とできる．この四配置で全て尽くされる．

\paragraph{三つの非相互配置．}
$(2,1)$ は二つの単位枝を交換すると $(3,0)$ になるので，
$f_0=3,f_2\in\{0,1\}$ としてよい．
単位枝 $01$ の共通負位置は $2$ であり，そこの外側 pair は負である．
従って $q_{02}=q_{12}=3$ である．唯一の外側正 pair は core $3$ にあるため，
補題~\ref{cm25:negative-common} より node 零位置は $3$ であり，
core $0$ が正低行側なので $q_{03}\ge7$ である．

$f_2=0$ なら，単位枝 $23$ の共通負位置 $1$ は外側正 pair である．
従って位置 $0$ の node は零または一で，そこの外側 pair は負である．
反対符号表における正低行側分母は $q_{02}=3$ である．
node 零の例外は正側分母 $4r+1\ge5$ なので不可能であり，
node 一の通常負族と三種類の例外はいずれも負側分母 $q_{03}\le5$ を与える．
これは $q_{03}\ge7$ に反する．
$f_2=1$ なら，単位枝 $23$ の共通負位置は $0$ で外側 pair も負である．
その node は四なので $q_{02}=q_{03}=3$ となり，再び矛盾する．

\paragraph{相互配置 $(f_0,f_2)=(2,0)$．}
この場合，単位枝 $01,23$ の唯一の外側正 pair はそれぞれ共通負位置 $3,1$ にある．
両方に補題~\ref{cm25:one-positive} を適用する．
単位枝 $01$ の node 零位置は $2$ または $4$，
単位枝 $23$ の node 零位置は $0$ または $4$ である．
共通高分母 $q_{13}$ は，前者ではそれぞれ $3$ または $2R+1\ge5$，
後者ではそれぞれ $3$ または $2R'+1\ge5$ となる．
従って混合選択は不可能であり，二つの場合だけが残る．

両 node 零位置が $2,0$ なら，単位枝 $01$ の正規形は $i=0,j=1,a=2,b=4,c=3$，
単位枝 $23$ の正規形は $i=2,j=3,a=0,b=4,c=1$ である．
$q_{02}=4r+1=4r'+1$ より $r=r'$．
式~\eqref{cm25:E0} と第二正規形の低 column により，core $2$ の ray $0$ の行は
\[
 p^2_{01}=-2r,\qquad p^2_{03}=2r,\qquad
 p^2_{04}=-q_{02}.
\]
元の pair equation
$p^2_{04}p^2_{40}=q_{02}q_{24}-1$ は
$q_{02}\mid1$ を要求し，$q_{02}>1$ に反する．

両 node 零位置が $4$ なら，正規形はそれぞれ
$(i,j,a,b,c)=(1,0,4,2,3)$ と $(3,2,4,0,1)$ である．
共通負位置の分母三から
\[
 q_{12}=k+1=3,\qquad q_{03}=k'+1=3,
 \qquad k=k'=2.
\]
第一正規形の \eqref{cm25:E3} は $p^2_{10}=-1$ を与える．
第二正規形の低 column は $p^2_{13}=-2$ を与える．
従って core $2$ の ray $1$ の行和は
\[
 -1-2+p^2_{14}=-q_{12}=-3,
\]
すなわち $p^2_{14}=0$ を強制する．
しかし $p^2_{14}p^2_{41}=q_{12}q_{24}-1>0$ であり矛盾する．
これで四符号配置と全 node 零選択が排除された．

四高 core の正 graph が完全マッチングであることを使うと，
$\{01,23\}$ 以外の二 matching は pair $01,23$ がともに負なので既に排除されている．
\end{proof}

この結論は中央 matching の二通りを除く構造還元である．
整合する中央 matching $\{01,23\}$ を持つ $G_{5,14}$ 全体の存在・非存在は未解決である．
\nolinkurl{verify_g514_central_matching_v25.py} は全符号配置の有限網羅，
正規形の原 clutching/pair 恒等式，二つの最終行和障害を独立に確認する．

% v23. Complete ordinary G5_15 obstruction; no scale/denominator cutoff.
\section{外側の pair 方程式による通常型 \texorpdfstring{$G_{5,15}$}{G5,15} の完全排除}
\label{g51523:ordinary-section}

この節では，前節までと同じ標準五次整数接着の仮定を用いる．
特に各辺の node 総数は五であり，局所行和，pair equations，正の basin weights，
奇数分母，marked clutching および対向 $D$ の相殺を課す．
支持を $K_5-\{23,24\}=G_{5,15}$ とする．
通常型とは，共有高辺 $34$ の共通低 puncture $2$ において
$p^3_{24}+p^4_{23}=0$ となる分岐である．

\begin{theorem}[通常型の無上界排除]
\label{g51523:no-ordinary}
上記の仮定の下で，$G_{5,15}$ の通常型は存在しない．
この結論には，共有分母 $Q=q_{34}$，他の分母，行列成分，尺度のいずれにも上界を仮定しない．
\end{theorem}

証明では，四高 cores $0,1$ を完成させる必要はない．
それぞれにおける一つの pair equation を，隣接する二辺の clutching と組み合わせる．
これにより，既知の四中心模型に付く正 reverse 側の二尺度を有限表へ還元する．
共有高 $D$ の非共鳴性や固定共有分母の二次復元を追加の仮定として用いないため，
すでに処理した共鳴枝や $T=1$ 境界もこの証明に含まれる．

\subsection{低辺に接する外側三角形の一次式}

\begin{lemma}[外側 pair による尺度の有限化]
\label{g51523:outer-pair}
core $2$ と core $p\in\{3,4\}$ の辺が低く，外側 core $i\in\{0,1\}$ との二辺が高いとする．
\[
 q_{2i}=a>1,\qquad p^2_{ip}=c,\qquad
 p^p_{2i}=-j,\qquad p^p_{i2}=-\sigma,
 \qquad j,\sigma\in\mathbb Z_{>0}
\]
と置く．低--高 pair により $q_{ip}=q=1+j\sigma$ である．
node 変数
\[
 u=1-n_{2i;p},\qquad v=1-n_{ip;2},\qquad
 u,v\in\{-4,-3,-2,-1,0,1\}
\]
と $E=au-c$ を用いると，外側 core $i$ の二 rays $2,p$ の pair equation は
\begin{equation}\label{g51523:outer-linear}
 \bigl\{E(jv+1)-aj\bigr\}\sigma=a-1-Ev
\end{equation}
に同値である．
特に $a-1>j$ なら，この一次式が $\sigma$ について恒等的に零となることはない．
したがって全候補は上記三十六組の $(u,v)$ から有限に復元できる．
さらに各候補は
\begin{equation}\label{g51523:outer-divisor}
 q\bigl\{jEv+E-aj\bigr\}=E-j\ne0,
 \qquad q\mid E-j,
 \qquad q\le\max_{-4\le u\le1}|au-c-j|
\end{equation}
を満たす．
\end{lemma}

\begin{proof}
二辺の clutching をそれぞれ punctures $p,2$ で読むと
\[
 p^i_{2p}=a u-c=E,\qquad
 p^i_{p2}=qv+\sigma
\]
である．core $i$ の pair equation は
\begin{equation}\label{g51523:outer-original}
 E(qv+\sigma)=aq-1
\end{equation}
となる．$q=1+j\sigma$ を代入すれば
式~\eqref{g51523:outer-linear} を得る．

一次係数と定数がともに零なら
\[
 E(jv+1)=aj,\qquad Ev=a-1.
\]
後者の $j$ 倍を前者から引いて $E=j$ を得る．
よって $v=(a-1)/j>1$ となり node 範囲に反する．
一次係数のみが零の場合は不整合であり，これも除外する．
非零の一次係数で割る場合は，得られる値が正整数で，$1+j\sigma$ が奇数であることを確認する．

最後に，式~\eqref{g51523:outer-linear} の係数を $A$，右辺を $B$ とすれば
$jB+A=E-j$ であるから，$(j\sigma+1)A=E-j$ を得る．
$E-j=0$ なら $A=0$，さらに $B=A\sigma=0$ となり，すでに排除した共鳴に戻る．
従って $E-j\ne0$ であり，式~\eqref{g51523:outer-divisor} が従う．
\end{proof}

この補題は外側 core の完全復元を仮定しない．
node 総数から得られる各変数の範囲だけを用い，異なる node 変数の残りの整合を課していない．
したがって表に残る値は必要候補であり，外側 core の存在を意味しない．
また $a-1\le j$ の別の局所問題では共鳴を再検討する必要がある．
補題の仮定を外して一般の自由尺度を除いてはならない．

\subsection{十四配置を尽くす六つの正側データ}

通常型の四中心および十四の低 row 配置は
表~\ref{g51521:center-table}，\ref{g51521:ordinary-row-table} で完全に定まっている．
共有高 ray への低 row 成分が正となる側を core $p$ とし，そのデータを
\[
 R^p=(-j_0,-j_1,m),\qquad
 (p^p_{02},p^p_{12},p^p_{\bar p2})=(-S,-T,U),
 \qquad S,T,U>0
\]
と書く．ここで $\bar p$ は $\{3,4\}\setminus\{p\}$ の元である．
低係数 $D^p_2=d$ により
\begin{equation}\label{g51523:positive-column}
 U=S+T-d-1,\qquad Q=q_{34}=1+mU
\end{equation}
である．外側補題に必要な中心データだけをまとめると，次の六行になる．
二つの負成分の交換を含めており，Schur パラメータの向きは選んでいない．

\begin{table}[htbp]
\centering\small
\begin{tabular}{c|c|c|r|c|c|r}
中心&$p$&$R^p$&$d$&$(q_{20},q_{21})$&$(p^2_{0p},p^2_{1p})$&配置数\\\hline
$M_1$&4&$(-1,-4,4)$&5&$(7,31)$&$(-2,6)$&1\\
$M_1$&3&$(-1,-1,1)$&$-9$&$(7,31)$&$(3,-10)$&1\\
$M_2$&4&$(-1,-2,2)$&2&$(9,11)$&$(-4,5)$&3\\
$M_2$&4&$(-2,-1,2)$&2&$(9,11)$&$(-4,5)$&3\\
$M_3$&4&$(-1,-2,2)$&$-7$&$(31,7)$&$(-10,2)$&3\\
$M_4$&4&$(-1,-2,2)$&9&$(19,43)$&$(-6,14)$&3
\end{tabular}
\caption{外側 pair に必要な正側データ．配置数の和は十四である．}
\label{g51523:positive-profiles}
\end{table}

全ての入力で $a\ge7$，$j\le4$ なので $a-1>j$ が成り立つ．
従って外側 pair の恒等共鳴はない．
正整数尺度と奇数分母を課して式~\eqref{g51523:outer-linear} を評価すると，
表~\ref{g51523:scale-table} を得る．各要素 $(\sigma;u,v)$ は尺度とその node 変数を表す．

\begin{table}[htbp]
\centering\small
\begin{tabular}{c|c|c|c}
中心，正側 row&$i=0$：$(S;u,v)$&$i=1$：$(T;u,v)$&帰結\\\hline
$M_1$, $(-1,-4,4)$&$(2;-1,-2)$&$(5;1,1)$&$U=1$\\
$M_1$, $(-1,-1,1)$&$(2;1,1)$&$\varnothing$&不整合\\
$M_2$, $(-1,-2,2)$&$(2;1,0)$&$(1;-1,-1)$&$U=0$\\
$M_2$, $(-2,-1,2)$&$\varnothing$&$(4;1,1)$&不整合\\
$M_3$, $(-1,-2,2)$&$\varnothing$&$(1;1,1)$&不整合\\
$M_4$, $(-1,-2,2)$&$\varnothing$&$(13;1,1)$&不整合
\end{tabular}
\caption{外側二 core の各一つの pair equation が許す全尺度候補．}
\label{g51523:scale-table}
\end{table}

この有限表は次の手順で手計算でも確認できる．
各行・各外側標識について $a,c,j$ を表~\ref{g51523:positive-profiles} から読み，
$u=-4,-3,-2,-1,0,1$ のそれぞれに対し $E=au-c$ とする．
$E=0$ は式~\eqref{g51523:outer-original} で不可能である．
$E=j$ は補題の共鳴議論で不可能である．それ以外では，
\[
 q\mid |E-j|,\quad q\ge3\text{ は奇数},\quad q\equiv1\pmod j,
 \qquad \sigma=(q-1)/j
\]
を列挙し，
\begin{equation}\label{g51523:node-from-divisor}
 v=\frac{(E-j)/q-E+aj}{jE}
 \in\{-4,-3,-2,-1,0,1\}
\end{equation}
を確認すればよい．例えば最初の行の $i=0$ では $E=7u+2$ であり，
残るのは $u=-1$，$E=-5$，$q=3$，$\sigma=2$，$v=-2$ だけである．
同じ手順を全ての入力に適用した結果が表であり，分母や尺度の任意の cutoff は使っていない．

\subsection{最後の一点の矛盾と証明の完了}

表の第二，第四，第五，第六行は一方の候補集合が空である．
第三行は式~\eqref{g51523:positive-column} により $U=0$ となり，高分母 $Q>1$ に反する．
残る第一行は
\[
 R^4=(-1,-4,4),\qquad (-S,-T,U)=(-2,-5,1),\qquad
 (q_{40},q_{41},q_{43})=(3,21,5)
\]
に限られる．

三高 matching 定理により，高 rays を $(a,b,c)=(0,1,3)$ とした core $4$ では
\[
 p^4_{ab}=H>0,\quad p^4_{ba}=K>0,\quad
 p^4_{ac}=-(H+1),\quad p^4_{ca}=-M,\quad p^4_{cb}=-N,
 \qquad M,N>0
\]
である．$p^4_{ac}=-(H+1)$ は ray $a$ の元の行和から従う．
二つの元の pair equations と ray $c$ の元の行和は
\begin{equation}\label{g51523:last-point}
 HK=62,\qquad (H+1)M=14,\qquad M+N=6
\end{equation}
を与える．$H$ は $62$ の正約数なので $H\in\{1,2,31,62\}$ である．
$H+1\mid14$ を満たすのは $H=1$ だけであり，そのとき $M=7$ となる．
これは $M,N>0$ と $M+N=6$ に反する．
全六行，従って全十四配置が排除され，定理~\ref{g51523:no-ordinary} が証明された．

\paragraph{結論の量化と検算．}
本定理は，五つの固定した $(\text{low row},D_0)$ に対する\emph{孤立した三高局所模型}を
全て分類したという主張ではない．既知の中心へ接着し，さらに外側 cores の必要 pair equations を
満たす通常型配置を全て排除するものである．
四高 core の完成や正 weights の新たな仮定を付け加えずに，その必要条件だけで矛盾した．
定理~\ref{g51521:reduction} と合わせると，この段階で残る $G_{5,15}$ は
共有分母 $Q=3,5$ の負例外 template に限られる．

検証器 \nolinkurl{verify_g5_15_ordinary_outer_pair_v23.py} は標準ライブラリのみを用い，
四中心の元の行和・pair・$D$，全 node triples から得る十四配置，二種類の独立した有限表計算，
最後の一点の整除矛盾を再確認する．
直接の三十六 node 列挙と，式~\eqref{g51523:outer-divisor} の全約数列挙の結果を照合する．
旧版の探索器は import せず，$Q$ の scan も実施しない．
これは証明を補助する厳密算術検算であり，証明支援系での形式化や外部査読ではない．

\section{\texorpdfstring{$G_{5,15}$}{G5,15} の負例外の全尺度排除}
\label{ne23:section}

共有高分母が三または五であることだけでは，三高 core の他の分母や尺度は
有限にならない．ここでは未完成の四高 cores における単一の pair equation を先に使い，
低 columns を中心模型の尺度の定数シフトへ限定する．最後は二つの三高局所方程式の
resultant を整数多項式として扱う．尺度の数値的な探索上限は用いない．

\begin{theorem}[二つの負例外 template の排除]
\label{ne23:no-negative-exception}
定理~\ref{g51521:reduction} と同じ標準五次の算術的接着条件の下で，
命題~\ref{g51521:negative-templates} の負例外は存在しない．
\end{theorem}

証明に必要な条件は，二高中心の完全尺度表示，三高 matching 性，奇数高分母，
低辺の対向 $D$，非負 node と各辺総数五，および外側 cores の一部の pair equations である．
外側四高 cores の全行列や正の weights を先に完成する必要はない．
また，最後の排除では共有高辺の対向 $D$ 条件さえ追加する必要がない．

\subsection{外側 pair から得る二つの尺度表}

以下の二つの表は，補題~\ref{oc23:linear-transfer} の整数符号を詳しく使うものである．
全ての外側 node について $u,v\in\{-4,-3,-2,-1,0,1\}$ を保つ．

\begin{lemma}[正の中心 entry を横切る尺度]
\label{ne23:positive-transfer}
$r\ge3,x\ge2$ とし，$a=rx+1$，$j=r+\nu-1$，$0\le\nu\le3$ と置く．
外側 pair が
\[
 E=au-x,\qquad
 E\{(jA+1)v+A\}=a(jA+1)-1,\qquad A\in\mathbb Z_{>0}
\]
を満たすなら，可能性は次の三つだけである：
\[
\begin{array}{c|c|c|c}
\nu&A&u&v\\\hline
0&x-1&1&1\\
1&1&-1&-1\\
2&x+1&-1&-1
\end{array}
\]
$r=2$ でも同じ表が成り立つが，$\nu=0$ に限り
$A=x,u=0,v=-3$ が追加される．
\end{lemma}

\begin{proof}
$E$ と反対 entry の符号が一致するため，$u=1$ なら $v=0,1$ である．
$v=0$ は不可能で，$v=1$ は
$A(1-\nu x)=x-1$ を与えるため表の第一行だけが残る．

残りでは $u=-h$，$v=-k$，$0\le h\le4$，$1\le k\le4$ と書ける．
尺度の方程式は
\[
 A\{a[j(hk-1)-h]+x(jk-1)\}=a(1-hk)-kx-1
\]
である．$h\ge1,hk\ge2$ では，$j\ge2$ により左辺の係数が正，右辺が負となる．
$h=k=1$ では
$A\{1+(2-\nu)x\}=x+1$ となり，表の第二・第三行を得る．
$h=0$ では
\[
 A\{xj(k-r)-x-j\}=x(r-k).
\]
$k\le r$ なら左辺の係数は負，右辺は非負である．
$k>r$ なら $x,j\ge2$ より左辺の係数は
$xj-x-j\ge0$，右辺は負となる．

$r=2$ で追加検討が必要なのは $\nu=0,j=1$ だけである．
$h=0,k=3$ が $A=x$ を与え，$k=4$ は符号で除かれる．
$h\ge1,k=1$ は
$A=h-1+(x+1)/(2x+1)$ を要求するため整数にならず，
$k\ge2$ は再び符号で除かれる．
\end{proof}

\begin{lemma}[負の中心 entry を横切る尺度と共鳴]
\label{ne23:negative-transfer}
$r\ge4,y\ge1$，$b=(r+1)y+1$ は奇数，$k=r+h$，$h\in\{-1,0,1,2\}$ とする．
外側 pair が
\[
 E=bu+y,\qquad
 E\{(kB+1)v-B\}=b(kB+1)-1,\qquad B\in\mathbb Z_{>0}
\]
を満たすなら，可能性は次の表に限られる：
\[
\begin{array}{c|c|c|c|l}
h&B&u&v&条件\\\hline
0&y-1&-1&-1&y\ge2\\
1&1&1&1&---\\
1&\text{任意の正整数}&-1&-1&y=1\\
2&y+1&1&1&---
\end{array}
\]
第三行は係数も定数も零になる共鳴であり，捨ててはならない．
また $r\ge2,y\ge2,h\in\{0,1,2\}$ に限定する場合にも，同じ非共鳴表が成り立つ．
\end{lemma}

\begin{proof}
$E>0$ なら $u=0,1$，$v=1$ であり，$u=0$ は符号で不可能である．
$u=1$ では
$B\{1+(2-h)y\}=y+1$ となり，表の第二・第四行だけが残る．
$u=-a\le-1,v=-c\le0$ とすると
\[
 B\{b[k(ac-1)+a]-y(kc+1)\}=b(1-ac)+cy-1.
\]
$ac\ge2$ では左辺の係数が正，右辺が負である．
$a=c=1$ では $B(1-hy)=y-1$ となり，第一行と $h=1,y=1$ の共鳴を得る．
$c=0$ では，正の尺度のために $a\ge k+1$ が必要である．
$r\ge4$ では $r=4,h=-1,a=4$ の場合だけを調べればよく，
$B=5y/(4y+1)$ が正整数になるのは $y=1$ のみだが，このとき $b=6$ は偶数である．

最後の追加範囲でも $c=0$ の場合だけを補えばよい．
$a=k+1$ では $(r+1)y/(ry+1)$ は一と二の間にあり，
$a\ge k+2$ では比は一未満である．従って $y\ge2$ の正整数尺度は生じない．
\end{proof}

\subsection{第一 template の四つの中心と排除}

中心の記号 $r,s,z,S,T,x,y,D_C,D_D$ は式~\eqref{g51521:central} に従う．
負枝では $S=T=w$，$r\ge4$，$x\ge2,y\ge1$ である．
第一 template は $Q=3,L=1$ であり，core $3$ の低 row と column を
\[
 (-k,k,-1),\quad (-A,B,-2),\qquad k=r+a_0-1,\quad0\le a_0\le3
\]
と書ける．低 $D$ 相殺は
\begin{equation}\label{ne23:first-low-D}
 A-B+1=-D_C=x-y
\end{equation}
である．二つの外側 pair に前二補題を適用する．
$a_0=0,3$ は不可能である．$a_0=1$ では $A=1,B=y-1$ なので $x=3$ となる．
$a_0=2$ では $A=x+1,B=y+2$ が必要であり，後者は第二補題の $y=1$ の共鳴に限られる．
共鳴の全尺度をまず保持し，式~\eqref{ne23:first-low-D} を加えている．

負枝の完全尺度表示では，$w=\ell(t+1)$ として
\[
\begin{array}{c|cc}
&x&y\\\hline
N_+&\ell t+1&wt\\
N_-&wt&\ell t-1
\end{array}
\]
である．$x=3$ は $N_+$ の $\ell t=2$ に限り，
$y=1$ は $N_-$ の $\ell t=2$ に限られる．従って四つの中心で尽くされる．
core $3$ の共有高 row を $(Y,-Y-1,-2)$ とし，その元の pair equations を整理すると，
次の原始整数多項式を得る：
\[
\begin{array}{c|cc|ccc|l}
\text{中心族}&\ell&t&k&A&B&\text{必要多項式}\\\hline
N_+&1&2&6&1&5&3Y-20\\
N_+&2&1&8&1&3&34Y-299\\
N_-&1&2&6&7&3&3Y^2-97Y+544\\
N_-&2&1&8&5&3&4Y^2-184Y+1403
\end{array}
\]
前二式は整数根を持たない．後二式の判別式は $2881$ と $11408$ であり，
それぞれ $53^2<2881<54^2$，$106^2<11408<107^2$ により整数根を持たない．
従って第一 template は全て排除された．

\subsection{第二 template の定数シフトと全尺度の parity}

第二 template では core $3$ を $J$，core $4$ を $K$ と呼ぶ．
高 ray の順は，$J$ では外部負・外部正・共有高の順 $(0,1,4)$，
$K$ では $(1,0,3)$ とする．その低 rows と columns は
\[
\begin{array}{c|c|c}
&\text{低 row}&\text{低 column}\\\hline
J&(-j,j+Q-2,-(Q-1))&(-A_J,B_J,-1)\\
K&(-k,k,-1)&(-A_K,B_K,-(Q-1))
\end{array}
\]
となる．ここで $j=r+a_0-1$，$k=S+1-b_0$ である．

\begin{lemma}[第二 template の残る全尺度形]
\label{ne23:shift-table}
第二 template が存在するなら，次のいずれかである：
\[
\begin{array}{c|c|ccc|ccc}
\text{中心}&Q&j&A_J&B_J&k&A_K&B_K\\\hline
N_\pm&3&r+\varepsilon&x+\varepsilon&y+\varepsilon&w+\eta&z+\eta&s+\eta\\
P_\pm&5&r-1&x-1&y+1&S+1&z-1&s+1
\end{array}
\]
第一行は $\varepsilon,\eta\in\{-1,1\}$ である．
表示された全ての columns は正でなければならない．
$N_\pm$ は負枝の $E$ 正負，$P_\pm$ は正 $L=3$ 枝の $F$ 正負を表す．
全四族で，さらに $\ell$ は正の奇数，$t$ は正の偶数である．
\end{lemma}

\begin{proof}
まず中心が正の $L=3$ 枝なら，$b_0=0$，$k=S+1$ である．
低辺 $24$ の二つの外側 pair は
$A_K=z-1$，$B_K=s+1$ を与える．最初は $T\ge5,z\ge2$ に対する
補題~\ref{ne23:positive-transfer} の $\nu=0$，後者は
補題~\ref{ne23:negative-transfer} の $r=S-1\ge2,y=s\ge2,h=2$ である．
低 $D$ は
\[
 A_K-B_K+Q-2=z-s+Q-4=-D_D=z-s+1
\]
なので $Q=5$ となる．$a_0=0$ であり，低辺 $23$ の外側 pair は
$A_J=x-1,B_J=y+1$ を与える．
$r=2$ の唯一の中心は $(r,s,z,S,T)=(2,2,6,3,5)$ で，$x=3$ である．
第一補題の追加候補 $A_J=x=3$ は $q_{30}=1+A_J=4$ が偶数なので除かれる．
$r=3$ は $P_-$ の共鳴 $\ell=0$，任意の $t\ge1$ であり，
この全族では $k=5,A_K=1$，従って高分母 $1+kA_K=6$ が偶数となる．
残る $r\ge4$ では表を直接適用できる．

中心が負枝の場合，$Q=5$ なら $a_0=0$ であり，
$A_J=x-1,B_J=y+1$ となる．これは低 $D$ 相殺
$A_J-B_J=x-y$ に反する．従って $Q=3$ である．
$a_0=0,2$ は第一行の $J$ 側を与える．
$a_0=1$ では $A_J=1$ である．非共鳴では $B_J=1$ なので $D_C=0$ となり，
負枝の唯一の模型 $w=2,r=4$ を要求するが，$J$ 側の高分母は六となる．
$y=1$ の共鳴を残した場合も，低 $D$ 相殺から $B_J=2-x\le0$ となる．

低辺 $24$ について，まず $w\ge3$ とする．$b_0=0,1,2$ に応じて
$A_K=z+1,1,z-1$ であり，外側のもう一つの pair は
$B_K=s+1,1,s-1$ を要求する．$b_0=1$ の低 $D$ 相殺は $s=z$，すなわち $D_D=-1$ を要求するが，
負枝の二つの完全尺度式はいずれもこの値を取らない．
$w=2$ の唯一の中心については $s=4,z=5$ である．
二つの通常シフト $(A_K,B_K)=(6,5),(4,3)$ は奇数性に反し，
$b_0=1$ では $B_K=0$ となる．第一補題の追加候補 $(5,4)$ も高分母が六となる．
従ってこの境界を含めて，残るのは第一行の $K$ 側だけである．

最後に parity を調べる．$N_+$ では $B_K=s+\eta$，$N_-$ では $A_K=z+\eta$ が奇数である．
高分母 $1+kB_K$ または $1+kA_K$ の奇数性は $k=w+\eta$ を偶数にし，$w$ を奇数にする．
$P_+$ では $B_K=s+1$，$P_-$ では $A_K=z-1$ が奇数であり，同様に
$k=S+1$ を偶数にする．それぞれ $S=3R$，$S=3R-2$ なので $R$ は奇数である．
$w=\ell(t+1)$，$R=\ell(t+1)-2$ または $\ell(t+1)+2$ より，全ての場合で
$\ell$ は奇数，$t$ は偶数となる．
\end{proof}

正 $L=4$ の七の共鳴族は，元の node 条件 $L+2\le5$ により本 template の対象外である．
正 $L=3$ の十三の共鳴族は，前証明の通り全ての自由尺度を保持してから奇数性で排除した．
固定した小分母だけから尺度有限性を主張した箇所はない．

\subsection{元の pair 残余と四十個の resultant}

低 row と column が $(-j,j+m-1,-m)$，$(-A,B,-U)$ である三高 core を考え，
$Q=mU+1$ とする．共有高 row は
$(X,-W,-U)$，$W=X+Q-U$ となる．pair 方程式で決まる正の量は
\[
 K=\frac{Q(jA+1)-1}{X},\qquad
 N=\frac{Q((j+m-1)B+1)-1}{W}.
\]
外部二 rays の row sums は，反対 entries の大きさを
$K+(j-1)A+1$ と $(j+m)B+1-N$ に決める．従って必要な元の pair 残余は
\begin{align}
\mathcal R_{j,m,U,A,B}(X)
={}&\{QjA+Q-1+X((j-1)A+1)\}\notag\\
 &\cdot\{(X+Q-U)((j+m)B+1)-Q(j+m-1)B-Q+1\}\notag\\
 &-X(X+Q-U)\{j(j+m-1)AB+jA+(j+m-1)B\}=0.
\label{ne23:residual}
\end{align}
これは $X$ の二次以下の整数多項式である．$X,W>0$ なので，払分母による必要性の喪失はない．
$U=1$ の側では，$m=Q-1$，$v=j-X$ と置けば
\[
\mathcal R=(X(X+m)-Qv(v-1))AB+XvA+(X+m)(1-v)B+X(X+m)
\]
とも書ける．この式は低 $D$ 相殺 $A-B=-D_C$ を直接使う場合にも有用である．

第二 template の共有高辺では puncture $2$ の node が二である．
残る外部二 punctures の node の和は三なので，$n=n_{34;0}\in\{0,1,2,3\}$ と置くと
\[
 X-Y=\delta_n:=Q+1-Qn
\]
となる．補題~\ref{ne23:shift-table} の各行について
\[
 f(X)=\mathcal R_J(X),\qquad g(X)=\mathcal R_K(X-\delta_n)
\]
を定義する．負中心は二族・四つの符号配置・四つの nodes で三十二場合，
正中心は二族・四つの nodes で八場合，合計四十場合である．

$f=aX^2+bX+c$，$g=dX^2+eX+f_0$ と書き，
\begin{equation}\label{ne23:resultant}
 \mathcal P=(af_0-cd)^2-(ae-bd)(bf_0-ce)
\end{equation}
と置く．共通根を持てば必ず $\mathcal P=0$ である．
実際，$P=ae-bd$，$Q_0=af_0-cd$，$W_0=bf_0-ce$ と書けば
\[
 Q_0^2-PW_0=(Q_0-PX)(ag-df)-P(bg-ef)
\]
という整数多項式恒等式がある．この必要性は $a$ や $d$ が零になる場合も変わらず，
先頭係数で割っていない．

\subsection{整数係数による全尺度の符号証明}

各中心族を次の完全な多項式表示で式~\eqref{ne23:residual} に代入する：
\[
\begin{array}{c|ccc|cccc}
&r&S&T&x&y&s&z\\\hline
N_+&2w&w&w&\ell t+1&wt&2(\ell t+1)&(2w+1)t\\
N_-&2w-1&w&w&wt&\ell t-1&(2w-1)t&2(\ell t-1)\\
P_+&2R&3R&3R+2&3(\ell t-1)&(3R+2)t&2(\ell t-1)&(2R+1)t\\
P_-&2R-1&3R-2&3R&(3R-2)t&3(\ell t+1)&(2R-1)t&2(\ell t+1)
\end{array}
\]
ここで $N_\pm$ は $w=\ell(t+1)$，$P_+$ は $R=\ell(t+1)-2$，
$P_-$ は $R=\ell(t+1)+2$ である．四十個の $\mathcal P(\ell,t)$ はこの表と
補題~\ref{ne23:shift-table}，式~\eqref{ne23:residual}--\eqref{ne23:resultant} によって一意に指定される．

全ての正の奇数 $\ell$ と正の偶数 $t$ は，次の三領域と四点で尽くされる：
\begin{align*}
\mathrm A&:\quad \ell=2v+3,\quad t=2u+4,\quad u,v\ge0,\\
\mathrm B&:\quad \ell=2v+5,\quad t=2,\quad v\ge0,\\
\mathrm C&:\quad \ell=1,\quad t=2u+8,\quad u\ge0,\\
\mathrm D&:\quad (\ell,t)=(1,2),(3,2),(1,4),(1,6).
\end{align*}
これらは整数尺度全域の分割であり，有界探索の範囲指定ではない．

\begin{lemma}[全 resultant の係数符号証明書]
\label{ne23:coefficient-certificate}
各行の四つの記号は $n=0,1,2,3$ の順である．領域 A，B，C のいずれでも，
表示した符号を掛けた $\mathcal P$ の全係数は非負，定数項は正である：
\[
\begin{array}{c|cc|c}
\text{中心族}&\varepsilon&\eta&n=0,1,2,3\text{ の符号}\\\hline
N_+&-1&-1&+,+,-,-\\
N_+&-1& 1&-,-,-,-\\
N_+& 1&-1&+,-,-,-\\
N_+& 1& 1&-,+,-,-\\
N_-&-1&-1&-,-,+,-\\
N_-&-1& 1&-,-,-,+\\
N_-& 1&-1&-,-,-,-\\
N_-& 1& 1&-,-,+,+\\
P_+&\multicolumn{2}{c|}{---}&+,+,+,+\\
P_-&\multicolumn{2}{c|}{---}&+,+,+,+
\end{array}
\]
領域 D の百六十個の評価で零となるのは
\[
 (N_-,\varepsilon,\eta,n,\ell,t)= (N_-,-1,-1,2,1,2)
\]
の一つだけである．この場合は $B_J=y-1=0$ であり，正値性に反する．
\end{lemma}

\begin{proof}
整数多項式の係数演算で確認する．再現可能な標準ライブラリ実装
\nolinkurl{verify_g5_15_negative_exception_v23.py} は，中心表から全多項式を作り，
三領域の百二十個のシフト多項式の全係数と，厳密に正の定数項を
\nolinkurl{g5_15_negative_exception_verification_v23.json} に記録する．
四点の百六十値も整数のまま記録する．例えば
$(N_+,-1,-1,n=0)$ の $(\ell,t)=(1,2)$ の値は $1458466816$，
$(P_+,n=3)$ の同じ点での値は $-958992384$ である．
後者の符号は尾部と異なるため，境界を一様な符号と取り違えない．

検査は，表示された有限個の多項式の全係数を調べるものであり，
$\ell,t$ の有限範囲を試す手順ではない．非負変数における非負係数と正の定数項は，
その領域全体で厳密な正値性を証明する．唯一の零評価については
$N_-$ の $\ell=1,t=2$ から $y=\ell t-1=1$，従って $B_J=0$ を直接確認する．
同実装は元の pair からの残余恒等式，Sylvester determinant と式~\eqref{ne23:resultant} の一致，
四族の中心方程式，parity の全剰余類，および外側尺度表の千三百六十八個の
記号的 node 場合分けも確認する．
\end{proof}

\begin{proof}[定理~\ref{ne23:no-negative-exception} の証明の完了]
第一 template は四つの必要局所多項式で排除された．
第二 template は定数シフト表と parity により，上記の四十場合と全尺度分割に含まれる．
各有効場合で $\mathcal P\ne0$ なので，二つの必要局所方程式は共通根を持たない．
唯一の零評価は正の column 条件を満たさない．従って第二 template も存在しない．
\end{proof}

本定理は内部の整数係数証明書を伴う算術的排除であり，proof assistant による形式化や
外部査読を終えたという意味ではない．指定された標準接着の結論であり，
任意の Lagrangian torus fibration の不存在定理へ拡張するものではない．

% v13: complete numerical leaf-to-degree-two transfer, with the continuation
% along the second edge imposed.  No nonnegative-node hypothesis is used.

\section{葉から二価 core を経由する係数の伝播}
\label{lt13:section}

命題~\ref{pe12:leaf-escalation} は，葉から二価 core へ進むと分母が増大することを示した．
次に，その二価 core のもう一方の高 branch を実際に第三の core へ接続する条件を用いる．
これにより，分母だけでなく係数 $D$ の値まで決まり，木の中心に向かう定量的な障害が得られる．

\begin{lemma}[負の補助変数の差は $-1$ にならない]
\label{lt13:no-negative-unit-gap}
正の二高分母標準形で $\alpha,\beta<0$ とする．このとき $\alpha-\beta=-1$ は不可能である．
\end{lemma}

\begin{proof}
$\alpha=-u$，$\beta=-v$ と置く．補題~\ref{ml12:dichotomy} により
$w=S=T$，$H=r-w\geq1$ と書ける．式~\eqref{ml12:E-equation} は
\[
 \left(1-\frac{w^2}{H(H+1)}\right)uv=v-u+1
\]
を与える．$H^2<H(H+1)<(H+1)^2$ なので左辺の括弧は零ではない．
$u=v+1$ なら右辺が零となり，$u,v>0$ に反する．
\end{proof}

\begin{proposition}[葉に接続した二高分母模型の三つの数値型]
\label{lt13:three-types}
高分母 support において，vertex $i$ の次数が一，その隣接 vertex $j$ の次数が二であるとする．
$j$ のもう一つの高隣接 vertex を $k$ とし，$q_{ij}=a>1$，$q_{jk}=b>1$ と書く．
正の局所 charts，辺 $ij$ の整数 marked clutching，および同辺の対向 $D$ の相殺を仮定する．
この段階では辺 $jk$ の clutching を仮定しない．

core $j$ の canonical な高 ray $A$ は必ず $ji$ に対応する．
canonical データと高係数は，次の三つの数値型のいずれかになる：
\begin{align*}
\textnormal{I}\quad
 &a=h(h+1)+1,\qquad b=(h+1)(h+2)+1,\qquad h\geq1,\\
 &(r,s,z,S,T)=\bigl(h,1,1,h(h+1),(h+1)(h+2)\bigr),\\
 &(D_A,D_B)=(a,-b);\tag{I}\label{lt13:type-I}\\[1mm]
\textnormal{II}\quad
 &a=2t+1,\qquad b=6t+1,\qquad t\geq1,\\
 &(r,s,z,S,T)=(1,1,1,2t,6t),\\
 &(D_A,D_B)=(3t,1-8t);\tag{II}\label{lt13:type-II}\\[1mm]
\textnormal{III}\quad
 &a=6t+1,\qquad b=20t-1,\qquad t\geq1,\\
 &(r,s,z,S,T)=(1,2,2,3t,10t-1),\\
 &(D_A,D_B)=(7t,3-18t).\tag{III}\label{lt13:type-III}
\end{align*}
型 I の $h=1$ と型 II の $t=1$ は同じ数値型 $(a,b)=(3,7)$ である．
この表は共有辺の条件から従う必要条件であり，任意の標識付き接着や大域 atlas の存在を主張しない．
node 数の非負性は用いない．
\end{proposition}

\begin{proof}
葉 $i$ の高 ray の outgoing entries は，命題~\ref{qh9:isolated-local} より
\[
 (-m,-(h+1),h),\qquad m=a-1=t h(h+1),\qquad h,t\geq1
\]
である．そのうち label $k$ へ向かう entry を $x$ とし，core $j$ で $ji$ から $jk$，
$jk$ から $ji$ へ向かう entries をそれぞれ $-\mu,-\nu$ と書く．
命題~\ref{pe12:leaf-escalation} の証明により
\[
 \mu=ca+x,\qquad \nu=ca+\frac{m}{x},\qquad
 b=c^2a+c\left(x+\frac{m}{x}\right)+1,
 \quad c=\begin{cases}1,&x>0,\\2,&x<0.\end{cases}
\]

\paragraph{二つの共通 punctures における entries．}
core $j$ の高 ray $ji$ から二つの低 rays へ向かう entries は一正一負である．
補題~\ref{ls11:low-puncture} と，二高分母側の負の entry が $-m$ にはならないという
補題~\ref{cc12:proper-negative} を用いると，それらの正負の値は
\[
\begin{array}{c|rr}
 \text{葉側の }x&\text{core }j\text{ 側の正の entry}&\text{同じ row の負の entry}\\\hline
 h&h+1&-1\\
 -m&h+1&-h\\
 -(h+1)&m&-h
\end{array}
\]
に限られる．第一行では，葉側の二つの負の entries のうち $-m$ が二高分母側の $-1$ と接着する．
残る二行で共通 punctures の符号を逆同士に合わせれば，表は直ちに従う．
同符号同士に合わせる場合には，葉側の正の entry $h$ が $m$ の真約数なので $h=1$ が必要である．
さらに負の組は葉側の $-m$ と二高分母側の $-1$ でなければならない．
$x=-m$ の行ではこれが $m=h+1=2$ を要求し，$x=-(h+1)$ の行では $-h=-1$ となる．
従って同符号の場合も同じ表に含まれる．

\paragraph{$x=h$ の場合．}
core $j$ で葉に向かう高 ray が canonical な $B$ だったと仮定すると，
その正負の entries は $z=h+1$，$-y=-1$ である．
もう一方の高分母 $b$ に対して
\[
 \alpha=\nu-b=-(h+1),\qquad \beta=y-z=-h
\]
となり，補題~\ref{lt13:no-negative-unit-gap} に反する．
従って葉に向かう ray が $A$ である．表から $s=1$，canonical な正の entry は $h+1$ なので
\[
 r=t h,\quad S=t h(h+1),\quad
 b-1=(h+1)\{t(h+1)+1\}.
\]
また $\nu=b-y+z$ と $y=(b-1)/(r+1)$ から
\[
 z=\frac{t(h+1)}{t h+1}
   =1+\frac{t-1}{t h+1}.
\]
右辺は $[1,2)$ に属するため，整数性から $t=1$，$z=1$ となり，型 I を得る．

\paragraph{$x=-m$ の場合．}
この場合は $b=2a+1$，$\mu=a+1$，$\nu=2a-1$ である．
葉に向かう ray が $B$ なら $\alpha=\nu-b=-2$，$\beta=h-(h+1)=-1$ なので，
再び補題~\ref{lt13:no-negative-unit-gap} に反する．
従って $A$ であり，表から
\[
 r=t h,\quad s=h,\quad S=t(h+1),\quad
 z=\frac{2t h^2}{t h+1}=2h-\frac{2h}{t h+1}
\]
となる．整数性は $t h+1\mid2h$ を要求する．$t h+1>h$ なので $t h+1=2h$，
従って $(2-t)h=1$，すなわち $h=t=1$ である．これは型 I の $h=1$ に含まれる．

\paragraph{$x=-(h+1)$ の場合．}
表の正の entry は最大約数 $m$ なので，補題~\ref{cf13:full-entry} により葉に向かう ray は $A$ である．
従って $r=1$，$s=h$ となる．ここでは
\[
 \nu=2a-t h,\qquad b=4a-2(h+1+t h)+1
\]
であるから，$y=(b-1)/2$ と $\nu=b-y+z$ より $z=h$ を得る．
さらに
\[
 T=\frac{b-1}{h}=2t(2h+1)-2+\frac2h
\]
の整数性から $h\mid2$，すなわち $h=1,2$ だけが残る．
$h=1$ が型 II，$h=2$ が型 III である．

最後に，標準形の係数公式
$D_A=-a+B-r+S$，$D_B=-b+A+r+1-T$ に各データを代入すると，表示した $D$ を得る．
以上は符号，整除性，row sums と共有辺 $ij$ の $D$ の相殺だけを用いている．
\end{proof}

\begin{theorem}[第三の core へ接続したときの完全な係数伝播]
\label{lt13:exact-transfer}
命題~\ref{lt13:three-types} の設定で，辺 $jk$ の marked puncture にも整数 clutching を仮定する．
このとき，ある整数 $h\geq1$ に対して
\[
 a=h(h+1)+1,\qquad b=(h+1)(h+2)+1,
\]
であり，特に $b-a=2(h+1)$ および $\gcd(a,b)=1$ が成立する．
core $j$ の係数は
\[
 D^j_i=a,\qquad D^j_k=-b
\]
となる．さらに core $k$ の ray $kj$ から低 ray $ki$ への entry は
\[
 p^k_{ji}=-(h+1),\qquad p^k_{ij}=-(h+2)
\]
に強制される．core $k$ の高分母次数や node 数の非負性は仮定しない．
\end{theorem}

\begin{proof}
$i$ は高分母 support の葉なので $q_{ki}=1$ である．
$u=p^k_{ji}$ と置くと，pair product と辺 $jk$ の puncture $i$ の clutching により
\begin{equation}\label{lt13:continuation-congruence}
 u\mid b-1,\qquad 0<|u|\leq b-1,\qquad u\equiv\nu\pmod b
\end{equation}
を得る．三つの型のいずれでも $0<\nu<b$ なので，可能な整数は $u=\nu$ または $u=\nu-b$ だけである．

型 II では $\nu=3t+2$，$b=6t+1$ である．
$u=3t+2$ が $6t$ を割れば $3t+2\mid4$ となって不可能である．
$u=-(3t-1)$ が $6t$ を割る場合には $3t-1\mid2$ なので $t=1$ となる．
これは型 I の $h=1$ と同じ数値型である．

型 III では $\nu=10t+2$，$b=20t-1$ である．
$u=10t+2$ の整除性は $10t+2\mid6$ を，$u=-(10t-3)$ の整除性は
$10t-3\mid4$ を要求する．いずれも $t\geq1$ では不可能なので，型 III は全て排除される．

従って型 I だけが残り，その係数は既に $(a,-b)$ である．
この場合 $\nu=(h+1)^2+1$，$b-1=(h+1)(h+2)$ である．
$u=\nu$ の整除性は $\nu\mid h$ を要求するが，$\nu>h$ なので不可能である．
従って $u=\nu-b=-(h+1)$ である．その reverse entry は
$(b-1)/u=-(h+2)$ となる．
また $a$ は奇数で，$a\equiv1\pmod{h+1}$ だから
\[
 \gcd(a,b)=\gcd\bigl(a,2(h+1)\bigr)=1
\]
である．
\end{proof}

この結論は残る数値型を指定するものであり，葉側の entry $x$ の符号を指定するものではない．
特に $(a,b)=(3,7)$ は，前命題の型 I と型 II の両方の導出に現れる．
型 III では $D_B=3-18t>-b$ であり，局所的な不等式 $D_B\leq-b$ そのものは偽である．
この型を除くのは，第三の core への整数 clutching である．
一般の二高分母模型に係数の大きさの伝播を拡張したり，第三の core の条件を省略したりしてはならない．

% v13: the forest obstruction uses quantitative leaf transfer, not an
% unproved sign assertion for an arbitrary three-high local model.
\subsection{高分母 support の木成分は存在しない}
\label{fo13:section}

本項の議論では，三高分母模型の全ての高係数が正となり得るかという
一般の局所問題を解く必要はない．葉に接する辺が要求する係数の
\emph{大きさ}と，分母一の row が与える整数上界を比較する．

\begin{lemma}[三高分母の係数和の上界]\label{fo13:center-bound}
分母が $(a,b,c,1)$，$a,b,c>1$ である整数局所模型において，
row sums と pair equations が成立すると仮定する．
このとき
\begin{equation}\label{fo13:center-inequality}
 D_a+D_b+D_c
 \leq a+b+c-\min\{a,b,c\}-2.
\end{equation}
正の weights，node の条件，高分母間の numerators の符号は仮定しない．
特に高分母が奇数なら，右辺は $a+b+c-5$ 以下である．
\end{lemma}

\begin{proof}
分母一の頂点を $0$ とし，$r_i=p_{0i}$，$x_i=p_{i0}$ と置く．
row sum と pair equations は
\[
 \sum_{i=a,b,c}r_i=-1,\qquad r_ix_i=q_i-1>0
\]
を与える．従って，ある $j$ について $r_j<0$，従って $x_j\leq-1$ である．
他の二つの添字については，$x_i$ が正なら $q_i-1$ の正の約数であり，
負なら同じ上界が自明に成り立つため，いずれも $x_i\leq q_i-1$ である．
全ての $D$ の和は零なので
\[
 D_a+D_b+D_c=-D_0=1+x_a+x_b+x_c
 \leq\sum_{i\ne j}q_i-2
 \leq a+b+c-\min\{a,b,c\}-2.
\]
これで示された．
\end{proof}

\begin{theorem}[木成分の排除]\label{fo13:no-tree}
全平面五次の既存の $K_5$ branch graph と標準 $K_4$ cones を保ち，
正の局所 charts，marked puncture の整数 clutching，
および高辺の対向端における $D$ の相殺が成立するとする．
高分母 support $H=\{ij:q_{ij}>1\}$ は，辺を持つ木の連結成分を持たない．
従って，$H$ が forest なら全ての分母は一である．
この結論に node 数の非負性は不要である．
\end{theorem}

\begin{proof}
木の連結成分は高々五頂点を持つ．全ての次数が二以下なら path なので，
定理~\ref{pc12:no-path} によって排除される．
残る形は，三本の葉を持つ star $K_{1,3}$，四本の葉を持つ star $K_{1,4}$，
または $K_{1,3}$ の一辺を一度だけ細分した五頂点の木である．

葉の唯一の高分母を $q$ とする．命題~\ref{qh9:isolated-local} によれば，
ある正整数 $h,t$ により $q-1=t h(h+1)$ かつ
\[
 D_{\mathrm{leaf}}=-(q-1)-t\leq-q.
\]
従って，その反対側の端では $D\geq q$ が必要である．
$K_{1,3}$ の中心における三つの分母を $a,b,c$ とすれば，
三つの葉が $D_a+D_b+D_c\geq a+b+c$ を要求する．
これは補題~\ref{fo13:center-bound} に反する．
$K_{1,4}$ では四つの係数全てが正となり，その和が零であることに反する．

最後に細分された star を考える．長さ二の枝を，葉 $i$，二価頂点 $j$，
三価中心 $k$ の順に書く．$j$ でのもう一方の高辺 $jk$ は実際に $k$ へ接着するので，
継続する葉辺の係数移送定理~\ref{lt13:exact-transfer} を適用できる．
従って $jk$ の $j$ 側の係数は $-q_{jk}$ 以下であり，
$k$ 側では $q_{jk}$ 以上である．残る二つの枝も葉に直結するので，
中心の三つの高係数はいずれも対応する分母以上となる．
再び補題~\ref{fo13:center-bound} に矛盾する．
\end{proof}

\begin{remark}[局所符号問題との区別]\label{fo13:scope}
定理~\ref{fo13:no-tree} は，一般の整数三高分母模型について
$D_a,D_b,D_c>0$ が不可能であるとは述べない．
三つ全ての正性の排除は，独立した局所問題として残る．
ここで用いたのは，葉とその接着条件が強制する
$D_i\geq q_i$ という，より強い下界である．
また，結論は木の\emph{連結成分}に関するものであり，
cycle を含む成分の内部にある木の部分 graph や pendant tree を全て排除するものではない．
\end{remark}

\begin{corollary}[未排除の support の大域的な形]\label{fo13:remaining-support}
定理~\ref{fo13:no-tree} の仮定に node 数の非負性を加える．
$H$ が空でなければ，辺を持つ連結成分はちょうど一つであり，
その成分は四または五頂点を持ち，cycle を含み，最大次数は三以上である．
\end{corollary}

\begin{proof}
辺を持つ成分は木ではないので，cycle を含む．
三頂点のそのような成分は $C_3$ だけだが，
定理~\ref{ls11:no-triangle-component} により排除される．
従ってその成分は四頂点以上を使い，$K_5$ 内に二つは存在しない．
最大次数が二以下なら $C_4$ または $C_5$ となるが，
系~\ref{cf13:degree-two-rigidity} に反する．
\end{proof}

% v23: complete G5_15 obstruction leaves three dense support types.
\subsection{三型への帰着と同じ core に接する二本の単位枝}
\label{sf23:section}\label{sf22:section}\label{sf21:section}\label{sf20:section}

\begin{corollary}[非空支持の最小次数と三点連結性]
\label{sf23:no-leaves}\label{sf22:no-leaves}\label{sf21:no-leaves}\label{sf20:no-leaves}
正の局所整数 charts，marked puncture の整数 clutching，対向 $D$ 相殺，
および非負 node と各辺の node 総数五を満たす標準五次 $K_5$ 接着データを考える．
高分母支持 $H=\{ij:q_{ij}>1\}$ が非空なら，五頂点全てを使い，
\[
 \min\deg(H)\ge3,\qquad |E(H)|\ge8,\qquad b_1(H)\ge4
\]
が成り立つ．さらに $H$ は三点連結であり，二つ以下の任意の頂点を削除しても連結である．
補 graph $K_5\setminus H$ は高々二辺の matching である．
\end{corollary}

\begin{proof}
全低 core の伝播，木と閉路の障害，二価二頂点の三角形，葉付き四角形，
$K_{2,3}$，diamond，二葉三角形，四使用頂点 $K_4$ の障害を適用する．
さらに定理~\ref{hs18:no-house}，\ref{gten18:no-support}，
\ref{lh18:no-K4-leaf}，\ref{g719:no-G5-7}，\ref{dl19:no-G5-9}，
\ref{g51120:no-support} により，必要候補は従前の五型
$G_{5,13},G_{5,14},G_{5,15},G_{5,16},G_{5,17}$ に限られる．
Book $G_{5,13}$ は定理~\ref{bc22:book-obstruction} により排除される．
$G_{5,15}$ は通常型の定理~\ref{g51523:no-ordinary}，
負例外の定理~\ref{ne23:no-negative-exception}，および既存の正例外排除
（定理~\ref{g51521:reduction}）を合わせて排除される．
残る三型は，$K_5$ から二辺以下の matching を除いたものである．
従って各頂点の次数は三以上であり，辺数と $b_1$ の主張も従う．

三点連結性を確認する．任意の二頂点を削除した後には，三頂点上の完全 graph から
高々一辺を除いた graph が残る．matching の二辺を三頂点上に同時に含めることはできないためである．
この graph は連結である．一頂点または零頂点の削除でも連結性は保たれる．
\end{proof}

\begin{corollary}[三型・二十六標識]
\label{sf23:remaining-support}\label{sf22:remaining-support}\label{sf21:remaining-support}\label{sf20:remaining-support}
前系の仮定の下で，非空高分母支持の必要候補は次の三型に限られる．
補 graph $\overline H=K_5\setminus H$ の表示では孤立頂点を省略する．
\begin{center}\small
\begin{tabular}{crcrcr}
型&$|E|$&降順次数列&$b_1$&$\overline H$&標識付き個数\\\hline
$G_{5,14}$&8&$(4,3,3,3,3)$&4&$2K_2$&15\\
$G_{5,16}$&9&$(4,4,4,3,3)$&5&$K_2$&10\\
$G_{5,17}$&10&$(4,4,4,4,4)$&6&$\varnothing$&1
\end{tabular}
\end{center}
識別子は従前のものを保ち，欠番を詰めない．これらは必要条件を通過した
graph の個数であり，整数接着，作用 atlas，幾何的非同型類の存在数ではない．
\end{corollary}

\begin{proof}
二辺の matching は，孤立する補頂点を五通りに選び，残る四頂点の三つの完全 matching を選ぶため
$5\cdot3=15$ 標識となる．一辺の matching は $\binom52=10$ 通り，空 matching は一通りであり，
和は $15+10+1=26$ である．
\nolinkurl{verify_support_agenda_v23.py} は，全 $2^{10}=1024$ 標識付き部分 graph を
旧 graph 実装を import せずに列挙する．現行三型の自己同型群，最小次数，
任意の二頂点削除後の連結性を確認し，第22版の五十六標識との差分が
$G_{5,15}$ の三十標識ちょうどであることを記録する．
算術障害は前系の証明を用い，graph 検査で算術証明を置き換えない．
\end{proof}

\begin{corollary}[隣接する二本の単位分母からの伝播]
\label{sf23:adjacent-unit-branches}\label{sf22:three-unit-branches}
同じ仮定の下で，一つの core に接する相異なる二本の branches がともに分母一なら，
全ての分母は一である．同値に，一つの core の高分母 ray が高々二本なら全分母一である．
また，十本の相異なる branches のうち三本以上が分母一なら，全分母一である．
\end{corollary}

\begin{proof}
非空高分母支持の最小次数は三以上である．一つの core に二本の単位枝が接すれば
その高次数は二以下となり，矛盾する．三本の単位枝があれば，五頂点上では
その中の二本が必ず頂点を共有するため，同じ結論になる．
二本の互いに素な単位枝だけでは，残る $G_{5,14}$ を排除していない．
\end{proof}

\paragraph{残る研究課題．}
Book と $G_{5,15}$ は解決済みであり，専用の大分母探索や四高 core の完成は課題から外す．
第三 core の単一 pair による尺度の線形復元と共鳴分類を，
二本の互いに素な低辺を持つ $G_{5,14}$ へ移すことが具体的な入口となる．
四高 core に三高 matching 定理を適用せず，必要な pair と node 条件から先に比較する．
三高係数の同時非負性には新たな無上界排除領域が得られたが，
一般の非対称問題は残る．三型の整数接着と幾何的実現，一般の五次分母剛性は未解決である．
指定した cores，cones，marked punctures の接着問題の結論を，
任意の Lagrangian torus fibration の不存在へ拡張してはいない．

% v11: general integer proof of the equal-two-high-denominator defect sign.
% The finite table below follows from unbounded inequalities; it is not a
% denominator-bounded search. Uses the general sign lemma ls11:two-high-signs.

\subsection{等二高分母の係数和の負性}
\label{uh11:two-section}

等二高分母の場合に残っていた符号問題は，整除条件を保った次の整数不等式で解決する．

\begin{proposition}[等二高分母の整数不等式]\label{uh11:two-integer-inequality}
$r,t\geq1$ を整数，$m=tr(r+1)$ とし，$s,z$ を $m$ の正の divisors とする．
\begin{equation}\label{uh11:conic}
 sz+(1+tr^2)s-\{1+t(r+1)^2\}z=t+1
\end{equation}
が成立すれば，
\begin{equation}\label{uh11:strict-inequality}
 (s-z)\left(1+\frac{m}{sz}\right)>t+1
\end{equation}
である．$r,t,m$ の上限は仮定しない．
\end{proposition}

\begin{proof}
$R=r(r+1)$，$n=2r+1$，$q=m+1$ と置き，
\[
 S=\frac ms,\qquad Z=\frac mz,\qquad H=Z-S,\qquad \delta=s-z
\]
とする．以下の全ての補助変数の整数性を保つことが重要である．

\paragraph{差の正値性と上限．}
まず $\delta>0$ を示す．式~\eqref{uh11:conic} は
\[
 \left(q+\frac mr-s\right)
 \left(q-\frac{m}{r+1}+z\right)=q^2-1
\]
と同値である．$s=m/S$，$z=m/Z$ を代入して分母を払うと
\begin{equation}\label{uh11:reciprocal-identity}
 m^2\{nS-r^2H-R\}+m(SZ-RH)+RSZ=0
\end{equation}
を得る．もし $H\leq-1$ なら
\[
 nS-r^2H-R\geq n+r^2-R=r+1>0
\]
であり，\eqref{uh11:reciprocal-identity} の左辺の全項が正となって矛盾する．
$H=0$ なら $s=z$ であり，\eqref{uh11:conic} は
\[
 s\{s-t(2r+1)\}=t+1
\]
となる．これは同時に $s\mid t+1$ と $s\geq t(2r+1)+1>t+1$ を要求するので不可能である．
従って $H>0$，従って $\delta>0$ である．

上の二因子を $A,B$ と書くと，$s\leq m$ より $A>0$，また $B>0$ であり，
$AB=q^2-1$ である．$A,B$ は整数なので $A+B\geq2q$ である．
一方 $A+B=2q+t-\delta$ だから
\begin{equation}\label{uh11:delta-range}
 1\leq\delta\leq t,\qquad e=t-\delta\geq0.
\end{equation}
目標~\eqref{uh11:strict-inequality} は
\begin{equation}\label{uh11:eta}
 \eta:=\delta+H-t-1=H-e-1>0
\end{equation}
と同値である．

\paragraph{二つの整数 $G,J$．}
式~\eqref{uh11:reciprocal-identity} と $m=tR$ から
\begin{equation}\label{uh11:G}
 G:=r^2H+R-nS=\frac{R(e+1)}{sz}>0
\end{equation}
を得る．左辺から $G$ は整数である．さらに
\[
 J:=H-\delta G
\]
と置くと，$H=m\delta/(sz)$ より
\begin{equation}\label{uh11:J-identities}
 tJ=(\delta-1)H,\qquad Jsz=\delta(\delta-1)R.
\end{equation}
また $\delta S=zH$ と $H=\delta G+J$ から
\begin{equation}\label{uh11:J-divisibility}
 \delta\mid Jz
\end{equation}
が従う．元の式は
\begin{equation}\label{uh11:K}
 tK=sz+\delta-1,
 \qquad K=nz-r^2\delta+1>0
\end{equation}
とも書ける．$K$ は正の整数なので $nz\geq r^2\delta$ である．

\paragraph{$\delta=1$ の場合．}
この場合は $s=z+1$，$sz=tK$，$H=R/K$ である．$H,K$ が正の整数なので $K\leq R$，従って
\[
 nz-r^2+1\leq R,\qquad nz\leq nr-1.
\]
よって $z<r$ であり，$s=z+1\leq r$ から $sz<R$ となる．従って
$H/t=R/(sz)>1$，すなわち $\eta=H-t>0$ である．

\paragraph{$\delta\geq2$ の場合の $J$ の制限．}
式~\eqref{uh11:J-identities} より $J>0$ である．$x=z/\delta$ と置くと
\[
 x\geq\frac{r^2}{n},\qquad
 Jx(x+1)=\frac{\delta-1}{\delta}R<R.
\]
一方
\[
 x(x+1)\geq\frac{r^2(r+1)^2}{n^2}=\frac{R^2}{n^2}
\]
なので，$J<n^2/R=4+1/R$ であり，$J\leq4$ である．
$J=4$ なら，\eqref{uh11:J-divisibility} により $4x$ は整数である．
\[
 x\geq\frac{r^2}{2r+1}>\frac r2-\frac14
\]
より $4x\geq2r$，すなわち $x\geq r/2$ となる．これは
$4x(x+1)\geq r(r+2)>R$ を与えて矛盾する．従って
\begin{equation}\label{uh11:J-small}
 J\in\{1,2,3\}.
\end{equation}

もし $J\geq\delta-1$ なら，\eqref{uh11:J-identities} より $t\leq H$ であり，
$\eta\geq\delta-1>0$ が従う．以下では残る可能性
\begin{equation}\label{uh11:bad-range}
 1\leq J\leq\delta-2
\end{equation}
が存在しないことを示す．

\paragraph{残存候補の一般的な係数評価．}
式~\eqref{uh11:G} と $H=\delta G+J$ は
\[
 nS=(\delta r^2-1)G+(J+1)r^2+r
\]
を与える．また
\[
 SZ=\frac{mH}{\delta}
   =\frac{(\delta-1)R(\delta G+J)^2}{\delta J}.
\]
前式と $Z=S+H$ を後式に代入して展開すると
\begin{equation}\label{uh11:polynomial-elimination}
 R^2L-RM-N=0
\end{equation}
となる．ここで
\begin{align*}
 L={}&\delta^2(\delta J-4\delta+4)G^2
      +\delta J(2\delta J-6\delta+8)G
      +J\{\delta(J-1)^2+4J\},\\
 M={}&\delta^2(2J+\delta-1)G^2
      +\delta J(\delta+2J)G-J^2,\\
 N={}&\delta J(\delta-1)G^2+\delta J^2G.
\end{align*}
$G\geq1$ と $\delta\geq J+2$ から $M,N>0$ なので，必要条件は $L>0$ である．

$J=1$，$\delta\geq3$ の場合は
\[
 L=\delta^2(4-3\delta)G^2+\delta(8-4\delta)G+4
   \leq-3\delta^3+8\delta+4<0.
\]
$J=2$，$\delta\geq4$ の場合も
\[
 L=\delta^2(4-2\delta)G^2+2\delta(8-2\delta)G+2\delta+16
   \leq-2\delta^3+18\delta+16<0.
\]
いずれも矛盾する．$J=3$，$\delta\geq5$ では
\[
 L=\delta^2(4-\delta)G^2+24\delta G+12\delta+36.
\]
$\delta\geq9$ ならこれは $G\geq1$ で減少し，
$L(1)=-\delta^3+4\delta^2+36\delta+36<0$ となる．
残る $\delta=5,6,7,8$ における必要条件 $L>0$ は，それぞれ
\[
 G\leq5,\qquad G\leq2,\qquad G\leq1,\qquad G\leq1
\]
を与える．さらに
\[
 t=\frac{(\delta-1)(\delta G+3)}{3}\in\mathbb Z
\]
を課すと，残る組は
\[
 (\delta,G)=(5,3),(6,1),(6,2),(7,1)
\]
だけである．

\paragraph{最後の四つの二次方程式．}
これらを \eqref{uh11:polynomial-elimination} に代入し，共通の正の整数因子を除くと，
次表の二次方程式が $R=r(r+1)$ に必要となる．
\begin{center}
\small
\begin{tabular}{ccrl}
$\delta$ & $G$ & 必要な二次式 & 判別式の厳密な区間 \\
\hline
$5$ & $3$ & $77R^2-912R-225=0$ & $949^2<901044<950^2$ \\
$6$ & $1$ & $20R^2-67R-16=0$ & $75^2<5769<76^2$ \\
$6$ & $2$ & $12R^2-223R-52=0$ & $228^2<52225<229^2$ \\
$7$ & $1$ & $47R^2-284R-63=0$ & $304^2<92500<305^2$
\end{tabular}
\end{center}
全ての判別式が非平方なので，いずれも整数 $R$ を持たない．
従って \eqref{uh11:bad-range} は排除される．
$J\geq\delta-1$ に戻れば $\eta>0$ が従い，証明が完了する．
\end{proof}

この証明の有限表は，分母の上限を置いた探索の結果ではない．
任意の整数解を補助整数 $J\in\{1,2,3\}$ と一般の係数評価によって四候補へ帰着したものである．

\begin{corollary}[等二高分母 core の負の総和]\label{uh11:two-negative}
正の局所四直線 chart の分母が $(q,q,1,1)$，$q>1$，であるとする．
二つの高分母頂点を $A,B$ と書けば
\[
 D_A+D_B<0.
\]
\end{corollary}

\begin{proof}
二高分母の表示と符号補題~\ref{ls11:two-high-signs} により，標識交換後に
$m=q-1=tr(r+1)$，$u=-s<0<v=z$ と書け，$s,z$ は $m$ の正の divisors で
\eqref{uh11:conic} を満たす．分母一の頂点を $C,D$ とすると
\[
 D_C=\frac mz-\frac ms-t,\qquad D_D=s-z-1.
\]
従って $D_C+D_D$ は \eqref{uh11:strict-inequality} の左辺から $t+1$ を引いた値であり，正である．
四頂点の係数総和は零なので，$D_A+D_B=-(D_C+D_D)<0$ となる．
\end{proof}

% Complete classification of the positive equal-two-high numerator data.
% Depends on uniform_two_high_negativity_v11.tex; no geometric moduli count.

\subsection{等二高分母の整数模型の完全分類}
\label{uc11:classification-section}

前の証明の補助整数は，符号だけでなく整数模型そのものを分類する．

\begin{theorem}[等二高分母の全 numerator data]\label{uc11:classification}
正の局所四直線 chart の分母を $(q,q,1,1)$，$q>1$ とする．
二高分母の標準表示において $u=-s<0<v=z$，$q-1=tr(r+1)$ と書く．
標識の交換を除いて，可能な整数データ $(r,t,s,z)$ は次表で尽きる．
最終行では $t\geq1$ は任意の整数である．
\begin{center}
\small
\begin{tabular}{ccccccl}
$r$ & $t$ & $s$ & $z$ & $q$ & $s-z$ & $D_C+D_D$ \\
\hline
$2$ & $1$ & $2$ & $1$ & $7$ & $1$ & $2$ \\
$1$ & $6$ & $4$ & $1$ & $13$ & $3$ & $5$ \\
$2$ & $5$ & $6$ & $3$ & $31$ & $3$ & $2$ \\
$6t-3$ & $t$ & $3t$ & $3t-1$ & $t(6t-3)(6t-2)+1$ & $1$ & $3t-2$
\end{tabular}
\end{center}
ここで $C,D$ は分母一の頂点である．各行は正の weights を持ち，局所の幾何学的な chart を実現する．
従って可能な $q$ の全体は
\[
 \{7,13,31\}\ \cup\
 \{t(6t-3)(6t-2)+1:t\geq1\}
\]
である．$q=13$ では表の第二行と最終行の $t=1$ が異なる整数データを与える．
この分類は整数 numerator systems の分類であり，weights の変形や幾何学的同型類の個数を主張しない．
\end{theorem}

\begin{proof}
命題~\ref{uh11:two-integer-inequality} の証明の記号
$\delta=s-z$，$R=r(r+1)$，$n=2r+1$，$H=Z-S$，$G,J,K$ を用いる．
$\delta\geq2$ の場合，既に
\[
 \delta-1\leq J\leq3
\]
を証明したので，$\delta\leq4$ である．また同じ証明で用いた不等式を少し正確に書けば
\begin{equation}\label{uc11:sharper-J}
 J\leq\frac{\delta-1}{\delta}\left(4+\frac1R\right)
\end{equation}
となる．これは $x=z/\delta\geq r^2/n$ と
$Jx(x+1)=(\delta-1)R/\delta$ から直接従う．

\paragraph{$\delta=2,4$ の排除．}
$\delta=2$ なら \eqref{uc11:sharper-J} より $J\leq2$ である．
$J=1$ の場合は，\eqref{uh11:polynomial-elimination} の必要条件 $L>0$ に対し
$L=-8G^2+4<0$ となり矛盾する．
$J=2$ の場合は $sz=R$，$s=z+2$ だから
\[
 (z+1)^2=r^2+r+1
\]
となる．右辺は $r^2$ と $(r+1)^2$ の間に厳密にあり，平方ではない．

$\delta=4$ なら $J=3$ であり，$\delta\mid Jz$ から $z=4k$，$s=4(k+1)$ と書ける．
$Jsz=\delta(\delta-1)R$ より $R=4k(k+1)$ となるので
\[
 (2r+1)^2=4(2k+1)^2-3.
\]
従って
\[
 \{2(2k+1)-(2r+1)\}\{2(2k+1)+(2r+1)\}=3.
\]
左の二因子は正の整数なので一と三でなければならないが，その和は
$4(2k+1)\geq12$ であり矛盾する．

\paragraph{$\delta=3$ の二つの模型．}
$J=2$ または $3$ である．$J=2$ の場合には
\[
 L=-18G^2+12G+22>0
\]
から $G=1$ だけが残る．式~\eqref{uh11:polynomial-elimination} は
\[
 4R^2-23R-6=(R-6)(4R+1)=0
\]
となるので $R=6$，$r=2$ である．
$sz=18$ と $s-z=3$ から $(s,z)=(6,3)$，
$t=(\delta-1)(\delta G+J)/J$ から $t=5$ となる．

$J=3$ の場合には \eqref{uc11:sharper-J} が $R\leq2$ を与えるので $r=1$，$R=2$ である．
$sz=4$ と $s-z=3$ から $(s,z)=(4,1)$ となり，
$tK=sz+\delta-1$ と $K=1$ から $t=6$ となる．

\paragraph{$\delta=1$ の無限族と一つの例外．}
この場合は $s=z+1$，$R=HK$，$sz=tK$ であり，先の証明から $z<r$ である．
$a=r-z\geq1$ と置くと
\[
 K=R-na+1,\qquad a^2=(t-1)K+1,\qquad na=(H-1)K+1.
\]
最後の式を二乗し，$n^2=4HK+1$ と二番目の式を代入して整理すると
\begin{equation}\label{uc11:LK}
 LK=t+2H+1,\qquad L=(H+1)^2-4Ht>0
\end{equation}
となる．$L=H(H-4t+2)+1>0$ と $H\geq1$ から $H\geq4t-2$ である．
そこで $H=4t-2+k$，$k\geq0$ と書く．
\eqref{uc11:LK} は
\[
 K=\frac{9t-3+2k}{1+(4t-2)k+k^2}
\]
を与える．$k\geq3$ なら分母と分子の差は
\[
 k^2+(4t-4)k-9t+4\geq3t+1>0
\]
となり，正の整数 $K$ は存在しない．
$k=2$ なら $K=(9t+1)/(8t+1)$ は一と二の間にあり，やはり整数ではない．
$k=1$ なら $K=2+(t-1)/(4t)$ なので $t=1$ だけが可能であり，
$H=3$，$K=2$，$R=6$ から $(r,t,s,z)=(2,1,2,1)$ を得る．

残る $k=0$ の場合は
\[
 H=4t-2,\qquad K=9t-3,\qquad
 R=HK=(6t-3)(6t-2).
\]
$r>0$ なので $r=6t-3$ であり，
$a^2=(t-1)K+1=(3t-2)^2$ と $a>0$ から $a=3t-2$ となる．
従って $z=r-a=3t-1$，$s=3t$ を得る．

最後に，各行は \eqref{uh11:conic} と $s,z\mid tr(r+1)$ を満たす．
無限族では $s=3t$ の整除性は $3\mid6t-3$，$z=3t-1$ の整除性は
$6t-2=2(3t-1)$ から従い，conic は多項式恒等式である．
$u=-s<0$ の標準形なので正の weights が存在し，局所 chart の十分条件を満たす．
実際，$x=m/r$，$Z=m/z$ と書けば，全ての行で
\[
 (w_A,w_B,w_C,w_D)
 =\left(q,q,x+\frac1{2r},z+\frac1{2Z}\right)
\]
と取れる．正の numerators は $A\leftrightarrow C$ と $B\leftrightarrow D$ のみであり，
対応する四つの basin areas は $1/(2r),1/2,1/(2Z),1/2$ である．
$D_C+D_D=\delta+H-t-1$ を代入すると表の最終列を得る．
\end{proof}

無限族の高分母二頂点の係数は，標準表示の順で
\[
 D_A=2(3t-2)(t-1),\qquad
 D_B=-(3t-2)(2t-1)
\]
である．従って $t\geq2$ では一方が正である一方，その総和は常に $-(3t-2)$ である．

% v10. General arithmetic sign theorem for q-pattern (q,q,q,1).
% No basin-positivity assumption is needed for the exclusion step.

\subsection{三つの等しい高分母を持つ core の係数和}
\label{uh10:three-section}

一般の最大分母についての局所符号原理は成立しないが，
三つの高分母が等しい場合には，その三頂点の係数和に符号がある．

\begin{proposition}[三つの等高分母の負の総和]\label{uh10:three-negative}
局所四直線模型の分母が $(q,q,q,1)$，$q\geq3$，であり，
全 row sums と pair products が
\[
 \sum_{j\ne i}p_{ij}=-q_i,
 \qquad p_{ij}p_{ji}=q_iq_j-1
\]
を満たすとする．分母一の頂点を $0$ と書けば
\[
 D_0>0,\qquad \sum_{i\ne0}D_i=-D_0<0.
\]
この結論には basin positivity を仮定する必要はない．
\end{proposition}

\begin{proof}
$m=q-1\geq2$ とする．分母一の row の三 entries は非零の signed divisors of $m$ で，
和は $-1$ である．従って，二つが負で一つが正か，二つが正で一つが負である．

まず $(-a,-b,c)$，$a,b,c>0$，$c=a+b-1$ と書ける場合を考える．
pair products から
\[
 D_0=-1+m\left(\frac1a+\frac1b-\frac1c\right).
\]
$a=1$ なら $c=b$ なので $D_0=m-1>0$ である．
$a>1$ なら $c>b$ なので括弧内は $1/a$ より大きい．
$a\mid m$，$a>0$ より $m/a\geq1$ であり，この場合も $D_0>0$ となる．

残る $(a,b,-c)$，$c=a+b+1$ の場合が存在しないことを示す．
\[
 r=\frac ma,\qquad s=\frac mb,\qquad t=\frac mc,
 \qquad N=q^2-1
\]
と置く．高分母の三頂点における，分母一の頂点へ向かう entries は $r,s,-t$ である．
従って高分母三頂点の内部 entries だけを取り，diagonal を
\[
 A=q+r,\qquad B=q+s,\qquad C=q-t
\]
とすれば，三行の和が全て零である $3\times3$ 行列を得る．
その三つの off-diagonal pair products は全て $N$ である．
二つの向きの三角 cycle products を $X,Y$ と書くと，determinant が零であることから
\[
 X+Y=-\{ABC-N(A+B+C)\},
 \qquad XY=N^3.
\]
従って，実数 $X,Y$ の存在の必要条件は
\begin{equation}\label{uh10:triangle-discriminant}
 \{ABC-N(A+B+C)\}^2\geq4N^3
\end{equation}
である．

一方，$c=a+b+1$ は
\[
 rs-rt-st=\frac{rst}{m}
\]
と同値である．これを用いて展開すると
\[
 ABC-N(A+B+C)=-L+U,
 \qquad L=2q^3-3q,
 \qquad U=r+s-t+\frac{rst}{m}
\]
となる．$c>a,b$ より $t<r,s$ なので $U>0$ であり，
$1\leq r,s,t\leq m$ から
\[
 U\leq2m+m^2=N<L.
\]
従って $(-L+U)^2<L^2$ である．さらに
\[
 4N^3-L^2=3q^2-4>0
\]
なので，\eqref{uh10:triangle-discriminant} に矛盾する．
以上で二つが正の row は排除される．
最後に，四頂点全体の係数和は row-sum 条件から零なので，
高分母三頂点の係数和は $-D_0<0$ である．
\end{proof}

\subsection{単色高分母成分と二値分母の剛性}
\label{uh11:monochromatic-section}

局所の係数和を高分母 support の連結成分上で足すと，
分母の上限を置かずに二値分母を排除できる．
ここで「単色」とは，その成分の全ての高分母が同じ正整数であることを意味する．

\begin{theorem}[単色高分母成分の不存在]\label{uh11:monochromatic-component}
本稿の全平面五次の SNC cores，標準 $K_4$ cones，既存の branch graph を保ち，
追加の branch singularities を入れない大域作用適合 atlas を考える．
高分母 support $H=\{ij:q_{ij}>1\}\subset K_5$ の，辺を持つ連結成分では，
少なくとも二種類の高分母が現れる．
この結論は分母の上限を仮定せず，node 数の非負性や全 node budget も用いない．
\end{theorem}

\begin{proof}
辺を持つ連結成分 $C$ の全ての高分母が同じ $q>1$ であると仮定する．
三価 branch の整数格子条件から $q$ は奇数であり，従って $q\geq3$ である．
core $i\in V(C)$ で高分母 branch の本数を $k_i$ と書く．
$C$ は $H$ の連結成分なので，その core の残る $4-k_i$ 個の分母は全て一である．
その core における高分母 rays の normal-action 係数の和を
\[
 S_i=\sum_{\substack{j:\,ij\in E(C)}}D^i_j
\]
と置く．$k_i=1$ なら命題~\ref{qh9:isolated-local}，
$k_i=2$ なら系~\ref{uh11:two-negative}，
$k_i=3$ なら命題~\ref{uh10:three-negative} により $S_i<0$ である．
$k_i=4$ なら全四 rays を足しているため，row-sum 条件から $S_i=0$ である．
従って全ての $S_i\leq0$ であり，次数一から三の vertex が一つでもあれば総和は負となる．

一方，各共有 branch では対向する normal actions の係数が相殺するので，
\[
 \sum_{i\in V(C)}S_i
 =\sum_{ij\in E(C)}(D^i_j+D^j_i)=0.
\]
従って $C$ の全 vertices は次数四でなければならない．
$C\subset K_5$ であるため，これは $C=K_5$ を意味する．
すると $q$ の任意の素因数 $\ell$ に対し $G_\ell=K_5$ となり，
定理~\ref{ps10:forbidden} の $K_4$ 禁止に反する．
\end{proof}

\begin{corollary}[二値分母の剛性]\label{uh11:two-value-rigidity}
ある正整数 $q>1$ について，全十本の branch denominators が
$q_{ij}\in\{1,q\}$ を満たすなら，全て $q_{ij}=1$ である．
\end{corollary}

\begin{proof}
高分母 branch があれば，それを含む $H$ の連結成分は単色となり，
定理~\ref{uh11:monochromatic-component} に反する．
\end{proof}

この剛性は，局所の等高分母 chart の不存在を述べるものではない．
局所には $(q,q,1,1)$ の正の模型が存在する．
その高分母係数の負の総和と，大域の対向端相殺を同時に要求すると矛盾するのである．
異なる高分母が同じ連結成分内に混在する場合には，この総和の符号は一般には保たれない．
従って任意分母に関する全剛性は，なお別の問題である．

% v10: independently audited modular consequences of the local equations
% and puncture clutching.  The modular results use no basin positivity or
% bounded search.  The last geometric corollary also uses qh9:isolated-edge.

\section{素数 support に対する五次の合同障害}
\label{ps10:section}

全平面五次の五つの cores の標識集合を $V=\{0,1,2,3,4\}$ とする．
共有 branch $ij$ の分母を $q_{ij}=q_{ji}$，core $i$ の branch $ij$ における
puncture $ijk$ の numerator を $p^i_{jk}$ と書く．
以下では $q_{ij}$ が正の奇整数で，$p^i_{jk}\in\ZZ$ が，相異なる $i,j,k$ に対し
\begin{align}
 \sum_{k\in V\setminus\{i,j\}}p^i_{jk}&=-q_{ij},
       \label{ps10:row}\\
 p^i_{jk}p^i_{kj}&=q_{ij}q_{ik}-1,
       \label{ps10:det}\\
 q_{ij}&\mid p^i_{jk}+p^j_{ik}
       \label{ps10:clutch}
\end{align}
を満たすと仮定する．最初の二式は局所の三 puncture period relation と
crossing unimodularity，第三式は共有 puncture cycles の整数 clutching から来る．
以下の合同障害はこの三条件だけを用い，basin 面積の正値性は用いない．
後の系~\ref{ps10:all-high-divisible-three} では，さらに既存の幾何的接着障害を用いる．

奇素数 $\ell$ ごとに，頂点集合 $V$ 上の単純 graph
\[
 E(G_\ell)=\bigl\{\{i,j\}:\ell\mid q_{ij}\bigr\}
\]
を定める．部分 graph を含むという表現は induced であることを要求しない．
以下の合同計算はすべて $\mathbf F_\ell$ 内で行う．
$ij\in E(G_\ell)$ なら，\eqref{ps10:det} により，すべての $k\ne i,j$ で
\begin{equation}\label{ps10:unit}
 p^i_{jk}\ne0,
 \qquad p^i_{kj}=-(p^i_{jk})^{-1}
\end{equation}
となることを繰り返し用いる．

\begin{lemma}[三角形の符号]\label{ps10:triangle}
$i,j,k$ が $G_\ell$ の三角形をなすなら，ある
$\varepsilon\in\{1,-1\}\subset\mathbf F_\ell$ に対し
\[
 p^i_{jk}=p^j_{ki}=p^k_{ij}=\varepsilon,
 \qquad
 p^i_{kj}=p^j_{ik}=p^k_{ji}=-\varepsilon
\]
である．特に，この三角形に属する各 numerator の二乗は一である．
\end{lemma}

\begin{proof}
$x=p^i_{jk}\ne0$ と置く．core $i$ の determinant と，辺 $ij,ik$ の
clutching を順に用いると
\[
 p^i_{kj}=-x^{-1},\qquad p^j_{ik}=-x,\qquad p^k_{ij}=x^{-1}
\]
を得る．core $j$ では $p^j_{ki}=x^{-1}$ であり，辺 $jk$ の clutching は
$p^k_{ji}=-x^{-1}$ を与える．core $k$ の determinant より
$-x^{-2}=-1$，すなわち $x^2=1$ となる．
$\ell$ は奇素数なので $x=\varepsilon=\pm1$ であり，表示した六式を得る．
\end{proof}

\begin{theorem}[素数 support の禁止配置]\label{ps10:forbidden}
\eqref{ps10:row}--\eqref{ps10:clutch} の下で，次が成り立つ．
\begin{enumerate}
\item 任意の奇素数 $\ell$ に対し，$G_\ell$ は $K_4$ を含まない．
\item $G_3$ は matching，すなわち最大次数が一以下の graph である．
\item 任意の奇素数 $\ell$ に対し，$G_\ell$ は book $B_3$ を含まない．
  ここで $B_3$ は一つの辺 $ij$ と，その両端に接続する三つの相異なる
  vertices $k,l,m$ からなる七辺の graph である．
\item $G_\ell$ が $K_1\vee P_4$ を含むなら，$\ell=5$ である．
  ここで $K_1\vee P_4$ は，path $a-b-c-d$ の四 vertices すべてに
  一つの vertex $o$ を接続した七辺の graph を表す．
\end{enumerate}
\end{theorem}

\begin{proof}
まず $K_4$ の vertices を $i,j,k,l$，残る vertex を $m$ とする．
core $i$ で，row $j$ の entries $p^i_{jk},p^i_{jl}$ は
補題~\ref{ps10:triangle} により $\pm1$ である．
これらが逆符号なら，row sum より $p^i_{jm}=0$ となり，
\eqref{ps10:unit} に反する．従って二つは同じ符号 $s_j$ を持つ．
同様に rows $k,l$ の二つの内部 entries も，各々同じ符号 $s_k,s_l$ を持つ．
各 pair の determinant は
\[
 s_js_k=s_js_l=s_ks_l=-1
\]
を要求するが，三つの $\pm1$ が互いに逆符号となることはできない．
これが第一項である．

第二項では $\ell=3$ とし，ある core に接する二本以上の edges が
$G_3$ に属すると仮定する．この core の四つの局所 branch labels を
$0,1,2,3$ と書き直し，$q_0=q_1=0$ in $\mathbf F_3$ としてよい．
rows $0,1$ の三つの entries はそれぞれ全て非零であり，row sum は零である．
$\mathbf F_3$ で三つの $\pm1$ の和が零であるためには，全てが同じ符号でなければならない．
従ってある $\sigma,\tau\in\{1,-1\}$ に対し
\[
 p_{01}=p_{02}=p_{03}=\sigma,
 \qquad p_{10}=p_{12}=p_{13}=\tau.
\]
pair $01$ の determinant は $\sigma\tau=-1$ を与えるので，$\tau=-\sigma$ である．
残る $j=2,3$ について，pairs $0j,1j$ の determinants より
\[
 p_{j0}=-\sigma,\qquad p_{j1}=\sigma.
\]
これらは各 row で打ち消し合うため，$p_{23}=-q_2$，$p_{32}=-q_3$ となる．
しかし pair $23$ の determinant は
\[
 q_2q_3=p_{23}p_{32}=q_2q_3-1
\]
を要求し，矛盾する．$q_2,q_3$ の剰余には追加条件を課していない．
従って一つの core に三で割れる分母は高々一つであり，$G_3$ は matching である．
この項では clutching 条件も面積正値性も使っていない．

第三項の $B_3$ があれば，core $i$ の row $j$ は三つの三角形 $ijk,ijl,ijm$ により
全 entries が $\pm1$ となる．その和が零であることから $\ell=3$ が必要である．
しかし $i$ は $G_3$ で次数四を持ち，第二項に反する．

最後に $K_1\vee P_4$ を仮定し，core $o$ の添字 $o$ を省略する．
三角形の補題から
\[
 p_{ab}=s,\qquad p_{bc}=t,\qquad p_{cd}=u,
 \qquad s,t,u\in\{1,-1\}
\]
と書ける．この core の四分母は全て零であり，すべての非対角 entries が非零である．
row $b$ より $p_{bd}=s-t\ne0$ なので，$t=-s$，$p_{bd}=2s$ となる．
row $c$ より $p_{ca}=t-u=-s-u\ne0$ なので，$u=s$，$p_{ca}=-2s$ となる．
determinant と row $a$ を用いると
\[
 p_{ac}=\frac{s}{2},\qquad p_{ad}=-\frac{3s}{2}.
\]
$p_{ad}\ne0$ より $\ell\ne3$ である．従って row $d$ は
\[
 p_{da}+p_{db}+p_{dc}
   =\frac{2s}{3}-\frac{s}{2}-s
   =-\frac{5s}{6}=0
\]
となり，$\ell=5$ が従う．
\end{proof}

\begin{proposition}[$K_{2,2,1}$ と残る二分母]\label{ps10:k221}
$G_\ell$ が parts $\{a,b\},\{c,d\},\{o\}$ の $K_{2,2,1}$ を含むなら，
$\ell=5$ であり，$G_5$ はこの $K_{2,2,1}$ に等しい．
さらに，ある $s\in\{1,-1\}\subset\mathbf F_5$ に対し
\begin{equation}\label{ps10:k221-residues}
 q_{ab}\equiv-s,\qquad q_{cd}\equiv s\pmod5.
\end{equation}
特に $q_{ab}+q_{cd}\equiv0$ および $q_{ab}^2\equiv q_{cd}^2\equiv1\pmod5$ である．
\end{proposition}

\begin{proof}
$K_{2,2,1}$ は $K_1\vee P_4$ を含むので，定理~\ref{ps10:forbidden} より
$\ell=5$ である．辺 $ab$ または $cd$ を加えると $K_4$ が生じるため，
$q_{ab},q_{cd}$ は五で割れず，support graph はちょうど $K_{2,2,1}$ となる．

まず core $o$ を考える．各 row の反対 part への二 entries は $\pm1$ であり，
残る一 entry が非零であるため，その二つは同符号である．
cross-part determinants から，part $\{a,b\}$ の二 rows は同じ符号 $s$，
part $\{c,d\}$ の二 rows は逆符号 $-s$ を持つ．
row sums より
\[
 p^o_{ab}=p^o_{ba}=-2s,
 \qquad p^o_{ac}=p^o_{ad}=p^o_{bc}=p^o_{bd}=s,
 \qquad p^o_{cd}=p^o_{dc}=2s.
\]

次に core $a$ を考える．辺 $ao$ の clutching と local determinants より
\[
 p^a_{ob}=2s,\qquad p^a_{oc}=p^a_{od}=-s,
 \qquad p^a_{bo}=2s,\qquad p^a_{co}=p^a_{do}=s
\]
となる．ここでは $-1/2=2$ in $\mathbf F_5$ を用いた．
$p^a_{cd}=s x$ と書くと，$x\ne0$ であり，rows $c,d$ と determinants は
\[
 p^a_{cb}=-s(1+x),\qquad
 p^a_{db}=s(-1+x^{-1}),\qquad
 p^a_{bc}=\frac{s}{1+x},\qquad
 p^a_{bd}=\frac{s}{1-x^{-1}}
\]
を与える．最初の二 entries が非零なので $x\ne\pm1$，従って $x=\pm2$ である．
row $b$ の和は
\[
 -q_{ab}
 =s\left(2+\frac1{1+x}+\frac1{1-x^{-1}}\right)
 =\begin{cases}s&x=2,\\0&x=-2.\end{cases}
\]
$q_{ab}\not\equiv0$ なので $x=2$ であり，$q_{ab}\equiv-s$ となる．
二つの parts を交換すると $s$ は $-s$ に変わるので，同じ議論から
$q_{cd}\equiv s$ を得る．
\end{proof}

\begin{corollary}[予備的な edge 数評価と共通因子]\label{ps10:counts-gcd}
\eqref{ps10:row}--\eqref{ps10:clutch} の下で
\[
 |E(G_\ell)|\le
 \begin{cases}
 2&\ell=3,\\
 8&\ell=5,\\
 7&\ell\ne3,5
 \end{cases}
\]
である．また，任意の四元部分集合 $S\subset V$ に対し
\begin{equation}\label{ps10:gcd}
 \gcd\{q_{ij}:i,j\in S,\ i<j\}=1.
\end{equation}
従って，全十本の共有 branches が同じ分母 $q>1$ を持つ配置は存在しない．
\end{corollary}

\begin{proof}
$G_3$ は五 vertices 上の matching なので，edge 数は二以下である．
五 vertices 上の $K_4$-free graph は八 edges 以下である．
実際，九 edges 以上なら欠ける edge は高々一本なので $K_4$ を含む．
八 edges の場合，欠ける二 edges が共通 vertex を持つと，残る四 vertices が
$K_4$ をなす．従って欠ける二 edges は disjoint であり，元の graph は
$K_{2,2,1}$ に等しい．この場合は命題~\ref{ps10:k221} により $\ell=5$ である．
これが三つの edge bounds を与える．

\eqref{ps10:gcd} の左辺が一より大きければ，その奇素因数 $\ell$ に対し
$S$ が $G_\ell$ 内の $K_4$ をなす．これは定理~\ref{ps10:forbidden} に反する．
全分母が同じ $q$ なら，任意の四元部分集合の gcd も $q$ なので最後の結論を得る．
\end{proof}

\begin{corollary}[全高分母が三で割れる場合]\label{ps10:all-high-divisible-three}
本稿の全平面五次の SNC core models，標準 $K_4$ cones，既存の branch graph を用い，
追加の branch singularities を入れない大域作用適合 atlas を考える．
すべての共有 branch に対して
\[
 q_{ij}>1\quad\Longrightarrow\quad 3\mid q_{ij}
\]
が成り立つなら，すべての $q_{ij}$ は一である．
同値に，高分母を持つ大域 atlas には，三で割れない高分母も必要である．
\end{corollary}

\begin{proof}
高分母 branches だけからなる support graph を $H$ とする．
仮定の下では $H=G_3$ であり，定理~\ref{ps10:forbidden} により $H$ は matching である．
$H$ が非空なら，その各 edge は単独で一つの連結成分をなす．
これは系~\ref{qh9:isolated-edge} に反する．
具体的には，その両端の core は分母 pattern $(q,1,1,1)$ を持ち，
実際の atlas の正の basin areas により，命題~\ref{qh9:isolated-local} が両端の
normal-action 分岐係数をともに負とする．
同じ neck の対向 normal circles から必要な係数の相殺は成立し得ない．
従って $H$ は空である．
\end{proof}

\begin{remark}[七辺の場合の残る graph]\label{ps10:seven-edge-shape}
$\ell\ne3,5$ で $|E(G_\ell)|=7$ なら，complement は $P_3\sqcup K_2$ に限られる．
同値に，$G_\ell$ は $K_{2,3}$ の三点側に一本の edge を加えた graph である．
実際，五 vertices 上の三辺 graph は，isolated vertices を省いて
$K_3,K_{1,3},P_4,P_3\sqcup K_2$ のいずれかである．
前の三つの complements は，順に $B_3$，$K_4$ を含む graph，$K_1\vee P_4$ なので排除される．
この段階の必要条件は，残る七辺 graph の実現可能性を主張しない．
次節の定理~\ref{ps11:dense} では，五つの cores の式を用いて七辺ちょうどの場合を全て排除する．
\end{remark}

以上は素数ごとの普遍的な合同障害であり，分母の大きさを制限した探索の結論ではない．
一方，異なる素因数を持つ分母の配置や，許された support graphs の整数 lifting，
basin areas と有限 neck の同時接着は，これらの必要条件だけでは決まらない．
特に，系~\ref{ps10:all-high-divisible-three} の追加仮定を外した任意の分母配置について，
全平面五次のすべての分母が一であることまでは結論しない．

% v11 addition after quintic_prime_support.tex.  No existing source is edited.
% This is an unbounded-prime proof; the companion program only checks examples
% and exhausts the small F_5 completion problem.

\subsection{七辺 support の排除と密な合同模型の完全記述}
\label{ps11:dense-section}

本節でも条件~\eqref{ps10:row}--\eqref{ps10:clutch} を仮定し，
合同計算には basin 面積の正値性を用いない．
前節の七辺 bound は，五つの cores の局所式を同時に使うと強化できる．

\begin{theorem}[密な素数 support]\label{ps11:dense}
任意の奇素数 $\ell$ に対し，$|E(G_\ell)|\ge7$ なら
\[
 \ell=5,\qquad G_\ell\cong K_{2,2,1}.
\]
特に，七辺ちょうどの $G_\ell$ は存在せず，
\begin{equation}\label{ps11:bounds}
 |E(G_\ell)|\le
 \begin{cases}
 2&\ell=3,\\
 8&\ell=5,\\
 6&\ell\ne3,5
 \end{cases}
\end{equation}
となる．$\ell=5$ でも，許される edge 数は $0,\ldots,6,8$ に限られる．

固定された標識付き $K_{2,2,1}$ に対し，対応する $\mathbf F_5$ 上の
row sums，determinants，零分母辺の clutching の解は，後述する
$s=\pm1$ の二つだけである．これらは normal-action 係数の対向相殺も満たす．
これは合同式の解の分類であり，正整数への lifting や大域作用 atlas の実現を含まない．
\end{theorem}

\begin{lemma}[二部 graph に一本を加えた場合]\label{ps11:seven-bipartite}
$G_\ell$ が，parts $\{a,b\},\{c,d,e\}$ の $K_{2,3}$ に
辺 $cd$ を加えた七辺 graph $H$ を含むとする．
このとき $\ell=5$ であり，$ce,de$ のちょうど一方も $G_5$ の辺である．
従って $G_5$ は $K_{2,2,1}$ であり，$H$ そのものを exact support とすることはできない．
\end{lemma}

\begin{proof}
三角形 $acd,bcd$ の符号を $s,t\in\{\pm1\}$ とする．
core $c$ の row $d$ は，最初の二 entries が $s,t$ であり，
第三 entry $-s-t$ も非零である．従って $t=s$ である．
この合同計算の間だけ，全 $p,q$ を同じ $s$ で割る．
$s^2=1$ なので row sums，determinants，零分母辺の clutching は保たれ，
二つの三角形の符号は一に正規化される．

\[
 u=p^c_{ab},\qquad v=p^d_{ab},\qquad
 x=q_{ab},\quad y=q_{ce},\quad z=q_{de}
\]
と置く．core $c$ の rows $a,b,d$ は，列順を $a,b,d,e$ として
\[
 \begin{pmatrix}
 0&u&-1&1-u\\
 -1/u&0&-1&1+1/u\\
 1&1&0&-2
 \end{pmatrix},
\]
core $d$ の対応する rows は，列順 $a,b,c,e$ で
\[
 \begin{pmatrix}
 0&v&1&-1-v\\
 -1/v&0&1&-1+1/v\\
 -1&-1&0&2
 \end{pmatrix}
\]
となる．高分母の局所 branch に接する entries は全て非零なので，
$u,v,1-u,1+u,1-v,1+v$ は非零である．

辺 $ac,ad,bc,bd$ の clutching と各 core の determinants を使うと，
core $a$ の row $b$ と core $b$ の row $a$ はそれぞれ
\begin{align}
 -x&=\frac1u+\frac1v+
       \frac{(1-u)(1+v)}{u+v},\label{ps11:x-a}\\
 -x&=-u-v+\frac{(1+u)(1-v)}{u+v}.
       \label{ps11:x-b}
\end{align}
ここで $u+v\ne0$ は，core $a$ の row $e$ の entry
\[
 p^a_{eb}=-\frac{u+v}{(1-u)(1+v)}
\]
が非零であることから従う．辺 $ae,be$ の clutching により，
core $e$ の rows $a,b$ は列順 $a,b,c,d$ で
\begin{equation}\label{ps11:e-high-rows}
 \begin{pmatrix}
 0&\dfrac{u+v}{(1-u)(1+v)}&-\dfrac1{1-u}&\dfrac1{1+v}\\[6pt]
 \dfrac{u+v}{(1+u)(1-v)}&0&-\dfrac u{1+u}&-\dfrac v{1-v}
 \end{pmatrix}
\end{equation}
となる．その pair $ab$ の determinant は
\[
 (u+v)^2+(1-u^2)(1-v^2)=(uv+1)^2=0
\]
を要求する．体上の計算なので $uv=-1$ である．
\eqref{ps11:x-a} から \eqref{ps11:x-b} を引くと
\[
 0=\frac{u+v}{uv}+(u+v)+\frac{2(v-u)}{u+v}
   =\frac{2(v-u)}{u+v}.
\]
$\ell$ は奇数であるため，$v=u$，従って $u^2=-1$ を得る．
core $c,d$ の残る row sums と \eqref{ps11:x-a} は
\begin{equation}\label{ps11:xyz}
 x=3u,\qquad y=u+\frac12,\qquad z=u-\frac12
\end{equation}
を与える．具体的に，$y,z$ を決める式は
\[
 y=\frac1{1-u}+\frac u{1+u}-\frac12,
 \qquad
 z=-\frac1{1+v}+\frac v{1-v}+\frac12
\]
である．

\eqref{ps11:e-high-rows} の reverse entries を determinants で求め，
残る二つの row sums を使うと，core $e$ の rows $c,d$ は
\[
 \begin{pmatrix}
 1-u&1-u&0&u-5/2\\
 -1-u&-1-u&u+5/2&0
 \end{pmatrix}
\]
となる．最後の pair $cd$ の determinant の左右の差は
\begin{equation}\label{ps11:five-residual}
 \left(u-\frac52\right)\left(u+\frac52\right)
 -\left[\left(u+\frac12\right)\left(u-\frac12\right)-1\right]
 =-5.
\end{equation}
従って $\ell=5$ が必要である．$\mathbf F_5$ 内の $u^2=-1$ の解は
$u=\pm2$ なので，\eqref{ps11:xyz} から $y,z$ のちょうど一方は零，
他方は非零となる．また $x=3u\ne0$ である．
元の符号 $s$ を戻しても零か否かは変わらず，結論を得る．
特に，ここで禁止しているのは exact 七辺 support であり，
$H$ を部分 graph として含む $K_{2,2,1}$ を排除していない．
\end{proof}

次の小さな局所計算を，もう一つの七辺 graph と八辺模型の記述に使う．

\begin{lemma}[三つの零分母を持つ局所合同模型]\label{ps11:three-zero}
$\mathbf F_5$ 上の局所四 branch を $o,b,c,d$ とし，
分母が $(0,r,0,0)$，$r\ne0$ であるとする．
row $o$ が $(p_{ob},p_{oc},p_{od})=(2s,-s,-s)$，$s=\pm1$ なら，
row sums と determinants は
\begin{equation}\label{ps11:outer-matrix}
 r=-s,\qquad
 P=s\begin{pmatrix}
 0&2&-1&-1\\
 2&0&2&2\\
 1&2&0&2\\
 1&2&2&0
 \end{pmatrix}
\end{equation}
を強制する．
\end{lemma}

\begin{proof}
$p_{cd}=st$ と置く．determinants と rows $c,d$ から
\[
 p_{cb}=-s(1+t),\quad p_{db}=s(-1+t^{-1}),\quad
 p_{bc}=\frac{s}{1+t},\quad p_{bd}=\frac{s}{1-t^{-1}}.
\]
全 entries は零分母 branch に接するので非零であり，$t\ne0,\pm1$ である．
従って $t=\pm2$ である．row $b$ の和は，$t=2$ なら $s$，
$t=-2$ なら零である．$r\ne0$ より前者に限られ，表示した値が全て従う．
\end{proof}

\begin{lemma}[$K_1\vee P_4$ の completion]\label{ps11:path-completion}
$G_\ell$ が，vertex $o$ と path $a-b-c-d$ の join を含むなら，
$\ell=5$ であり，さらに辺 $ad$ を含む．従ってその support は
parts $\{a,c\},\{b,d\},\{o\}$ の $K_{2,2,1}$ に等しい．
\end{lemma}

\begin{proof}
定理~\ref{ps10:forbidden} により $\ell=5$ であり，$K_4$ の禁止により
$q_{ac},q_{bd}$ は非零である．同定理の証明で得た core $o$ の行列を
$\mathbf F_5$ で書くと，列・行の順 $a,b,c,d$ に対し
\[
 P^o=s\begin{pmatrix}
 0&1&3&1\\
 4&0&4&2\\
 3&1&0&1\\
 4&2&4&0
 \end{pmatrix},\qquad s\in\{\pm1\}.
\]
core $c$ と core $b$ に補題~\ref{ps11:three-zero} を適用すると，それぞれ
\[
 q_{ac}=-s,\qquad q_{bd}=s
\]
となる．core $b$ の row $a$ は，列 $o,c,d$ に対して
$(-s,-2s,-2s)$ である．辺 $ab$ の clutching を使うと，
core $a$ の row $b$ は $(s,2s,2s)$ となる．
さらに辺 $ao$ の clutching と determinants により，core $a$ では
\[
 p^a_{co}=2s,\quad p^a_{do}=s,\quad
 p^a_{cb}=p^a_{db}=2s.
\]
$y=q_{ad}$ と置くと，rows $c,d$ より
\[
 p^a_{cd}=2s,\qquad p^a_{dc}=-y-3s.
\]
pair $cd$ の determinant は
\[
 2s(-y-3s)=(-s)y-1
\]
を要求する．$\mathbf F_5$ 内ではこの式は $sy=0$ と同値であるから，$y=0$ である．
これが追加辺 $ad$ を与える．他の二辺を加えると $K_4$ となるので，
support は表示した $K_{2,2,1}$ に等しい．
\end{proof}

\begin{proof}[定理~\ref{ps11:dense} の証明]
八辺以上の場合は系~\ref{ps10:counts-gcd} と命題~\ref{ps10:k221} による．
七辺の場合，complement は三辺 graph である．
孤立 vertices を省くと，それは $K_3,K_{1,3},P_4,P_3\sqcup K_2$ のいずれかである．
最初の二つの complements はそれぞれ $B_3$ と $K_4$ を含むため排除される．
後の二つはそれぞれ $K_1\vee P_4$ と補題~\ref{ps11:seven-bipartite} の $H$ なので，
いずれも補題により八本目の辺を強制する．従って七辺ちょうどの場合はない．
三に関する bound は既存の matching 定理から従う．

残る八辺の場合を具体的に記述する．parts を $\{a,b\},\{c,d\},\{o\}$ と固定する．
命題~\ref{ps10:k221} より，ある $s=\pm1$ に対し
\[
 q_{ab}=-s,\qquad q_{cd}=s,\qquad
 q_{ij}=0\quad(ij\ne ab,cd)
\]
となる．core $o$ の行列は，行・列順 $a,b,c,d$ で
\begin{equation}\label{ps11:central-matrix}
 P^o=s\begin{pmatrix}
 0&-2&1&1\\
 -2&0&1&1\\
 -1&-1&0&2\\
 -1&-1&2&0
 \end{pmatrix}.
\end{equation}
core $a,b$ は，他方を $b'$ と書き，行・列順 $(o,b',c,d)$ で
\eqref{ps11:outer-matrix} の行列である．core $c,d$ は同様に，
行・列順 $(o,d',a,b)$ に対し同じ行列の符号 $s$ を $-s$ に置き換えたものである．
補題~\ref{ps11:three-zero} により，これら以外の局所 completion はない．

逆に，表示した二組のデータを代入すると，全ての row sums，determinants，
零分母辺の clutching が成り立つ．normal-action 係数
\[
 D^i_j=-q_{ij}-\sum_{k\ne i,j}p^i_{kj}
\]
は core $o$ で $s(-1,-1,1,1)$，core $a,b$ の上記順序で
$s(1,0,2,2)$，core $c,d$ の上記順序で $-s(1,0,2,2)$ となる．
従って全十辺で $D^i_j+D^j_i=0$ も満たす．
これは，この八辺型を同じ mod~$5$ 条件だけでさらに排除できないことを示す．
\end{proof}

\begin{corollary}[七分母の共通因子の指数]\label{ps11:dense-gcd}
任意に選んだ七本の異なる branches の分母の最大公約数は五を割る．
特に
\[
 \#\{ij:25\mid q_{ij}\}\le6.
\]
$G_5=K_{2,2,1}$ なら，その八辺の分母の最大公約数は正確に五である．
従って，その八辺のうち少なくとも二本では $v_5(q_{ij})=1$ となる．
\end{corollary}

\begin{proof}
七分母に共通する素因数は，定理~\ref{ps11:dense} により五に限られる．
七分母の全てが二十五で割れると仮定する．その七辺は $G_5=K_{2,2,1}$ の
部分 graph であるから，$K_1\vee P_4$ または補題~\ref{ps11:seven-bipartite} の $H$ となる．
以下では $R=\ZZ/25\ZZ$ 上で計算する．
高分母 branch に接する全 entries は determinant が $-1$ なので units である．
三角形の計算から $x^2=1$ を得るが，この環でもその解は $x=\pm1$ に限られる．
実際，$(x-1)(x+1)=0$ で両因子の差は unit の二であるから，
片方は五で割れず，他方は二十五で割れなければならない．

$K_1\vee P_4$ の場合，定理~\ref{ps10:forbidden} の path 計算がそのまま成り立ち，
\[
 p_{ab}=s,\quad p_{bc}=-s,\quad p_{cd}=s,
 \qquad p_{da}=\frac{2s}{3},\quad
 p_{db}=-\frac{s}{2},\quad p_{dc}=-s
\]
となる．分母二と三はこの環の units であり，row $d$ の和は
$-5s/6\ne0$ in $R$ となる．これは row sum に反する．

$H$ の場合も，補題~\ref{ps11:seven-bipartite} の有理式で現れた分母は
全て units であり，\eqref{ps11:x-a}--\eqref{ps11:e-high-rows} は $R$ 上で有効である．
ただし $r=uv+1$ と置いたとき，この環で $r^2=0$ から $r=0$ とは結論しない．
代わりに，二つの $x$ の式の差から
\[
 2uv(v-u)=-(u+v)^2r
\]
を使う．従って $v-u$ は ideal $(r)$ に属し，
\[
 u^2+1=r-u(v-u),\qquad v^2+1=r+v(v-u)
\]
も $(r)$ に属する．よってこれらの平方は零である．
\[
 A=\frac{2u}{1-u^2},\quad B=\frac{2v}{1-v^2},\quad
 C=u-u^{-1}-A,\quad D=v-v^{-1}-B
\]
と置くと
\[
 C-A=\frac{(u^2+1)^2}{u(u^2-1)}=0,\qquad D-B=0,
 \qquad
 A-B=\frac{2(u-v)r}{(1-u^2)(1-v^2)}=0.
\]
core $c,d$ の残る row sums は $y=A+1/2,z=B-1/2$ を与え，
core $e$ の最後の二 entries は $p^e_{cd}=C-5/2$，$p^e_{dc}=D+5/2$ となる．
その determinant の左右の差は
\[
 CD-AB+\frac52(C-D)+\frac12(A-B)-5=-5\ne0
 \quad\text{in }R.
\]
これも矛盾である．従って七分母の gcd は二十五では割れず，五を割る．

二十五で割れる枝が七本以上あれば，そのうち七本を選んで直前の結論に反する．
最後に $G_5=K_{2,2,1}$ の八分母は全て五で割れ，その gcd は五を割るので，
gcd は正確に五である．二十五で割れるものは高々六本であるため，
残る少なくとも二本の五進付値は一となる．
\end{proof}

\begin{remark}[合同分類と整数接着の区別]\label{ps11:modular-scope}
零でない $q_{ij}$ に対する整数整除
$q_{ij}\mid p^i_{jk}+p^j_{ik}$ は，mod~$5$ に落とすとそれだけでは
追加の制約を課さない．上の二解は，この意味での reduced system の全解である．
素数ごとの residue data が存在しても，正の奇整数 $q$，整数 $p$，
全ての整除条件，正の basin 面積，node 配分，有限 neck の作用接着を
同時に持つとは限らない．この定理は任意分母の大域存在も不存在も結論しない．

標識付き $K_{2,2,1}$ は $5\cdot3=15$ 個であり，各々に二つの符号解がある．
付属の標準 Python による検算は，この三十個の合同データについて全式を確認する．
さらに $K_1\vee P_4$ を含む mod~$5$ データの全百六十通りの
分母・符号 assignments を，numerator の全 residue pairs を含めて列挙し，
八辺 completion の二解のみが残ることを確認する．
一般の素数に対する結論はこの有限計算ではなく，
\eqref{ps11:five-residual} を含む上の恒等式の証明による．
\end{remark}

\subsection{高分母全体の原始性と単位枝による合同条件の強化}
\label{gp23:section}

高分母 support の大きさと密な素数 support の分類を合わせると，
単色成分の排除より強い共通因子の制約が得られる．
以下も指定した標準五次接着データに関する主張である．

\begin{corollary}[高分母全体は原始的である]
\label{gp23:primitive}
高分母 support $H$ が非空なら，
\[
 \gcd\{q_{ij}:ij\in E(H)\}=1.
\]
従って非空の高分母全体が一つの素数の倍数からなることはない．
\end{corollary}

\begin{proof}
系~\ref{sf22:no-leaves} より $|E(H)|\ge8$ である．
全ての高分母を割る素数 $\ell$ が存在すると仮定する．分母の奇数性から $\ell$ も奇数である．
単位分母は $\ell$ で割れないから，素数 support はちょうど $G_\ell=H$ となる．
定理~\ref{ps11:dense} は
$\ell=5$，$H\cong K_{2,2,1}$ を強制する．
その二つの非辺を $ab,cd$ とすると，実際の分母は両方とも一である．
一方，同定理の合同分類は，ある $\varepsilon\in\{1,-1\}$ に対し
\[
 q_{ab}\equiv-\varepsilon,\qquad q_{cd}\equiv\varepsilon\pmod5
\]
を与える．従って $2=q_{ab}+q_{cd}\equiv0\pmod5$ となり矛盾する．
\end{proof}

\begin{corollary}[二本の単位枝と七分母の互いに素性]
\label{gp23:two-units}
相異なる二本以上の branches の分母が一であるなら，任意に選んだ
相異なる七本の branches の分母の最大公約数は一である．
特に残る唯一の八高辺型 $G_{5,14}$ の場合，五で割れる分母は高々六本である．
\end{corollary}

\begin{proof}
系~\ref{ps11:dense-gcd} により，七分母の最大公約数は五を割る．
もし五が共通因子なら，定理~\ref{ps11:dense} により
$G_5\cong K_{2,2,1}$ は八辺を持つ．
二本の単位枝はその外にあるため，補 graph の二辺に一致しなければならない．
前証明と同じ $q_{ab}+q_{cd}\equiv0\pmod5$ の矛盾が生じる．
また五で割れる辺が七本あれば八本目が強制されるので，上限六を得る．
\end{proof}

\begin{corollary}[一本の単位枝がある密な五進配置]
\label{gp23:one-unit-dense-five}
$G_5$ が八辺を持ち，ある branch の分母が一であるとする．
単位枝はちょうど一本であり，実際の高分母 support は $G_{5,16}$ となる．
$G_5$ の外のもう一本の branch は単位枝と頂点を共有せず，その分母 $q$ は
\[
 q\equiv-1\pmod5,\qquad q\equiv9\pmod{10},\qquad q\ge9
\]
を満たす．
\end{corollary}

\begin{proof}
$G_5$ の補 graph は二つの互いに素な辺である．一方が単位枝なら，
合同分類により他方の分母は $-1$ modulo five となる．
それは一ではなく，さらに局所方程式から分母は正の奇数であるため表示式を得る．
八本の五の倍数とこの一本が高分母で，残る一本だけが単位分母となる．
従って高分母 support は $K_5$ から一辺を除いた $G_{5,16}$ である．
\end{proof}

これらは既に証明した二つの support 分類を組み合わせた系であり，
合同条件だけから実際の整数接着の存在を結論するものではない．
局所模型単独には共通因子を持つ例があるため，原始性を局所分類の仮定として
先に課してはいけない．

% v10: general triangle-completion test and a symbolic infinite obstruction family.
% Requires q9:four-line-criterion, mg8:end-cancellation and the standard macros.
\section{三角形の第三 corner による整数接着の障害}
\label{qt10:section}

一つの共有 branch の両端で分母，分岐係数，puncture clutching が整合しても，
その三つの punctures が属する三角形全体で整合するとは限らない．
全平面五次では各共有直線上の meridian node 数が五であることを用い，
第三 corner の必要条件を分母の大きさによらない有限個の整数式にできる．

\begin{lemma}[第三 corner の有限な必要条件]\label{qt10:third-corner}
三つの異なる core labels を $i,j,k$ とする．
core $i,j$ の整数作用データが既に与えられているとし，
\[
 A=q_{ik},\qquad B=q_{jk},\qquad
 r=p_{i;kj},\qquad s=p_{j;ki}
\]
と置く．これらを全平面五次の大域 atlas に接続できるためには，
\begin{equation}\label{qt10:completion}
 (Au-r)(Bv-s)=AB-1
\end{equation}
を満たす $u,v\in\{-4,-3,-2,-1,0,1\}$ が必要である．
\end{lemma}

\begin{proof}
第三 core $k$ で未定の numerators を
$e=p_{k;ij}$，$f=p_{k;ji}$ と書く．
共有直線 $ik,jk$ の puncture clutching は
\[
 e=A(1-n_{ik;j})-r,\qquad
 f=B(1-n_{jk;i})-s
\]
を要求する．各直線の全 meridian node 数は五で，その各 basin の個数は非負なので，
$0\le n_{ik;j},n_{jk;i}\le5$ である．
$u=1-n_{ik;j}$，$v=1-n_{jk;i}$ と置けば指定した六値の範囲となる．
core $k$ の crossing determinant $AB-ef=1$ が \eqref{qt10:completion} を与える．
\end{proof}

本補題は第三 core の他の二つの分母，weights，作用切片を用いない．
従って候補となる二つの端点について，三つの punctures のそれぞれに適用できる．
逆に，三つ全てがこの検査を通っても，第三 core の全データや大域 atlas の存在は従わない．

\subsection{単独の共有辺では整合する無限族}

\begin{proposition}[隣接する二高分母の局所族]\label{qt10:ladder-family}
整数 $n\ge1$ に対し
\[
 Q_n=n^2+n+1,\qquad R_n=n^2-n+1=Q_{n-1}
\]
と置く．次のデータは正の局所四直線作用 chart $L_n$ を与える：
\[
 q=(Q_n,1,1,R_n),\qquad
 P_n=\begin{pmatrix}
 0&-(n+1)&1&-(n^2+1)\\
 -n&0&0&n-1\\
 n(n+1)&-(2n+1)&0&-n(n-1)\\
 -n^2&n&-1&0
 \end{pmatrix},
\]
\[
 w=\left(Q_n-\frac12,\ 1,\ 1,\ n-1+\frac1{2n}\right).
\]
その分岐係数は $(-Q_n,2n+1,-1,R_n)$ である．
\end{proposition}

\begin{proof}
$Q_n,R_n$ は正の奇数である．全 row sums は $-q_i$ であり，
六つの pair products は $q_iq_j-1$ となる．
例えば $Q_nR_n-1=n^2(n^2+1)$ である．
正の numerators に対応する basin areas は
\[
 B_{02}=\frac1{2Q_n},\qquad B_{20}=\frac12,
 \qquad B_{13}=\frac1{2n},\qquad B_{31}=\frac1{2R_n}
\]
である．$n=1$ では $p_{13}=0$ だが，同じ正値の公式を用いてよい．
残る numerators は非正なので，その basin areas も正である．
命題~\ref{q9:four-line-criterion} を適用する．
各 column sum から記した分岐係数を得る．
\end{proof}

$L_n$ の分母 $Q_n$ の端と，$L_{n+1}$ の分母 $R_{n+1}=Q_n$ の端を比較する．
その分岐係数は $-Q_n,+Q_n$ なので相殺する．
両端の三つの outgoing numerators はそれぞれ
\[
 (-(n+1),\ 1,\ -(n^2+1)),
 \qquad (-(n+1)^2,\ n+1,\ -1).
\]
対応を，最初の列の順に第二列の第二，第三，第一 entry へ取ると，
両端の和は $(0,0,-2Q_n)$ となる．従って node 配分は $(1,1,3)$ であり，
全て非負で合計五となる．incoming numerators に対しても，同じ順の和は
$(0,0,-2Q_n)$ であり，三枚の face の star shears は $(0,0,2)$ となる．
その和二は $K_4$ の ambient star relation と整合する．
ここで確認したのは整数データの接着条件であり，actual weights や切片を一致させた
二つの幾何的 neck germs の接着を主張するものではない．

\begin{theorem}[この無限族の隣接端は三角形へ延長しない]\label{qt10:ladder-no-triangle}
$n\ge2$ とする．$L_n$ と $L_{n+1}$ の上記の分母 $Q_n$ の端は，
単独の branch の整数条件を満たすが，現在の全平面五次の大域作用 atlas には同時に現れない．
\end{theorem}

\begin{proof}
$L_n$ 側の outgoing residues は $Q_n$ を法として $\{1,n,n^2\}$ である．
$n\ge2$ では三つは相異なるので，puncture clutching に適合する対応は上記のものに限られる．
従って $L_n$ のもう一つの高分母 $R_n$ に向かう puncture と，
$L_{n+1}$ のもう一つの高分母 $Q_{n+1}$ に向かう puncture は同じ第三 core $k$ を指す．

この triangle について
\[
 R=n^2-n+1,\quad S=n^2+3n+3,\quad C=n^2+2n+2
\]
と置くと，補題~\ref{qt10:third-corner} は
\begin{equation}\label{qt10:ladder-third-corner}
 (Ru+n^2)(Sv+C)=RS-1=:M=n^2(n+1)^2+2
\end{equation}
を要求する．$u,v\in\{-4,\ldots,1\}$ の全場合を次の不等式で排除する．

$u,v\ge0$ なら左辺は $n^2C$ 以上であり，$n^2C-M=n^2-2>0$ となる．
$u=-1,v\ge0$ なら，左辺は正で $(n-1)(S+C)$ 以下であり，
$M-(n-1)(S+C)=n^4-2n^2+7>0$ となる．
$u\ge-1,v\le-1$ または $u\le-2,v\ge0$ なら，二因子の符号が異なるので左辺は負である．

残る場合には $a=-u\in\{2,3,4\}$，$b=-v\in\{1,2,3,4\}$ と書く．
左辺は正の二因子の積 $(aR-n^2)(bS-C)$ であり，$a,b$ に関して厳密に増加する．
$b=1$ ならその最大値でも
\[
 M-(4R-n^2)(n+1)=n^4-n^3+2n^2-2>0.
\]
$(a,b)=(2,2)$ でも
\[
 M-(2R-n^2)(2S-C)=3n^2-6>0.
\]
他の全ての組は $(3,2)$ または $(2,3)$ 以上であり，
\begin{align*}
 (3R-n^2)(2S-C)-M&=n^4+3n^3-2n^2+10>0,\\
 (2R-n^2)(3S-C)-M&=n^4+n^3-4n^2+12>0.
\end{align*}
最後の二つも $n\ge2$ で厳密に正である．
従って \eqref{qt10:ladder-third-corner} は解を持たず，矛盾する．
\end{proof}

この無限族は，分母を大きくすると局所の高分母 chart と単独 branch の整数整合を
いくらでも作れることを示す．同時に，それらから大域存在を推測する際には，
同じ puncture が三つの cores に属することと，全 node budget を保持する必要がある．

% v10 computer-assisted finite exclusion, with complete enumeration proof.
\section{分母百一以下における大域排除の厳密証明書}
\label{ex10:section}

一般の高分母に対する大域存在問題は残るが，固定された有限の分母集合については，
numerators に探索上限を置かずに，全ての必要な局所候補を列挙できる．
以下の結論は明示した有限範囲のものであり，任意分母への外挿を含まない．

\begin{theorem}[有限分母範囲の排除]\label{ex10:bounded-exclusion}
現在の全平面五次の SNC core models，標準 $K_4$ cones，既存の branch graph を用い，
追加の branch singularities を入れない大域作用適合 atlas を考える．
全ての branch denominators が百一以下なら，それらは全て一である．
従って高分母を持つ大域 atlas が存在するとすれば，少なくとも一つの分母は百三以上である．
\end{theorem}

\begin{proof}
各分母は奇数なので，調べる集合は
$\mathcal Q=\{1,3,5,\ldots,101\}$ である．
まず二つ以上の高分母を持つ一つの core を考える．
標識を替えて $q_0,q_3>1$ とし，$N_{ij}=q_iq_j-1$ と置く．
三つの signed divisors
\[
 a\mid N_{01},\qquad b\mid N_{02},\qquad d\mid N_{13}
\]
を選ぶだけで，次の順に全 numerators が決まる．
\begin{align*}
 p_{01}&=a,&p_{10}&=N_{01}/a,&
 p_{02}&=b,&p_{20}&=N_{02}/b,\\
 p_{03}&=-q_0-a-b,&p_{30}&=N_{03}/p_{03},&
 p_{13}&=d,&p_{31}&=N_{13}/d,\\
 p_{32}&=-q_3-p_{30}-p_{31},&p_{23}&=N_{23}/p_{32},\\
 p_{12}&=-q_1-p_{10}-p_{13},&p_{21}&=-q_2-p_{20}-p_{23}.
\end{align*}
途中の分母が零または整除性を満たさない場合と，最後の
$p_{12}p_{21}=N_{12}$ が成立しない場合を除く．
全 row sums と pair products を満たす行列は必ずこの順序で再現されるので，
これは固定された $q$ に対して完全な有限列挙である．
特に $N_{12}=0$ の場合にも，$p_{12},p_{21}$ は他の row sums から既に決まり，
自由な整数を任意の範囲で打ち切る必要はない．

正の weights の存在は補題~\ref{kl10:cycle-positive} の全 simple directed cycles の
有理数積で厳密に判定する．正値性を満たす場合には，同補題の方法で有理 weights を構成し，
全 basin inequalities を再確認する．

高分母が一つしかない core では，低分母 vertices の numerators は有界とは限らない．
そこで全行列を列挙せず，命題~\ref{qh9:isolated-local} が与える高分母端の
全ての必要な normal forms を候補として残す．
これは候補集合を広げる方向の操作であり，大域排除に対して安全である．

各高分母端には，分母 $q$，係数 $D$，三 numerators の residue multiset
$\{p_1,p_2,p_3\}\bmod q$ を付ける．
対向する端は，同じ $q$，係数 $-D$，負の residue multiset を持たなければならない．
候補の中にその相手がない端を一つでも持つ core candidate を除く操作を繰り返す．
punctures の全ての置換を許し，共有 weights や三角形の incidence の条件もこの段階では課さないので，
実際の大域 atlas から来る候補が誤って除かれることはない．

収録した標準 Python の厳密検算は，二つ以上の高分母を持つ $316\,200$ 個の
分母多重集合を扱い，$2\,973$ 個の整数行列候補，そのうち $1\,864$ 個の正の局所候補を得る．
さらに一高分母の必要な端点 normal forms $87$ 個を加えた $1\,951$ 個の候補は，
上の相手の不存在による削除で
\[
1951\longrightarrow118\longrightarrow45\longrightarrow28\longrightarrow17
\longrightarrow11\longrightarrow7\longrightarrow4\longrightarrow1\longrightarrow0
\]
となる．従って有限範囲内では，高分母を含む大域 atlas の局所データは存在しない．
\end{proof}

この列挙の件数は，分母の標識を正規化した有限な整数データの件数である．
非同型な幾何的 charts や atlases の個数ではない．
検算は \texttt{verification/verify\_quintic\_denominators\_v10.py} で再生成でき，
候補データ，削除の理由，実行範囲とコード hash を対応する JSON に記録する．
proof の整数式と正値性判定はすべて厳密であり，浮動小数点 tolerance や numerators の上限を用いない．

定理の証明に使う削除は，node 数の非負性や全 node budget を必要としない．
同じ候補にそれらと補題~\ref{qt10:third-corner} を追加すると，
単独 branch の整数条件を満たす $835$ 個の標識付き端点の非順序対のうち，
三つ全ての第三 corner の検査を通るものは $12$ 個に減り，
core candidates は二回の削除で $1951\to4\to0$ となる．
この追加検査は一般の高分母候補を調べる際にも使えるが，
百一を超える全ての分母を排除する一般証明は与えていない．

\begin{thebibliography}{99}

\bibitem{EM}
J. D. Evans and M. Mauri,
\emph{Constructing local models for Lagrangian torus fibrations},
Annales Henri Lebesgue \textbf{4} (2021), 537--570.
Lemmas~2.16, 2.22, 3.1; Remarks~2.7, 2.19, 2.23; \S3.1;
Assumption~6.6, Proposition~6.13 and Theorem~6.15.
\url{https://doi.org/10.5802/ahl.80};
\url{https://arxiv.org/abs/1905.09229v3}.

\bibitem{FernandezBobadillaPelka}
Javier Fern\'andez de Bobadilla and Tomasz Pe\l ka,
\emph{Fibrations by Lagrangian tori for maximal Calabi--Yau degenerations and beyond},
Communications on Pure and Applied Mathematics (2026), e70060.
Published online 8 July 2026.
Theorems~1.1--1.2 and~6.14; \S8.
The linked arXiv version is v4, 9 July 2026.
\url{https://doi.org/10.1002/cpa.70060};
\url{https://arxiv.org/abs/2312.13248v4};
\url{https://arxiv.org/html/2312.13248v4}.

\bibitem{GrossSiebert}
M. Gross and B. Siebert,
\emph{From real affine geometry to complex geometry},
Annals of Mathematics (2) \textbf{174} (2011), no.~3, 1301--1428.
\url{https://doi.org/10.4007/annals.2011.174.3.1}.

\bibitem{LeungWan}
N. C. Leung and T. Y. H. Wan,
\emph{Calabi--Yau components in general type hypersurfaces},
Journal of Differential Geometry \textbf{83} (2009), no.~1, 43--74.
\S\S1 and 3.1.
\url{https://doi.org/10.4310/jdg/1253804351};
\url{https://arxiv.org/abs/0807.1784}.

\bibitem{Mikhalkin}
G. Mikhalkin,
\emph{Decomposition into pairs-of-pants for complex algebraic hypersurfaces},
Topology \textbf{43} (2004), no.~5, 1035--1065.
Theorems~1 and~1$'$; \S4.2.
\url{https://doi.org/10.1016/j.top.2003.11.006};
\url{https://arxiv.org/abs/math/0205011v3}.

\bibitem{Zhang}
Y. Zhang,
\emph{Pair-of-Pants Decompositions of 4-Manifolds Diffeomorphic to General Type Hypersurfaces},
Journal of Mathematical Study \textbf{58} (2025), no.~2, 211--252.
Theorem~1.1; the quintic example in \S1.
\url{https://doi.org/10.4208/jms.v58n2.25.06};
\url{https://arxiv.org/abs/2102.10037}.

\bibitem{FeichtnerSturmfels}
Eva Maria Feichtner and Bernd Sturmfels,
\emph{Matroid polytopes, nested sets and Bergman fans},
Portugaliae Mathematica (N.S.) \textbf{62} (2005), no.~4, 437--468.
\href{https://arxiv.org/abs/math/0411260}{arXiv:math/0411260}.

\bibitem{ArdilaKlivans}
Federico Ardila and Caroline J. Klivans,
\emph{The Bergman complex of a matroid and phylogenetic trees},
Journal of Combinatorial Theory, Series B \textbf{96} (2006), no.~1, 38--49.
\href{https://doi.org/10.1016/j.jctb.2005.06.004}{doi:10.1016/j.jctb.2005.06.004}.
\href{https://arxiv.org/abs/math/0311370}{arXiv:math/0311370}.

\bibitem{ShawSurfaces}
Kristin Shaw,
\emph{Tropical surfaces},
arXiv:1506.07407v1 (2015), 43 pp.
Definition~3.1, pp.~15--16 (ambient integer affine transitions);
\S3.2 (boundary), \S4.2 (tropical summation).
\url{https://arxiv.org/abs/1506.07407v1};
\url{https://arxiv.org/pdf/1506.07407}.

\bibitem{JellRauShaw}
Philipp Jell, Johannes Rau, and Kristin Shaw,
\emph{Lefschetz $(1,1)$-theorem in tropical geometry},
\'Epijournal de G\'eom\'etrie Alg\'ebrique \textbf{2} (2018), Article~11, 27 pp.
Definitions~2.1, 3.1 and 5.1 (rational polyhedral spaces, the affine sheaf, and matroidal charts).
\url{https://doi.org/10.46298/epiga.2018.volume2.4126};
\url{https://epiga.episciences.org/4981/pdf};
\url{https://arxiv.org/abs/1711.07900}.

\bibitem{CartwrightComplexes}
Dustin Cartwright,
\emph{Tropical complexes},
Manuscripta Mathematica \textbf{163} (2020), 227--261.
Proposition~2.1; Lemma~2.2; Definitions~2.5, 2.8;
Proposition~2.10; Construction~3.2 and Remark~3.4.
The linked arXiv version is v4 (2019).
\url{https://doi.org/10.1007/s00229-019-01148-3};
\url{https://arxiv.org/abs/1308.3813v4};
\url{https://arxiv.org/pdf/1308.3813}.

\bibitem{SergeevSchneiderButkovic}
S. Sergeev, H. Schneider and P. Butkovi\v c,
\emph{On visualization scaling, subeigenvectors and Kleene stars in max algebra},
Linear Algebra Appl. \textbf{431} (2009), no.~12, 2395--2406.
\url{https://doi.org/10.1016/j.laa.2009.03.040};
\url{https://arxiv.org/abs/0808.1992};
\url{https://web.mat.bham.ac.uk/P.Butkovic/My%20papers/VisSSB15hps.pdf}.
Theorem~2.2 and Proposition~2.5 in the linked author manuscript (24 March 2009).
% Theorem 2.2 and Proposition 3.4 checked in arXiv:0808.1992v3.

\bibitem{CorralesValencia}
H. Corrales and C. E. Valencia,
\emph{Arithmetical structures on graphs},
Linear Algebra Appl. \textbf{536} (2018), 120--151.
\url{https://doi.org/10.1016/j.laa.2017.09.018};
\url{https://arxiv.org/abs/1604.02502}.
% Theorems 3.2 and 3.7 checked in arXiv:1604.02502v3.
% Published numbering differs at the finiteness result; cite section or version.

\bibitem{CuntzHolmFriezes}
M. Cuntz and T. Holm,
\emph{Frieze patterns over integers and other subsets of the complex numbers},
Journal of Combinatorial Algebra \textbf{3} (2019), no.~2, 153--188.
\S\S2--3 and 7 (tame friezes, quiddity cycles, finiteness and the integral model).
\url{https://doi.org/10.4171/JCA/29};
\url{https://arxiv.org/abs/1711.03724v2}.

\bibitem{CuntzHolmPagano}
M. Cuntz, T. Holm and C. Pagano,
\emph{Frieze patterns over algebraic numbers},
Bulletin of the London Mathematical Society \textbf{56} (2024), no.~4, 1417--1432.
\S\S3--4 (finiteness over algebraic numbers and signs of integral friezes).
\url{https://doi.org/10.1112/blms.13003};
\url{https://arxiv.org/abs/2306.12148v2}.

\bibitem{LinSturmfels}
S. Lin and B. Sturmfels,
\emph{Polynomial relations among principal minors of a $4\times4$-matrix},
J. Algebra \textbf{322} (2009), 4121--4131.
\href{https://doi.org/10.1016/j.jalgebra.2009.06.026}{doi:10.1016/j.jalgebra.2009.06.026}.
\href{https://arxiv.org/abs/0812.0601}{arXiv:0812.0601v3}.

\bibitem{SurfaceGeometry}
\emph{一般型射影曲面のLagrangian torus fibrationと作用適合matroidal atlas：相対SNC構成、次数分割と分岐底空間、大域作用座標の接着}．未発表草稿，第25稿，2026年10月28日．

\end{thebibliography}

\end{document}
